\documentclass[11pt,a4paper]{article}

\usepackage[margin=2.0cm]{geometry}
\usepackage{amsmath,amssymb,amsthm}
\usepackage{booktabs}
\usepackage{array}
\usepackage{tabularx}
\usepackage[shortlabels]{enumitem}
\usepackage{xcolor}
\usepackage{tikz}
\usepackage{fancyhdr}
\usepackage{setspace}
\usepackage{microtype}
\usepackage{titlesec}
\usepackage{hyperref}
\usepackage{listings}
\usepackage[skins,breakable]{tcolorbox}

\usetikzlibrary{positioning,arrows.meta,calc,fit,shapes.geometric,matrix}

\definecolor{ink}{HTML}{1E1E1E}
\definecolor{rulegray}{HTML}{444444}
\definecolor{boxline}{HTML}{555555}
\definecolor{boxfill}{HTML}{F5F5F5}
\definecolor{boxtitle}{HTML}{E6E6E6}
\definecolor{solline}{HTML}{888888}
\definecolor{soltitle}{HTML}{EFEFEF}
\definecolor{g1}{HTML}{D0D0D0}
\definecolor{g2}{HTML}{8A8A8A}
\definecolor{g3}{HTML}{555555}

\hypersetup{
  colorlinks=true,
  linkcolor=ink,
  urlcolor=ink,
  pdftitle={MTL 458 Major Exam Question Bank (CS-537 \& ISI Kolkata)},
  pdfauthor={Compiled for MTL 458, IIT Delhi}
}

\setstretch{1.10}
\raggedbottom
\setlist[itemize]{leftmargin=1.3em,itemsep=0.22em,topsep=0.3em}
\setlist[enumerate]{leftmargin=1.5em,itemsep=0.25em,topsep=0.32em}

\titleformat{\part}[display]
  {\centering\Huge\bfseries\color{ink}}
  {\vspace{0.5em}\Large\MakeUppercase{\partname~\thepart}}
  {0.4em}
  {\rule{\textwidth}{0.8pt}\\[0.5em]}
  [\vspace{0.3em}\rule{\textwidth}{0.8pt}]

\titleformat{\section}
  {\Large\bfseries\color{ink}}{\thesection.}{0.55em}{}
  [{\color{rulegray}\titlerule[0.5pt]}]
\titleformat{\subsection}
  {\large\bfseries\color{ink}}{\thesubsection}{0.5em}{}
\titleformat{\subsubsection}
  {\normalsize\bfseries\color{ink}}{\thesubsubsection}{0.45em}{}

\titlespacing*{\section}{0pt}{1.3em}{0.65em}
\titlespacing*{\subsection}{0pt}{1.0em}{0.4em}

\pagestyle{fancy}
\fancyhf{}
\renewcommand{\headrulewidth}{0.4pt}
\renewcommand{\footrulewidth}{0.3pt}
\fancyhead[L]{\small\color{ink}MTL 458 \textbullet\ Operating Systems}
\fancyhead[R]{\small\color{ink}Major Exam Question Bank}
\fancyfoot[C]{\small\thepage}
\setlength{\headheight}{14pt}

% ---------------------------------------------------------------------------
% Boxes & Code Styling
% ---------------------------------------------------------------------------
\tcbset{
  qboxstyle/.style={
    colback=boxfill, colframe=boxline, boxrule=0.55pt,
    left=6pt, right=6pt, top=5pt, bottom=5pt,
    fonttitle=\bfseries\small, coltitle=black, colbacktitle=boxtitle,
    enhanced, sharp corners=south, breakable
  },
  sboxstyle/.style={
    colback=white, colframe=solline, boxrule=0.45pt,
    left=6pt, right=6pt, top=5pt, bottom=5pt,
    fonttitle=\bfseries\small, coltitle=ink, colbacktitle=soltitle,
    enhanced, sharp corners=south, breakable
  }
}

\newtcolorbox{keybox}[1][]{qboxstyle, title={#1}}

\lstset{
  language=C,
  basicstyle=\ttfamily\footnotesize,
  keywordstyle=\bfseries,
  commentstyle=\itshape\color{g3},
  numbers=none,
  frame=single, rulecolor=\color{g1},
  backgroundcolor=\color{white},
  xleftmargin=0.6em, framexleftmargin=0.4em,
  columns=fullflexible, keepspaces=true,
  showstringspaces=false, breaklines=true,
  aboveskip=0.5em, belowskip=0.5em,
  morekeywords={cond_t,mutex_t,sem_t,lock_t,rwlock_t,pthread_t,pthread_cond_t,pthread_mutex_t,uint32_t,size_t,off_t,ssize_t,pid_t}
}

% Keep at least #1 space before starting a new question
\makeatletter
\newcommand{\keepspace}[1]{\par\penalty-100\begingroup
  \setlength{\dimen@}{#1}\dimen@ii\pagegoal\advance\dimen@ii-\pagetotal
  \ifdim\dimen@>\dimen@ii\ifdim\dimen@ii>\z@\vfil\fi\break\fi\endgroup}
\makeatother

\newcounter{qn}[section]

% \begin{Qn}{Difficulty}{Lecture}{Source}{Title}
\newenvironment{Qn}[4]{%
  \refstepcounter{qn}\keepspace{7\baselineskip}%
  \begin{tcolorbox}[qboxstyle, title={Q\thesection.\theqn\ \ [#1 \textbar\ #2]\ \ #4}]%
  {\footnotesize\color{g3}\textit{Source: #3}}\par\smallskip%
}{%
  \end{tcolorbox}%
}

\newenvironment{Sol}{%
  \begin{tcolorbox}[sboxstyle, title={Solution}]%
}{%
  \end{tcolorbox}\vspace{0.35em}%
}

\begin{document}

% ---------------------------------------------------------------------------
% Title Header
% ---------------------------------------------------------------------------
\begin{center}
{\color{rulegray}\rule{\textwidth}{1.2pt}}\\[0.5em]
{\LARGE\bfseries Operating Systems (MTL 458)}\\[0.25em]
{\Large\bfseries Major Exam Comprehensive Question Bank}\\[0.35em]
{\normalsize Compiled Strictly from the Two Matched Minor Exam Sources:}\\
{\normalsize\textbf{UW--Madison CS-537 Old Exams (2001--2022)} \textperiodcentered\ \textbf{ISI Kolkata OS Past Papers (Prof.\ Mandar Mitra, 2010--2024)}}\\[0.35em]
{\color{rulegray}\rule{\textwidth}{1.2pt}}
\end{center}

\vspace{0.35em}
\begin{keybox}[How to Use This Question Bank]
\begin{itemize}
  \item \textbf{Priority-Ordered Structure:}
    \begin{itemize}
      \item \textbf{Part I (50 Problems --- Essential Core, Cover First):} The 50 highest-yield, completely \emph{non-overlapping} problems spanning every core skill (\textbf{44 Post-Minor [23 Concurrency + 10 VM + 11 Persistence]} + \textbf{6 Pre-Minor}). Ideal when time is tight.
      \item \textbf{Part II (53 Problems --- Extended Challenging Set, Cover Next):} The remaining 53 verified high-priority analytical, code-writing, and tricky problems (\textbf{45 Post-Minor [25 Concurrency + 9 VM + 11 Persistence]} + \textbf{8 Pre-Minor}).
      \item \textbf{Part III (240 Problems --- Additional Practice, MCQs \& Drills):} All remaining deduplicated syllabus-relevant problems, MCQs, true/false sets, and conceptual drills (\textbf{137 Post-Minor} + \textbf{103 Pre-Minor}) for extra practice if time permits.
    \end{itemize}
  \item \textbf{Strict Source \& Syllabus Fidelity:} Every single question (343 total) is lifted from an actual \textbf{UW--Madison CS-537} or \textbf{ISI Kolkata (Mandar Mitra)} exam paper, with exact source citations. Out-of-syllabus topics (Lottery/Stride/CFS, Linux 2.4/2.6 scheduler internals, Multi-CPU scheduling, Working Set/WSClock/PFF, Banker's algorithm, FFS cylinder groups, Journaling/FSCK/LFS, Flash/SSDs, NFS/AFS, SysV IPC, real-time scheduling) have been strictly excluded.
  \item \textbf{Format \& Solutions:} Written problems are preserved as written problems, and MCQs/T-F blocks retain their original options. A step-by-step \textbf{Solution} box (verified against official CS-537 keys where published and derived from first principles otherwise) appears immediately below each question.
\end{itemize}
\end{keybox}

\vspace{0.3em}
\begin{center}
\small
\begin{tabularx}{\textwidth}{@{}lXrrrr@{}}
\toprule
\textbf{Block} & \textbf{Lectures \& Topics} & \textbf{Part I} & \textbf{Part II} & \textbf{Part III} & \textbf{Total} \\
\midrule
1 & \textbf{Lect 13--15:} Virtual Memory (TLB, Adv.\ Paging, Swapping) & 10 & 9 & 49 & 68 \\
2 & \textbf{Lect 16--20:} Concurrency (Threads, Locks, CVs, Semaphores, Deadlock) & 23 & 25 & 51 & 99 \\
3 & \textbf{Lect 21--24 (incl.\ 22A):} Persistence (I/O, HDDs, RAID, File Systems) & 11 & 11 & 37 & 59 \\
\midrule
\multicolumn{2}{@{}l}{\textbf{Post-Minor Subtotal (Lectures 13--24)}} & \textbf{44} & \textbf{45} & \textbf{137} & \textbf{226} \\
\midrule
4 & \textbf{Lect 1--7 (Pre-Minor):} CPU Virtualization (Processes, Traps, Scheduling) & 3 & 5 & 61 & 69 \\
5 & \textbf{Lect 8--12 (Pre-Minor):} Early Memory Virt.\ (B\&B, Seg., Free Space, Paging) & 3 & 3 & 42 & 48 \\
\midrule
\multicolumn{2}{@{}l}{\textbf{Pre-Minor Subtotal (Lectures 1--12)}} & \textbf{6} & \textbf{8} & \textbf{103} & \textbf{117} \\
\midrule
\multicolumn{2}{@{}l}{\textbf{Grand Total (Full Syllabus, Deduplicated)}} & \textbf{50} & \textbf{53} & \textbf{240} & \textbf{343} \\
\bottomrule
\end{tabularx}
\end{center}

\clearpage
\setcounter{tocdepth}{2}
\tableofcontents
\clearpage


% ===========================================================================

% ===========================================================================

% ===========================================================================
\part{Essential Core Problems --- Non-Overlapping Must-Do Set (Cover First)}
% ===========================================================================

\section{Lectures 13--15: Virtual Memory --- TLB, Advanced Paging \& Swapping}

\subsection{Lectures 13--14: Paging with TLB \& Advanced Paging}

\begin{Qn}{Hard}{Lect 13}{CS-537 Sp10 Final, P4}{Forcing a TLB miss on (almost) every access}
TLB misses can be nasty. The following code can cause a lot of TLB misses,
depending on the values of \texttt{STRIDE} and \texttt{MAX}. Assume that the
system has a \textbf{32-entry} TLB and an \textbf{8\,KB} page size.
\begin{lstlisting}[numbers=none]
int value = 0;
int data[MAX];   /* a big array */

for (int j = 0; j < 1000; j++) {
    for (int i = 0; i < MAX; i += STRIDE) {
        value = value + data[i];
    }
}
\end{lstlisting}
Assume \texttt{sizeof(int) = 4} and that \texttt{data} is page-aligned.
There is enough physical memory that a modestly sized array need not be
swapped.
\begin{enumerate}[(a)]
\item What should you set \texttt{MAX} and \texttt{STRIDE} to so that you
      achieve a TLB miss (but \emph{not} a page fault) upon pretty much
      every access to the array \texttt{data}? Describe the constraints.
\item What happens if \texttt{MAX} is too high? Too low?
\item What happens if \texttt{STRIDE} is too high? Too low?
\end{enumerate}
\end{Qn}
\begin{Sol}
An access misses in the TLB when its VPN is not among the 32 cached
translations. Because the outer loop repeats the inner scan 1000 times,
any working set of \textbf{at most 32 pages} will, after the first pass,
produce almost all hits. To keep missing we need a working set of
\textbf{strictly more than 32 distinct pages}, while still remaining
resident in DRAM (no page faults).

Each page holds $8192/4 = 2048$ integers. One access per page is obtained
by striding a full page of integers:
\[
\texttt{STRIDE} \ge 2048.
\]
The number of distinct pages touched is
\[
\left\lceil\frac{\texttt{MAX}}{\texttt{STRIDE}}\right\rceil
\]
(exactly $\texttt{MAX}/\texttt{STRIDE}$ when both are multiples of the
page). We therefore need
\[
\frac{\texttt{MAX}}{\texttt{STRIDE}} \ge 33.
\]
A canonical choice is $\texttt{STRIDE}=2048$ and $\texttt{MAX}=33\times 2048$
(or any larger multiple that still fits in physical memory). LRU then
thrashes: the 33rd page evicts the 1st, and the next pass again misses on
every page. Random replacement would miss on most, but not necessarily
all, accesses; LRU on a cyclic working set of size $\mathrm{TLB}+1$ misses
on \emph{every} access.

\begin{enumerate}[(a)]
\item \texttt{STRIDE} equal to (at least) the number of \texttt{int}s per
      page, and $\texttt{MAX}/\texttt{STRIDE} > 32$, with the array still
      fitting in RAM.
\item Too high: the array no longer fits in physical memory, so most
      accesses become \textbf{page faults} (disk) rather than mere TLB
      misses. Too low: $\le 32$ pages are touched, so after the first
      outer-loop iteration the TLB holds the whole working set and the
      remaining $\approx 999$ scans are hits.
\item Too high: for a fixed \texttt{MAX} you touch \emph{fewer} pages
      (you skip whole pages), so the working set may drop to $\le 32$ and
      the TLB starts hitting. Too low: several accesses land in the same
      page, so only the first of each run of about
      $2048/\texttt{STRIDE}$ references ($\lceil\cdot\rceil$ or
      $\lfloor\cdot\rfloor$ depending on alignment) misses;
      the rest are spatial-locality hits.
\end{enumerate}
This is the lecture's LRU-thrash example (5 pages, 4 TLB slots, cyclic
access $\Rightarrow$ 100\% miss rate) scaled to a 32-entry TLB, and it is
also why \textbf{TLB reach} $=(\#\textrm{entries})\times(\textrm{page size})$
matters: here reach is $32\times 8\,\mathrm{KB}=256\,\mathrm{KB}$, so any
array scan larger than that, at one-touch-per-page, thrashes the TLB.
\end{Sol}

\begin{Qn}{Hard}{Lect 14}{CS-537 Sp02 Final, Q5}{MiniTable: a na\"ive inverted page table}
TinyOS on a handheld: 1\,KB pages, 128\,KB of physical memory, 20-bit
virtual and physical addresses.
\begin{enumerate}[(a)]
\item Linear page table. Each PTE stores the translation plus 2
      protection bits and 4 miscellaneous bits. How much memory does
      \emph{one process's} page table occupy?
\item \textbf{MiniTable}: a linear array with one entry per physical
      page. Each entry stores a 6-bit PID and the VPN that maps to that
      frame. How much space does MiniTable occupy?
\item Time overhead of a VA$\to$PA lookup in MiniTable versus the linear
      array of (a). State your assumptions.
\item MiniTable does not easily allow two processes to share a page.
      Why, and how would you fix it?
\end{enumerate}
\end{Qn}
\begin{Sol}
Page size $2^{10}$, so a 20-bit VA has a 10-bit VPN ($2^{10}=1024$ virtual
pages) and a 10-bit offset. Physical memory is $128\,\mathrm{KB}/1\,\mathrm{KB}
= 128 = 2^7$ frames, so a PPN is 7 bits.
\begin{enumerate}[(a)]
\item PTE payload: 7-bit PPN $+$ 2 protection $+$ 4 misc $= 13$ bits,
      stored in 2 bytes (byte-aligned; bit-packing would take
      $1024\times 13\text{ bits}=1664\,\mathrm{B}$, and reserving a full
      10-bit PPN for 20-bit physical addresses gives $10+2+4=16$ bits,
      still 2 bytes). Linear table:
      $1024\times 2\,\mathrm{B} = 2\,\mathrm{KB}$
      \emph{per process}. This table must be contiguous.
\item One entry per frame: 128 entries. PID (6) $+$ VPN (10) $= 16$ bits
      $= 2$ bytes, so MiniTable is
      $128\times 2\,\mathrm{B} = 256\,\mathrm{B}$
      \emph{for the whole machine} --- independent of the number of
      processes.       This is a na\"ive \textbf{inverted page table}:
      $\mathrm{PPN}\mapsto(\mathrm{VPN},\mathrm{ASID/PID})$ (with a
      reserved PID marking free frames).
\item Linear array: $O(1)$ (index by VPN). MiniTable: the hardware/OS
      has the VPN and must find the matching $(\mathrm{VPN},\mathrm{PID})$
      row --- a scan of all 128 frames, $O(\#\textrm{frames})$. A
      \textbf{hashed inverted table} (PowerPC-style) reduces this to
      expected $O(1)$, which is why inverted tables pair naturally with
      a \textbf{software-managed TLB} (the OS may use any data structure
      on a miss). A hardware walker wants a tree it can parse in MMU
      logic.
\item Each frame has a \emph{single} $(\mathrm{PID},\mathrm{VPN})$ slot,
      so it cannot record that process A \emph{and} process B both map
      that frame (shared library, copy-on-write, \ldots). Fixes: store a
      list/chain of $(\mathrm{PID},\mathrm{VPN})$ pairs per frame; keep a
      separate sharing table; or hold a reference count plus one
      distinguished owner and look up sharers elsewhere.
\end{enumerate}
\end{Sol}

\begin{Qn}{Hard}{Lect 14}{ISI Mandar 2017 End-sem, Q3(a); cf.\ 2011 End-sem, Q3(b)}{Min/max page size for exactly two paging levels}
An OS uses paging. Logical addresses are $n$ bits; each PTE is 4 bytes.
In a 2-level scheme the second-level page table must fit in a single
page. Compute the \textbf{minimum} and \textbf{maximum} page sizes for
which \emph{exactly} 2 levels of paging are required.
\begin{enumerate}[(a)]
\item $n=28$ (2017).
\item $n=24$ (2011; same template).
\end{enumerate}
\end{Qn}
\begin{Sol}
Let the page size be $2^{p}$ bytes, $p\ge 2$ (a page must hold at least
one 4-byte PTE). VPN bits $= n-p$. A leaf page-table page holds
$2^{p}/4 = 2^{p-2}$ PTEs, so it consumes $p-2$ VPN bits.

\textbf{One level suffices} iff the whole linear table fits in one page:
\[
2^{n-p}\cdot 4 \le 2^{p} \iff n+2 \le 2p \iff p \ge \left\lceil\frac{n+2}{2}\right\rceil.
\]
\textbf{Two levels suffice} iff the outer directory, after paging the
linear table into page-sized leaves, itself fits in one page. The outer
table has $2^{(n-p)-(p-2)} = 2^{n-2p+2}$ PDEs, size
$2^{n-2p+2}\cdot 4 = 2^{n-2p+4}$, so
\[
n-2p+4 \le p \iff p \ge \left\lceil\frac{n+4}{3}\right\rceil.
\]
Exactly two levels means one level is \emph{not} enough and two levels
\emph{are}:
\[
\left\lceil\frac{n+4}{3}\right\rceil \le p \le
\left\lceil\frac{n+2}{2}\right\rceil - 1.
\]

(a) $n=28$: $\lceil 32/3\rceil=11$ and $\lceil 30/2\rceil=15$, so
$11\le p\le 14$.
\[
p_{\min}=11\ \Rightarrow\ 2\,\mathrm{KB},\qquad
p_{\max}=14\ \Rightarrow\ 16\,\mathrm{KB}.
\]
Check $p=11$: VPN $=17$, leaf index $=9$, directory index $=8$,
directory $=256\times 4=1\,\mathrm{KB} < 2\,\mathrm{KB}$. Linear table
$=2^{17}\times 4=512\,\mathrm{KB} \not\le 2\,\mathrm{KB}$.
Check $p=14$: linear table $=2^{14}\times 4=64\,\mathrm{KB} > 16\,\mathrm{KB}$;
directory index $=2$ bits, directory $=16\,\mathrm{B}$.
Check $p=15$: linear table $=2^{13}\times 4=32\,\mathrm{KB}=$ one page
(only 1 level). Check $p=10$: directory $=2^{10}\times 4=4\,\mathrm{KB}
> 1\,\mathrm{KB}$ page (needs 3 levels).

(b) $n=24$: $\lceil 28/3\rceil=10$ and $\lceil 26/2\rceil=13$, so
$10\le p\le 12$.
\[
p_{\min}=10\ \Rightarrow\ 1\,\mathrm{KB},\qquad
p_{\max}=12\ \Rightarrow\ 4\,\mathrm{KB}.
\]
\end{Sol}

\begin{Qn}{Hard}{Lect 14}{CS-537 Sp12 Mid, Q10}{Two-level walk on a 15-bit address space}
15-bit virtual address, 32-byte pages, two-level page table. PDE and PTE
format: a valid bit followed by a 7-bit page-frame number (1 byte).
PDBR $= 73$ (decimal). Relevant physical pages (bytes 0..31):
\begin{lstlisting}[numbers=none]
pg   6: 0a 1c 01 14 0b 1a 19 0a 0a 1a 0c 14 02 0c 1c 0c
        15 04 0e 13 17 11 08 05 08 07 04 13 0f 1d 0f 1e
pg  73: a2 d2 97 96 d9 7f 87 b4 b7 f2 f4 82 bf 7f be 93
        e8 9d 99 9e f1 7f 7f b0 d8 da eb b1 81 c3 c2 f6
pg 114: 7f 7f 7f 7f 82 7f 7f 7f 7f 7f 7f 7f 99 7f 7f 7f
        7f 7f 7f 86 7f 7f 7f 7f 7f 7f 8f 7f 7f 7f 7f 7f
\end{lstlisting}
The split is 5 bits (PD index) $+$ 5 bits (PT index) $+$ 5 bits (offset).
Valid bit is the high bit of each entry; PFN is the low 7 bits.
\begin{enumerate}[(a)]
\item Translate VA \texttt{0x1787}. If valid, what byte is loaded?
\item Translate VA \texttt{0x2665}. If valid, what byte is loaded?
\end{enumerate}
\end{Qn}
\begin{Sol}
\textbf{(a)} \texttt{0x1787} in 15 bits is \texttt{001\,0111\,1000\,0111}:
\[
\mathrm{PD}=00101_2=5,\quad
\mathrm{PT}=11100_2=28,\quad
\mathrm{off}=00111_2=7.
\]
PDE $=$ byte 5 of page 73 $= \texttt{0x7f} = 0111\,1111_2$:
\textbf{valid bit 0} $\Rightarrow$ fault. No inner table, no load.
(McFlub-style: an invalid PDE means ``this 32-VPN range is unmapped'',
which is the whole point of multi-level tables.)

\textbf{(b)} \texttt{0x2665} $= \texttt{010\,0110\,0110\,0101}$:
\[
\mathrm{PD}=01001_2=9,\quad
\mathrm{PT}=10011_2=19,\quad
\mathrm{off}=00101_2=5.
\]
PDE $=$ byte 9 of page 73 $= \texttt{0xf2} = 1111\,0010_2$: valid, PFN
$= 1110010_2 = 114$. PTE $=$ byte 19 of page 114 $= \texttt{0x86}
= 1000\,0110_2$: valid, PFN $= 0000110_2 = 6$. Data $=$ byte 5 of
page 6 $= \texttt{0x1a}$. Physical address
$(6\times 32)+5 = 197 = \texttt{0xc5}$. The load returns \texttt{0x1a}.

Memory references on this TLB-less walk: PDE, PTE, data --- three
physical accesses (the lecture's ``TLB miss costs 2 PT refs $+$ data
$= 3$'').
\end{Sol}

\begin{Qn}{Hard}{Lect 13--14}{CS-537 Fa11 Final, Q1}{Worst/best memory-access counts: linear, TLB, two-level, swapping}
A simple load \texttt{mov VirtualAddress, register}. Each memory access
costs $M$; a disk access costs $D$. Count instruction fetch \emph{and}
the data reference (both are translated).
\begin{enumerate}[(a)]
\item Linear page table, \textbf{no} TLB, no swapping. Time?
\item Same system \textbf{with} a TLB. Fastest (best-case) time,
      assuming TLB hits?
\item Two-level page table with a page directory. Every TLB reference
      misses; all pages are in memory. Slowest (worst-case) time?
\item Same as (c), but now some of the referenced pages may be on disk.
      Worst-case time?
\end{enumerate}
\end{Qn}
\begin{Sol}
The load is two virtual references: fetch the instruction, then fetch
the data.
\begin{enumerate}[(a)]
\item Each virtual reference: walk one PTE, then access the byte.
      Two references $\Rightarrow 4M$.
\item Both translations hit in the TLB, then one instruction fetch and
      one data fetch: $2M$. (TLB lookup itself is not a memory access.)
\item Each virtual reference walks PDE then PTE then the byte: $3M$ per
      reference, two references $\Rightarrow 6M$.
      (Lecture: ``TLB miss, 2-level, present: 2 PT refs $+$ data $= 3$''
      per reference.)
\item Assume the page directory is pinned in RAM (since PDBR holds a
      physical address), while leaf page-table pages and the instruction
      and data pages can be swapped out, the TLB always misses, and the
      walk restarts from the page directory after each fault. For one
      virtual reference: read PDE ($M$), fault on leaf PT ($D$); restart
      and read PDE$+$PTE ($2M$), fault on code/data page ($D$); restart
      and read PDE$+$PTE$+$byte ($3M$), totaling $6M+2D$ per reference.
      Across both instruction fetch and data load this gives
      $\mathbf{12M+4D}$ (or $\mathbf{12M+8D}$ if every evicted victim is
      dirty and requires a write-back first). (In a simplified model
      that charges no re-walk after a fault and allows the directory
      itself to swap, one gets $6M+6D$ without write-backs or $6M+12D$
      with dirty write-backs; always state your assumptions explicitly.)
\end{enumerate}
A hardware-managed two-level walk never traps on a present PTE; a
software-managed TLB traps on every miss, even when the page is present
--- the $M$-counts above are the \emph{memory} operations, not the trap
overhead.
\end{Sol}

\subsection{Lecture 15: Swapping \& Page Replacement Policies}

\begin{Qn}{Hard}{Lect 15}{CS-537 Fa09 Mid, Q4}{Valid bit vs present bit, and a dumped PTE}
LostOS uses a linear page table. Each PTE is 64 bits and contains only four
fields: a \textbf{valid} bit, a \textbf{present} bit, a 32-bit timestamp
(seconds since 1970; reused as a \textbf{disk address} when the page is
swapped out), and a 30-bit PFN. When valid$\,=0$, the remaining bits are
filled with all 0s or all 1s (never mixed). The current time is
\texttt{0x4ADCCD3F}.

First PTEs of process 1 (hex bytes, one PTE per line):
\begin{lstlisting}[numbers=none]
C0 00 01 00 4A DC 70 04
80 00 00 00 00 00 10 00
7F FF FF FF FF FF FF FF
7F FF FF FF FF FF FF FF
\end{lstlisting}
First PTEs of process 2:
\begin{lstlisting}[numbers=none]
C0 00 00 10 4A DC 90 CB
C0 00 00 11 4A DC 90 C0
00 00 00 00 00 00 00 00
00 00 00 00 00 00 00 00
\end{lstlisting}
\begin{enumerate}
  \item[(a)] What is the valid bit used for?
  \item[(b)] What is the present bit used for?
  \item[(c)] Infer the bit layout of the PTE.
  \item[(d)] For each of process 1, VPN $0,1,2$ and process 2, VPN $0,1,2$,
        state what happens on an access (translate / page fault / segfault).
\end{enumerate}
\end{Qn}
\begin{Sol}
(a) Valid$\,=1$ means this VPN is part of the process's address space
(allocated). Valid$\,=0$ means the address is illegal: a
\textbf{segmentation fault}, not a page fault.

(b) Present is consulted only on a valid page. Present$\,=1$: the PFN field
is a frame number and the page is in RAM. Present$\,=0$: the timestamp
field is a \textbf{swap-block number}; the access traps as a
\textbf{page fault}.

(c) Top bit valid, next bit present, then 30-bit PFN, then 32-bit
timestamp/disk address (big-endian in the dump). Check: first byte
\texttt{C0}$\,=\,$\texttt{1100\,0000} (valid and present);
\texttt{80}$\,=\,$\texttt{1000\,0000} (valid, not present);
\texttt{7F}/\texttt{00} have the top bit clear (invalid). The low 32 bits
of a present PTE are near \texttt{0x4ADC....}, matching the given clock.

(d) Process 1, VPN 0: \texttt{C0 00 01 00 | 4A DC 70 04} --- valid, present,
PFN \texttt{0x0100}: ordinary translation.
Process 1, VPN 1: \texttt{80 ... 00 00 10 00} --- valid, not present, disk
address \texttt{0x1000}: \textbf{page fault}.
Process 1, VPN 2: \texttt{7F FF...}: valid$\,=0$: \textbf{segfault}.
Process 2, VPN 0: valid, present, PFN \texttt{0x10}: translate.
Process 2, VPN 1: valid, present, PFN \texttt{0x11}: translate.
Process 2, VPN 2: all zeros: \textbf{segfault}.
\end{Sol}

\begin{Qn}{Hard}{Lect 15}{CS-537 Sp18 Final, Q77--80}{Which policy left $\{3,4,5\}$ in a 3-page memory?}
A 3-page memory currently holds pages 3, 4 and 5. The last 10 references
(oldest to newest, memory empty at the start of the 10) were
\[
8,\;7,\;4,\;2,\;5,\;4,\;7,\;3,\;4,\;5.
\]
\begin{enumerate}
  \item[77.] Which replacement policy was used?
  \begin{enumerate}
    \item[(a)] FIFO \quad (b) LRU \quad (c) MRU \quad (d) LFU
    \item[(e)] Cannot tell from the given information
  \end{enumerate}
  \item[78.] If the policy had been FIFO, how many misses occurred among
        those 10 accesses? (a) 7 (b) 8 (c) 9 (d) 10 (e) none of these
  \item[79.] Same question for LRU.
  \item[80.] Same question for OPT (Belady's MIN).
\end{enumerate}
\end{Qn}
\begin{Sol}
Simulate 3 frames, empty start. (Contents shown after the access; $*$ =
miss.)

\textbf{FIFO} (9 misses, final $\{3,4,5\}$):
\begin{center}
\small
\begin{tabular}{c|cccccccccc}
ref & 8 & 7 & 4 & 2 & 5 & 4 & 7 & 3 & 4 & 5 \\
\hline
    & 8 & 87& 874& 742& 425& 425& 257& 573& 734& 345 \\
    & $*$& $*$& $*$& $*$& $*$&  & $*$& $*$& $*$& $*$ \\
\end{tabular}
\end{center}

\textbf{LRU} (8 misses, final $\{3,4,5\}$):
\begin{center}
\small
\begin{tabular}{c|cccccccccc}
ref & 8 & 7 & 4 & 2 & 5 & 4 & 7 & 3 & 4 & 5 \\
\hline
    & 8 & 87& 874& 742& 425& 254& 547& 473& 734& 345 \\
    & $*$& $*$& $*$& $*$& $*$&  & $*$& $*$&  & $*$ \\
\end{tabular}
\end{center}

\textbf{OPT} (6 misses, final $\{3,4,5\}$). Evict the resident whose
next use is furthest (or never):
$8*$ $\{8\}$;
$7*$ $\{8,7\}$;
$4*$ $\{8,7,4\}$;
$2*$ evict 8 (never reused) $\{7,4,2\}$;
$5*$ evict 2 (never reused) $\{7,4,5\}$;
4 hit; 7 hit;
$3*$ evict 7 (never reused) $\{4,5,3\}$;
4 hit; 5 hit.

\textbf{MRU} ends with $\{8,3,5\}$. \textbf{LFU} with history-wide
counts ends with $\{4,5,7\}$, whereas LFU with counts reset on eviction
(FIFO/LRU tie-break) ends with $\{3,4,5\}$ as well. Either way, FIFO,
LRU, and OPT already all produce $\{3,4,5\}$.

77. FIFO, LRU \emph{and} OPT all finish with $\{3,4,5\}$.
Answer: \textbf{(e) cannot tell}.

78. FIFO misses $=9$ $\to$ \textbf{(c)}.

79. LRU misses $=8$ $\to$ \textbf{(b)}.

80. OPT misses $=6$, which is not among 7, 8, 9, 10 $\to$
\textbf{(e) none of the above}.
\end{Sol}

\begin{Qn}{Hard}{Lect 15}{CS-537 Fa09 Mid, Q5}{Perfect LRU via timestamps, then a second-chance \texttt{replace()}}
LostOS first implemented perfect LRU with hardware timestamps: on every
reference the MMU writes the current time into that PTE. Linear page
tables, as in Q4 of the same exam.
\begin{enumerate}
  \item[(a)] Given that hardware, what must the OS do to pick a victim,
        and what is wrong with the approach?
  \item[(b)] Hardware is then reduced to a single reference bit (set to 1
        on each access). Write \texttt{replace(int *process, int *VPN)}
        that selects a not-recently-referenced resident page. The
        tables are
\end{enumerate}
\begin{lstlisting}[numbers=none]
typedef struct {
    int valid, present, referenced;
    unsigned pfn;
} pte_t;
typedef struct {
    pte_t pagetable[PAGE_TABLE_ENTRIES];
} page_table_t;
page_table_t *p[MAX_PROCESSES];  /* NULL if process does not exist */
\end{lstlisting}
Consider: where the scan starts; whether the reference bit is ever
cleared; that only valid \emph{and} present pages are candidates.
\end{Qn}
\begin{Sol}
(a) The hardware already records recency. To find the LRU page the OS
must \textbf{scan every valid, present PTE of every process} and take
the one with the smallest timestamp. On a 32-bit system with 4\,KB
pages that is up to $2^{20}$ PTEs \emph{per process} --- far too slow
on the page-fault path (while a process is blocked waiting for a frame).
Even in hardware, writing a 32-bit timestamp to the PTE in memory on
\emph{every} memory reference adds huge memory-traffic overhead.

(b) This is \textbf{clock / second chance} over the concatenated page
tables. Clear referenced bits as you skip them so that a later sweep
can tell cold from hot; do not always start at process 0 / VPN 0
(that repeatedly victimizes low-numbered pages instead of approximating LRU). One acceptable implementation:

\begin{lstlisting}[numbers=none]
int cur_p = 0, cur_v = 0;   /* persistent hand */

void replace(int *process, int *VPN) {
    while (1) {
        if (p[cur_p] != NULL) {
            pte_t *e = &p[cur_p]->pagetable[cur_v];
            if (e->valid && e->present) {
                if (e->referenced == 0) {
                    *process = cur_p;
                    *VPN     = cur_v;
                    cur_v    = (cur_v + 1) % PAGE_TABLE_ENTRIES;
                    if (cur_v == 0)
                        cur_p = (cur_p + 1) % MAX_PROCESSES;
                    return;
                }
                e->referenced = 0;   /* second chance */
            }
        }
        cur_v = (cur_v + 1) % PAGE_TABLE_ENTRIES;
        if (cur_v == 0)
            cur_p = (cur_p + 1) % MAX_PROCESSES;
        /* A full sweep clears all referenced bits so the next pass
           succeeds (FIFO degeneration). */
    }
}
\end{lstlisting}
A random start instead of a stored hand is also accepted. Returning a
page with \texttt{referenced==1} without a clearing pass is not.
\end{Sol}

\begin{Qn}{Medium-Hard}{Lect 15}{ISI Mandar 2016 End-sem, Q2(a)}{Maximum page-fault rate for a 200\,ns EAT, with dirty writebacks}
Demand paging. A memory access is 100\,ns. The page table has 8 entries
and is held in \textbf{registers} (no extra memory reference for the
PTE). A page fault takes 10\,ms if a free frame exists or the victim is
clean, and 20\,ms if the victim is dirty. The victim is dirty 80\% of
the time. Find the maximum page-fault rate $p$ such that the effective
access time stays $\le 200$\,ns.
\end{Qn}
\begin{Sol}
Because the PTEs live in registers, a hit costs just the 100\,ns memory
access. A fault costs
\[
  0.2\times 10\,\mathrm{ms} + 0.8\times 20\,\mathrm{ms}
  = 18\,\mathrm{ms} = 1.8\times 10^{7}\,\mathrm{ns}.
\]
\[
  \mathrm{EAT}
  = (1-p)\cdot 100 + p\cdot 1.8\times 10^{7}
  \le 200.
\]
\[
  p\bigl(1.8\times 10^{7}-100\bigr) \le 100
  \quad\Rightarrow\quad
  p \le \frac{100}{1.8\times 10^{7}-100}
  = \frac{1}{179999}
  \approx 5.556\times 10^{-6}.
\]
(Dropping the 100\,ns inside the fault service --- legitimate because
$18\,\mathrm{ms}\gg 100\,\mathrm{ns}$ --- gives the same three-digit
value $p\le 100/(1.8\times 10^{7})=1/180000$.) At most about five
faults per million references.
\end{Sol}

\begin{Qn}{Hard}{Lect 15}{ISI Mandar 2018 End-sem, Q4(a)}{FIFO page-fault rate of an $N\times N$ transpose}
A program computes \texttt{for (i = 0; i < N; i++) for (j = 0; j < N; j++) B[i,j] = A[j,i];} with $A,B$ stored row-major.
Parameters: page size $2^{12}$ bytes; 4 frames, FIFO; $N=2^{15}$; 8-byte
entries; $A$ and $B$ each begin on a page boundary; no $A$/$B$ page is
initially resident. Count only references to matrix entries.
\begin{enumerate}
  \item[(i)] How many pages do $A$ and $B$ occupy?
  \item[(ii)] Page-fault rate of the nested loops.
\end{enumerate}
\end{Qn}
\begin{Sol}
(i) Each matrix is $N^{2}$ entries $\times$ 8 bytes
$= 2^{15}\cdot 2^{15}\cdot 2^{3} = 2^{33}$ bytes, over $2^{12}$-byte
pages: $2^{21}$ pages each. Equivalently, 512 entries/page so 64 pages
per row and $2^{15}$ rows.

(ii) Inner index $j$ walks \textbf{across} a row of $B$ (sequential,
512 elements per page) and \textbf{down} a column of $A$.
$A[j,i]$ for successive $j$ lives 64 pages apart, so every $A$
reference in an inner loop hits a different page. Although $A[j,i+1]$
lies on the \emph{same} page as $A[j,i]$ (for 511 out of 512 column
steps), it is not touched again until $N = 2^{15}$ other $A$ pages have
been accessed, long after the 4-frame FIFO has evicted it. Thus every
one of the $N^{2}$ $A$ references faults.

$B$ is not free. The access stream is $A,B,A,B,\ldots$ with a new $A$
page every time and the same $B$ page for 512 consecutive $j$. FIFO-4
on that stream:
\begin{center}
\small
\begin{tabular}{c|cccccc}
$j \bmod 4$ & 0 & 1 & 2 & 3 & 0 & $\cdots$ \\
\hline
$A$ & miss & miss & miss & miss & miss & \\
$B$ & miss & hit  & hit  & hit  & miss & \\
\end{tabular}
\end{center}
After a 4-step fill, $B$ is evicted every 4th $A$ miss, so each 512-long
$B$-page window has $512/4=128$ $B$ faults (and 384 $B$ hits), and every
$A$ faults. There are $N^{2}/512 = 2^{21}$ such windows.

Total references $2N^{2}=2^{31}$. Total faults
$N^{2} + 2^{21}\cdot 128 = 2^{30}+2^{28}$. Rate
\[
  \frac{2^{30}+2^{28}}{2^{31}} = \frac{5}{8} = 0.625.
\]
\end{Sol}

\clearpage

\section{Lectures 16--20: Concurrency --- Threads, Locks, CVs, Semaphores \& Deadlock}

\subsection{Lectures 16--17: Concurrency, Threads \& Locks}

\begin{Qn}{Hard}{Lect 17}{Mandar ISI 2012 Mid Q3; repeated Mandar ISI 2022 Mid Q3}{Concurrent \texttt{pop} on a shared linked-list stack}
Consider a stack implemented as a singly linked list. The \texttt{pop} function is:

\begin{lstlisting}[numbers=none]
NODE *pop(STACK *s)
{ NODE *tmp;
  if (s->top == NULL) return NULL;
  tmp = s->top;
  s->top = tmp->next;
  tmp->next = NULL;
  return tmp;
}
\end{lstlisting}

(a)~Suppose two processes $P_1$ and $P_2$ share a stack. Show how inconsistent results may be obtained if they execute \texttt{pop} concurrently. You must show the initial stack, the intended results of the two pops, the interleaved execution, and the actual results.

(b)~[2012 Mid only] Redesign \texttt{pop} to prevent the race if the hardware provides an atomic \texttt{TestAndSet} operation.
\end{Qn}
\begin{Sol}
(a)~Let the stack be $A \to B \to C \to \bot$, with \texttt{s->top} pointing at $A$. Intended result: $P_1$ pops $A$, $P_2$ pops $B$, leftover $C$.

A bad interleaving (lost update / duplicate pop):

\begin{center}
\begin{tabular}{c|l|l|l}
Step & $P_1$ & $P_2$ & Shared state \\
\hline
0 & & & \texttt{top}$=A \to B \to C$ \\
1 & \texttt{tmp = s->top} \quad ($A$) & & \\
2 & & \texttt{tmp = s->top} \quad ($A$) & \\
3 & \texttt{s->top = tmp->next} \quad ($B$) & & \texttt{top}$=B$ \\
4 & & \texttt{s->top = tmp->next} \quad ($B$) & \texttt{top}$=B$ \\
5 & \texttt{tmp->next = NULL}; return $A$ & & \\
6 & & \texttt{tmp->next = NULL}; return $A$ & \\
\end{tabular}
\end{center}

Both processes return a pointer to $A$ (double-free / use-after-pop). Node $B$ is still the top, so $B$ was never popped. Node $A$ is unlinked twice, and if each caller later \texttt{free}s it, the heap is corrupted. Symmetric lost-node variants exist if the two stores to \texttt{s->top} use different snapshots of \texttt{tmp->next}.

(b)~Protect the check-and-unlink as a critical section with a TAS spinlock:

\begin{lstlisting}[numbers=none]
shared int lock = 0;   /* 0 = free, 1 = held */

NODE *pop(STACK *s)
{ NODE *tmp;
  while (TestAndSet(&lock)) ;   /* spin */
  if (s->top == NULL) { lock = 0; return NULL; }
  tmp = s->top;
  s->top = tmp->next;
  tmp->next = NULL;
  lock = 0;
  return tmp;
}
\end{lstlisting}

\texttt{TestAndSet(\&lock)} atomically writes 1 and returns the old value, so only one caller enters the body. Unlock is a plain store of 0 (safe here: the lock owner is the only writer of the free value). The empty-stack path must release the lock on every return.
\end{Sol}

\begin{Qn}{Hard}{Lect 17}{Mandar ISI 2012 End Q2(a)}{Modified Peterson: exhibit a mutual-exclusion failure}
Consider this modified 2-process CSP scheme:

\begin{lstlisting}[numbers=none]
shared char want[2] = {0,0};
shared int turn = 0;
/* P_i, i in {0,1}, j = 1-i */
while (1) {
    turn = j;
    want[i] = 1;
    while (want[j] && turn != i) ;
    critical_section();
    want[i] = 0;
    remainder_section();
}
\end{lstlisting}

Show that the scheme does \emph{not} guarantee mutual exclusion, by constructing an interleaved sequence of instructions of $P_0$ and $P_1$.
\end{Qn}
\begin{Sol}
Peterson's real entry section writes the flag \emph{before} writing \texttt{turn}. Here \texttt{turn = j} happens first, so a process can publish ``you go'' before it has published ``I want in''. Counterexample (both start with \texttt{want[0]=want[1]=0}):

\begin{center}
\small
\begin{tabularx}{\linewidth}{@{}c|l|l|X@{}}
Step & $P_0$ & $P_1$ & \texttt{want}, \texttt{turn} \\
\hline
1 & \texttt{turn = 1} & & $(0,0),\ 1$ \\
2 & & \texttt{turn = 0} & $(0,0),\ 0$ \\
3 & & \texttt{want[1] = 1} & $(0,1),\ 0$ \\
4 & & \texttt{while(want[0] \&\& turn!=1)} & \texttt{want[0]=0}, enter CS \\
5 & \texttt{want[0] = 1} & & $(1,1),\ 0$ \\
6 & \texttt{while(want[1] \&\& turn!=0)} & &
    \texttt{want[1]=1} but \texttt{turn==0==i}, wait is false: enter CS \\
\end{tabularx}
\end{center}

Both are in the critical section. At step~6 the wait condition is \texttt{want[1] \&\& turn != 0}. Because $P_1$ last wrote \texttt{turn = 0} \emph{before} $P_0$ set \texttt{want[0]}, $P_0$ sees \texttt{turn} already equal to its own id and does not spin.
\end{Sol}

\begin{Qn}{Medium-Hard}{Lect 17}{CS-537 Spring 2018 Mid Problem XXI}{Possible \texttt{(ticket, turn)} pairs}
\begin{lstlisting}[numbers=none]
typedef struct __lock_t {
    int ticket, turn;
} lock_t;
void lock_init(lock_t *lock) {
    lock->ticket = 0;
    lock->turn   = 0;
}
void lock(lock_t *lock) {
    int myturn = FetchAndAdd(&lock->ticket);
    while (lock->turn != myturn)
        ; /* spin */
}
void unlock(lock_t *lock) {
    lock->turn = lock->turn + 1;
}
\end{lstlisting}

At most 5 threads, lock used properly (acquire/release paired). Which simultaneous values are possible? A~=~possible, B~=~not.

\begin{enumerate}
\item[101.] \texttt{ticket=0}, \texttt{turn=0}
\item[102.] \texttt{ticket=0}, \texttt{turn=1}
\item[103.] \texttt{ticket=1}, \texttt{turn=0}
\item[104.] \texttt{ticket=16}, \texttt{turn=5}
\item[105.] \texttt{ticket=1000}, \texttt{turn=999}
\end{enumerate}
\end{Qn}
\begin{Sol}
Invariant: \texttt{ticket} is the number of \texttt{lock()} calls so far; \texttt{turn} is the number of completed \texttt{unlock()}s. Always $0 \le \texttt{turn} \le \texttt{ticket}$ and $\texttt{ticket}-\texttt{turn} \le \#\text{threads}$ (here 5), because that difference is the number of threads currently in \texttt{lock()} (holding or spinning).

\begin{itemize}
\item 101: initial state. A.
\item 102: \texttt{turn} advanced without a matching ticket. Impossible. B.
\item 103: one thread has taken a ticket and not yet unlocked (it holds the lock, spinning waiters would also increase ticket). A.
\item 104: $\texttt{ticket}-\texttt{turn}=11>5$. Impossible. B.
\item 105: difference 1, after many acquire/release cycles plus one in-progress acquire. A.
\end{itemize}

Official key: 101~A, 102~B, 103~A, 104~B, 105~A.
\end{Sol}

\begin{Qn}{Hard}{Lect 17}{CS-537 Spring 2012 Mid Q14}{TAS versus test-and-test-and-set}
A spin-lock acquire with atomic exchange:

\begin{lstlisting}[numbers=none]
while (xchg(&lock, 1) == 1)
    ; /* spin */
\end{lstlisting}

A subtly different acquire:

\begin{lstlisting}[numbers=none]
while (1) {
    while (lock > 0)
        ; /* spin */
    if (xchg(&lock, 1) == 0)
        return;
}
\end{lstlisting}

Does the new lock work? How does the behaviour change?
\end{Qn}
\begin{Sol}
It works (still a correct mutex) assuming \texttt{lock} is 0 when free and 1 when held, and unlock stores 0. This is \textbf{test-and-test-and-set}. Waiters spin on an ordinary load, which hits in their cache, until they observe the lock look free; only then do they fire the expensive \texttt{xchg}.

Differences:
\begin{itemize}
\item Correctness: same ME. Two waiters can both see \texttt{lock==0} and both \texttt{xchg}; one gets 0 and wins, the other gets 1 and returns to the outer loop. No double-enter.
\item Performance: far less interconnect / cache-line bouncing while the lock is held, because spinning loads are local. TAS hammers the line on every iteration.
\item Fairness: still unfair (no FIFO order; a waiter can lose every \texttt{xchg} race); not a ticket lock. (In C, \texttt{lock} must also be \texttt{volatile}/atomic so the compiler does not hoist the inner load out of the loop.)
\item Uniprocessor: still disastrous if the holder is preempted --- waiters spin a whole timeslice. Lecture~17: spinlocks need preemption on 1~CPU and are for short CS on $N$~CPUs.
\end{itemize}

(Near-duplicate skill: CS-537 Spring 2009 Final Q3 uses the same TTAS loop with \texttt{TestAndSet}.)
\end{Sol}

\begin{Qn}{Hard}{Lect 17}{CS-537 Fall 2021 Midterm 2 Q30}{CAS increment: how many times can L2 run?}
\begin{lstlisting}[numbers=none]
void increment(int *p) {
     int v, n;
     do {
           v = *p;                             /* L1 */
           n = v + 1;                          /* L2 */
     } while (compare_and_swap(p, v, n) == 0); /* L3 */
}
\end{lstlisting}

Three threads call \texttt{increment} on the same variable at roughly the same time. How many times \emph{might} L2 execute (worst case among the options)?
(a)~3 (b)~4 (c)~5 (d)~6 (e)~none of the above
\end{Qn}
\begin{Sol}
CAS succeeds only if \texttt{*p} still equals \texttt{v}. Worst case: all three execute L1--L2 on the same snapshot (3 executions of L2); two CAS fail; those two retry L1--L2 (2 more); one of those two fails again and retries (1 more). Total 6.

Trace: $p=0$. T1,T2,T3 all read $v=0$, all compute $n=1$ (L2 $\times 3$). T1 CAS 0$\to$1 succeeds. T2,T3 CAS fail. T2,T3 read $v=1$, compute $n=2$ (L2 $\times 2$). T2 CAS 1$\to$2 succeeds. T3 fails, reads $v=2$, computes $n=3$ (L2 $\times 1$). T3 CAS succeeds. L2 ran $3+2+1=6$ times. Answer (d).
\end{Sol}

\begin{Qn}{Hard}{Lect 17}{Mandar ISI 2022 Mid Q4(b); cf.\ Mandar ISI 2015 Mid Q5}{Waiting-array TAS: ME / progress / bounded wait}
Each $P_i$ runs (all shared state initialized to 0):

\begin{lstlisting}[numbers=none]
shared char lock, waiting[n];   /* local int key per process */
waiting[i] = 1;
key = 1;
while (waiting[i] && key) { key = TestAndSet(&lock); }
waiting[i] = 0;
critical_section();
for (j = 0; j < n; j++)
    if (waiting[j] == 1) { waiting[j] = 0; break; }
if (j == n) lock = 0;
\end{lstlisting}

For each of the three CSP properties, does this solution possess it? If yes, argue; if no, give a violating scenario.

[2015 Mid Q5 is the simpler lock \texttt{while(TestAndSet(\&lock)); CS; lock=0;} --- same three properties, different answers; treat as the warm-up form of this question.]
\end{Qn}
\begin{Sol}
Recall TAS returns the old value and writes 1.

\textbf{Mutual exclusion: yes.} A thread leaves the while-loop only if \texttt{key==0} (TAS saw a free lock and consumed it) or \texttt{waiting[i]==0} (handed off by an exiting thread). The hand-off clears exactly one waiter's flag and does \emph{not} set \texttt{lock=0}, so the token moves to that waiter. Only one token exists.

\textbf{Progress: yes.} If the CS is free and some thread is in its entry section, either TAS succeeds (\texttt{lock} was 0) or an exiting thread finds a waiter and clears its \texttt{waiting[j]}.

\textbf{Bounded waiting: no (for $n \ge 3$).} The exit scan is \texttt{for (j = 0; j < n; j++)}, always starting from index 0 rather than $(i+1)\bmod n$. With only 2 processes, the exiting process has already cleared its own \texttt{waiting[i]=0}, so it must hand off to the other process. With $n \ge 3$ (say $P_0, P_1, P_2$), low-index processes can shut out $P_2$ forever:
\begin{itemize}
  \item $P_0$ is in the CS (\texttt{waiting[0]=0}); both $P_1$ and $P_2$ are spinning with \texttt{waiting[1]=waiting[2]=1} and \texttt{lock=1}.
  \item $P_0$ exits: its scan from $j=0$ skips $j=0$, hits $j=1$, clears \texttt{waiting[1]=0}, and breaks (leaving \texttt{lock=1}). $P_1$ enters the CS.
  \item Before $P_1$ exits, $P_0$ re-enters its entry section and sets \texttt{waiting[0]=1}.
  \item $P_1$ exits: its scan from $j=0$ immediately hits $j=0$, clears \texttt{waiting[0]=0}, and breaks. $P_0$ re-enters the CS.
\end{itemize}
Repeating this ping-pong between $P_0$ and $P_1$, \texttt{lock} never returns to 0 and \texttt{waiting[2]} is never cleared, so $P_2$ starves indefinitely. (A cyclic scan starting at $(i+1)\bmod n$ restores bounded waiting.)

Warm-up (2015 Mid Q5) \texttt{while(TAS(\&lock)); CS; lock=0;}:
\begin{itemize}
\item ME: yes (TAS).
\item Progress: yes (when lock is 0 the next TAS succeeds).
\item Bounded waiting: \textbf{no} --- a thread can lose every TAS race (unfair spinlock; Lecture~17).
\end{itemize}
\end{Sol}

\begin{Qn}{Hard}{Lect 16}{Mandar ISI 2010 Mid Q3}{Possible finals of racy \texttt{x++} / \texttt{x--}}
Shared \texttt{int x = 0}. Two processes, $i\in\{0,1\}$:

\begin{lstlisting}[numbers=none]
for (j = 0; j < 10; j++) {
    if (i % 2) x++;
    else       x--;
}
\end{lstlisting}

After both finish, what set of values can \texttt{x} take? Justify carefully. (Compile \texttt{x++}/\texttt{x--} as load, add/sub, store.)
\end{Qn}
\begin{Sol}
$P_1$ performs 10 increments; $P_0$ performs 10 decrements. Each operation is three memory-relevant steps: \texttt{r = x; r = r $\pm$ 1; x = r}. Lost updates occur when a store is based on a stale load.

The possible finals are all integers in $[-10,10]$.

\textbf{Extremes.} For $+10$: $P_1$ load-computes its first increment from 0 but \emph{delays the store}; $P_0$ then runs all 10 decrements to completion (\texttt{x=-10}); $P_1$ stores $1$ (overwriting $-10$) and then chains its remaining 9 increments $1\to 2\to\cdots\to 10$. Symmetric delay of $P_0$'s first store yields $-10$.

\textbf{Interior values.} Assuming loop counter \texttt{j} is private to each process, any $k\in[0,10]$ is obtained explicitly as follows: $P_1$ loads $0$ for its first increment and pauses before storing; $P_0$ completes $k$ decrements (making \texttt{x = -k}); $P_1$ stores $1$ (overwriting $-k$) and executes its remaining 9 increments to completion (\texttt{x = 10}); finally $P_0$ runs its remaining $10-k$ decrements sequentially, yielding $\texttt{x} = 10 - (10-k) = k$. Any negative value $-k\in[-10,0]$ follows symmetrically by swapping the roles of $P_0$ and $P_1$.

\textbf{Impossible:} $|x|>10$, because each thread writes at most 10 times and each committed chain can contribute at most one $\pm 1$ per operation of that thread.
\end{Sol}

\begin{Qn}{Hard}{Lect 17}{CS-537 Fall 2011 Final Q2}{Best- and worst-case time of a spinlock on one CPU}
A spinlock acquire/unlock each cost $A$ when the lock is free. A context switch costs $C$. A timeslice is $T$. Two threads on \emph{one} processor, at roughly the same time, run

\begin{lstlisting}[numbers=none]
mutex_lock();
do_something(); /* time 0 */
mutex_unlock();
exit();         /* time 0 */
\end{lstlisting}

(a)~Best-case time for both to finish.

(b)~Worst-case time. At most three context switches.

(c)~If the spinlock is replaced by a queue-based lock (\texttt{park}/\texttt{unpark}), how does the worst case change?
\end{Qn}
\begin{Sol}
(a)~No contention, no extra switches: T1 locks ($A$), unlocks ($A$), exits; OS context-switches to T2 ($C$); T2 locks ($A$), unlocks ($A$). Total: $\mathbf{4A + C}$.

(b)~Holder is preempted inside the CS (here the CS is empty, so: T1 acquires, timeslice ends before unlock). T2 runs and \emph{spins for up to $T$} until T1 is rescheduled, then T1 unlocks, then T2 acquires. With $\le 3$ switches:
T1 runs, acquires, preempted ($C$);
T2 spins almost a full slice $T$, preempted ($C$);
T1 unlocks, exits, switch to T2 ($C$);
T2 acquires and finishes.
Worst-case $\approx \mathbf{4A + 3C + T}$ (the wasted spin). Lecture~17: on 1~CPU a spinner wastes the rest of the timeslice.

(c)~A queue lock \texttt{park}s T2 immediately on contention instead of spinning for $T$: T1 acquires and is preempted ($C$); T2 tries to acquire, calls \texttt{park()}, and the OS switches back to T1 ($C$); T1 unlocks, \texttt{unpark}s T2, and exits ($C$); T2 acquires and unlocks. The worst-case time drops to $\approx \mathbf{4A + 3C}$ (plus small \texttt{park}/\texttt{unpark} queue overhead), completely eliminating the timeslice term $T$.
\end{Sol}

\begin{Qn}{Hard}{Lect 16}{CS-537 Spring 2018 Mid Problems XXIV--XXV}{Printer threads: heap char vs.\ stack char}
\begin{lstlisting}[numbers=none]
void *printer(void *arg) {
  char *p = (char *) arg;
  printf("%c", *p);
  return NULL;
}
\end{lstlisting}

\textbf{XXIV.} \texttt{main} loops $i=0..4$: \texttt{malloc} a \texttt{char}, set it to \texttt{'a'+i}, \texttt{pthread\_create(..., printer, c)}, then joins all. Possible outputs? A/B:

116.~\texttt{abcde} \quad 117.~\texttt{edcba} \quad 118.~\texttt{cccde} \quad 119.~\texttt{eeeee} \quad 120.~\texttt{aaaaa}

\textbf{XXV.} Same, but \texttt{char c = 'a'+i; pthread\_create(..., printer, \&c);} (no malloc). Possible:

121.~\texttt{abcde} \quad 122.~\texttt{edcba} \quad 123.~\texttt{cccde} \quad 124.~\texttt{eeeee} \quad 125.~\texttt{aaaaa}
\end{Qn}
\begin{Sol}
\textbf{XXIV.} Each thread gets a \emph{distinct heap cell} with a unique letter. Any permutation of \texttt{a}--\texttt{e} is possible; repeats are not. 116~A, 117~A, 118~B, 119~B, 120~B.

\textbf{XXV.} All five threads receive the address of the \emph{same stack slot} \texttt{c} (assuming the compiler reuses one stack slot for \texttt{c} across iterations and before the joins; strictly in C, reading a loop-scoped local after its iteration ends is UB). Output depends on when each thread dereferences vs.\ when \texttt{main} writes.

\begin{itemize}
\item 121 \texttt{abcde}: each child runs to completion before the next \texttt{create}. A.
\item 122 \texttt{edcba}: each child evaluates its argument \texttt{*p} in creation order before entering \texttt{printf}, but the actual output of earlier threads is delayed until later threads print. A (official key).
\item 123 \texttt{cccde}: first three creates happen before those threads run, \texttt{c} is already \texttt{'c'}; then \texttt{d} and \texttt{e} print in order. A.
\item 124 \texttt{eeeee}: all five created before any runs; \texttt{c=='e'}. A.
\item 125 \texttt{aaaaa}: after the second create, \texttt{c} is at least \texttt{'b'}. Impossible to have five \texttt{a}'s. B.
\end{itemize}

Official: 121--124 A, 125 B.
\end{Sol}

\subsection{Lecture 18: Condition Variables}

\begin{Qn}{Hard}{Lect 18}{CS-537 Fa21 Mid2, Q16; cf.\ Fa11 Mid, Q21}{Join with a CV: missing mutex, missing state, lost wakeup}
Bug tries to solve the fork/join problem (a parent must not print
\texttt{done} until the child has printed \texttt{child}) as follows.
\texttt{cond\_t} is a condition variable; \texttt{cond\_wait}/\texttt{cond\_signal}
are the usual primitives. CVs and mutexes, if used, are properly initialised.
\begin{lstlisting}[numbers=none]
cond_t c;

void child(void *arg) {
    printf("child\n");
    cond_signal(&c);
}

int main() {
    thread_t p1;
    thread_create(&p1, child, NULL);
    cond_wait(&c);          /* no mutex! */
    printf("done\n");
    return 0;
}
\end{lstlisting}
\begin{enumerate}[(a)]
\item Give \emph{two} schedules: one in which the program prints
      \texttt{child} then \texttt{done}, and one in which it prints only
      \texttt{child} and the parent waits forever.
\item The call \texttt{cond\_wait(\&c)} does not even take a mutex.
      Why does a correct \texttt{wait} require the caller to hold a mutex,
      and what does \texttt{wait} do with that mutex?
\item Attempt~2 adds a state variable but the child does not take the
      lock:
\begin{lstlisting}[numbers=none]
int done = 0;
void thread_join() {                 void thread_exit() {
    Mutex_lock(&m);          /* w */     done = 1;            /* a */
    if (done == 0)           /* x */     Cond_signal(&c);     /* b */
        Cond_wait(&c, &m);   /* y */ }
    Mutex_unlock(&m);        /* z */
}
\end{lstlisting}
      Give a schedule in which the parent waits forever.
\item Write a correct \texttt{thread\_join}/\texttt{thread\_exit} using one
      mutex, one CV, and \texttt{done}. Why is a spin on \texttt{done}
      also correct, and why is it a bad idea?
\end{enumerate}
\end{Qn}
\begin{Sol}
(a)~\emph{Works:} parent creates the child, then waits; child prints
\texttt{child}, signals, parent wakes, prints \texttt{done}.
\emph{Lost wakeup:} child runs to completion first, prints
\texttt{child}, and signals while \emph{nobody is waiting}. A Mesa
\texttt{signal} with an empty queue is a no-op (it is not remembered).
The parent then calls \texttt{wait} and sleeps forever. Both outcomes
are possible, so the official MCQ key is
``just \texttt{child}, \emph{or} \texttt{child} then \texttt{done}''.
(One can also argue the snippet would not compile because POSIX
\texttt{pthread\_cond\_wait} requires a mutex; that is option~(e) on
the Scantron. The intended concurrency bug is the lost wakeup.)

(b)~The caller must hold the mutex. \texttt{wait} \textbf{atomically}
releases the mutex and puts the caller to sleep; after a wakeup it
\textbf{reacquires} the mutex before returning. The atomic
release-and-sleep is what closes the window between ``I have seen that
I should wait'' and ``I am on the queue''. Without the mutex, a
signaller can change the shared state and signal in that window.

(c)~Parent \texttt{w}, \texttt{x} (sees \texttt{done == 0}, about to
wait). Child \texttt{a} (\texttt{done = 1}), \texttt{b} (signal; nobody
is waiting yet, so it is lost). Parent \texttt{y}: sleeps forever.
The child updates \texttt{done} and signals \emph{outside} the lock, so
the update can fall between the parent's check and its wait.

(d)~Correct join (state + lock held around both the check/wait and the
update/signal):
\begin{lstlisting}[numbers=none]
int done = 0;
mutex_t m;  cond_t c;

void thread_join() {                 void thread_exit() {
    Mutex_lock(&m);                      Mutex_lock(&m);
    while (done == 0)                    done = 1;
        Cond_wait(&c, &m);               Cond_signal(&c);
    Mutex_unlock(&m);                    Mutex_unlock(&m);
}
\end{lstlisting}
With one parent and one child, \texttt{if (done == 0)} happens to
suffice; \texttt{while} is always correct (Mesa, and spurious wakeups).
A spin \texttt{while (done == 0) ;} is also correct --- the child
stores \texttt{done = 1}, the parent eventually observes it --- but it
\textbf{wastes a CPU} for the whole wait, which is the reason CVs exist:
ordering without spinning.

Rules of thumb used: (1)~keep state in addition to the CV;
(2)~wait/signal with the lock held; (3)~recheck with \texttt{while}.
\end{Sol}

\begin{Qn}{Hard}{Lect 18}{CS-537 Sp12 Mid, Q13}{Hoare lock-handoff: does \texttt{if} suddenly work?}
A bounded buffer is shared between producer and consumer threads. The
solution uses one lock \texttt{m} and two CVs \texttt{fill} and
\texttt{empty}.
\begin{lstlisting}[numbers=none]
/* Consumer */                         /* Producer */
while (1) {                              for (i = 0; i < 10; i++) {
    mutex_lock(m);                           mutex_lock(m);
    if (numfull == 0)                        if (numfull == MAX)
        wait(fill, m);                           wait(empty, m);
    tmp = get();                             put();
    signal(empty);                           signal(fill);
    mutex_unlock(m);                         mutex_unlock(m);
}                                        }
\end{lstlisting}
The CV library is changed in two ways: (1)~\texttt{wait()} never wakes
unless signalled (no spurious wakeups); (2)~\texttt{signal()}
\textbf{immediately} transfers control to a waiting thread (if there is
one), and \textbf{implicitly passes the lock} to that waiter. Does this
subtle change affect the correctness of the code above?
\end{Qn}
\begin{Sol}
\textbf{Yes --- under these two assumptions, the \texttt{if}-based code becomes correct.}
This change converts the CV library from Mesa to \textbf{Hoare} semantics (immediate transfer of both CPU and lock) with no spurious wakeups. Because there are \emph{two} separate CVs (\texttt{fill} and \texttt{empty}), a producer only signals \texttt{fill} immediately after \texttt{put()} (so \texttt{numfull >= 1} holds) and immediately hands the lock and CPU to the woken consumer before any other consumer can enter and steal the item. Symmetrically, a consumer only signals \texttt{empty} immediately after \texttt{get()} (so \texttt{numfull < MAX} holds) and hands the lock directly to a waiting producer. Thus the woken thread knows the condition is still true without re-checking.

(Official handwritten key: ``Code above usually doesn't work because of \texttt{if} (need to re-check condition). [The change] makes it work, as now when the thread wakes it knows the condition is true (no need to recheck). See Mesa vs.\ Hoare semantics.'')

In real POSIX/Mesa systems (where \texttt{signal} is only a hint and spurious wakeups can occur), \texttt{while} is always mandatory. (Compare the Mesa \texttt{if}-vs-\texttt{while} stolen-wakeup problem immediately below.)
\end{Sol}

\begin{Qn}{Hard}{Lect 18}{CS-537 Fa11 Mid, Q22; Sp16 Mid, Q6 trace 5; Fa04 Final, Q8}{\texttt{if} instead of \texttt{while}: stolen wakeup}
Consider the following two-CV bounded-buffer code under standard \textbf{Mesa} semantics, where \texttt{while} has been replaced with \texttt{if}:
\begin{lstlisting}[numbers=none]
/* Producer */                            /* Consumer */
p1: mutex_lock(&m);                       c1: mutex_lock(&m);
p2: if (count == max)                     c2: if (count == 0)
p3:     cond_wait(&empty, &m);            c3:     cond_wait(&fill, &m);
p4: put(i);                               c4: tmp = get();
p5: cond_signal(&fill);                   c5: cond_signal(&empty);
p6: mutex_unlock(&m);                     c6: mutex_unlock(&m);
\end{lstlisting}
Assume two consumers and one producer (and any \texttt{max} you like). Give a concrete interleaving that breaks, and say what goes wrong.
\end{Qn}
\begin{Sol}
Take \texttt{max = 1}, \texttt{count = 0}. Consumers $C_a,C_b$,
producer $P$.

\begin{center}
\small
\begin{tabular}{@{}llll@{}}
  \textbf{Thread} & \textbf{Lines} & \textbf{count} & \textbf{Comment} \\
  \midrule
  $C_a$ & c1, c2, c3 & 0 & empty: $C_a$ sleeps on \texttt{fill} \\
  $P$   & p1, p2, p4 & 1 & puts the one item \\
  $P$   & p5, p6     & 1 & signals \texttt{fill}: $C_a$ is \emph{ready} \\
  $P$   & p1, p2, p3 & 1 & buffer full: $P$ sleeps on \texttt{empty} \\
  $C_b$ & c1, c2     & 1 & \texttt{if}: does \emph{not} wait \\
  $C_b$ & c4         & 0 & sneaks in and \textbf{steals} the item \\
  $C_b$ & c5, c6     & 0 & signals \texttt{empty}: wakes $P$ \\
  $C_a$ & c4         & 0 & returns from wait, \textbf{no re-check}: \texttt{get()} on empty \\
\end{tabular}
\end{center}
$C_a$ was signalled when the buffer was full, but $C_b$ ran first and
emptied it. Mesa does not run $C_a$ immediately, so the hint is stale.
\texttt{while} would send $C_a$ back to sleep. (This is the lecture
trace with one CV and \texttt{if}; with two CVs the same steal still
happens, because the bug is the missing re-check, not the number of
CVs.) Fa11 Mid Q22 asked for exactly this 2-consumer/1-producer
interleaving.
\end{Sol}

\begin{Qn}{Hard}{Lect 18}{CS-537 Sp16 Mid, Q6 trace 6; Fa08 Final, Q3(e); Fa21 Mid2, Q31}{One CV, even with \texttt{while}: everyone asleep}
Now restore \texttt{while}, but collapse the two CVs onto a
\textbf{single} condition variable \texttt{cond}. Both producer and
consumer wait on \texttt{cond} and signal \texttt{cond}. Assume one
producer and two consumers, no spurious wakeups, and threads deschedule
only on wait. Show a trace in which every thread is asleep and nobody
will ever signal. Why does this \emph{not} happen with two CVs?
\end{Qn}
\begin{Sol}
\texttt{max = 1}, \texttt{count = 0}. $C_a$ and $C_b$ both wait on
\texttt{cond}. $P$ puts (\texttt{count = 1}) and signals: suppose it
wakes $C_a$. $P$ loops, sees full, waits on the same \texttt{cond}.
$C_a$ re-checks (\texttt{while}), gets (\texttt{count = 0}), and
signals \texttt{cond}. The waiters on \texttt{cond} are $C_b$
\emph{and} $P$; the signal may wake $C_b$. $C_a$ loops and waits
(empty). $C_b$ re-checks, sees empty, waits again. Now $C_a$, $C_b$,
and $P$ are all asleep: a consumer woke a consumer, and the producer
has nobody left to wake it.

Flow (one of several):
\[
\begin{aligned}
C_a:&\; 1,2,3 \quad\text{(wait)}\\
C_b:&\; 1,2,3 \quad\text{(wait)}\\
P:&\; 1,2,4,5,6,1,2,3 \quad\text{(put, wake $C_a$, then wait)}\\
C_a:&\; 2,4,5 \quad\text{(get, signal --- wakes $C_b$)}\\
C_a:&\; 6,1,2,3 \quad\text{(wait)}\\
C_b:&\; 2,3 \quad\text{(re-check, wait)}.
\end{aligned}
\]
Stuck.

\textbf{Why two CVs fix it.} Producers wait only on \texttt{empty} and
signal only \texttt{fill}; consumers wait only on \texttt{fill} and
signal only \texttt{empty}. A consumer cannot wake a consumer. The
Fa21 Mid2 Q31 distractors ``producer can incorrectly wake a consumer''
/ ``consumer can incorrectly wake a producer'' are the \emph{healthy}
direction; the actual bug is same-kind wakeup. Official key: (e)~None
of the above --- the listed options never name the single-CV bug.

\textbf{Broadcast alternative.} \texttt{broadcast} on the single CV
would also unstick the trace (both $C_b$ \emph{and} $P$ wake; $P$
re-checks, sees empty, produces). That is a covering condition: correct,
wasteful. Two CVs are the clean fix.
\end{Sol}

\begin{Qn}{Medium-Hard}{Lect 18}{CS-537 Sp18 Mid, Problem XXVI (Q126--130)}{Allocator: \texttt{signal} may wake the wrong waiter}
A multi-threaded heap, sketched as follows:
\begin{lstlisting}[numbers=none]
int bytes_left = MAX_HEAP_SIZE;
pthread_cond_t c;
pthread_mutex_t m;

void *allocate(int size) {
    pthread_mutex_lock(&m);
    while (bytes_left < size)
        pthread_cond_wait(&c, &m);
    void *ptr = ...;           /* get mem from internal data structs */
    bytes_left -= size;
    pthread_mutex_unlock(&m);
    return ptr;
}
void free(void *ptr, int size) {
    pthread_mutex_lock(&m);
    bytes_left += size;
    pthread_cond_signal(&c);
    pthread_mutex_unlock(&m);
}
\end{lstlisting}
All of memory is used up (\texttt{bytes\_left = 0}).
Then: T1 calls \texttt{allocate(100)}; later T2 calls
\texttt{allocate(1000)}; later T3 calls \texttt{free(200)}.
All library calls work as expected. For each statement, answer
\textbf{A}~=~possible or \textbf{B}~=~not possible, just after this
sequence.

\begin{enumerate}[(a)]
\item 126. T1 and T2 remain blocked inside \texttt{allocate()}.
\item 127. T1 becomes unblocked, gets 100 bytes, and returns.
\item 128. T2 becomes unblocked, gets 1000 bytes, and returns.
\item 129. T3 becomes blocked inside \texttt{free()}.
\item 130. T1, T2, and T3 become deadlocked.
\end{enumerate}
Then: what one-line change makes the allocator correct, and why is
that change a \emph{covering condition}?
\end{Qn}
\begin{Sol}
\texttt{signal} wakes \emph{one} waiter, and the signaller cannot name
which. 200 free bytes suffice for T1 (wants 100) but not for T2 (wants
1000). If the signal wakes T2, T2 re-checks \texttt{200 < 1000}, goes
back to sleep, and T1 is never woken --- progress is lost even though
an allocation is possible.

\begin{enumerate}[(a)]
\item 126. \textbf{A (possible).} Signal may wake T2; T2 re-waits; T1
      sleeps forever. Official key.
\item 127. \textbf{A (possible).} If instead T1 is awoken, $200 \ge 100$,
      T1 allocates and returns.
\item 128. \textbf{B (not possible).} Only 200 bytes are free; T2's
      \texttt{while} will not let it allocate 1000.
\item 129. \textbf{B (not possible).} \texttt{free} never waits on
      the CV; the only lock is held and released in one straight-line
      section. (A brief wait to \emph{acquire} the mutex is not
      ``blocked inside \texttt{free}'' in the exam's sense.)
\item 130. \textbf{B (not possible).} One lock, no lock cycle.
\end{enumerate}

\textbf{Fix:} replace \texttt{pthread\_cond\_signal(\&c)} by
\texttt{pthread\_cond\_broadcast(\&c)}. Both waiters wake. Whichever
runs first re-checks: T1 takes 100 bytes (\texttt{bytes\_left = 100});
T2 sees $100 < 1000$ and sleeps again. The broadcast
\textbf{covers} every waiter that might now be able to run, because
the signaller cannot know the sizes. Cost: extra wakeups of threads
that will immediately wait again (Lampson--Redell covering condition;
OSTEP Ch.~30 / Lect 18 allocator).
\end{Sol}

\subsection{Lecture 19: Semaphores}

\begin{Qn}{Hard}{Lect 19}{Mandar Mitra, ISI Kolkata, 2014 End-sem Q1(b)}{Elephants and lions at the lake}
Once upon a time, somewhere in Africa, there was a big lake. A large herd of
elephants and a pride of lions would both come to the lake to drink, but the
two species always avoided each other.

Model each animal by a process that executes \texttt{elephant()} or
\texttt{lion()}. You may assume a function \texttt{drink()}. Write
pseudo-code for the two functions, using only semaphores (and ordinary
integers), so that two different species are never simultaneously executing
\texttt{drink()}. Full credit requires a \emph{starvation-free} solution:
a waiting lion must eventually drink even if elephants keep arriving, and
vice versa.

State every semaphore's initial value and argue mutual exclusion of the two
species, deadlock freedom, and starvation freedom.
\end{Qn}
\begin{Sol}
Same-species animals may drink together; mixed species may not. This is the
two-tribe analogue of a reader/writer lock in which \emph{each} tribe plays
the role of ``readers''. The first-readers-writers lightswitch (count +
group lock) gives species exclusion but lets a continuous stream of
elephants starve lions. Fairness is restored with a \emph{turnstile} that a
waiter of the other species can hold.

\begin{lstlisting}[numbers=none]
sem_t mutexE = 1, mutexL = 1;   /* protect the two counts   */
sem_t empty    = 1;             /* lake free / empty        */
sem_t turnstile = 1;            /* FIFO door                */
int elephants = 0, lions = 0;

void elephant(void) {
    sem_wait(&turnstile);
    sem_wait(&mutexE);
    elephants++;
    if (elephants == 1)
        sem_wait(&empty);       /* first elephant claims the lake */
    sem_post(&mutexE);
    sem_post(&turnstile);

    drink();

    sem_wait(&mutexE);
    elephants--;
    if (elephants == 0)
        sem_post(&empty);       /* last elephant frees the lake  */
    sem_post(&mutexE);
}
\end{lstlisting}
\texttt{lion()} is symmetric (\texttt{mutexL}, \texttt{lions}, same
\texttt{empty} and \texttt{turnstile}).

\emph{Species exclusion.} \texttt{empty} is a binary semaphore. The first
elephant to arrive takes it; later elephants only increment the count. A
lion's first arrival also needs \texttt{empty}, so no lion can
\texttt{drink} until the last elephant has posted it. Symmetric for lions.

\emph{Deadlock.} When the first elephant blocks on \texttt{empty} while
holding \texttt{mutexE}, \texttt{elephants} was previously 0, so no
elephant is currently drinking and thus no exiting elephant needs
\texttt{mutexE}; meanwhile \texttt{turnstile} blocks any newly arriving
elephant. Only lions (which use \texttt{mutexL}, not \texttt{mutexE})
are drinking, so the last lion will acquire \texttt{mutexL} and post
\texttt{empty} without blocking (and symmetrically for lions). Hence no
deadlock can occur.

\emph{Starvation freedom} (assuming FIFO semaphore queues). Suppose elephants
are drinking and a lion arrives. The lion takes \texttt{turnstile} and, being
first of its kind, waits on \texttt{empty} \emph{while still holding
\texttt{turnstile}}. No further elephant can pass the door. Currently
drinking elephants finish, the last posts \texttt{empty}, the lion enters
and releases the turnstile. Bound: at most the elephants already past the
door drink ahead of a waiting lion. Symmetric for a waiting elephant.
\end{Sol}

\begin{Qn}{Hard}{Lect 19}{Mandar Mitra, ISI Kolkata, 2010 Mid-sem Q4; also CS-537 Fall 2004 Midterm 1 Q8, Fall 2013 Midterm Q15}{First readers--writers, and writer starvation}
Consider the \emph{first} readers--writers problem (readers have priority).
Readers and writers synchronise with semaphores.

\begin{enumerate}
  \item[(a)] Write the reader and writer code (including initial values).
  \item[(b)] Does the solution satisfy bounded waiting? If not, exhibit a
        schedule in which a writer waits forever.
  \item[(c)] \textbf{[CS-537 Fall 2013 Midterm Q15]} The following code
        was presented with the comment ``pretty sure this is useful for
        something''. Identify the protocol and the fairness bug.
\end{enumerate}
\begin{lstlisting}[numbers=none]
typedef struct _foo_t { sem_t lock1, lock2; int count; } foo_t;
void foo_init(foo_t *f) {
    f->count = 0;
    sem_init(&f->lock1, 1); sem_init(&f->lock2, 1);
}
void do1(foo_t *f) {
    sem_wait(&f->lock1); f->count += 1;
    if (f->count == 1) sem_wait(&f->lock2);
    sem_post(&f->lock1);
}
void undo1(foo_t *f) {
    sem_wait(&f->lock1); f->count -= 1;
    if (f->count == 0) sem_post(&f->lock2);
    sem_post(&f->lock1);
}
void do2(foo_t *f)   { sem_wait(&f->lock2); }
void undo2(foo_t *f) { sem_post(&f->lock2); }
\end{lstlisting}
\end{Qn}
\begin{Sol}
(a)~OSTEP / Lect 19 reader/writer lock. \texttt{lock} (binary, init 1)
protects the integer \texttt{readers}; \texttt{writelock} (binary, init 1)
is the exclusive resource. First reader takes \texttt{writelock} on behalf
of all readers; last reader releases it. A writer just waits on
\texttt{writelock}.
\begin{lstlisting}[numbers=none]
sem_t lock = 1, writelock = 1;  int readers = 0;

void reader(void) {                  void writer(void) {
  sem_wait(&lock);                     sem_wait(&writelock);
  readers++;                           /* write */
  if (readers == 1)                    sem_post(&writelock);
      sem_wait(&writelock);          }
  sem_post(&lock);
  /* read */
  sem_wait(&lock);
  readers--;
  if (readers == 0) sem_post(&writelock);
  sem_post(&lock);
}
\end{lstlisting}

(b)~No. Bounded waiting fails for writers. Schedule: reader $R_1$ enters first, acquires \texttt{writelock} (\texttt{readers = 1}), and begins reading. Writer $W$ then calls \texttt{sem\_wait(\&writelock)} and blocks. While \texttt{readers} stays positive, new readers $R_2, R_3, \ldots$ keep entering (they only take \texttt{lock}, see \texttt{readers >= 1}, and skip \texttt{writelock}) before earlier readers leave. \texttt{readers} never hits 0, so nobody posts \texttt{writelock} and $W$ starves. (By contrast, the ``second'' readers--writers problem gives
priority to writers, which can starve readers.)

(c)~Official CS-537 Fall 2013 answer: this \emph{is} a reader/writer lock.
\texttt{do1}/\texttt{undo1} are read-acquire/read-release
(\texttt{lock1} $=$ mutex, \texttt{lock2} $=$ writelock,
\texttt{count} $=$ readers). \texttt{do2}/\texttt{undo2} are
write-acquire/write-release. Successfully calling \texttt{do1()}
guarantees that no \texttt{do2()} succeeds until every shared-mode user
has finished --- the same writer-starvation problem as (b).
\end{Sol}

\begin{Qn}{Hard}{Lect 19}{Mandar Mitra, ISI Kolkata, 2012 End-sem Q2(b)--(c); also 2017 End-sem Q1, CS-537 Fall 2004 Midterm 1 Q7}{Dining philosophers: deadlock, asymmetry, footman, respected guest}
Five philosophers, five chopsticks, semaphore \texttt{chopstick[5]} each
init 1. Philosopher $i$ runs
\begin{lstlisting}[numbers=none]
while (1) {
    think();
    wait(chopstick[i]); wait(chopstick[(i+1)%5]);
    eat();
    signal(chopstick[i]); signal(chopstick[(i+1)%5]);
}
\end{lstlisting}

\begin{enumerate}
  \item[(a)] Exhibit a deadlocking interleaving.
  \item[(b)] The odd/even (lefty/righty) fix is one way to avoid deadlock.
        Give an \emph{alternate} deadlock-free semaphore solution. (Hint,
        Mandar 2012 Q1(b): prevent circular wait.) Also state whether the
        CS-537 / Lect 19 code below deadlocks, and why.
\begin{lstlisting}[numbers=none]
void acquire(int i) {
    if (i < 4) { sem_wait(chop[i]); sem_wait(chop[i+1]); }
    else       { sem_wait(chop[0]); sem_wait(chop[4]); }
}
\end{lstlisting}
  \item[(c)] Variant: one philosopher is highly respected. When he
        announces that he is hungry, his two neighbours give him their
        forks. Can deadlock still occur? If yes, give a deadlocked
        situation; if not, which Coffman condition is violated?
  \item[(d)] Write a deadlock-free POSIX-style \texttt{philosopher(int i)}
        (Mandar 2017 End Q1). Show creation/initialisation and use. Do
        \emph{not} use System V \texttt{semget}/\texttt{semop}.
\end{enumerate}
\end{Qn}
\begin{Sol}
(a)~Each $P_i$ takes \texttt{chopstick[i]} (left) before anyone takes a
right chopstick. Then $P_i$ waits for \texttt{chopstick[(i+1)\%5]}, held by
$P_{i+1}$. Cycle $P_0 \to P_1 \to P_2 \to P_3 \to P_4 \to P_0$. All five
sleep holding one chopstick.

(b)~Two standard alternate fixes, both destroying circular wait.

\emph{Global lock order} (the 2012 hint). Every philosopher picks up the
lower-numbered chopstick first:
\begin{lstlisting}[numbers=none]
lo = min(i, (i+1)%5);  hi = max(i, (i+1)%5);
wait(chopstick[lo]); wait(chopstick[hi]);
\end{lstlisting}
No cycle can form in the wait-for graph.

\emph{Footman / multiplex.} A sixth semaphore \texttt{room} init 4.
\texttt{wait(room)} before picking up any chopstick,
\texttt{signal(room)} after putting both down. At most four philosophers
are at the table, so at least one of those four can obtain two chopsticks;
the hold-and-wait cycle of all five cannot close.

\emph{CS-537 / Lect 19 reverse-one.} Official: this does \emph{not}
deadlock. Philosophers $0$--$3$ take left then right; philosopher 4 takes
\texttt{chop[0]} then \texttt{chop[4]} (right then left). If $0$--$3$ each
hold their left chopstick, $P_4$ waits on \texttt{chop[0]} \emph{without}
holding \texttt{chop[4]}, so $P_3$ can take \texttt{chop[4]}, eat, and
release. The circular wait is impossible. (If every philosopher is a
``righty'', deadlock is again possible by symmetry; at least one lefty and
one righty is enough.)

(c)~\textbf{No --- deadlock cannot occur.} Let the respected philosopher be $P_0$ (whose two forks under the given code are \texttt{chopstick[0]} and \texttt{chopstick[1]}). There are two exhaustive cases:
\begin{itemize}
  \item If $P_0$ is \emph{not} hungry (thinking), he holds zero forks. Among the remaining four philosophers $P_1,\ldots,P_4$, only $P_4$ ever requests \texttt{chopstick[0]}, so any wait-for relation among $P_1,\ldots,P_4$ forms a chain ending at $P_4$ (who can take \texttt{chopstick[0]} and eat), never a cycle.
  \item If $P_0$ \emph{is} hungry, he announces hunger and both neighbours surrender their forks to him. Thus $P_0$ acquires both forks, eats, and releases them, breaking any potential cycle.
\end{itemize}
The Coffman condition violated is \textbf{No Preemption} (neighbours yield their held forks on demand when $P_0$ is hungry; when $P_0$ is thinking, \textbf{Circular Wait} cannot form because $P_0$ holds nothing).

(d)~Asymmetric POSIX solution (SysV APIs are out of syllabus):
\begin{lstlisting}[numbers=none]
sem_t chop[5];
void init(void) {
    int i;
    for (i = 0; i < 5; i++) sem_init(&chop[i], 0, 1);
}
void philosopher(int i) {
    int left = i, right = (i + 1) % 5;
    while (1) {
        think();
        if (i == 4) {                 /* reverse one philosopher */
            sem_wait(&chop[right]); sem_wait(&chop[left]);
        } else {
            sem_wait(&chop[left]);  sem_wait(&chop[right]);
        }
        eat();
        sem_post(&chop[left]); sem_post(&chop[right]);
    }
}
\end{lstlisting}
At most $\lfloor 5/2 \rfloor = 2$ philosophers eat at once (each eater
uses two of five forks).
\end{Sol}

\begin{Qn}{Hard}{Lect 19}{Mandar Mitra, ISI Kolkata, 2011 Mid-sem Q4; 2018 End-sem Q1(b)}{Counting semaphore from binary, and the bounded-waiting failure}
\begin{enumerate}
  \item[(a)] Let \texttt{wait\_b}/\texttt{signal\_b} be the operations on
        \emph{binary} semaphores. Implement a counting semaphore $C$
        (operations \texttt{wait}/\texttt{signal}) using only binary
        semaphores and an integer. How many binary semaphores does your
        implementation need, and why?
  \item[(b)] \textbf{[2018 End Q1(b)]} Consider this attempt
        (\texttt{S1 = 1}, \texttt{S2 = 0}, integer \texttt{C = 1}):
\begin{lstlisting}[numbers=none]
wait_c(C) {                    signal_c(C) {
    wait(S1);                      wait(S1);
    C--;                           C++;
    if (C < 0) {                   if (C <= 0) signal(S2);
        signal(S1);                signal(S1);
        wait(S2);              }
    } else signal(S1);
}
\end{lstlisting}
        Two processes $P_1,P_2$ each run
        \texttt{wait\_c(C); critical\_section(); signal\_c(C);}.
        Complete a schedule showing that the number of times $P_2$ can
        re-enter the critical section before $P_1$ enters is
        \emph{not} bounded by a constant independent of the scheduler.
\end{enumerate}
\end{Qn}
\begin{Sol}
(a)~Two binary semaphores plus an integer. \texttt{S1} (init 1) is a mutex
protecting $C$; \texttt{S2} (init 0) is a delay semaphore on which waiters
sleep. The integer $C$ holds the counting-semaphore value (and, when
negative, the number of waiters --- matching Lect 19).

\emph{Important subtlety (baton-passing):} if \texttt{signal} released
\texttt{S1} after posting \texttt{S2}, two back-to-back \texttt{signal}s
(when two or more waiters have decremented $C$ and released \texttt{S1}
but not yet reached \texttt{wait\_b(S2)}) could post the \emph{binary}
semaphore \texttt{S2} twice before any waiter consumes the first post,
losing the second wakeup and deadlocking a waiter. The correct construction passes the \texttt{S1}
mutex ``baton'' from the signaller directly to the woken waiter:
\begin{lstlisting}[numbers=none]
wait(C):
    wait_b(S1);
    C--;
    if (C < 0) {
        signal_b(S1);
        wait_b(S2);   /* woken waiter inherits S1 from signal(C) */
    }
    signal_b(S1);

signal(C):
    wait_b(S1);
    C++;
    if (C <= 0)
        signal_b(S2); /* pass S1 baton to one waiter; do NOT signal S1 */
    else
        signal_b(S1);
\end{lstlisting}
You cannot fold the mutex and the delay into one binary semaphore: the
mutex must be releasable by a thread that is about to sleep, while the
delay must be posted by a \emph{different} thread that is incrementing $C$.
Baton-passing ensures at most one unconsumed \texttt{signal\_b(S2)} is ever
in flight, and also prevents outsiders from stealing the wakeup before the
woken waiter passes \texttt{S2}.

(b)~In the \emph{broken} attempt given in part~(b) (which releases
\texttt{S1} unconditionally in \texttt{signal\_c} instead of baton-passing),
the binary delay \texttt{S2} remembers a wakeup that a re-entering $P_2$ can
steal before $P_1$ reaches \texttt{wait(S2)}:

\begin{center}
\small
\begin{tabularx}{\linewidth}{@{}lX@{}}
  \toprule
  Process & Completes \\
  \midrule
  $P_2$ & \texttt{wait\_c} (C: $1\to 0$), critical section \\
  $P_1$ & \texttt{wait\_c} through \texttt{C--} (C: $0\to -1$), then
          \texttt{signal(S1)}, but \emph{not yet} \texttt{wait(S2)} \\
  $P_2$ & \texttt{signal\_c}: \texttt{C++} (C: $-1\to 0$),
          \texttt{C <= 0} so \texttt{signal(S2)} ($S2 = 1$),
          \texttt{signal(S1)} \\
  $P_2$ & \texttt{wait\_c} again: \texttt{C--} (C: $0\to -1$),
          \texttt{C < 0}, \texttt{signal(S1)}, \texttt{wait(S2)}
          --- consumes the remembered $S2=1$ and \emph{re-enters} CS \\
  $P_1$ & (after $P_2$ repeats the previous two rows $k$ times) now executes \texttt{wait(S2)} and blocks \\
  \bottomrule
\end{tabularx}
\end{center}
$P_2$ can repeat ``\texttt{signal\_c} then \texttt{wait\_c} that steals
\texttt{S2}'' arbitrarily often, depending on how long $P_1$ is delayed
before its \texttt{wait(S2)}. The waiting bound is not a fixed constant;
it is a function of the scheduler. (With the baton-passing fix of part~(a),
$P_2$'s second \texttt{wait\_c} blocks on \texttt{S1} until $P_1$ has
consumed \texttt{S2} and released \texttt{S1}; full $n$-process fairness
still also requires FIFO queues on \texttt{S1}/\texttt{S2}.)
\end{Sol}

\begin{Qn}{Medium-Hard}{Lect 19}{CS-537 Spring 2010 Final Q8}{Why a condition variable built from a semaphore is buggy}
A co-worker implements CVs using a semaphore:
\begin{lstlisting}[numbers=none]
typedef struct __cond_t { sem_t s; } cond_t;
void cond_init(cond_t *c) { sem_init(&c->s, 0); }
void cond_wait(cond_t *c, mutex_t *m) {
    mutex_unlock(m);          /* release lock, then sleep */
    sem_wait(&c->s);
    mutex_lock(m);
}
void cond_signal(cond_t *c) {
    sem_post(&c->s);          /* wake one waiter, if any */
}
\end{lstlisting}
Why doesn't this work properly? (Assume Mesa CV semantics as in Lect 18.)
\end{Qn}
\begin{Sol}
Two major defects break standard Mesa CV semantics:

\emph{1.~Remembered signals (accumulation of stale credits).} A Mesa
\texttt{cond\_signal} with nobody waiting is a no-op. By contrast,
\texttt{sem\_post} is \emph{never} lost (Lect 19): it increments the
semaphore counter every time. Repeated \texttt{cond\_signal}s when no
thread is waiting pile up positive credits in \texttt{c->s}, so future
\texttt{cond\_wait} calls return immediately without sleeping even when
no signaller just established the condition. With \texttt{while} this
causes useless wakeups/spinning; with any caller relying on blocking
until a \emph{subsequent} signal (e.g.\ barriers or \texttt{if}-style
waits), it breaks correctness outright.

\emph{2.~Cannot implement \texttt{cond\_broadcast} (and non-atomic
unlock/sleep).} Because the implementation tracks no waiter count under
the mutex, it cannot know how many threads are sleeping to wake them
all in a broadcast. (Note: unlike a real CV, the non-atomic gap between
\texttt{mutex\_unlock} and \texttt{sem\_wait} does \emph{not} cause a
lost wakeup here because \texttt{sem\_post} is counted, though a signal
fired in that window may wake a different waiter than expected.)

A correct CV built from semaphores must track a waiter count \emph{under
the caller's mutex} so that \texttt{signal} posts only when waiters are
present (no remembered credits) and \texttt{broadcast} knows how many
posts to issue --- which is why Lecture~19 notes that building CVs from
semaphores is surprisingly tricky.
\end{Sol}

\subsection{Lecture 20: Concurrency Bugs \& Deadlock}

\begin{Qn}{Hard}{Lect 20}{CS-537 Spring 2010 Final, Problem 7}{\texttt{vector\_add}: lock-order inversion on two vectors}
You are given the following code, which adds two vectors together, and does
so in a multithread-safe way.
\begin{lstlisting}[numbers=none]
void
vector_add(vector *v1, vector *v2) {
    mutex_lock(v1->lock);
    mutex_lock(v2->lock);
    for (i = 0; i < v1->size; i++) {
        v1[i] = v1[i] + v2[i];
    }
    mutex_unlock(v1->lock);
    mutex_unlock(v2->lock);
}
\end{lstlisting}
Then you are told that two different concurrently-executing threads, 1 and
2, call this code as follows:
\begin{lstlisting}[numbers=none]
Thread 1:                          Thread 2:
  vector_add(&vectorA, &vectorB);    vector_add(&vectorB, &vectorA);
\end{lstlisting}
Unfortunately, this can lead to a deadlock, in which the program gets stuck,
with each thread waiting for the other to make progress.
\begin{enumerate}
  \item[(a)] Why does this happen? (Describe, or show with a picture.)
  \item[(b)] How could you write \texttt{vector\_add()} so that this
        deadlock never happens?
\end{enumerate}
\end{Qn}
\begin{Sol}
(a)~Lock-order inversion. Thread~1 locks \texttt{vectorA->lock} then tries
\texttt{vectorB->lock}. Thread~2 locks \texttt{vectorB->lock} then tries
\texttt{vectorA->lock}. Interleaving:
\begin{lstlisting}[numbers=none]
T1: mutex_lock(vectorA->lock)     /* holds A */
T2: mutex_lock(vectorB->lock)     /* holds B */
T1: mutex_lock(vectorB->lock)     /* waits for B */
T2: mutex_lock(vectorA->lock)     /* waits for A */
\end{lstlisting}
A cycle $T1 \to B \to T2 \to A \to T1$. All four Coffman conditions hold:
the locks are exclusive; each thread holds one and waits for the other;
neither lock is forcibly taken; the wait-for graph has a cycle.

(b)~Official / Lect~20 fix: \emph{total lock order} (break circular wait).
Acquire by lock address (or any stable lock ID), so every call grabs the
``lower'' lock first:
\begin{lstlisting}[numbers=none]
void vector_add(vector *v1, vector *v2) {
    mutex_t *L = v1->lock, *H = v2->lock;
    if (L == H) {                  /* same vector */
        mutex_lock(L);
        /* ... add ... */
        mutex_unlock(L);
        return;
    }
    if ((uintptr_t)L > (uintptr_t)H) { mutex_t *t = L; L = H; H = t; }
    mutex_lock(L);
    mutex_lock(H);
    /* ... add ... */
    mutex_unlock(H);
    mutex_unlock(L);
}
\end{lstlisting}
Both threads then acquire \texttt{vectorA->lock} before \texttt{vectorB->lock} (assuming
\texttt{vectorA.lock < vectorB.lock}), so a cycle is impossible.

Other Lect~20 preventions that also work here:
\begin{itemize}
  \item \texttt{trylock} on the second lock; on failure, drop the first and
        retry (break no-preemption). Watch livelock; add exponential
        back-off.
  \item Grab both locks under a process-wide meta-lock (break
        hold-and-wait). Degenerates to one big lock.
\end{itemize}
Lock ordering is the usual encapsulation-friendly fix for this pattern
(OSTEP ``transfer/swap two accounts'').
\end{Sol}

\begin{Qn}{Hard}{Lect 20}{CS-537 Spring 2016 Midterm, Q2}{Deadlock or not deadlock: seven lock-order arrangements}
Some of the following arrangements of threads, and the locks that they
will grab, can lead to deadlock. For each, answer: could this deadlock?
(Yes / No.)

\begin{enumerate}
  \item[(i)] Thread~1 will try to grab locks 1 and 2 in an \emph{arbitrary}
        order. Thread~2 will try to grab locks 1 and 2 in \emph{fixed}
        order, 1 then 2.
  \item[(ii)] Thread~1: locks 1 and 2 in fixed order, 1 then 2.
        Thread~2: locks 1 and 2 in fixed order, 2 then 1.
  \item[(iii)] Thread~1: locks 1 and 2 in fixed order, 1 then 2.
        Thread~2: locks 1, 2 and 3 in fixed order, 1 then 2 then 3.
  \item[(iv)] Thread~1: locks 1 and 2 in arbitrary order.
        Thread~2: locks 2 and 3 in arbitrary order.
        Thread~3: locks 1 and 3 in arbitrary order.
  \item[(v)] Thread~1: locks 1 and 2 in arbitrary order.
        Thread~2: locks 2 and 3 in arbitrary order.
        Thread~3: will try to grab lock 1.
  \item[(vi)] Thread~1: locks 1 and 2 in fixed order, 1 then 2.
        Thread~2: locks 2 and 3 in fixed order, 2 then 3.
        Thread~3: locks 1 and 3 in fixed order, 1 then 3.
  \item[(vii)] Thread~1: locks 1 and 2 in fixed order, 1 then 2.
        Thread~2: locks 2 and 3 in fixed order, 2 then 3.
        Thread~3: locks 1 and 2 in arbitrary order.
\end{enumerate}
\end{Qn}
\begin{Sol}
Official CS-537 Spring 2016 answers (from the answer sheet), with the
cycle test: deadlock is possible iff some interleaving of the allowed
orders produces a circular wait.

(i)~\textbf{Yes.} Thread~1 may grab 2 then 1, while Thread~2 grabs 1 then
2. Same inversion as \texttt{vector\_add}.

(ii)~\textbf{Yes.} The two threads use opposite total orders on $\{1,2\}$.
Classic lock-order inversion; a cycle is forced as soon as each takes its
first lock.

(iii)~\textbf{No.} (Official: ``no cycle''.) Both threads acquire in
increasing lock-ID order. Wait-for edges only point toward a
higher-numbered lock still held by someone who will release it after
finishing a prefix of $1 < 2 < 3$. The wait-for graph is acyclic.

(iv)~\textbf{Yes.} Each adjacent pair shares two locks with no agreed
order, so
\[
T1\ \text{holds 1 waits 2},\quad
T2\ \text{holds 2 waits 3},\quad
T3\ \text{holds 3 waits 1}
\]
is allowed. Three-cycle.

(v)~\textbf{No.} Thread~3 takes only lock~1. The only shared locks are
$\{1\}$ (T1--T3) and $\{2\}$ (T1--T2). One shared lock cannot close a
cycle: the holder of that lock is not waiting for the other thread. T2
and T3 share nothing.

(vi)~\textbf{No.} (Official: ``ordered acquire''.) All three use
increasing IDs. Thread~3 in particular takes 1 then 3, so it can never
\emph{hold} 3 while waiting for 1. The would-be cycle
$T1(1{\to}2), T2(2{\to}3), T3(3{\to}1)$ is illegal under T3's order.

(vii)~\textbf{Yes.} Thread~3 may grab 2 then 1 while Thread~1 holds 1 and
waits for 2. T2 is irrelevant; T1 and T3 alone deadlock on $\{1,2\}$.
\end{Sol}

\begin{Qn}{Medium-Hard}{Lect 20}{CS-537 Fall 2021 Midterm 2, Q24}{Unlocking \texttt{outer} inside \texttt{function2}: deadlock or livelock?}
Bug was told to speed up \texttt{function2()}. There were two locks
(\texttt{outer} and \texttt{inner}), so Bug unlocked the outer lock so as
to allow more concurrency. Could this lead to a problem?
\begin{lstlisting}[numbers=none]
function1() {
    mutex_lock(&outer);
    mutex_lock(&inner);
    function2();
    mutex_unlock(&inner);
    mutex_unlock(&outer);
}
function2() {
    mutex_unlock(&outer);
    // do some stuff
    mutex_lock(&outer);
}
\end{lstlisting}
Assuming two (or more) threads call into \texttt{function1()}, this
code\ldots
\begin{enumerate}
  \item[(a)] will always run to completion
  \item[(b)] can deadlock
  \item[(c)] will deadlock
  \item[(d)] will livelock
  \item[(e)] None of the above
\end{enumerate}
\end{Qn}
\begin{Sol}
Official CS-537 Fall 2021 Midterm~2 answer: \textbf{(b) can deadlock}.

If one thread runs \texttt{function1} to completion before the other
really does anything, both finish. So (a) and (c) are false: deadlock is
possible, not guaranteed, and the code does not \emph{always} complete.

Deadlocking schedule (official):
\begin{lstlisting}[numbers=none]
T1: lock outer, lock inner, call function2, unlock outer
T2: lock outer                  /* T1 holds inner; T2 holds outer */
T1: lock outer                  /* waits for outer, still holds inner */
T2: lock inner                  /* waits for inner, still holds outer */
\end{lstlisting}
Cycle $T1 \to \texttt{outer} \to T2 \to \texttt{inner} \to T1$.

This is \emph{not} livelock (d). Livelock (Lect~20, \texttt{trylock} +
retry with no back-off) is threads that keep running, fail, release,
retry, and make no progress. Here both threads are \emph{blocked} on a
mutex --- classic deadlock, not a retry loop. Hence (b), not (d).
\end{Sol}

\begin{Qn}{Hard}{Lect 20}{Mandar Mitra, ISI Kolkata, 2010 End-sem Q4}{RAG: construct a cycle that deadlocks, and a cycle that does not}
\begin{enumerate}
  \item[(a)] Describe how a resource-allocation graph (RAG) is constructed
        for a given set of processes and resources.
  \item[(b)] Let $P = \{P_1,P_2,P_3,P_4\}$ be a set of processes and
        $R = \{R_1,R_2,R_3,R_4\}$ a set of resources. Each resource $R_i$
        ($1 \le i \le 4$) has \emph{exactly two} instances, and each
        process $P_i$ requires \emph{at most one} instance of any
        resource. Construct a RAG $G = (P \cup R, E)$ such that
        \begin{enumerate}
          \item[(i)] $G$ has a cycle \emph{and} $P_1,P_2,P_3,P_4$ are
                deadlocked;
          \item[(ii)] $G$ has a cycle but $P_1,P_2,P_3,P_4$ are
                \emph{not} deadlocked.
        \end{enumerate}
\end{enumerate}
\end{Qn}
\begin{Sol}
(a)~Bipartite directed graph. Vertices: one per process, one per
\emph{resource type}. Edges:
\begin{itemize}
  \item \emph{Assignment} (allocation) $R_j \to P_i$ if $P_i$ currently
        holds one instance of $R_j$ (draw one such edge per held
        instance).
  \item \emph{Request} $P_i \to R_j$ if $P_i$ is waiting for an instance
        of $R_j$.
\end{itemize}
A cycle is Coffman circular wait. With a \emph{single} instance per type,
a cycle is necessary \emph{and} sufficient for deadlock. With \emph{multiple} instances, a cycle is necessary but not sufficient: if
another process holding an extra instance of a resource on the cycle is
\emph{not} blocked and can finish and release it, the circular wait unwinds.

(b)~Two instances each, at most one instance of a given type per process.

(i)~\textbf{Cycle and deadlock.} Allocate \emph{both} instances of every
type to the four waiters, so no request on the cycle can be granted:
\begin{lstlisting}[numbers=none]
holds:  P1:{R1,R4}  P2:{R1,R2}  P3:{R2,R3}  P4:{R3,R4}
wants:  P1:R2       P2:R3       P3:R4       P4:R1
\end{lstlisting}
Each type's two instances are held; each $P_i$ holds instances of two types
and waits for a third whose both instances are held by other blocked processes. Cycle
\[
P_1 \to R_2 \to P_2 \to R_3 \to P_3 \to R_4 \to P_4 \to R_1 \to P_1.
\]
Nothing can run; all four are deadlocked.

(ii)~\textbf{Cycle, not deadlocked.} Keep all requested resource types
fully allocated (so requests legitimately block), but let processes
\emph{not} on the cycle hold the second instances:
\begin{lstlisting}[numbers=none]
holds:  P1:{R1}  P2:{R2}  P3:{R2}  P4:{R1}
wants:  P1:R2    P2:R1    P3:none  P4:none
\end{lstlisting}
Both instances of $R_1$ and $R_2$ are held, and there is a cycle
$P_1 \to R_2 \to P_2 \to R_1 \to P_1$. However, $P_3$ and $P_4$ are not
waiting for any resource: $P_3$ finishes and releases $R_2$, which is
granted to $P_1$, allowing $P_1$ to finish and release $R_1$ for $P_2$.
Thus no process is deadlocked. Circular wait is not sufficient when
multi-instance resources have non-blocked holders. That is the point of
the two-instance constraint, and it is \emph{not} Banker's algorithm: we
never compute a safe sequence or a Need matrix.
\end{Sol}

\clearpage

\section{Lectures 21--24 (incl.\ 22A): Persistence --- I/O Devices, HDDs, RAID \& File Systems}

\subsection{Lectures 21, 22 \& 22A: I/O Devices, Hard Disk Drives \& RAID}

\begin{Qn}{Hard}{Lect 22}{CS-537 Fa13 Final, Q1}{Worst-case and average HDD I/O time}
Assume a hard-disk drive with a transfer rate of $200$\,MB/s, a
\textbf{maximum} seek time of $9$\,ms, and a rotation rate of
$10{,}000$\,RPM.
\begin{enumerate}[(a)]
\item When reading a $4$\,KB block, what is the worst-case time?
\item What would you expect the average time to read a $4$\,KB block to be?
\item When reading a $2$\,MB block, what is the worst-case time?
\item What would you expect the average time to read a $2$\,MB block to be?
\item You have a series of ten $4$\,KB reads. What series of ten reads would
      you send to the drive to get the \textbf{worst-case} performance?
      Assume the disk scheduler is FIFO.
\item How long will these ten requests take to complete?
\item Now assume a shortest-access-time-first (SATF) scheduler. What ten
      requests could you send to get the worst-case performance, given that
      the disk will re-order requests as per SATF?
\item How long will these ten SATF requests take to complete?
\end{enumerate}
\end{Qn}
\begin{Sol}
I/O time is $\mathrm{seek}+\mathrm{rotation}+\mathrm{transfer}$.
Full rotation $= 60/10{,}000 = 6$\,ms, so average rotation $= 3$\,ms.
Average seek $\approx \tfrac13$ of the maximum, i.e.\ $3$\,ms.
Transfer of $4$\,KB at $200$\,MB/s:
\[
\frac{4\times 10^{3}}{200\times 10^{6}} = 2\times 10^{-5}\,\mathrm{s}
= 0.02\,\mathrm{ms}.
\]
Transfer of $2$\,MB: $2/200 = 10$\,ms.

(a) Worst $4$\,KB: $9 + 6 + 0.02 = \mathbf{15.02}$\,ms.

(b) Average $4$\,KB: $3 + 3 + 0.02 = \mathbf{6.02}$\,ms.

(c) Worst $2$\,MB: $9 + 6 + 10 = \mathbf{25}$\,ms.

(d) Average $2$\,MB: $3 + 3 + 10 = \mathbf{16}$\,ms.

(e) Alternate distinct sectors on the innermost and outermost cylinders
(so the drive's track cache cannot satisfy repeats) so that every FIFO
request pays a \textbf{full-stroke seek} plus nearly a full rotation.

(f) $10\times 15.02 = \mathbf{150.2}$\,ms.

(g) SATF reorders a closed batch, so FIFO ping-ponging no longer works.
To maximize time under SATF, place the ten requests at the \emph{same}
angular position (just behind the head after each access) on ten
distinct tracks across the disk whose inter-track seek is smaller than a
rotation. Whichever remaining request SATF selects next, it must pay a
seek plus nearly a full rotation. (A simpler two-extreme split --- five
innermost, five outermost --- forces only two full-stroke seeks plus two
rotations.)

(h) For the ten-track construction: one full-stroke seek ($9$\,ms)
spread across the sweep plus up to ten rotations ($10\times 6$\,ms) and
ten transfers ($0.2$\,ms) gives $\approx 9 + 60 + 0.2 = \mathbf{69.2}$\,ms.
(For the simpler 5-innermost $+$ 5-outermost batch, two full positionings
suffice: $2\times(9+6)+10\times 0.02 = \mathbf{30.2}$\,ms.)
\end{Sol}

\begin{Qn}{Medium-Hard}{Lect 21}{CS-537 Fa13 Final, Q8(c)--(d)}{Polling versus interrupts, and the hybrid}
\begin{enumerate}[label=(\alph*), start=3]
\item You are trying to decide whether to \textbf{poll} a particular
      device or to use \textbf{interrupts} to get notified of device
      completion. When is one approach better than the other?
\item Could there be a situation where both polling and interrupts
      should be used? Describe.
\end{enumerate}
\end{Qn}
\begin{Sol}
(c) \textbf{Poll} when the device is fast (completion is expected in
well under the cost of a context switch / interrupt): spinning wastes
less CPU than taking a trap, saving state, and restoring it.
\textbf{Interrupt} when the device is slow (disk, network round-trip):
the CPU should run another process rather than busy-wait on STATUS.

(d) Yes --- two standard hybrids from the canonical
STATUS/COMMAND/DATA protocol:
\begin{itemize}
\item \textbf{Spin-then-interrupt.} If completion time is unknown, poll
      for a short bound; if the device is still BUSY, sleep and take the
      interrupt. Combines low latency of polling with the CPU efficiency
      of interrupts.
\item \textbf{Interrupt livelock / coalescing.} A flood of completions
      (e.g.\ a fast NIC) can livelock the CPU in the handler. The driver
      then \textbf{switches to polling}, or the device uses
      \textbf{interrupt coalescing} (one interrupt per batch of
      completions).
\end{itemize}
PIO versus DMA and explicit \texttt{in}/\texttt{out} versus MMIO are
orthogonal axes of the same design space (status $=$ poll$|$interrupt;
data $=$ PIO$|$DMA; control $=$ port-mapped$|$MMIO).
\end{Sol}

\begin{Qn}{Hard}{Lect 22A}{CS-537 Fa11 Final, Q4}{Twelve random writes on RAID-0/1/4/5, two disk models}
A workload of 12 writes to random locations, spread evenly across the
disks of a \textbf{4-disk} RAID.
\begin{enumerate}
\item Simple model: each read or write takes $D$ time units. Time to
      complete the 12 writes on RAID-0 (striping)?
\item On RAID-1 (mirroring)?
\item On RAID-4 (parity)?
\item On RAID-5 (rotated parity)?
\item Better model: random seek costs $S$, full rotation costs $R$,
      transfer is free. Time for the 12 random writes on RAID-0?
\item On RAID-1?
\item On RAID-4?
\item On RAID-5?
\end{enumerate}
\end{Qn}
\begin{Sol}
Official key (CS-537 Fa11 Final answers, p.~7). Layout sketches:
RAID-0 stripes \texttt{0 1 2 3} / \texttt{4 5 6 7}; RAID-1 two pairs;
RAID-4 dedicated P; RAID-5 rotating P.

\textit{Model $D$.}
\begin{itemize}
\item RAID-0: $12/4 = 3$ writes per disk, all parallel:
      $\boldsymbol{\sim 3D}$.
\item RAID-1: each logical write hits one mirrored pair (2 disks in
      parallel). Two pairs $\Rightarrow$ $12/2 = 6$ sequential writes
      per pair: $\boldsymbol{\sim 6D}$.
\item RAID-4 small-write: 2 reads then 2 writes, and every write hits
      the parity disk. 12 writes $\times$ 2 parity I/Os:
      $\boldsymbol{\sim 12D \times 2 = 24D}$.
\item RAID-5: same 4 I/Os per small write, but parity rotates, so
      effective random-write bandwidth is $N/4 \cdot (1/D)$.
      Time $\boldsymbol{\sim 6D \times 2 = 12D}$.
\end{itemize}

\textit{Model $S,R$.} Average rotation $R/2$. A small write's read
phase costs $S+R/2$; the subsequent write phase, now on-track, waits
about a full rotation $R$ (the sector has just passed). Net per small
write: $S + 3R/2$.
\begin{itemize}
\item RAID-0: $\boldsymbol{\sim 3\bigl(S + \tfrac12 R\bigr)}$.
\item RAID-1: $\boldsymbol{\sim 6\bigl(S + \tfrac12 R\bigr)}$
      (slightly longer in practice --- two seeks/rotates to keep the
      mirrors consistent).
\item RAID-4: $\boldsymbol{\sim 12\bigl(S + \tfrac32 R\bigr)}$.
\item RAID-5: $\boldsymbol{\sim 6\bigl(S + \tfrac32 R\bigr)}$.
\end{itemize}
These match the OSTEP table (writing $R_{\mathrm{bw}}$ for single-disk random bandwidth, distinct from rotation time $R$ above): RAID-0 rand-write $NR_{\mathrm{bw}}$; RAID-1 $(N/2)R_{\mathrm{bw}}$; RAID-4 $R_{\mathrm{bw}}/2$; RAID-5 $(N/4)R_{\mathrm{bw}}$.
\end{Sol}

\begin{Qn}{Hard}{Lect 22A}{CS-537 Sp10 Final, P6}{RAID-4 SMALLWRITE with subtractive parity}
RAID-4, $4$\,KB chunk, 5 disks (4 data $+$ 1 parity):
\begin{lstlisting}[numbers=none]
DISK-0     DISK-1     DISK-2     DISK-3     DISK-4
block0     block1     block2     block3     parity(0..3)
block4     block5     block6     block7     parity(4..7)
\end{lstlisting}
Primitives: \texttt{READ(disk, offset, data)},
\texttt{WRITE(disk, offset, data)}, \texttt{XOR(d1, d2)}
(the last XORs one block into another). Fill in
\texttt{SMALLWRITE(int block, char *data)} for a single-block write.
\end{Qn}
\begin{Sol}
Subtractive (RMW) update:
\[
P_{\mathrm{new}} = P_{\mathrm{old}} \oplus D_{\mathrm{old}} \oplus D_{\mathrm{new}}.
\]
Disk $= \texttt{block} \bmod 4$, offset $= \lfloor\texttt{block}/4\rfloor$,
parity always on disk 4. Assuming \texttt{XOR(dst, src)} computes
\texttt{dst \textasciicircum= src} in place (so we never overwrite the caller's
\texttt{data} buffer!):
\begin{lstlisting}[numbers=none]
void SMALLWRITE(int block, char *data) {
    int disk = block % 4;
    int off  = block / 4;
    char olddata[4096], oldparity[4096];
    READ(disk, off, olddata);     /* parallel */
    READ(4,    off, oldparity);
    XOR(oldparity, olddata);      /* oldparity = P_old XOR D_old */
    XOR(oldparity, data);         /* oldparity = P_old XOR D_old XOR D_new */
    WRITE(disk, off, data);       /* parallel */
    WRITE(4,    off, oldparity);
}
\end{lstlisting}
The two READs may issue together; the two WRITEs may issue together
after the new parity is computed. This is the 2-read / 2-write small-write
path and is why RAID-4 random-write bandwidth collapses to $R/2$
(the parity disk is the bottleneck).
\end{Sol}

\begin{Qn}{Hard}{Lect 22A}{CS-537 Sp18 Final, Q1--10}{Forensic identification of RAID-0/1/4/5 from timing and physical I/Os}
All arrays have \textbf{8 disks} and a $4$\,KB chunk; the logical block
size is $4$\,KB. Requests are chunk-aligned unless stated. Each question
is about a \emph{different} RAID. Options for Q1--5:
(a) RAID-1\quad (b) RAID-4\quad (c) RAID-5\quad
(d) RAID-4 or RAID-5\quad (e) all of RAID-1/4/5.

\begin{enumerate}
\item[1.] RAID A: a block read takes $\approx 10$\,ms and a block write
      also takes $\approx 10$\,ms. Likely?
\item[2.] RAID B: two small random writes usually take about twice as
      long as one. Likely?
\item[3.] RAID C: a large write of 7 blocks usually takes about as long
      as a small write of 1 block. Likely?
\item[4.] RAID D: many 1-block random writes deliver $\approx 4$\,MB/s;
      a comparable RAID-0 delivered $\approx 8$\,MB/s. Likely?
\item[5.] RAID E: large sequential writes $\approx 700$\,MB/s; RAID-0
      sequential writes $\approx 800$\,MB/s. Likely?
\item[6.] RAID F is RAID-4 or RAID-5. One perfectly aligned block write.
      Number of physical I/Os?
      (a) 0 (b) 1 (c) 2 (d) 3 (e) 4
\item[7.] RAID G is RAID-1. One aligned 1-block write. Physical I/Os?
      (a) 0 (b) 1 (c) 2 (d) sometimes 2, sometimes $>2$ (e) always $>2$
\item[8.] RAID H is RAID-1. One small write, \emph{not} necessarily
      aligned. Physical I/Os?
      (same options as Q7)
\item[9.] RAID I is RAID-5. One small 1-block read that is
      \emph{never} aligned. Physical I/Os?
      (a) 0 (b) 1 (c) 2 (d) sometimes 2, sometimes $>2$ (e) always $>2$
\item[10.] RAID J, read-only workload: number of physical I/Os equals
      the number of logical reads, and \textbf{one disk is never
      accessed}. The RAID is:
      (a) RAID-0\quad (b) RAID-1\quad (c) RAID-4\quad
      (d) RAID-0 or RAID-1\quad (e) RAID-0 or RAID-1 or RAID-4
\end{enumerate}
\end{Qn}
\begin{Sol}
Official key (CS-537 Sp18 Final answers.txt):
1\,a, 2\,b, 3\,d, 4\,a, 5\,d, 6\,e, 7\,c, 8\,d, 9\,c, 10\,c.

1. \textbf{(a) RAID-1.} A RAID-4/5 small write is two sequential phases
(RMW), so write latency $\approx 2\times$ read latency. RAID-1 writes
the two copies in parallel, so read and write both cost one positioning
($\approx 10$\,ms).

2. \textbf{(b) RAID-4.} Two small writes serialize on the dedicated
parity disk, doubling the time. RAID-5 (and RAID-1) can often issue
two writes in parallel on distinct disks.

3. \textbf{(d) RAID-4 or RAID-5.} A 7-block write on 8 disks is a
full-stripe write: 7 data $+$ 1 parity issued in parallel, latency
$\approx$ one disk write --- comparable to a small write's dominant
positioning. RAID-1 would write 14 physical copies and not match.

4. \textbf{(a) RAID-1.} Random-write throughput is $N/2\cdot R$ versus
RAID-0's $NR$: half of $8$\,MB/s is $4$\,MB/s. RAID-5 would be
$N/4\cdot R = 2$\,MB/s; RAID-4 would be $R/2 \approx 0.5$\,MB/s.

5. \textbf{(d) RAID-4 or RAID-5.} Sequential write uses $N-1$ disks:
$(7/8)\times 800 = 700$\,MB/s. RAID-1 sequential write is
$(N/2)/N = 1/2$ of RAID-0, i.e.\ $400$\,MB/s, which does not match.

6. \textbf{(e) 4.} Subtractive small write: 2 reads $+$ 2 writes.

7. \textbf{(c) 2.} One aligned RAID-1 write updates both mirrors.

8. \textbf{(d) sometimes $=2$, sometimes $>2$.} If the unaligned 4\,KB
spans two chunks, each chunk is mirrored $\Rightarrow$ 4 physical
writes; if it happens to sit in one chunk, 2 writes.

9. \textbf{(c) 2.} A never-aligned 4\,KB read with 4\,KB chunks always
spans two disks.

10. \textbf{(c) RAID-4.} Only RAID-4 \emph{structurally} never accesses
the dedicated parity disk on reads, and each logical read is one
physical read. RAID-0 stripes across every disk, and a RAID-1 that
load-balances reads across mirrors touches every disk as well.
\end{Sol}

\subsection{Lectures 23--24: File System API \& Implementation}

\begin{Qn}{Hard}{Lect 24}{CS-537 Fa13 Final, Q4}{VSFS access-path I/O counts for read, append, create, unlink}
The very simple file system (VSFS) uses a superblock, bitmaps, inodes, and
data blocks. Each inode is a multi-level index with 12 direct pointers, one
indirect pointer, and one double-indirect pointer. Disk blocks are 4\,KB.
Assume no relevant structures are in the file cache unless a part explicitly
says otherwise, and count both reads and writes as disk accesses.
\begin{enumerate}[(a)]
\item You need to read a 4\,KB (one-block) file, given its inode number.
      How many disk accesses will the read take? Describe them.
\item If this same file is read frequently, how many disk accesses will
      typically occur (with caching)?
\item Append a 4\,KB block to a \emph{small} file that already exists.
      How many disk accesses will the append take? Describe them.
\item If this small file is appended to frequently, how many disk accesses
      will typically occur (with caching)?
\item Append a 4\,KB block to a \emph{large} file that already exists
      (already using its indirect pointer). How many disk accesses will
      the append take? Describe them.
\item Create a file in the root directory (e.g.\ \texttt{/foo.txt}).
      How many disk \emph{writes} take place in the worst case?
\item Delete a file in the root directory (e.g.\ \texttt{/foo.txt}).
      How many disk \emph{writes} take place in the worst case?
\end{enumerate}
\end{Qn}
\begin{Sol}
Nothing is cached, the disk reads and writes whole blocks, and path
walks start at a well-known root inode. Given an inode number, the inode's
on-disk block is
\[
\text{inode-table start} + \bigl\lfloor n / (\text{block size}/\text{inode size})\bigr\rfloor;
\]
no directory walk is required in (a).
\begin{enumerate}[(a)]
\item \textbf{2 accesses:} read the inode block, then read the single
      data block named by a direct pointer. Bitmaps are not consulted on
      a pure read. (If atime is updated, a third access writes the inode
      back; many exam answers omit atime and report 2.)
\item \textbf{0.} After the first read, the inode and data block sit in
      the buffer cache; subsequent reads hit memory.
\item \textbf{5 accesses (2 reads $+$ 3 writes)}, matching the lecture
      write-append trace: read the inode (where does the file end?);
      read the data bitmap; write the data bitmap (allocate); write the
      new data block; write the inode (new pointer, size, mtime).
\item \textbf{0 immediate disk I/Os} under write buffering: the bitmap,
      data, and inode updates remain dirty in memory and may even be
      avoided (create-then-delete of a temp file). They are forced only
      by \texttt{fsync}, a periodic flush, or memory pressure.
\item The file already has an indirect block. Extra I/O: read and write
      that indirect block as well.
      \textbf{7 accesses (3 reads $+$ 4 writes):}
      inode R, data-bitmap R, indirect R;
      data-bitmap W, data W, indirect W, inode W.
      (If the append \emph{allocates} a brand-new indirect or
      double-indirect, add one data-bitmap allocation and one write of
      the new index block.)
\item Common case (root directory already has a free dirent slot):
      write the inode bitmap (allocate), the new inode, the root
      directory data block, and the root inode (updating mtime/ctime;
      nlink is unchanged for a regular file) $\Rightarrow$ 4 writes.
      \textbf{Worst case} (root directory is full and must allocate a
      new data block): also write the data bitmap and the new directory
      data block $\Rightarrow$ \textbf{5 writes} (or 6 if a new indirect
      block must also be allocated for the directory).
\item Unlink of a one-block file (common/worst case): write the
      directory data (remove dirent), the parent inode (mtime), the
      file's own inode (decrement nlink to 0 / mark free), the inode
      bitmap (free inode), and the data bitmap (free data block)
      $\Rightarrow$ \textbf{5 writes} (4 if the file inode itself is not
      written back when its bitmap bit is cleared; 3--4 if the file had
      no data blocks).
\end{enumerate}
These are the lecture's uncached traces: create \texttt{/foo/bar} is
$6\mathrm{R}+4\mathrm{W}$; open is $5\mathrm{R}$; write-append is
$2\mathrm{R}+3\mathrm{W}$; read is $2\mathrm{R}+1\mathrm{W}$ (atime);
\texttt{close} touches \textbf{no} disk.
\end{Sol}

\begin{Qn}{Hard}{Lect 24}{CS-537 Sp18 Final, Q31--38}{Reconstruct a VSFS disk image}
Empty VSFS (8 inodes, 8 data blocks):
\begin{lstlisting}[numbers=none]
inode bitmap  10000000
inodes        [d a:0 r:2] [] [] [] [] [] [] []
data bitmap   10000000
data          [(.,0) (..,0)] [] [] [] [] [] [] []
\end{lstlisting}
Notation \texttt{[d a:0 r:2]}: type directory/file, data-block address,
reference count. Root is inode 0.
\begin{enumerate}[(31)]
\item \texttt{mkdir("/n")} on the empty root. Inodes become
      \texttt{[d a:0 r:3] [d a:1 r:2] []...}, data bitmap
      \texttt{11000000}, data
      \texttt{[(.,0) (..,0) (n,1)] [(.,1) (..,0)] []...}.
      Missing inode bitmap?
      \begin{enumerate}[(a)]
      \item \texttt{10000000}
      \item \texttt{11000000}
      \item \texttt{11100000}
      \item \texttt{10100000}
      \item None of the above
      \end{enumerate}
\item Instead, \texttt{creat("/z")} on the empty file system. Inodes
      \texttt{[d a:0 r:2] [f a:- r:1] []...}, data bitmap
      \texttt{10000000}, root data block corrupt. What belongs there?
      \begin{enumerate}[(a)]
      \item \texttt{[(.,0) (..,0)]}
      \item \texttt{[(.,0) (..,0) (z,1)]}
      \item \texttt{[(.,0) (..,0) (z,2)]}
      \item \texttt{[(.,0) (..,0) (/z,1)]}
      \item \texttt{[(.,0) (..,0) (/z,2)]}
      \end{enumerate}
\item After two operations the image is
\begin{lstlisting}[numbers=none]
inode bitmap 11100000
inodes       [d a:0 r:4] [d a:1 r:2] [d a:2 r:2] [] ...
data bitmap  11100000
data         [(.,0) (..,0) (d,1) (w,2)] [????] [(.,2) (..,0)] [] ...
\end{lstlisting}
      Contents of missing data block 1?
      \begin{enumerate}[(a)]
      \item \texttt{(.,0) (..,0)}
      \item \texttt{(.,1) (..,0)}
      \item \texttt{(.,2) (..,0)}
      \item \texttt{foofoofoo}
      \item None of the above
      \end{enumerate}
\item Which two operations produced that state?
      \begin{enumerate}[(a)]
      \item \texttt{creat("/d"); creat("/w");}
      \item \texttt{creat("/d"); link("/d","/w");}
      \item \texttt{creat("/w"); unlink("/d");}
      \item \texttt{mkdir("/d"); mkdir("/w");}
      \item None of the above
      \end{enumerate}
\item Corrupt image:
\begin{lstlisting}[numbers=none]
inode bitmap 11000000
inodes       [d a:0 r:2] [f a:1 r:2] [] ...
data bitmap  ????????
data         [(.,0) (..,0) (c,1) (m,1)] [foofoofoo] [] ...
\end{lstlisting}
      Which ONE is \emph{not} true?
      \begin{enumerate}[(a)]
      \item There is a proper root directory
      \item The file \texttt{/c} exists
      \item The file \texttt{/m} exists
      \item Reading the first block of \texttt{/c} yields \texttt{foofoofoo}
      \item If you unlink \texttt{/c}, you can no longer read \texttt{foofoofoo}
      \end{enumerate}
\item Which ONE \emph{is} true of that same image?
      \begin{enumerate}[(a)]
      \item File \texttt{/m} is a hard link to \texttt{/c}
      \item File \texttt{/c} is a hard link to \texttt{/m}
      \item Both \texttt{/c} and \texttt{/m} are links to the same file
      \item The root directory has many data blocks
      \item The file \texttt{/m} has many data blocks
      \end{enumerate}
\item Final image:
\begin{lstlisting}[numbers=none]
inode bitmap 11111100
inodes       [d a:0 r:3] [f a:1 r:1] [f a:2 r:1] [d a:3 r:2]
             [f a:- r:1] [f a:4 r:1] [] []
data bitmap  11111000
data         [(.,0) (..,0) (d,1) (g,2) (y,3) (j,5)]
             [d] [v] [(.,3) (..,0) (n,4)] [/u] [] [] []
\end{lstlisting}
      Which regular files exist?
      \begin{enumerate}[(a)]
      \item Only \texttt{/d}, \texttt{/g}, \texttt{/j}
      \item Only \texttt{/d}, \texttt{/g}, \texttt{/j}, \texttt{/n}
      \item Only \texttt{/d}, \texttt{/g}, \texttt{/j}, \texttt{/y/n}
      \item Only \texttt{/d}, \texttt{/g}, \texttt{/j}, \texttt{/y}
      \item None of the above
      \end{enumerate}
\item In this 8-inode VSFS, what is the largest number of
      \emph{regular-file} inodes that can be allocated?
      \begin{enumerate}[(a)]
      \item 1 \quad (b) 6 \quad (c) 7 \quad (d) 8 \quad
      (e) as many as needed
      \end{enumerate}
\end{enumerate}
\end{Qn}
\begin{Sol}
Official key: 31\,b, 32\,b, 33\,b, 34\,d, 35\,e, 36\,c, 37\,c, 38\,c.
\begin{enumerate}[(31)]
\item \texttt{mkdir} allocates inode 1 \emph{and} a directory data block.
      Inode bitmap must show inodes 0 and 1 live: \texttt{11000000}.
      Root nlink $2\to 3$ because of the child's \texttt{..}.
      \textbf{(b)}
\item Empty regular file: no data block (\texttt{a:-}), nlink 1, dirent
      \texttt{(z,1)} using the \emph{component} name, not the full path.
      \textbf{(b)}
\item Inode 1 is a directory (type \texttt{d}) whose \texttt{.} must name
      itself: \texttt{(.,1) (..,0)}. \textbf{(b)}
\item Two new directories, root nlink 4. \textbf{(d)}
\item \texttt{/c} and \texttt{/m} are two names for inode 1 (nlink 2).
      Unlinking \texttt{/c} leaves \texttt{/m}, so the data remain
      readable. \textbf{(e)} is the statement that is \emph{not} true.
\item Hard links are symmetric; (c) is the precise claim. Key: \textbf{(c)}.
\item Inode 3 is directory \texttt{/y} containing file \texttt{n}
      (inode 4, empty). Regular files: \texttt{/d}, \texttt{/g},
      \texttt{/j}, \texttt{/y/n}. \textbf{(c)}
\item One inode is reserved for the root directory, leaving 7 regular-file
      inodes. \textbf{(c)}
\end{enumerate}
\end{Sol}

\begin{Qn}{Hard}{Lect 24}{ISI Mandar End-sem 2015 Q3(a) [cf.\ 2016E Q3(a)i, 2022E Q6(a), 2010E Q2]}{Max size with 10$+$2$+$2$+$1 TOC}
A variant of the SVR2 inode TOC has 15 entries: 10 direct, 2 single-indirect,
2 double-indirect, and 1 triple-indirect. Block size 4096 bytes; block
numbers stored in 32 bits. Calculate the maximum possible file size.
\end{Qn}
\begin{Sol}
Pointers per block $N = 4096/4 = 1024$. Addressable data blocks:
\begin{align*}
  10 + 2N + 2N^{2} + N^{3}
  &= 10 + 2048 + 2\,097\,152 + 1\,073\,741\,824 \\
  &= 1\,075\,841\,034.
\end{align*}
\[
  1\,075\,841\,034 \times 4096 = 4\,406\,644\,875\,264\ \text{bytes}
  \approx 4.008\,\mathrm{TiB}.
\]
If the exam also imposes a 32-bit file offset, the TOC is larger than
$2^{32}-1$ and the size would cap at $4\,\mathrm{GiB}-1$; 2015 does not
state an offset width, so the TOC figure is the intended answer.

\textbf{Cited numeric clones (same skill, standard 10$+$1$+$1$+$1 TOC):}
\begin{itemize}
\item 2016 End Q3(a)i, 512\,B blocks: $N=128$,
      $(10+128+128^{2}+128^{3})\times 512 = 1\,082\,201\,088$ bytes
      $\approx 1.008\,\mathrm{GiB}$ (below the 32-bit offset cap).
\item 2022 End Q6(a), 1024\,B blocks \emph{and} 32-bit file size /
      unsigned offset: TOC yields $(10+256+256^{2}+256^{3})\times 1024 =
      17\,247\,250\,432$ bytes $\approx 16.06\,\mathrm{GiB}$, capped by a
      32-bit file-size field at $2^{32}-1 = 4\,294\,967\,295\ \text{bytes}$
      (or $2^{32}$ bytes if only the byte offset is 32-bit unsigned).
\item 2010 End Q2(b): SVR2 file size $=$ one plus the highest byte
      offset that has ever been written (holes allowed).
\end{itemize}
\end{Sol}

\begin{Qn}{Hard}{Lect 24}{ISI Mandar End-sem 2016 Q3(a)ii [cf.\ 2022E Q6(b), 2010E Q2(c--d)]}{Min and max blocks for a given size (holes)}
SVR2 filesystem, 512-byte blocks, 32-bit block numbers, standard TOC
(10 direct $+$ 1 single $+$ 1 double $+$ 1 triple). Compute the
\textbf{minimum} and \textbf{maximum} number of disk blocks occupied by
a file of size (A)~4000 bytes and (B)~16\,000 bytes.
\end{Qn}
\begin{Sol}
$N=512/4=128$ pointers per index block. Directs cover the first 10
logical blocks (offsets $0..5119$). File size is the highest written
offset plus one, so a file may contain \emph{holes}: logical blocks
never written are not allocated.
\begin{center}
\begin{tabular}{lccc}
\hline
Size & last logical blk & min (hole) & max (dense) \\
\hline
4000  & $\lfloor 3999/512\rfloor=7$  (direct) & 1 data &
  $\lceil 4000/512\rceil=8$ data \\
16000 & $\lfloor 15999/512\rfloor=31$ (needs SI) &
  1 data $+$ 1 SI $=2$ &
  $\lceil 16000/512\rceil=32$ data $+$ 1 SI $=33$ \\
\hline
\end{tabular}
\end{center}
A size-4000 file can live in a single direct block by writing one byte
at offset 3999 (the rest is a hole). Size 16000's last block is past
the 10 directs, so even the hole case must allocate the single-indirect
block.

\textbf{Cited clones:}
\begin{itemize}
\item 2022 End Q6(b), size $2^{20}$, 1024\,B blocks: dense uses 1024
      data blocks of which 10 are direct, 256 via SI, 758 via double
      ($1$ DIB $+\;\lceil 758/256\rceil=3$ IBs) $\Rightarrow$
      \textbf{max 1029}; hole at last block 1023 (double region)
      $\Rightarrow$ \textbf{min 3} (data $+$ DIB $+$ IB).
\item 2010 End Q2(c): size 10\,000 in \emph{one} 2048\,B block --- write
      at offset 9999 (logical block 4); the first four logical blocks
      are holes.
\item 2010 End Q2(d): min blocks for size 25\,000 (2048\,B): last
      logical block $12$ needs the SI $\Rightarrow$ \textbf{2}.
\end{itemize}
\end{Sol}

\begin{Qn}{Hard}{Lect 23}{ISI Mandar End-sem 2015 Q4}{ProZilla: \texttt{fork} shares the GFT offset}
ProZilla downloads a file in \texttt{numParts} pieces by forking children
that each \texttt{fseek}/\texttt{fwrite} into a file opened once in the
parent:
\begin{lstlisting}[numbers=none]
partSize = totalDownloadSize / numParts;
fp = fopen(targetFile, "w");
for (count = 0; count < numParts; count++) {
    if (0 == (p = fork())) {
        socketfd = socket(...);
        getPartialData(socketfd, count, buffer, partSize);
        fseek(fp, count * partSize, SEEK_SET);
        fwrite(buffer, sizeof(char), partSize, fp);
        return;
    } else {
        count++;
        printf("Forked child (pid %d) to take care of fragment %d\n",
               p, count);
    }
}
\end{lstlisting}
\begin{enumerate}[(a)]
\item Draw the relationship between the inode cache, the global file
      table (GFT / open-file table), and the per-process file-descriptor
      tables. Show the relevant GFT fields. Why is this incorrect?
\item Give a simple modification that makes the scheme work.
\end{enumerate}
\end{Qn}
\begin{Sol}
\begin{enumerate}[(a)]
\item \texttt{fopen} creates \emph{one} GFT entry:
      $\{\text{offset},\;\text{inode ptr},\;\text{refcount},\;\text{mode}\}$.
      Each child's fd table (and the parent's) points at that \emph{same}
      GFT entry --- \texttt{fork} duplicates the fd table, not the
      open-file object. All children therefore share one file offset.
      Concurrent \texttt{fseek}/\texttt{fwrite} pairs race: child $A$'s
      seek is overwritten by child $B$'s seek, and pieces clobber each
      other. The in-core inode is also shared (one copy), which is
      correct; the bug is the shared offset.
      (A second, independent bug: the parent does \texttt{count++} in
      the \texttt{else} \emph{and} in the \texttt{for}-header, skipping
      parts.)
\item Any of: have the parent create/truncate the file once and then
      let each child re-open it independently with
      \texttt{fopen(targetFile, "r+")} or
      \texttt{open(targetFile, O\_WRONLY)} (never \texttt{"w"} /
      \texttt{O\_TRUNC}, which would truncate other children's data!),
      giving each child its own GFT entry and private offset; use
      \texttt{pwrite} (offset is an argument, not GFT state); or have
      the parent open $N$ distinct fds \emph{before} forking, one per
      part, and let each child inherit only its own. \texttt{dup} would
      \emph{not} fix it: \texttt{dup} shares the GFT entry on purpose.
\end{enumerate}
This is the lecture picture: per-process fd table $\to$ system-wide
open-file table (offset $+$ inode ptr) $\to$ inode.
\texttt{dup}/\texttt{dup2} and \texttt{fork} share the middle table;
a second \texttt{open} does not.
\end{Sol}

\begin{Qn}{Hard}{Lect 23}{CS-537 Fa13 Final, Q3}{Path walks, hard links, symlinks, and a directory loop}
A regular file is referred to by \texttt{/a/b/c/orig.txt}. Directories
are one block unless stated otherwise.
\begin{enumerate}[(a)]
\item How many directories are accessed when opening this file?
\item Hard link: \texttt{ln /a/b/c/orig.txt /hard.txt}. How many
      directories are accessed when opening \texttt{/hard.txt}?
\item Is \texttt{hard.txt} a special type of file? Explain.
\item What happens to \texttt{hard.txt} when we delete \texttt{orig.txt}?
      Can we still access it?
\item Soft link: \texttt{ln -s /a/b/c/orig.txt /a/b/c/soft.txt}. How
      many directories are accessed when opening \texttt{/a/b/c/soft.txt}?
\item Is \texttt{/a/b/c/soft.txt} a special kind of file? Explain.
\item What happens to \texttt{soft.txt} when we delete \texttt{orig.txt}?
\item \texttt{ln -s /a/b/c /a/b/c/loop}. How many different pathnames
      can refer to \texttt{orig.txt} in \texttt{/a/b/c}?
\end{enumerate}
\end{Qn}
\begin{Sol}
\begin{enumerate}[(a)]
\item Four directories: \texttt{/}, \texttt{a}, \texttt{b}, \texttt{c}
      (then the file's inode).
\item \textbf{One} directory: only the root directory \texttt{/} is
      accessed to read the \texttt{hard.txt} dirent and obtain the
      file's inode number. (The long path \texttt{/a/b/c} is not walked.)
\item No. A hard link is just another \texttt{(name, inumber)} pair for
      the \emph{same} inode. The file type remains regular.
\item \texttt{unlink(orig.txt)} decrements \texttt{nlink}. While
      \texttt{nlink}$\ge 1$ (the \texttt{/hard.txt} name), the inode and
      data stay. \texttt{/hard.txt} still works. Data are freed only
      when \texttt{nlink}$=0$ \emph{and} no fd is open.
\item Opening a symlink \emph{follows} the stored pathname: walk
      \texttt{/a/b/c/soft.txt} (directories \texttt{/}, \texttt{a},
      \texttt{b}, \texttt{c}) then walk the target
      \texttt{/a/b/c/orig.txt} (those four again). If components are
      not cached, \textbf{4 $+$ 4 $= 8$} directory accesses; with a
      cache of the second walk's prefixes, 4 unique directories, but
      the second walk still \emph{looks up} four dirents. Standard
      exam answer: 4 on the first walk plus 4 on the target $=8$.
\item Yes: a symlink is a \emph{new} file (its own inode) whose data
      are a pathname. It can dangle.
\item \texttt{orig.txt}'s inode is freed (if that was the last hard
      link). \texttt{soft.txt} still exists but is \textbf{dangling}:
      \texttt{open} of the symlink fails with \texttt{ENOENT}.
\item Infinitely many: \texttt{/a/b/c/orig.txt},
      \texttt{/a/b/c/loop/orig.txt},
      \texttt{/a/b/c/loop/loop/orig.txt}, \ldots
      (kernels cap symlink depth, e.g.\ 40, so ``infinite'' in theory,
      finite in practice).
\end{enumerate}
\end{Sol}

\clearpage

\section{Lectures 1--7 (Pre-Minor): CPU Virtualization (Processes, Traps \& Scheduling)}

\begin{Qn}{Hard}{Lect 3}{Mandar ISI 2018 Mid Q2}{\texttt{fork} then three \texttt{fprintf}s: how many distinct outputs?}
Explain how many distinct output(s) the following code fragment can generate, if \texttt{buffer} is a character array containing the string \texttt{"abcdef"}. Assume for this question that \texttt{rand()} generates a truly random, uniformly distributed positive integer.
\begin{lstlisting}[numbers=none]
fork();
for (i = 0; i < 3; i++) {
    fprintf(stderr, "%c", buffer[i]);
    sleep(rand() % 5);
}
\end{lstlisting}
\textit{Note: a close variant of this interleaving question appeared as MTL458 Minor Q3.}
\end{Qn}
\begin{Sol}
\texttt{fork()} duplicates the caller, so \emph{two} processes each execute the loop. They do not share \texttt{i} (separate address spaces), and each prints \texttt{buffer[0..2]}, i.e.\ the three characters \texttt{a},~\texttt{b},~\texttt{c} \emph{in that order}. The random \texttt{sleep}s only change the relative speed of the two processes; they cannot reorder a single process's own three prints. Output is therefore a \emph{shuffle} of the two identical strings \texttt{abc} and \texttt{abc}.

A 6-character string arises this way if and only if it contains two \texttt{a}'s, two \texttt{b}'s and two \texttt{c}'s, and every prefix satisfies
\[
\#a \ge \#b \ge \#c
\]
(each process emits \texttt{a} then \texttt{b} then \texttt{c}, so you can never have printed more \texttt{b}'s than \texttt{a}'s, nor more \texttt{c}'s than \texttt{b}'s). Equivalently, these are the standard Young tableaux of shape $(2,2,2)$; by the hook-length formula there are
\[
\frac{6!}{4\cdot 3\cdot 2\cdot 3\cdot 2\cdot 1} = 5
\]
such words. They are:
\begin{center}
\texttt{aabbcc},\quad
\texttt{aabcbc},\quad
\texttt{ababcc},\quad
\texttt{abacbc},\quad
\texttt{abcabc}.
\end{center}
Strings such as \texttt{abcbac} (a \texttt{b} appears before the second \texttt{a}) or \texttt{aabccb} (a \texttt{c} appears before the second \texttt{b}) cannot occur. So there are \textbf{exactly five} distinct outputs.
\end{Sol}

\begin{Qn}{Hard}{Lect 5--6}{Mandar ISI 2024 Mid Q3}{FCFS vs.\ RR($q=0.5$): Gantt, total TAT, utilisation with switch cost $\Delta$}
Consider $n$ processes $P_1,P_2,\ldots,P_n$ that are created in that order at roughly the same time. Each process has a single CPU burst and no I/O bursts. The duration of the CPU burst for process $P_i$ ($i\in\{1,2,\ldots,n\}$) is $n-i+1$ units. Neglect the time taken to switch from one process to another. For each of the following scheduling algorithms --- FCFS, and Round-robin with a time quantum of $0.5$ units ---
\begin{itemize}
  \item[(a)] draw the Gantt chart;
  \item[(b)] compute the total turnaround time.
\end{itemize}
Now assume that the kernel takes $\Delta$ units of time to switch from one process to another. Compute the CPU utilisation for each of the above scheduling algorithms.
\end{Qn}
\begin{Sol}
Bursts: $P_1{:}\,n$, $P_2{:}\,n-1$, \ldots, $P_n{:}\,1$. Total work $W = n(n+1)/2$. Arrivals $\approx 0$, so TAT $=$ completion time.

\textbf{FCFS.} Gantt: $P_1$ for $n$ units, then $P_2$ for $n-1$, \ldots, then $P_n$ for 1.
\[
C_k = \sum_{j=1}^{k}(n-j+1) = \frac{k(2n-k+1)}{2}.
\]
Total TAT:
\[
\sum_{k=1}^{n} C_k = \frac{n(n+1)(2n+1)}{6}.
\]
With switch cost $\Delta$ there are $n-1$ switches (plus, if one counts the initial dispatch, $n$; the paper's ``switch from one process to another'' is $n-1$):
\[
U_{\mathrm{FCFS}} = \frac{W}{W+(n-1)\Delta} = \frac{n(n+1)/2}{n(n+1)/2 + (n-1)\Delta}.
\]

\textbf{RR, $q=0.5$.} Each $P_i$ needs $2(n-i+1)$ slices of length $0.5$, issued in order $P_1,P_2,\ldots,P_n,P_1,\ldots$. Because every burst is a multiple of the quantum, the schedule proceeds in $n$ phases: phase $j$ still has $n-j+1$ alive jobs, each of which receives $1$ unit ($2\times 0.5$), so phase $j$ lasts $n-j+1$ units. The shortest job $P_n$ dies at the end of phase 1; $P_i$ (length $k=n-i+1$) dies at the end of phase $k$:
\[
C_i = \sum_{j=1}^{n-i+1}(n-j+1) = \frac{(n-i+1)(n+i)}{2}.
\]
Check $n=3$ (RR Gantt, each slot $=0.5$):
\begin{center}
\footnotesize
\begin{tabular}{c|cccccccccccc}
$t$ & 0.0 & 0.5 & 1.0 & 1.5 & 2.0 & 2.5 & 3.0 & 3.5 & 4.0 & 4.5 & 5.0 & 5.5--6.0 \\
\hline
run & $P_1$ & $P_2$ & $P_3$ & $P_1$ & $P_2$ & $P_3$ & $P_1$ & $P_2$ & $P_1$ & $P_2$ & $P_1$ & $P_1$ \\
\end{tabular}
\end{center}
$C_3=3$, $C_2=5$, $C_1=6$. Total TAT:
\[
\sum_{i=1}^{n} C_i = \frac{n(n+1)(2n+1)}{6},
\]
the same \emph{sum} as FCFS on this workload (FCFS here is longest-first, which is pessimal; RR merely reverses who pays --- short jobs finish first). The Gantt is not the FCFS Gantt.

CPU utilisation with $\Delta$: RR executes $n(n+1)$ slices in total. In the final phase only $P_1$ remains alive (taking its last two $0.5$-unit slices back-to-back with no switch to another process), so for $n\ge 2$ there are $n(n+1)-2$ actual process-to-process switches (0 for $n=1$):
\[
U_{\mathrm{RR}} = \frac{W}{W + \bigl(n(n+1)-2\bigr)\Delta}.
\]
(If every timer tick at quantum expiry still charges $\Delta$ even when $P_1$ continues, the denominator uses $n(n+1)-1$ switches.) In either case, RR utilisation collapses as $n$ grows because switch cost dominates the $0.5$ quantum --- the Lect~5--6 quantum tradeoff.
\end{Sol}

\begin{Qn}{Hard}{Lect 7}{CS-537 Fall 2004 Final Q11}{Broken MLFQ: gaming vs.\ starvation vs.\ works fine}
The Broken Multi-Level Feedback Scheduler doesn't seem to be working as expected. In our BMLFQ, we have many prioritised queues. Jobs in a higher-priority queue always run before jobs in a lower-priority queue. In each of the following examples, circle whether the scheduler is susceptible to \emph{gaming} by the user, whether \emph{starvation} can occur, or whether the scheduler \emph{works just fine}. (You can circle both gaming and starvation if appropriate.)
\begin{itemize}
  \item[(a)] Starts new jobs at the topmost priority. When the job has used up its quantum, it is moved downward in priority. If the job instead does an I/O before the quantum is up, when it runs again it is restarted at the same priority with a \emph{new} quantum. Jobs within a queue are scheduled in round-robin fashion.
  \item[(b)] Starts new jobs at the topmost priority. When the job has used up its quantum (\emph{regardless of I/Os issued during the quantum}), it moves downward in priority. Jobs within a queue are scheduled in round-robin fashion.
  \item[(c)] Same as (b), but the scheduler periodically moves a job \emph{up one queue}.
  \item[(d)] Same as (b), but the scheduler periodically moves a job \emph{up to the topmost queue}.
  \item[(e)] Same as (c), but jobs within a queue are scheduled in \emph{random} order.
\end{itemize}
\end{Qn}
\begin{Sol}
This is Lect~7's ``problems of basic MLFQ'' plus the two fixes (allotment accounting, boost $S$).
\begin{itemize}
  \item[(a)] \textbf{Gaming} and \textbf{starvation}. Gaming: run just under one quantum, issue a dummy I/O, get a fresh quantum at the same (high) priority --- the exact attack Rule~4's rewrite is designed to stop. Starvation: no boost, so a stream of interactive/high jobs can hold the CPU forever.
  \item[(b)] \textbf{Starvation} only. Allotment is charged whether or not the job did I/O, so the gaming trick fails. Still no boost, so long jobs on the bottom queue never run if high queues stay busy.
  \item[(c)] \textbf{Works just fine} (no gaming, no starvation). Because promotion happens periodically on a timer (not only when the CPU is idle), a job sitting in the bottom queue is boosted up one level every period; after at most $L-1$ boost periods ($L$ queues), it reaches the topmost queue and receives a round-robin turn (since it can only be demoted after running). It recovers from starvation more slowly than (d), but cannot starve.
  \item[(d)] \textbf{Works just fine.} This \emph{is} the final MLFQ: Rule~3 (enter at top), Rule~4 (allotment, no gaming), Rule~5 (boost \emph{to the top}, anti-starvation and behaviour change), Rules~1--2 (priority + RR).
  \item[(e)] \textbf{Starvation} possible (in the sense of no deterministic waiting-time bound; not gameable). Even after a job is boosted to the topmost queue, \emph{random} selection within a queue offers no finite deterministic bound (wait is unbounded, though the probability of never being picked tends to $0$).
\end{itemize}
\end{Sol}

\clearpage

\section{Lectures 8--12 (Pre-Minor): Early Memory Virtualization (B\&B, Seg., Free Space \& Paging)}

\begin{Qn}{Hard}{Lect 10}{CS-537 Fall 2009 Mid, Q2}{LostOS: reverse-engineer a two-segment (backward stack) translation}
An old version of LostOS used segmentation (not paging). The hardware
supported two segments: segment 0 for code and a growing heap, and
segment 1 for a negatively-growing stack. The virtual address space was
only 128 bytes, and physical memory was 1\,KB. The top bit of the
virtual address selects the segment.

The base of a grow-negative stack points to the physical address
\emph{one past} the last byte of that stack. The bounds (limit) register
holds the \emph{size} of the segment.

Programs were allowed to keep running after a segmentation violation, so
the trace below continues after faults. Recover the four register
values. If a value cannot be recovered exactly, give the tightest range
consistent with the trace.

\begin{lstlisting}[numbers=none]
VA 0x6c (dec 108) --> VALID in SEG1: PA 0x3ec (dec 1004)
VA 0x1d (dec  29) --> VALID in SEG0: PA 0x21d (dec  541)
VA 0x50 (dec  80) --> SEGMENTATION VIOLATION (SEG1)
VA 0x1e (dec  30) --> SEGMENTATION VIOLATION (SEG0)
VA 0x58 (dec  88) --> VALID in SEG1: PA 0x3d8 (dec  984)
VA 0x61 (dec  97) --> VALID in SEG1: PA 0x3e1 (dec  993)
VA 0x35 (dec  53) --> SEGMENTATION VIOLATION (SEG0)
VA 0x21 (dec  33) --> SEGMENTATION VIOLATION (SEG0)
VA 0x64 (dec 100) --> VALID in SEG1: PA 0x3e4 (dec  996)
VA 0x3d (dec  61) --> SEGMENTATION VIOLATION (SEG0)
VA 0x0c (dec  12) --> VALID in SEG0: PA 0x20c (dec  524)
VA 0x05 (dec   5) --> VALID in SEG0: PA 0x205 (dec  517)
VA 0x2f (dec  47) --> SEGMENTATION VIOLATION (SEG0)
\end{lstlisting}
\begin{enumerate}
  \item[(a)] Segment 0 base?
  \item[(b)] Segment 0 bounds (size)?
  \item[(c)] Segment 1 base?
  \item[(d)] Segment 1 bounds (size)?
  \item[(e)] LostOS later dropped segmentation for paging. Give one
        positive of the switch, and one reason one might still prefer
        segmentation.
\end{enumerate}
\end{Qn}
\begin{Sol}
VAS size $128$ means VAs $0..63$ are segment 0 (top bit $0$) and
VAs $64..127$ are segment 1 (top bit $1$). Translation for a
grow-positive segment is $\mathrm{PA}=\mathrm{base}+\mathrm{VA}$. For the
grow-negative stack, VA $127$ maps to $\mathrm{base}-1$, VA $126$ to
$\mathrm{base}-2$, and in general
$\mathrm{PA}=\mathrm{base}-(128-\mathrm{VA})$.

(a) A valid SEG0 pair: VA $5\to$ PA $517$. Then
$\mathrm{base}_0=517-5=512$ (\texttt{0x200}). Check: VA $12\to 524$
and VA $29\to 541$ both equal $512+\mathrm{VA}$.

(b) VA $29$ is valid and VA $30$ faults, so the legal SEG0 range is
exactly $0..29$. The size-encoded bounds register is therefore
\textbf{30}. (The corresponding physical limit would be $512+30=542$.)

(c) VA $108\to$ PA $1004$. The last virtual byte is $127$, which is
$19$ bytes ``up'' the stack from $108$, so it maps to $1004+19=1023$.
The base, one past that last byte, is \textbf{1024}. Check:
VA $88\to 1024-(128-88)=984$; VA $97\to 993$; VA $100\to 996$.

(d) VA $88$ is valid, so the stack must cover at least the $40$ bytes
$88..127$. VA $80$ faults, so the stack does not extend down to $80$;
the lowest legal VA could be $81$, giving maximum size $128-81=47$.
Hence the SEG1 size is in \textbf{$[40,47]$} (inclusive). The
trace does not pin it to a single integer.

(e) Paging eliminates \textbf{external fragmentation} (every frame is
the same size, so the OS free list is trivial). One might still prefer
segmentation because a linear page table is far larger than a handful
of segment registers, and a translation now costs an extra memory
reference to the PTE (Lect 12).
\end{Sol}

\begin{Qn}{Hard}{Lect 11}{CS-537 Spring 2011 Mid, Problem 3}{Allocator headers, splitting, coalescing, and best-fit}
A $50$-byte managed heap is tracked with a simple free list. A
\textbf{$4$-byte header is required before every allocated chunk}.
\texttt{H} is a header byte, \texttt{A} an allocated payload byte,
\texttt{F} a free byte. After one \texttt{allocate(10)} ($4+10=14$ bytes used, $36$ free) the heap looks like:
\begin{lstlisting}[numbers=none]
HHHHAAAAAAAAAA FFFFFFFFFFFFFFFFFFFFFFFFFFFFFFFFFFFF
\end{lstlisting}
The policy is \textbf{lowest-address-first} (a first-fit that always
restarts at the low end). Returned chunks are \textbf{coalesced}.
\begin{enumerate}
  \item[(a)] Start from a $50$-byte free heap. Draw the heap after
        \texttt{allocate(10)}, \texttt{allocate(10)},
        \texttt{allocate(10)}.
  \item[(b)] Start from a $50$-byte free heap. Draw it after
        \texttt{allocate(10)}, \texttt{allocate(20)},
        \texttt{allocate(20)}.
  \item[(c)] Start from a $50$-byte free heap. Sequence:
        \texttt{x=allocate(10)}, \texttt{y=allocate(10)},
        \texttt{z=allocate(10)}, \texttt{free(y)},
        \texttt{w=allocate(24)}.
  \item[(d)] Now there is \textbf{no coalescing}, and the policy is
        \textbf{best-fit}. Start from a $99$-byte free heap. Sequence:
        \texttt{x=allocate(10)}, \texttt{y=allocate(10)},
        \texttt{z=allocate(10)}, \texttt{free(y)}, \texttt{free(z)},
        \texttt{w=allocate(4)}.
\end{enumerate}
\end{Qn}
\begin{Sol}
A request of $n$ payload bytes occupies $n+4$ bytes on the heap.
\texttt{free(p)} finds the header just before the user pointer (that
is why \texttt{free} takes no size).

(a) Three $14$-byte allocated blocks ($42$ bytes), leftover $8$ free:
\begin{lstlisting}[numbers=none]
HHHHAAAAAAAAAA HHHHAAAAAAAAAA HHHHAAAAAAAAAA FFFFFFFF
\end{lstlisting}

(b) \texttt{allocate(10)} takes $14$; \texttt{allocate(20)} takes
$24$; leftover $12$. The third request needs $24$ bytes and only $12$
remain, so it \textbf{fails} (or the heap must grow via
\texttt{sbrk}/\texttt{mmap}, which this problem does not provide):
\begin{lstlisting}[numbers=none]
HHHHAAAAAAAAAA HHHHAAAAAAAAAAAAAAAAAAAA FFFFFFFFFFFF
\end{lstlisting}

(c) After three \texttt{allocate(10)}s the picture is (a).
\texttt{free(y)} returns the middle $14$ bytes. The leftover $8$ at
the high end is \emph{not} adjacent to that hole (z sits in between),
so coalescing does not help. \texttt{allocate(24)} needs $28$ bytes;
neither hole is large enough. \textbf{Fails.} This is external
fragmentation: $14+8=22$ free bytes exist, but not contiguously, and
the header tax makes the request even larger.

(d) No coalescing, best-fit, $99$ bytes.
Three \texttt{allocate(10)}s consume $42$ bytes and leave a $57$-byte
tail:
\texttt{[H10][H10][H10][57F]}.
\texttt{free(y)} then \texttt{free(z)} produce two \emph{separate}
$14$-byte holes (z is not merged with the tail):
\texttt{[H10][14F][14F][57F]}.
\texttt{allocate(4)} needs $8$ bytes. Best-fit picks a smallest
adequate chunk: either $14$-byte hole, remainder $6$:
\texttt{[H10][H4][6F][14F][57F]}.
\end{Sol}

\begin{Qn}{Hard}{Lect 12}{CS-537 Spring 2018 Final, Q71--Q76 (Q74 and Q76 later appeared as MTL458 Minor Q4--Q5)}{Linear page-table forensics: tiny translations and PT-size arithmetic}
A system has $64$ bytes of physical memory, $4$-byte pages, and
$16$-byte virtual address spaces. Each PTE stores a valid bit in the
high bit; the rest is the PFN.

\textbf{Q71.} A trace accesses virtual address \texttt{0x7}, which
translates to \texttt{0x33}. The page table is
\begin{lstlisting}[numbers=none]
[0] 0x00000000
[1] 0x800000??   <- missing two hex digits
[2] 0x00000000
[3] 0x00000000
\end{lstlisting}
What two hex digits are missing from PTE 1?
(a) \texttt{0x0a}\quad (b) \texttt{0x0b}\quad (c) \texttt{0x0c}\quad
(d) \texttt{0x0d}\quad (e) None of the above

\textbf{Q72.} Same page size / VAS, now the table is
\begin{lstlisting}[numbers=none]
[0] 0x00000000
[1] 0x80000005
[2] 0x8000000a
[3] 0x00000000
\end{lstlisting}
Virtual address $4$ translates to which \emph{decimal} physical
address?
(a) $5$\quad (b) $10$\quad (c) $20$\quad (d) $54$\quad (e) $45$

\textbf{Q73.} Table
\begin{lstlisting}[numbers=none]
[0] 0x00000000
[1] 0x8000000e
[2] 0x00000000
[3] 0x80000009
\end{lstlisting}
Which ranges of virtual address are valid?
(a) $0\ldots 64$\quad (b) $0\ldots 16$\quad
(c) $4\ldots 7$ and $12\ldots 15$\quad
(d) $0\ldots 3$ and $8\ldots 11$\quad
(e) $80\ldots 89$

\textbf{Q74.} A system has a $30$-bit virtual address and a $4$\,KB
page size. PTE size is $4$ bytes. How big is a linear page table for
one process?
(a) $1$\,MB\quad (b) $2$\,MB\quad (c) $4$\,MB\quad (d) $8$\,MB\quad
(e) None of the above

\textbf{Q75.} A system has a $10$-bit virtual address and a $256$-byte
page size. PTE size is $4$ bytes. How big is a linear page table?
(a) $16$ bytes\quad (b) $16$\,KB\quad (c) $16$\,MB\quad
(d) $16$\,GB\quad (e) None of the above

\textbf{Q76.} A system has a $1$\,MB linear page table per process and
a $1$\,KB page size. PTE size is $4$ bytes. How many bits are in the
VPN?
(a) $28$\quad (b) $18$\quad (c) $8$\quad (d) $32$\quad
(e) None of the above
\end{Qn}
\begin{Sol}
$16$-byte VAS, $4$-byte pages: $2$-bit VPN (4 PTEs) and $2$-bit
offset. $\mathrm{PA}=\mathrm{PFN}\times 4+\mathrm{offset}$.

\textbf{Q71. (c) \texttt{0x0c}.}
\texttt{0x7}$=0111_2$: VPN $01$ (PTE 1), offset $11_2=3$.
$\mathrm{PA}=0\mathrm{x}33=51$, so
$\mathrm{PFN}=(51-3)/4=12=0\mathrm{x}0\mathrm{c}$.
PTE 1 is \texttt{0x8000000c}.

\textbf{Q72. (c) $20$.}
VA $4=0100_2$: VPN $1$, offset $0$. PTE 1 is valid with PFN $5$.
$\mathrm{PA}=5\times 4+0=20$.

\textbf{Q73. (c).}
Valid PTEs are VPN $1$ and VPN $3$, i.e.\ bytes $4..7$ and
$12..15$.

\textbf{Q74. (a) $1$\,MB.}
$4$\,KB page $\Rightarrow$ $12$-bit offset, so a $30$-bit VA has an
$18$-bit VPN. Linear PT size
$=2^{18}\times 4=2^{20}=1$\,MB.

\textbf{Q75. (a) $16$ bytes.}
$256=2^8$ $\Rightarrow$ $8$-bit offset, $2$-bit VPN, $4$ PTEs,
$4\times 4=16$ bytes.

\textbf{Q76. (b) $18$.}
Number of PTEs $=1\,\mathrm{MB}/4=2^{18}$, so the VPN is $18$ bits.
(Page size $1$\,KB is not needed once the table size and PTE size are
known.)
\end{Sol}

\clearpage

% ===========================================================================
\part{Extended High-Priority \& Challenging Problems (Cover Next)}
% ===========================================================================

\section{Lectures 13--15: Virtual Memory --- TLB, Advanced Paging \& Swapping}

\subsection{Lectures 13--14: Paging with TLB \& Advanced Paging}

\begin{Qn}{Medium-Hard}{Lect 13}{CS-537 Fa13 Mid, Q10}{Array-walk TLB misses, including instruction fetches}
McFlub writes:
\begin{lstlisting}[numbers=none]
int a[4096];
int b[4096];
int c = 0;
for (i = 0; i < 4096; i++) {
    c = a[i] + b[i];
}
\end{lstlisting}
and claims: ``Without TLBs, on a system with 4-byte integers and a 4\,KB
page size, and assuming a linear page table, this code would access
memory an additional 8192 times (due to page-table accesses) --- this is
a worst-case analysis. With TLBs, the page tables are only accessed 8
times in the worst case.'' Then: ``Oh, hold on: I think I might have
forgotten to think about instruction references\ldots''

Identify what is right, what is wrong, and what instruction fetches do
to the accounting. State your layout assumptions.
\end{Qn}
\begin{Sol}
Each array is $4096\times 4 = 16\,\mathrm{KB} = 4$ pages, so $a$ and $b$
together occupy 8 pages if both are page-aligned, up to 9 if contiguous
but unaligned, and up to 10 if separate and unaligned (plus possibly 1
more page for the scalar \texttt{c}). The loop body does two data loads per iteration ($8192$ data loads), plus
one store to \texttt{c} per iteration ($4096$ stores) if \texttt{c} lives in
memory rather than a register. With a linear page table and no TLB, each
data load walks one PTE, so McFlub's ``additional 8192 page-table accesses''
is only the right worst-case count for \emph{loads} (a memory-resident
\texttt{c} adds 4096 more PTE walks, for 12{,}288 total data PTE walks).

With a TLB large enough to hold the data pages (8--10 entries) plus
the pages of the loop code, each data page contributes one compulsory
miss and then hits: \textbf{8--10 page-table walks for data} (plus up
to 1 for \texttt{c}), not
8192. McFlub's ``8'' is therefore the aligned, data-only, best-case
TLB-miss count --- it is not a worst case, and it ignores code.

Instruction fetches: a typical unoptimized loop is about 7--8 instructions
per iteration (load $a[i]$, load $b[i]$, add, store $c$, increment,
compare, branch), on the order of $7\text{--}8\times 4096 \approx 28\text{--}32\,\mathrm{K}$
fetches. Without a TLB each fetch walks the page table as well. The
code itself usually occupies one or two pages, so a TLB that also
caches those translations reduces instruction PTE walks to 1--2. A
clever compiler would further rewrite the loop as
\texttt{c += a[i] + b[i]} and keep $c$ in a register; McFlub's store to
\texttt{c} each iteration is itself a stream of data writes (possibly
yet another page).

Lecture connection: this is the array-walk miss-rate example
($2/2048 = 0.1\%$ for a 2048-int scan on 4\,KB pages) with the extra
requirement that instruction-side translations be counted.
\end{Sol}

\begin{Qn}{Medium-Hard}{Lect 13}{CS-537 Fa08 Mid, P3}{TLB hit/miss with a PID (ASID) tag}
32-bit virtual addresses, 20-bit VPN, 12-bit offset. Each TLB entry stores
a 20-bit VPN, a 20-bit PPN, and an 8-bit PID. The four entries are
(all hex):
\begin{center}
\begin{tabular}{ccc}
VPN    & PPN    & PID \\
\hline
00000  & 00FFF  & 00 \\
00000  & 00AAB  & 01 \\
00010  & F000A  & 00 \\
010FF  & 00ABC  & 01 \\
\end{tabular}
\end{center}
For each access, state hit or miss; on a hit give the physical address
in hex. The first five hex digits of a VA are the VPN; the last three are
the offset.
\begin{enumerate}[(a)]
\item PID 00, VA \texttt{00000000}
\item PID 01, VA \texttt{00000000}
\item PID 00, VA \texttt{FF00FFAA}
\item PID 00, VA \texttt{0010FFAA}
\item PID 01, VA \texttt{0010FFAA}
\item PID 00, VA \texttt{000000FF}
\item PID 01, VA \texttt{00000FAB}
\item PID 00, VA \texttt{010FFFFF}
\item PID 01, binary VA \texttt{00000001000011111111010100001111}
\item PID 00, same binary VA as (i)
\item PID 02, VA \texttt{00000000}
\end{enumerate}
\end{Qn}
\begin{Sol}
A hit requires \textbf{both} VPN match \textbf{and} PID match. The PA is
$\mathrm{PPN}\Vert\mathrm{offset}$.
\begin{enumerate}[(a)]
\item Hit. VPN \texttt{00000}, PID 00 $\to$ \texttt{00FFF000}.
\item Hit. VPN \texttt{00000}, PID 01 $\to$ \texttt{00AAB000}.
\item Miss. VPN \texttt{FF00F} is not in the TLB.
\item Miss. VPN \texttt{0010F} is not in the TLB.
\item Miss. Same VPN, still absent.
\item Hit. VPN \texttt{00000}, PID 00 $\to$ \texttt{00FFF0FF}.
\item Hit. VPN \texttt{00000}, PID 01 $\to$ \texttt{00AABFAB}.
\item Miss. VPN \texttt{010FF} is present, but only for PID 01.
\item Binary groups as \texttt{0000\,0001\,0000\,1111\,1111\,0101\,0000\,1111}
      $= \texttt{010FF50F}$. Hit for PID 01: $\texttt{00ABC50F}$.
\item Same VA, PID 00: miss (the \texttt{010FF} row is tagged 01).
\item Miss. No PID-02 row exists. (Two processes may share a VPN; the
      PID/ASID tag is what keeps their translations from aliasing.)
\end{enumerate}
Without a PID/ASID field the OS would have to \textbf{flush} the TLB on
every context switch; with it, stale entries from PID 00 can remain
while PID 01 runs.
\end{Sol}

\begin{Qn}{Hard}{Lect 14}{CS-537 Sp09 Mid, Q3}{Two-level footprint vs.\ a 16\,MB linear table}
32-bit virtual addresses, 1\,KB pages, 4-byte PTEs.
\begin{enumerate}[(a)]
\item Offset bits? VPN bits?
\item Linear table: number of entries? total size? memory occupied in a
      live system?
\item Two-level table: PTEs per page? bits in $\mathrm{VPN}_{\mathrm{chunk}}$?
      bits in $\mathrm{VPN}_{\mathrm{dir}}$? size of the page directory
      (assume 4-byte PDEs)?
\item For each workload, give multi-level size \emph{and} linear size:
      \begin{enumerate}[(i)]
      \item Code at address 0 (100 4-byte instructions); heap starts at
            page 1 and uses 3 pages; stack at the high end uses 3 pages.
      \item Same code; heap 1000 pages; stack 1000 pages.
      \item Every virtual page is valid.
      \end{enumerate}
\end{enumerate}
\end{Qn}
\begin{Sol}
\begin{enumerate}[(a)]
\item Offset: $10$ bits ($1\,\mathrm{KB}=2^{10}$). VPN: $32-10=22$ bits.
\item $2^{22}$ entries; $2^{22}\times 4\,\mathrm{B} = 16\,\mathrm{MB}$
      per process, and it must be contiguous. Live system: $16\,\mathrm{MB}
      \times(\#\textrm{processes})$ if every process has a full linear
      table (100 processes $\Rightarrow 1.6\,\mathrm{GB}$ of page tables
      alone).
\item PTEs per 1\,KB page: $1024/4 = 256 = 2^8$, so
      $\mathrm{VPN}_{\mathrm{chunk}}$ is 8 bits and
      $\mathrm{VPN}_{\mathrm{dir}}$ is $22-8=14$ bits. Page directory:
      $2^{14}$ PDEs $\times 4\,\mathrm{B} = 64\,\mathrm{KB}$ (also
      contiguous --- this is why two levels are not always enough).
\item Linear size is always 16\,MB.
      \begin{enumerate}[(i)]
      \item Code is 400 bytes $\subset$ page 0; heap is pages 1--3.
            Those four VPNs sit in the \emph{same} 256-page chunk, so
            one leaf page-table page. The stack at the other end of the
            22-bit VPN space needs a second leaf. Plus the 64\,KB
            directory: $64\,\mathrm{KB}+2\,\mathrm{KB}=66\,\mathrm{KB}$.
      \item First $1+1000=1001$ pages need
            $\lceil 1001/256\rceil = 4$ leaf pages; last 1000 pages
            need another 4. Total $64\,\mathrm{KB}+8\,\mathrm{KB}=72\,\mathrm{KB}$.
      \item All $2^{22}$ PTEs are valid $\Rightarrow$ all
            $2^{14}$ leaves exist ($16\,\mathrm{MB}$) plus the 64\,KB
            directory $= 16\,\mathrm{MB}+64\,\mathrm{KB}$. Multi-level
            is then \emph{larger} than linear.
      \end{enumerate}
\end{enumerate}
PDE valid bit: ``this inner page-table page exists''. PTE valid bit:
``this virtual page is mapped''. An invalid PDE means you must not even
look for an inner table (and you just saved 256 PTEs of storage).
\end{Sol}

\begin{Qn}{Medium-Hard}{Lect 14}{ISI Mandar 2010 End-sem, Q1(b)}{64-bit linear table, and how many paging levels}
A 64-bit machine has a page/frame size of $2^{16}$ bytes. Each PTE is
8 bytes long.
\begin{enumerate}[(i)]
\item What is the size of the complete first-level (primary) page table
      for a process?
\item If page tables are themselves stored using paging, how many levels
      of paging are required? (If $n$ levels are required, the level-$n$
      table must fit in a single page.)
\end{enumerate}
\end{Qn}
\begin{Sol}
VPN bits $= 64-16 = 48$.
\begin{enumerate}[(i)]
\item Linear/primary table:
      $2^{48}\times 8\,\mathrm{B} = 2^{51}\,\mathrm{B} = 2\,\mathrm{PB}$
      (contiguous!). This is the lecture's ``64-bit / 4\,KB / 4\,B
      $= 16\,\mathrm{PB}$'' disaster with a larger page and an 8-byte
      PTE --- still impossible as a contiguous array.
\item Each page holds $2^{16}/8 = 2^{13} = 8192$ PTEs, so each level
      consumes 13 VPN bits. We must cover 48 VPN bits, and the outermost
      table must fit in one page:
      \[
      \left\lceil\frac{48}{13}\right\rceil = 4
      \]
      levels. Explicitly: Level~1 (leaf table) has $2^{48}\times 8 = 2^{51}\,\mathrm{B} = 2^{35}$ pages; Level~2 has $2^{35}\times 8 = 2^{38}\,\mathrm{B} = 2^{22}$ pages; Level~3 has $2^{22}\times 8 = 2^{25}\,\mathrm{B} = 2^{9}$ pages; Level~4 (outermost directory) has $2^{9}\times 8 = 4\,\mathrm{KB} \le 64\,\mathrm{KB}$, which fits inside a single page. The 48-bit VPN is thus split $9+13+13+13$, requiring \textbf{4 levels}.
\end{enumerate}
A 4-level x86-64 table with 4\,KB pages and 8-byte PTEs similarly uses
9 bits per level ($512$ entries) and a 48-bit canonical VA
($4\times 9 + 12 = 48$).
\end{Sol}

\begin{Qn}{Medium-Hard}{Lect 14}{CS-537 Fa21 Mid, Q18--20}{PTEs per page; min valid pages for a max-size two-level table}
Minux uses a two-level page table. 32-bit VAs, 4\,KB pages, 4-byte PTEs.
\begin{enumerate}[(a)]
\item How many PTEs fit on one page?
\begin{enumerate}[(A)]
\item 1 \quad (B) 32 \quad (C) 1024 \quad (D) 4096 \quad (E) none of these
\end{enumerate}
\item What is the \textbf{minimum} number of valid virtual pages such
      that the multi-level page table attains its \textbf{maximum} size?
\begin{enumerate}[(A)]
\item 512 \quad (B) 1024 \quad (C) 2048 \quad (D) 4096 \quad (E) none
\end{enumerate}
\item What is the minimum number of pages needed for the multi-level
      table (including the page directory) when there are \textbf{1025
      contiguous} valid pages somewhere in the address space?
\begin{enumerate}[(A)]
\item 3 \quad (B) 4 \quad (C) 5 \quad (D) 1025 \quad (E) none
\end{enumerate}
\end{enumerate}
\end{Qn}
\begin{Sol}
VPN bits $=20$. Leaf page holds $4096/4=1024=2^{10}$ PTEs, so the split
is $10+10+12$. The directory has 1024 PDEs and occupies one page. Full
linear table is $2^{20}$ PTEs $= 4\,\mathrm{MB} = 1024$ leaf pages, plus
1 directory page.
\begin{enumerate}[(a)]
\item \textbf{(C) 1024.}
\item \textbf{(B) 1024.} To force every leaf page-table page to exist,
      each of the 1024 leaves needs at least one valid PTE. Place one
      valid virtual page in each of the 1024 directory slots (e.g.\ VPNs
      $0, 1024, 2048, \ldots$). That is the minimum number of valid
      pages that still allocates the entire $1024+1$ page skeleton.
      (1023 valid pages can always avoid at least one leaf.)
\item \textbf{(A) 3.} A leaf covers 1024 consecutive VPNs, so 1025
      contiguous valid pages always span \textbf{two} leaves (at least
      two because $1025>1024$, and at most two because even starting at
      the last slot of a leaf covers $1+1024=1025$ VPNs across two
      leaves). Plus the one directory page: 3 pages total.
\end{enumerate}
\end{Sol}

\subsection{Lecture 15: Swapping \& Page Replacement Policies}

\begin{Qn}{Hard}{Lect 15}{CS-537 Sp11 Final, P1}{Black-box VM: infer TLB size and the OS replacement policy}
A reference is observed as a TLB hit, a TLB miss (page present), or a
page fault. Memory starts empty.
\begin{enumerate}
  \item[(a)] Trace:
  \begin{center}
  \begin{tabular}{ll}
    load VPN 0,1,2,3 & page fault each \\
    load VPN 0,1,2,3 & TLB miss each
  \end{tabular}
  \end{center}
  Assuming an LRU TLB, what is the TLB size?
  \item[(b)] On a different machine the second round is four TLB hits.
        What can we conclude about the TLB size?
  \item[(c)] On a third machine the TLB has 2 entries and physical memory
        holds 3 pages. Trace:
  \begin{center}
  \begin{tabular}{ll}
    load VPN 0,1,2,3 & page fault each \\
    load VPN 0        & TLB hit \\
    load VPN 1        & TLB hit \\
    load VPN 2        & page fault \\
    load VPN 3        & TLB miss
  \end{tabular}
  \end{center}
  What can we conclude about the OS page-replacement policy?
\end{enumerate}
\end{Qn}
\begin{Sol}
(a) The four pages remain in memory (second round is not a fault), so
RAM holds at least 4 pages. Under an LRU TLB with $k \ge 4$ entries, all
four pages from the first pass remain in the TLB and the second round
would be four hits. Under \emph{any} LRU TLB with $k \in \{1,2,3\}$
entries, a cyclic scan over 4 distinct pages thrashes on every access
(e.g., for $k=3$, after $0,1,2,3$ the TLB holds $\{1,2,3\}$, so $0$
misses and evicts $1$, then $1$ misses and evicts $2$, then $2$ misses
and evicts $3$, then $3$ misses and evicts $0$). Therefore we conclude
that the LRU TLB has \textbf{at most 3 entries} ($1$, $2$, or $3$).

(b) All four translations survived, so the TLB has \textbf{at least 4}
entries.

(c) After four compulsory faults a 3-frame cache must have evicted one of
$\{0,1,2\}$. FIFO and LRU both evict 0, so the second access to 0 would
be a \textbf{page fault}. Observing a TLB hit on 0 (and then on 1) while
2 faults means the victim was \textbf{page 2}: memory holds
$\{0,1,3\}$, and the later TLB miss on 3 (not a fault) confirms 3 is
still resident. Evicting 2 is exactly OPT on this continuation (among
$\{0,1,2\}$, page 2 is used furthest) and is also what MRU does when
installing 3. Therefore the OS is \textbf{not} using FIFO or LRU; OPT,
MRU, or Random is consistent. (Note: any demand-filled 2-entry TLB must
hold 3 immediately after accessing 3, so two subsequent TLB hits on 0
and 1 are only possible if the fault handler does not install 3 into the
TLB, e.g.\ under software-managed TLB filtering; the inference about RAM
replacement is unaffected.)
\end{Sol}

\begin{Qn}{Hard}{Lect 15}{CS-537 Sp09 Mid, Q4}{TLB hits, misses, and page faults under LRU vs MRU, and the resulting EAT}
32-bit VAs, 4\,KB pages, a 32-entry LRU TLB, LRU page replacement, and
only 1024 physical pages. Ignore instruction fetches and stores to
\texttt{x}. Initialization (run once) zeros the first word of each of
\texttt{NUM\_PAGES} pages; the \textbf{test} then loads that word from
each page in order:
\begin{lstlisting}[numbers=none]
ptr = orig;
for (i = 0; i < NUM_PAGES; i++) {
    int x = *(int *)ptr;
    ptr += PAGE_SIZE;
}
\end{lstlisting}
A memory reference costs $M$, a disk access costs $D$. Count TLB hits,
TLB misses, and page faults during the \textbf{test}, then the running
time, for \texttt{NUM\_PAGES} $\in \{16,32,2048\}$. Repeat the running
time under \textbf{MRU} replacement (TLB and RAM). Finally, explain how
you would use this loop as a microbenchmark to discover the TLB size
and the RAM size of an unknown machine.
\end{Qn}
\begin{Sol}
Initialization already walked the same pages, so the test sees a warm
TLB/RAM except when the working set exceeds the cache.

\textbf{LRU.}
\begin{center}
\begin{tabular}{rccc}
\hline
\texttt{NUM\_PAGES} & TLB hits & TLB misses & Page faults \\
\hline
16   & 16   & 0    & 0 \\
32   & 32   & 0    & 0 \\
2048 & 0    & 2048 & 2048 \\
\hline
\end{tabular}
\end{center}
16 and 32 both fit in the 32-entry TLB (and in 1024 frames), so the
test is all hits, time $16M$ and $32M$. With 2048 pages LRU keeps the
1024 most recently initialized pages $\{1024,\ldots,2047\}$ and a TLB of
$\{2016,\ldots,2047\}$. The test then scans $0,1,\ldots,2047$: every
access is a new page, so every access is a TLB miss \emph{and} a page
fault. Each TLB miss costs two memory references (PTE then data) once
the page is in RAM, and each fault costs a disk I/O (ignoring both the
extra PTE re-walk on instruction restart and write-back of pages dirtied
during initialization --- otherwise LRU pays $+1024D$ and MRU pays $+1D$
in write-backs):
$2048\times 2M + 2048D$.

\textbf{MRU.} For 16 and 32 nothing changes ($16M$, $32M$). For 2048,
MRU \emph{protects the oldest} entries: after init, RAM holds
$\{0,\ldots,1022,2047\}$ and the TLB holds $\{0,\ldots,30,2047\}$. The
test therefore gets \textbf{32 TLB hits} (pages $0$--$30$ and $2047$),
\textbf{2016 TLB misses}, and \textbf{1024 page faults} (everything
except the 1024 retained frames). Time:
$32M + 2016\times 2M + 1024D$.

\textbf{Microbenchmark.} Sweep \texttt{NUM\_PAGES} and time one
iteration. A jump from $\approx N\cdot M$ to $\approx 2N\cdot M$ is the
TLB-capacity cliff (page-table walks appear). A second, much larger jump
(milliseconds, i.e.\ $D$) is the RAM-capacity cliff (page faults).
\end{Sol}

\begin{Qn}{Hard}{Lect 15}{CS-537 Fa04 Mid2, Q10; ISI Mandar 2012 End-sem, Q3(c)--(d)}{Belady's anomaly, and why FIFO is not a stack algorithm}
Reference string
\[
1,2,3,4,1,2,5,1,2,3,4,5
\]
(equivalently $A,B,C,D,A,B,E,A,B,C,D,E$).
\begin{enumerate}
  \item[(a)] FIFO misses with 3 frames and with 4 frames (memory empty at
        the start). Does the comparison surprise you?
  \item[(b)] What is a \textbf{stack algorithm}? Using the definition and
        this string, show that FIFO is \emph{not} a stack algorithm.
\end{enumerate}
Appears as CS-537 Fa04 Mid2 Q10 (parts (a) only; letters $A$--$E$) and
as ISI Mandar 2012 End-sem Q3(c)--(d) (numeric string plus the stack
definition).
\end{Qn}
\begin{Sol}
(a) FIFO, 3 frames (queue, oldest on the left). Misses marked $*$:
\begin{center}
\small
\begin{tabular}{c|cccccccccccc}
ref    & 1 & 2 & 3 & 4 & 1 & 2 & 5 & 1 & 2 & 3 & 4 & 5 \\
\hline
frames & 1 & 12& 123& 234& 341& 412& 125& 125& 125& 253& 534& 534 \\
       & $*$& $*$& $*$& $*$& $*$& $*$& $*$&   &   & $*$& $*$&     \\
\end{tabular}
\end{center}
\textbf{9} misses. FIFO, 4 frames:
\begin{center}
\small
\begin{tabular}{c|cccccccccccc}
ref    & 1 & 2 & 3 & 4 & 1 & 2 & 5 & 1 & 2 & 3 & 4 & 5 \\
\hline
frames & 1 & 12& 123&1234&1234&1234&2345&3451&4512&5123&1234&2345 \\
       & $*$& $*$& $*$& $*$&   &   & $*$& $*$& $*$& $*$& $*$& $*$ \\
\end{tabular}
\end{center}
\textbf{10} misses. More frames, \textbf{more} faults: \textbf{Belady's
anomaly}. FIFO's resident set is an arrival-order artefact, not ``the
$k$ most useful pages''.

(b) A replacement policy is a \textbf{stack algorithm} if, for every
prefix of every reference string, the resident set with $k$ frames is a
\textbf{subset} of the resident set with $k+1$ frames. Stack algorithms
(such as LRU, OPT, and history-wide LFU with consistent tie-breaking;
note that Clock/second-chance is \emph{not} a stack algorithm) cannot
exhibit Belady's anomaly: every $k$-frame hit is a $(k+1)$-frame hit, so
faults are non-increasing in $k$.

After the first 5 (the first access to 5), FIFO-3 holds $\{1,2,5\}$ and
FIFO-4 holds $\{2,3,4,5\}$. Since $1\in$ FIFO-3 but $1\notin$ FIFO-4, the
inclusion fails, so FIFO is not a stack algorithm.

(For comparison on the same string: LRU has 10 misses with 3 frames and
8 with 4 frames --- monotone, as the stack property requires.)
\end{Sol}

\begin{Qn}{Hard}{Lect 15}{ISI Mandar 2022 End-sem, Q4}{Protecting a swapped-out page, and a proof of FIFO Belady}
\begin{enumerate}
  \item[(a)] In a paging system a page that belongs to the logical address
        space need not be in physical memory. Explain how memory
        protection is still provided.
  \item[(b)] Prove that if the number of frames allocated to a process
        increases from $n$ to $n+1$ ($n\ge 3$), and memory starts empty,
        the reference string
        \[
          1,2,\ldots,n+1,\;
          1,2,\ldots,n-1,\;
          n+2,\;
          1,2,\ldots,n+2
        \]
        always exhibits Belady's anomaly for FIFO.
\end{enumerate}
\end{Qn}
\begin{Sol}
(a) The PTE still distinguishes three states. Valid$\,=0$: the VPN is
not in the address space $\to$ protection/segmentation fault, process
killed (or sent \texttt{SIGSEGV}). Valid$\,=1$, present$\,=0$: a
\emph{legal} page that lives on disk $\to$ page fault, OS fills a frame,
then the access is replayed. Protection bits (R/W/X) travel with the
PTE (and into the TLB once present$\,=1$), so a write to an
\texttt{r-x} page still traps after the page is resident. Swapping a
page out does not punch a hole in the protection check: it only delays
the translation.

(b) Write $S$ for the string; $|S|=(n+1)+(n-1)+1+(n+2)=3n+3$.
(For $n=3$ this is exactly $1,2,3,4,1,2,5,1,2,3,4,5$.)

\textbf{FIFO with $n$ frames.}
\begin{itemize}
  \item $1,\ldots,n+1$: $n+1$ compulsory misses, then frames
        $=[2,3,\ldots,n+1]$ (1 was the oldest).
  \item $1,2,\ldots,n-1$: access $1$ misses and evicts $2$, leaving
        $[3,\ldots,n+1,1]$; access $2$ then misses and evicts $3$; by
        induction each $k \in \{1,\ldots,n-1\}$ was just evicted at the
        previous step, so all $n-1$ accesses miss and the queue ends as
        $[n+1,1,2,\ldots,n-1]$.
  \item $n+2$: miss, frames $=[1,2,\ldots,n-1,n+2]$.
  \item $1,\ldots,n+2$: hits on $1,\ldots,n-1$; misses on $n$ and
        $n+1$ (whose victims are 1 and 2, both $\le n-1$ because
        $n \ge 3$, so $n+2$ survives in the queue); hit on $n+2$. Two
        misses.
\end{itemize}
Total: $(n+1)+(n-1)+1+2=2n+3$ misses.

\textbf{FIFO with $n+1$ frames.}
\begin{itemize}
  \item $1,\ldots,n+1$: $n+1$ misses, and now \emph{all} of them fit:
        $[1,2,\ldots,n+1]$.
  \item $1,\ldots,n-1$: all hits.
  \item $n+2$: miss, evict 1, frames $=[2,\ldots,n+1,n+2]$.
  \item $1,\ldots,n+2$: a rotating FIFO of $n+1$ frames over $n+2$
        distinct pages, so \emph{every} access misses ($n+2$ misses).
\end{itemize}
Total: $(n+1)+0+1+(n+2)=2n+4$ misses.

Since $2n+4 > 2n+3$, FIFO with $n+1$ frames faults more often than FIFO
with $n$ frames on $S$. That is Belady's anomaly.

(Check: $n=3$ gives 9 vs.\ 10, matching the previous question; $n=4$
gives 11 vs.\ 12, verified by direct simulation.)
\end{Sol}

\clearpage

\section{Lectures 16--20: Concurrency --- Threads, Locks, CVs, Semaphores \& Deadlock}

\subsection{Lectures 16--17: Concurrency, Threads \& Locks}

\begin{Qn}{Hard}{Lect 16--17}{Mandar ISI 2015 Mid Q4}{Concurrent BST \texttt{insert} corrupts the tree}
Consider

\begin{lstlisting}[numbers=none]
void insert(NODE *p, NODE **root)
{
    NODE *r;
    if (root == NULL) return;
    if (*root == NULL) { *root = p; return; }
    r = *root;
    if (p->value == r->value) return;
    if (p->value < r->value) { insert(p, &(r->left)); return; }
    if (p->value > r->value) { insert(p, &(r->right)); return; }
}
\end{lstlisting}

Two processes $P_1$ and $P_2$ share a binary search tree. Using an example, show how the tree may be corrupted if they execute \texttt{insert} concurrently. Show the initial tree, the two calls, the intended final tree, the interleaving, and the actual final tree.
\end{Qn}
\begin{Sol}
Corruption occurs whenever both threads reach the same \texttt{NULL} child pointer simultaneously. Take an initially empty tree (\texttt{*root == NULL}), with $P_1$ calling \texttt{insert(p5, \&root)} (key $5$) and $P_2$ calling \texttt{insert(p15, \&root)} (key $15$). Intended final tree: both keys $5$ and $15$ present (either $5$ at root with right child $15$, or $15$ at root with left child $5$).

\begin{center}
\begin{tabular}{c|l|l|l}
Step & $P_1$ (insert 5) & $P_2$ (insert 15) & \texttt{*root} \\
\hline
0 & & & \texttt{NULL} \\
1 & \texttt{*root == NULL} (about to write) & & \\
2 & & \texttt{*root == NULL} (about to write) & \\
3 & \texttt{*root = p5} & & $5$ \\
4 & & \texttt{*root = p15} & $15$ \\
\end{tabular}
\end{center}

The store of $P_2$ overwrites $P_1$'s root pointer. Node $5$ is leaked (orphaned); the tree contains only $15$. Intended tree had both keys. Same lost-update pattern as the racy linked-list insert in Lecture~17: two threads read a \texttt{NULL} child / \texttt{NULL} head, each link their node, and the later store wins.

A second (non-empty tree) example: initial tree has root $10$ with empty left child; $P_1$ inserts $4$ and $P_2$ inserts $6$ (both belong on the left). Both recurse to \texttt{\&(10->left)}, both see \texttt{NULL}, both write. One of $\{4,6\}$ is silently lost.
\end{Sol}

\begin{Qn}{Hard}{Lect 16--17}{CS-537 Spring 2018 Mid Problem XXII}{How long can a racy list become?}
Assume the following list insertion, which inserts into a list pointed to by shared global \texttt{head}:

\begin{lstlisting}[numbers=none]
int List_Insert(int key) {
  node_t *n = malloc(sizeof(node_t));
  if (n == NULL) { return -1; }
  n->key = key;
  n->next = head;
  head = n;
  return 0;
}
\end{lstlisting}

This code is executed by each of three threads exactly once, with no locks. Assume \texttt{malloc()} is thread-safe and succeeds, and the list starts empty. How long might the list be when the three threads finish? Fill in A for possible, B for not possible.

\begin{enumerate}
\item[106.] 0
\item[107.] 1
\item[108.] 2
\item[109.] 3
\item[110.] 4
\end{enumerate}
\end{Qn}
\begin{Sol}
The swing \texttt{n->next = head; head = n;} is not atomic. A node is lost whenever two (or three) threads both read the same \texttt{head} and the later \texttt{head = n} overwrites the earlier insert.

\begin{itemize}
\item Length 3 (serialized inserts): possible. A.
\item Length 2: two inserts succeed, one is overwritten. A.
\item Length 1: all three read \texttt{head==NULL} and the last store wins; two nodes leaked. A.
\item Length 0: impossible --- at least one store to \texttt{head} commits a non-\texttt{NULL} node. B.
\item Length 4: only three nodes are allocated. B.
\end{itemize}

Official key: 106~B, 107~A, 108~A, 109~A, 110~B.
\end{Sol}

\begin{Qn}{Hard}{Lect 17}{CS-537 Spring 2002 Mid Q5}{Peterson, but \texttt{turn} is a coin flip}
Classic Peterson pseudocode for a two-thread \texttt{deposit} (source exam syntax):

\begin{lstlisting}[numbers=none]
boolean lock[0] = lock[1] = false;
int turn = 0;
void deposit(int amount) {
    lock[pid] = true;
    turn = 1 - pid;
    while (lock[1-pid] && (turn == (1 - pid)))
        ; /* spin */
    balance = balance + amount;
    lock[pid] = false;
}
\end{lstlisting}

Replace \texttt{turn = 1 - pid} with \texttt{turn = BinaryRandom()} (returns 0 or 1 uniformly). Does the code still function properly? If not, what can go wrong?
\end{Qn}
\begin{Sol}
\textbf{No --- mutual exclusion fails.} Peterson's invariant relies on every thread setting \texttt{turn = 1 - pid} (``after you''), so a thread can \emph{never} grant itself \texttt{turn} while entering. Replacing that store with \texttt{BinaryRandom()} allows each thread to write its \emph{own} id into \texttt{turn}, letting both pass the spin-loop simultaneously:

\begin{center}
\small
\begin{tabularx}{\linewidth}{@{}c|X|X@{}}
Step & T0 & T1 \\
\hline
1 & \texttt{lock[0]=true} & \\
2 & & \texttt{lock[1]=true} \\
3 & \texttt{turn = 0} (random) & \\
4 & \texttt{while (lock[1] \&\& turn==1)} $\to$ false (\texttt{turn==0}), \textbf{enters CS} & \\
5 & & \texttt{turn = 1} (random) \\
6 & & \texttt{while (lock[0] \&\& turn==0)} $\to$ false (\texttt{turn==1}), \textbf{enters CS} \\
\end{tabularx}
\end{center}

Both threads are now inside the critical section simultaneously, violating mutual exclusion. In addition, bounded waiting / fairness is lost because a thread can repeatedly draw its own id on re-entry while the other thread waits.
\end{Sol}

\begin{Qn}{Hard}{Lect 17}{CS-537 Spring 2018 Mid Problem XX}{Broken flag lock: how many threads get in?}
\begin{lstlisting}[numbers=none]
void init(lock_t *mutex) {
  mutex->flag = 0; /* 0 -> lock is available, 1 -> held */
}
void lock(lock_t *mutex) {
  while (mutex->flag == 1) /* L1 */
    ;                     /* L2 */
  mutex->flag = 1;        /* L3 */
}
void unlock(lock_t *mutex) {
  mutex->flag = 0;        /* L4 */
}
\end{lstlisting}

Assume 5 threads compete for this lock. How many threads can possibly hold the lock \emph{simultaneously}? Fill in A for possible, B for not possible.

\begin{enumerate}
\item[96.] 1 \qquad 97.~2 \qquad 98.~3 \qquad 99.~4 \qquad 100.~5
\end{enumerate}
\end{Qn}
\begin{Sol}
The test (L1) and the set (L3) are not atomic --- this is Lecture~17's broken flag lock. All five threads can pass L1 while \texttt{flag} is still 0 and then all execute L3.

Explicit 2-thread trace, then scale:

\begin{center}
\begin{tabular}{c|l|l|l}
Step & T1 & T2 & \texttt{flag} \\
\hline
1 & L1: \texttt{flag==1}? no, fall through & & 0 \\
2 & & L1: \texttt{flag==1}? no, fall through & 0 \\
3 & L3: \texttt{flag=1} & & 1 \\
4 & & L3: \texttt{flag=1} & 1 \\
\end{tabular}
\end{center}

Both ``acquired''. The same pattern with five threads: all five read 0, all five write 1. So 1 through 5 concurrent holders are all possible (including the lucky case of only one). Official key: 96--100 all A.
\end{Sol}

\begin{Qn}{Hard}{Lect 17}{CS-537 Fall 2013 Mid Q12}{Build a spinlock from Load-and-Store-Zero}
A new instruction, in C pseudo-code:

\begin{lstlisting}[numbers=none]
int LoadAndStoreZero(int *addr) {
  int old = *addr;
  *addr   = 0;
  return old;
}
\end{lstlisting}

Write \texttt{init} / \texttt{acquire} / \texttt{release} for a working spinlock that uses this instruction.
\end{Qn}
\begin{Sol}
\texttt{LoadAndStoreZero} (abbreviated \texttt{LASZ}) can only \emph{force} the cell to 0, so 0 must mean \emph{held} and a non-zero value (say 1) must mean \emph{free}. Acquire spins until it observes a non-zero old value (it just stole the token by zeroing it). Release stores 1.

\begin{lstlisting}[numbers=none]
void init(int *lock)    { *lock = 1; }              /* 1 = free, 0 = held */
void acquire(int *lock) { while (LASZ(lock) == 0) ; } /* spin while already 0 */
void release(int *lock) { *lock = 1; }
\end{lstlisting}

If you invert the convention (free=0) you can never mark the lock held: \texttt{LASZ} cannot write a 1. Official solution uses exactly this polarity.
\end{Sol}

\begin{Qn}{Hard}{Lect 17}{CS-537 Fall 2013 Mid Q13}{Fetch-and-add built from \texttt{xchg} --- why it is wrong}
McFlub claims this implements fetch-and-add using only atomic exchange:

\begin{lstlisting}[numbers=none]
int FetchAndAdd(int *addr, int amt) {
  int should;
  do {
    should = *addr + amt;
  } while (atomic_exchange(addr, should) != should);
}
\end{lstlisting}

What is wrong? How should one build fetch-and-add given only \texttt{xchg}?
\end{Qn}
\begin{Sol}
Two fatal bugs (plus a missing \texttt{return}):
\begin{enumerate}
  \item \textbf{Infinite loop even with one thread:} \texttt{atomic\_exchange(addr, should)} returns the \emph{old} value of \texttt{*addr}. For any $\texttt{amt}\ne 0$, the returned old value is strictly different from $\texttt{should}=\texttt{old}+\texttt{amt}$, so \texttt{atomic\_exchange(addr, should) != should} is \emph{true} and the \texttt{do--while} loop spins forever, repeatedly adding \texttt{amt} to \texttt{*addr}.
  \item \textbf{Blind overwrite (lost updates) with multiple threads:} \texttt{atomic\_exchange} \emph{unconditionally} overwrites \texttt{*addr} with \texttt{should}, even if another thread modified \texttt{*addr} between the load and the exchange. Unlike compare-and-swap, \texttt{xchg} cannot test ``write only if unchanged''.
\end{enumerate}

Correct approach: use \texttt{xchg} to build a lock, then

\begin{lstlisting}[numbers=none]
int FetchAndAdd(int *addr, int amt) {
  int old;
  acquire(&faa_lock);
  old = *addr;
  *addr = old + amt;
  release(&faa_lock);
  return old;
}
\end{lstlisting}

Alternatively, if compare-and-swap (or LL/SC) is available, the textbook retry loop is
\texttt{do \{ old = *addr; \} while (CAS(addr, old, old+amt) fails);}.
\texttt{xchg} is not CAS.
\end{Sol}

\begin{Qn}{Hard}{Lect 17}{CS-537 Spring 2010 Final Problem 9}{Build a lock from \texttt{FetchAndSubtract}}
\begin{lstlisting}[numbers=none]
int FetchAndSubtract(int *location) {
    int value = *location; /* read */
    *location = value - 1; /* decrement, store */
    return value;          /* old value */
}
\end{lstlisting}

Write \texttt{lock\_t}, \texttt{lock\_init}, \texttt{lock}, and \texttt{unlock} that use this primitive.
\end{Qn}
\begin{Sol}
Two working constructions. Both use that FAS returns the \emph{old} value and writes $\mathrm{old}-1$.

\textbf{Construction A --- TAS analogue, free $=1$.} Only a thread that observes old $=1$ has the token. Failed FAS's can drive the cell negative, so waiters must \emph{not} FAS while the lock looks held (test-and-test-and-set), and unlock \emph{stores} 1 (which also repairs any leftover negative):

\begin{lstlisting}[numbers=none]
typedef struct { int v; } lock_t;          /* 1 = free, else held */

void lock_init(lock_t *l) { l->v = 1; }

void lock(lock_t *l) {
    while (1) {
        while (l->v != 1) ;                /* spin on a plain load */
        if (FetchAndSubtract(&l->v) == 1)
            return;                        /* v was 1; we decremented it to 0 */
        /* lost the race: another FAS won; retry after unlock stores 1 */
    }
}

void unlock(lock_t *l) { l->v = 1; }
\end{lstlisting}

\textbf{Construction B --- ticket lock.} FAS on a counter that starts at 0 issues tickets $0,-1,-2,\ldots$; negate to get $0,1,2,\ldots$:

\begin{lstlisting}[numbers=none]
typedef struct { int ticket; int turn; } lock_t;

void lock_init(lock_t *l) { l->ticket = 0; l->turn = 0; }

void lock(lock_t *l) {
    int myturn = -FetchAndSubtract(&l->ticket);  /* 0, then 1, then 2, ... */
    while (l->turn != myturn)
        ;
}

void unlock(lock_t *l) { l->turn = l->turn + 1; }
\end{lstlisting}

B is FIFO (Lecture~17 ticket lock; like any finite-width ticket lock, \texttt{ticket} eventually wraps after $2^{31}$--$2^{32}$ acquisitions). A is unfair, like TAS. Naive ``\texttt{while (FetchAndSubtract(\&l->v) != 1);}'' without the inner load-spin causes heavy cache-line invalidation traffic while spinning and --- more dangerously --- can suffer 32-bit signed integer underflow after $\sim 2^{32}$ failed decrements, wrapping \texttt{v} past \texttt{INT\_MIN} back to $1$ while the lock is still held and breaking mutual exclusion.
\end{Sol}

\begin{Qn}{Hard}{Lect 16}{Mandar ISI 2016 End Q1(b)}{Two threads, \texttt{putchar(A[i])}: four locality combinations}
Two threads of one process each execute

\begin{lstlisting}[numbers=none]
for (i = 0; i < 100; i++) putchar(A[i]);
\end{lstlisting}

Each of \texttt{A} and \texttt{i} may be local or global (four cases). For each, give the min and max number of characters printed. Justify.
\end{Qn}
\begin{Sol}
\texttt{putchar} of a single character is treated as atomic. The loop bound is 100 per thread if each has a private index.

\begin{center}
\small
\begin{tabularx}{\linewidth}{@{}l|l|c|X@{}}
\texttt{A} & \texttt{i} & min chars & max chars \\
\hline
local  & local  & 200 & 200 \\
global & local  & 200 & 200 \\
local  & global & 100 & 201 (atomic \texttt{i++}); $\ge 5149$ ($\Theta(N^2)$, non-atomic \texttt{i++}) \\
global & global & 100 & 201 (atomic \texttt{i++}); $\ge 5149$ ($\Theta(N^2)$, non-atomic \texttt{i++}) \\
\end{tabularx}
\end{center}

\textbf{\texttt{i} local (\texttt{A} local or global).} Each thread has its own private loop counter on its stack, so each runs 100 iterations $\Rightarrow$ always \textbf{200} characters. Sharing \texttt{A} affects only which characters appear, not the count.

\textbf{\texttt{i} global (\texttt{A} local or global).} Both threads share the single loop variable \texttt{i} (including the initializer \texttt{i = 0}!).
\begin{itemize}
  \item \textbf{Minimum = 100:} Both threads execute \texttt{i = 0} before either runs the loop body; thereafter they take turns without lost increments. Even under non-atomic \texttt{i++} (where a stale store can jump \texttt{i} forward after a reset), reaching \texttt{i = 100} starting from 0 still requires a causal chain of 100 increments each reading the previous value ($0\to 1\to\cdots\to 100$), and each increment executes after a distinct \texttt{putchar}, so at least 100 characters must be printed.
  \item \textbf{Maximum (atomic \texttt{i++}):} If T1 runs to completion (\texttt{i} goes $0\to 100$, printing 100 chars) and \emph{then} T2 begins and executes \texttt{i = 0}, T2 prints another 100 chars $\Rightarrow$ 200; and if T1 pauses right before checking \texttt{i < 100} when \texttt{i=100}, T2 resets \texttt{i=0}, and both threads later pass \texttt{i < 100} at \texttt{i=99} before either increments, they achieve \textbf{201}.
  \item \textbf{Maximum (non-atomic \texttt{i++}, load/add/store):} Lost updates allow $\Theta(N^2)$ prints ($\ge 5149$). With T1 paused just before its \texttt{i < 100} check after completing $0..99$ (100 chars) and T2 resetting \texttt{i = 0}, T1 executes \texttt{putchar} at \texttt{i = 0} (+1 char), loads \texttt{0} inside \texttt{i++}, pauses while T2 runs $0..98$ (99 chars), and stores \texttt{1}; next T2 prints at \texttt{i = 1} (+1), loads \texttt{1}, pauses while T1 runs $1..98$ (98 chars), and stores \texttt{2}, and so on ($\approx 100 + \sum_{k=1}^{99} (k+1) = 5149$ characters).
\end{itemize}
In all cases, whether \texttt{A} is local or global does not change the \emph{count} of printed characters.
\end{Sol}

\begin{Qn}{Medium-Hard}{Lect 17}{CS-537 Spring 2012 Mid Q15}{Locked list insert that still deadlocks itself}
\begin{lstlisting}[numbers=none]
mutex_t m    = PTHREAD_MUTEX_INITIALIZER;
node_t *head = NULL;

int List_Insert(int key) {
    mutex_lock(&m);
    node_t *n = malloc(sizeof(node_t));
    if (n == NULL) { return -1; } /* failed to insert */
    n->key = key;
    n->next = head;
    head    = n;
    mutex_unlock(&m);
    return 0;
}
\end{lstlisting}

Insert a subtle change that makes the code correct.
\end{Qn}
\begin{Sol}
The \texttt{return -1} path leaves \texttt{m} held, so every subsequent caller spins/blocks forever. Unlock on \emph{every} return path:

\begin{lstlisting}[numbers=none]
if (n == NULL) { mutex_unlock(&m); return -1; }
\end{lstlisting}

A better (still correct) rewrite also shrinks the critical section: \texttt{malloc} and the \texttt{key} assignment do not touch shared list state, so they may sit outside the lock. The two pointer writes \texttt{n->next = head; head = n;} must remain inside. Lecture~17 emphasises both properties: unlock on every path, and keep the CS as small as is safe.
\end{Sol}

\begin{Qn}{Hard}{Lect 17}{Mandar ISI 2016 Mid Q4(a)}{Another broken Peterson: flag first, \texttt{turn} only after the CS}
\begin{lstlisting}[numbers=none]
shared char want[2] = {0,0};
shared int turn = 0;
P_i() { while (1) {
    want[i] = 1;
    while (want[j] && turn != i) ;
    critical_section();
    want[i] = 0;
    turn = j;
    remainder_section();
}}
\end{lstlisting}

Show that mutual exclusion fails. (This is \emph{not} the same broken variant as 2012 End Q2(a).)
\end{Qn}
\begin{Sol}
\texttt{turn} is never written in the entry section; it is written only on exit, and it starts at 0. So on the first attempt $P_0$ treats \texttt{turn==0} as ``already my turn'' even if $P_1$ is inside.

\begin{center}
\small
\begin{tabularx}{\linewidth}{@{}c|l|l|X@{}}
Step & $P_0$ & $P_1$ & \texttt{want}, \texttt{turn} \\
\hline
1 & & \texttt{want[1] = 1} & $(0,1),\ 0$ \\
2 & & \texttt{while(want[0] \&\& turn!=1)} & \texttt{want[0]=0}, enter CS \\
3 & \texttt{want[0] = 1} & & $(1,1),\ 0$ \\
4 & \texttt{while(want[1] \&\& turn!=0)} & &
    \texttt{turn==0==i} $\Rightarrow$ wait false, enter CS \\
\end{tabularx}
\end{center}

Both processes are now in the critical section simultaneously, violating mutual exclusion. (Once the other process exits and clears \texttt{want[j]=0}, a lone process can still re-enter regardless of \texttt{turn}, so this is not strict alternation; the fatal flaw is the mutual-exclusion violation above.)
\end{Sol}

\begin{Qn}{Hard}{Lect 17}{CS-537 Fall 2021 Midterm 2 Q18--Q19}{A TAS spinlock that takes the flag by value}
Bug writes:

\begin{lstlisting}[numbers=none]
int m = 0; /* global */

void lock_acquire(int m) {
    while (xchg(&m, 1) == 0)
        ;
}
void lock_release(int m) {
    xchg(&m, 0);
}
\end{lstlisting}

Q18. What is wrong?
\begin{enumerate}
\item[(a)] Nothing; it works!
\item[(b)] It will spin forever
\item[(c)] The release does not free the lock
\item[(d)] The code will crash
\item[(e)] None of the above
\end{enumerate}

Q19. Same, but the wait is \texttt{while (xchg(\&m, 1) == 1)}. What is wrong? (same options.)
\end{Qn}
\begin{Sol}
\textbf{Q18.} Two independent bugs. (i)~The parameter \texttt{m} \emph{shadows} the global: \texttt{xchg} mutates the private stack copy, so every caller exits after two \texttt{xchg}s on its private integer (first returns 0, loops; second returns 1, exits) with zero mutual exclusion. (Option~(c) is not the sole defect because \texttt{lock\_acquire} never acquires the global lock either; the fundamental bug is total lack of mutual exclusion, which is not among (a)--(d).) (ii)~Even if the parameter were removed, the condition is inverted: \texttt{xchg} returns the \emph{old} value; spinning while it equals 0 means ``loop while the lock was free'' --- the first thread's first \texttt{xchg} on 0 returns 0 (loop continues), its second \texttt{xchg} on 1 returns 1 (exits leaving \texttt{m=1}), and a second thread's first \texttt{xchg} immediately sees 1 and also enters. Thus mutual exclusion fails even on a shared global. Official answer: \textbf{(e)}.

\textbf{Q19.} As written, the parameter \texttt{int m} still shadows the global, so each thread runs \texttt{xchg} on its own stack copy (which holds 0), exits the loop immediately, and enters the CS without mutual exclusion (again, (c) alone misses that \texttt{acquire} never locks the global either) $\Rightarrow$ \textbf{(e)}. (Only if the shadowed parameter is ignored as an exam typo does the code become a standard working TAS spinlock, which would give (a).)

The exam-grade takeaway: \texttt{xchg} must operate on the \emph{shared} flag, and the spin condition must be ``old value was held''.
\end{Sol}

\begin{Qn}{Hard}{Lect 16}{CS-537 Fall 2021 Midterm 2 Q3--Q4}{Split lock around \texttt{slot++} and \texttt{array[slot]}}
\begin{lstlisting}[numbers=none]
mutex_t m;
int slot = -1;
int array[2] = {0, 0};

function1() {
    mutex_acquire(&m);
    slot++;
    mutex_release(&m);
    int tid = get_counter();   /* 1, then 2, internally locked */
    mutex_acquire(&m);
    array[slot] = tid;
    mutex_release(&m);
}
\end{lstlisting}

Two threads call \texttt{function1}. Final \texttt{array}?
(a)~$[0,0]$ (b)~$[0,1]$ (c)~$[1,2]$ (d)~$[0,2]$ (e)~indeterminate

Q4: one lock around all three of \texttt{slot++}, \texttt{get\_counter()}, \texttt{array[slot]=tid}. Now?
\end{Qn}
\begin{Sol}
Q3. Because the lock is released between \texttt{slot++} and \texttt{array[slot] = tid}, several outcomes are possible:
\begin{itemize}
  \item If T0 completes the entire function before T1 starts: \texttt{array = [1,2]}.
  \item If both threads complete \texttt{slot++} before either writes \texttt{array[slot]}, then \texttt{slot == 1} for both writes, leaving \texttt{array[0] == 0} and yielding either \texttt{[0,1]} or \texttt{[0,2]} depending on which write lands last.
  \item (Note: \texttt{[2,1]} is \emph{not} possible: writing \texttt{array[0]} requires the other thread not to have run \texttt{slot++} yet, which means it cannot have called \texttt{get\_counter()} yet, so the writer of \texttt{array[0]} must have \texttt{tid = 1}.)
\end{itemize}
Because \texttt{[1,2]}, \texttt{[0,1]}, and \texttt{[0,2]} are all reachable, the answer is \textbf{(e) indeterminate}.

Trace for $[0,2]$:
\begin{center}
\begin{tabular}{c|l|l|l}
Step & T0 & T1 & slot, array \\
\hline
1 & \texttt{slot++} $\to 0$ & & 0 \\
2 & & \texttt{slot++} $\to 1$ & 1 \\
3 & & \texttt{tid=get\_counter()}$\to1$ & \\
4 & \texttt{tid=get\_counter()}$\to2$ & & \\
5 & & \texttt{array[1]=1} & $[0,1]$ \\
6 & \texttt{array[1]=2} & & $[0,2]$ \\
\end{tabular}
\end{center}
and swapping steps 5--6 yields $[0,1]$, while running T0 to completion before T1 yields $[1,2]$. Indeterminate.

Q4. First thread through holds the lock, gets \texttt{slot=0}, \texttt{tid=1}, writes \texttt{array[0]=1}. Second gets \texttt{slot=1}, \texttt{tid=2}, writes \texttt{array[1]=2}. (c)~$[1,2]$.
\end{Sol}

\begin{Qn}{Medium-Hard}{Lect 16}{CS-537 Spring 2018 Mid Problem XXIII}{Background \texttt{calloc} without a join}
\begin{lstlisting}[numbers=none]
void *background_malloc(void *arg) {
  int **int_ptr = (int **) arg;
  *int_ptr = calloc(1, sizeof(int));
  **int_ptr = 10;
  return NULL;
}
int main() {
  pthread_t p1;
  int *result = NULL;
  pthread_create(&p1, NULL, background_malloc, &result);
  printf("%d\n", *result);
  return 0;
}
\end{lstlisting}

Possible outcomes? A/B: 111~prints 0; 112~prints 10; 113~prints 100; 114~crashes; 115~hangs forever.
\end{Qn}
\begin{Sol}
No \texttt{pthread\_join}, so \texttt{main} races with the child on \texttt{result}.

\begin{itemize}
\item 111: child \texttt{calloc}s (zeroed) and publishes \texttt{result}, but has not yet stored 10; \texttt{main} prints 0. A.
\item 112: child runs to completion first. A.
\item 113: 100 is never written. B.
\item 114: \texttt{main} dereferences \texttt{NULL} \texttt{result} before the child publishes. A.
\item 115: nobody waits on anybody; no infinite lock. B.
\end{itemize}

Official: 111~A, 112~A, 113~B, 114~A, 115~B.
\end{Sol}

\subsection{Lecture 18: Condition Variables}

\begin{Qn}{Hard}{Lect 18}{CS-537 Fa04 Mid1, Q9; Fa04 Final, Q8}{Mesa semantics: two things wrong in a monitor producer/consumer}
The following producer/consumer is written inside a monitor (the monitor
lock is acquired on entry and released on exit / \texttt{wait}). It uses
\textbf{Mesa} semantics.
\begin{lstlisting}[numbers=none]
producer() {
    if (used == MAX_BUFFERS)
        condition.wait();
    /* produce a buffer */
    used++;
    condition.signal();
}
consumer() {
    if (used == 0)
        condition.wait();
    /* consume a buffer */
    used--;
    condition.signal();
}
\end{lstlisting}
Fix the \textbf{two major} things wrong with this code so that it works
with any number of producers and consumers. Then explain, independently,
why Mesa semantics demand \texttt{while} rather than \texttt{if} even
after the other fix.
\end{Qn}
\begin{Sol}
Two independent bugs.

\textbf{Bug 1 --- \texttt{if} instead of \texttt{while}.}
Mesa \texttt{signal} is a \emph{hint}: it moves one waiter to ready; the
waiter does \emph{not} run immediately, and the signaller keeps the
monitor. Between the signal and the waiter reacquiring the monitor,
another thread can sneak in and change \texttt{used}. The waiter must
re-check the predicate. Replace both \texttt{if}s by
\texttt{while}.

A concrete stolen-wakeup (single slot, two consumers $C_1,C_2$, one
producer $P$): $C_1$ waits (\texttt{used == 0}). $P$ produces, signals
$C_1$ (now ready, not running), and the producer may itself wait on a
full buffer. $C_2$ enters, sees \texttt{used == 1}, skips the
\texttt{if}, consumes, \texttt{used = 0}. $C_1$ finally re-enters and,
because of the \texttt{if}, does not re-check: it consumes an empty
buffer.

\textbf{Bug 2 --- one CV for two kinds of waiters.}
Producers wait for ``not full'' and consumers wait for ``not empty'',
but both sleep on the same \texttt{condition}. A consumer's
\texttt{signal} may wake another consumer; a producer's \texttt{signal}
may wake another producer. With \texttt{while} the wrongly-woken thread
goes back to sleep, and it is possible that \emph{everyone} is asleep
(the thread that should have been woken is still sleeping, and the
waker has itself gone to wait). Fix: two CVs, \texttt{empty} and
\texttt{fill} (producers wait on \texttt{empty} and signal
\texttt{fill}; consumers wait on \texttt{fill} and signal
\texttt{empty}). The covering-condition alternative is
\texttt{broadcast} on the single CV; that is correct but wakes extra
threads (see the single-CV \texttt{broadcast} question below).

\textbf{Why Mesa forces \texttt{while} even with two CVs.}
(i)~Another thread of the \emph{same} kind can acquire the monitor
after the signal and consume the state (two consumers, one item).
(ii)~Libraries may deliver \textbf{spurious wakeups}. Hoare semantics
run the waiter immediately and would make \texttt{if} safe; real
systems are Mesa, so always write \texttt{while}.
\end{Sol}

\begin{Qn}{Hard}{Lect 18}{CS-537 Sp16 Mid, Q6 traces 1--4; cf.\ Fa08 Final, Q3(a)--(d)}{Bounded-buffer traces with two CVs}
Assume the textbook two-CV bounded buffer. Lines are labelled as
follows (no spurious wakeups). A thread stops running only when it
blocks in a CV or a lock --- no untimely preemption.
\begin{lstlisting}[numbers=none]
void *producer(void *arg) {            void *consumer(void *arg) {
    for (i = 0; i < loops; i++) {          while (1) {
        Pthread_mutex_lock(&mutex);  /*p1*/    Pthread_mutex_lock(&mutex);  /*c1*/
        while (count == max)         /*p2*/    while (count == 0)           /*c2*/
            Pthread_cond_wait(&empty, &mutex); /*p3*/
                                               Pthread_cond_wait(&fill, &mutex); /*c3*/
        put(i);                      /*p4*/    int tmp = get();             /*c4*/
        Pthread_cond_signal(&fill);  /*p5*/    Pthread_cond_signal(&empty); /*c5*/
        Pthread_mutex_unlock(&mutex);/*p6*/    Pthread_mutex_unlock(&mutex);/*c6*/
    }                                      printf("%d\n", tmp);
}                                          }
                                       }
\end{lstlisting}
Write the line-number flow of each thread. Stop as specified.
\texttt{count} starts at 0.

\begin{enumerate}[(a)]
\item Trace 1: one producer \texttt{Pa}, one consumer \texttt{Ca},
      \texttt{max = 1}. \texttt{Pa} runs first. Stop when the consumer
      has consumed one entry.
\item Trace 2: same threads, \texttt{max = 3}. \texttt{Pa} runs first.
      Stop when the consumer has consumed one entry.
\item Trace 3: \texttt{max = 1}. \texttt{Ca} runs first. Stop when the
      consumer has consumed one entry.
\item Trace 4: one producer, two consumers \texttt{Ca},\texttt{Cb},
      \texttt{max = 1}. \texttt{Ca} runs first, then \texttt{Cb}, then
      \texttt{Pa}. Stop when \texttt{Pa} has produced one entry.
\item (Fa08 Final Q3(c), Hoare-style extra.) Same code, \texttt{max = 1},
      \texttt{Pa} first, but now \emph{every signal also context-switches}
      to a waiting thread if one exists. Give the flow until the
      consumer has consumed two entries.
\end{enumerate}
\end{Qn}
\begin{Sol}
Because a thread is descheduled only on \texttt{wait} (or lock block),
the running thread continues through lock/unlock/signal.

(a)~\texttt{max = 1}, producer first. \texttt{Pa} puts one item
(\texttt{count = 1}), loops, sees \texttt{count == max}, waits on
\texttt{empty}. Then \texttt{Ca} runs, sees \texttt{count == 1}, gets.
Official stop-at-consume:
\[
\texttt{Pa: } 1,2,4,5,6,1,2,3
\qquad
\texttt{Ca: } 1,2,4.
\]

(b)~\texttt{max = 3}. \texttt{Pa} can put three times before waiting:
\[
\texttt{Pa: } (1,2,4,5,6)^3,\; 1,2,3
\qquad
\texttt{Ca: } 1,2,4.
\]
(Three full produce loops, then the wait on \texttt{empty}.) Extra
slots mean fewer context switches: with \texttt{max = 1} the producer
blocks after every put; with \texttt{max = 3} it puts three times
before the first wait.

(c)~\texttt{Ca} first: it waits immediately. \texttt{Pa} puts, signals
\texttt{fill} (wakes \texttt{Ca}), loops, waits on the now-full buffer.
\texttt{Ca} returns from wait, \textbf{re-checks} \texttt{c2} (the
\texttt{while}), finds \texttt{count == 1}, gets.
\[
\texttt{Pa: } 1,2,4,5,6,1,2,3
\qquad
\texttt{Ca: } 1,2,3,\;2,4.
\]
The extra \texttt{c2} after wakeup is exactly Mesa: \texttt{wait}
returns holding the mutex, and the \texttt{while} predicate is tested
again.

(d)~Both consumers wait (\texttt{count = 0}). Then \texttt{Pa} puts
one item. Stop at the put:
\[
\texttt{Ca: } 1,2,3
\qquad
\texttt{Cb: } 1,2,3
\qquad
\texttt{Pa: } 1,2,4.
\]
\texttt{Pa}'s later \texttt{p5} would signal \texttt{fill} and wake
exactly one consumer (the right kind of waiter).

(e)~Signal-causes-switch (Hoare-like, assuming the woken thread also receives the mutex). \texttt{Pa} puts the first item
(\texttt{p5} signals \texttt{fill} with nobody waiting: no switch),
then waits on \texttt{empty}. \texttt{Ca} consumes and at \texttt{c5}
signals \texttt{empty}: \texttt{Pa} is waiting, so control passes to
\texttt{Pa} immediately. \texttt{Pa} resumes from \texttt{p3} into the
\texttt{while} test \texttt{p2}, produces the second item, signals
\texttt{fill} (consumer is not on the CV --- it is still inside
\texttt{consume}), then waits. \texttt{Ca} finishes the first consume
and takes the second item:
\[
\begin{aligned}
\texttt{Pa: }& 1,2,4,5,6,1,2,3,\; 2,4,5,6,1,2,3\\
\texttt{Ca: }& 1,2,4,5,\; 6,1,2,4.
\end{aligned}
\]
(The exact placement of \texttt{c6} relative to the second produce
depends on whether the switch happens \emph{during} \texttt{c5} or
after it; either way two items are produced and two consumed, and
the two CVs keep each signal aimed at the opposite kind of thread.)
\end{Sol}

\begin{Qn}{Medium-Hard}{Lect 18}{CS-537 Fa13 Mid, Q14}{Single CV plus \texttt{broadcast}: does it work?}
McFlub offers this producer/consumer, using a \emph{single} CV and
\texttt{Pthread\_cond\_broadcast} (wakes \emph{all} waiters):
\begin{lstlisting}[numbers=none]
void *producer(void *arg) {
    int i;
    for (i = 0; i < loops; i++) {
        Pthread_mutex_lock(&mutex);
        while (count == MAX)
            Pthread_cond_wait(&cv, &mutex);
        put(i);
        Pthread_cond_broadcast(&cv);   /* here! */
        Pthread_mutex_unlock(&mutex);
    }
}
void *consumer(void *arg) {
    int i;
    for (i = 0; i < loops; i++) {
        Pthread_mutex_lock(&mutex);
        while (count == 0)
            Pthread_cond_wait(&cv, &mutex);
        int tmp = get();
        Pthread_cond_broadcast(&cv);   /* here! */
        Pthread_mutex_unlock(&mutex);
        printf("%d\n", tmp);
    }
}
\end{lstlisting}
``Now I think that this code doesn't work. Or does it?''
\end{Qn}
\begin{Sol}
It \textbf{works}. The classic single-CV bug is that \texttt{signal}
wakes one waiter, possibly of the wrong kind, and everyone can end up
asleep. \texttt{broadcast} wakes \emph{all} waiters. Each of them
reacquires the mutex in turn, re-checks its \texttt{while} predicate,
and either proceeds or goes back to sleep. The thread that \emph{can}
make progress (a producer if \texttt{count < MAX}, a consumer if
\texttt{count > 0}) will find the predicate false and run; the others
re-wait. No lost wakeup, no everyone-asleep.

The cost is extra context switches: every signal-site wakes threads
that will immediately sleep again. That is why the textbook uses two
CVs and \texttt{signal}. Official Fa13 key: ``this code seemingly makes
the mistake of only using one condition variable \ldots but solves it
by using broadcast \ldots Because each thread re-checks the condition
when awoken, only those who should be able to make progress will do
so. Of course, this can be inefficient.''

Rule of thumb: if a program only works after changing
\texttt{signal} to \texttt{broadcast} and you do not think it should
need to, you probably have a bug (wrong CV, or \texttt{if} instead of
\texttt{while}). Here broadcast is a legitimate covering condition.
\end{Sol}

\subsection{Lecture 19: Semaphores}

\begin{Qn}{Medium-Hard}{Lect 19}{CS-537 Fall 2008 Final Q2}{Semaphore flows: binary and generalised, from lock + CV}
You are given these numbered chunks:
\begin{center}
\footnotesize
\begin{tabular}{@{}lll@{}}
  1. \texttt{while (s <= 0)} &
  2. \texttt{while (s >= 0)} &
  3. \texttt{while (s > 0)}  \\
  4. \texttt{while (s < 0)}  &
  5. \texttt{while (s == 0)} &
  6. \texttt{s--;}           \\
  7. \texttt{s++;}           &
  8. \texttt{cond\_wait(\&cond, \&lock);} &
  9. \texttt{cond\_signal(\&cond);} \\
  10. \texttt{mutex\_lock(\&lock);} &
  11. \texttt{mutex\_unlock(\&lock);} &
\end{tabular}
\end{center}
A flow is a comma-separated list of chunk numbers. \texttt{sem\_init}
sets \texttt{s = value}. Give a \emph{minimal} flow for each of:
(a)~binary \texttt{sem\_wait}; (b)~binary \texttt{sem\_post};
(c)~generalised \texttt{sem\_wait}; (d)~generalised \texttt{sem\_post}.
\end{Qn}
\begin{Sol}
This is the Zemaphore (Lect 19) written as a flow. Binary and generalised
\texttt{wait}/\texttt{post} share the same bodies; binary is the special
case whose value stays in $\{0,1\}$ when used as a lock.

(a)~Binary wait: \texttt{10, 1, 8, 6, 11}
\quad (lock; while $s\le 0$; wait; $s$\texttt{--}; unlock).
Equally valid: \texttt{10, 5, 8, 6, 11} (while $s==0$), since a correct
binary semaphore never stores a negative $s$ in this implementation.

(b)~Binary post: \texttt{10, 7, 9, 11}
\quad (lock; $s$\texttt{++}; signal; unlock).
A \emph{strict} binary post that refuses to increment past 1 can set
$s=1$ instead of $s$\texttt{++}; with the given chunks the increment is
what is available, and is correct if waiters pair with posts.

(c)~Generalised wait: \texttt{10, 1, 8, 6, 11}
\quad (must use \texttt{while (s <= 0)}, not \texttt{s == 0}, because
$s$ may be any non-negative integer).

(d)~Generalised post: \texttt{10, 7, 9, 11}.

\texttt{while} (not a one-shot \texttt{if}) is required under Mesa
semantics; omitting chunk 1/5, or signalling without the lock held, is
incorrect. \texttt{sem\_post} never waits, so chunks 2--4 do not appear
in post.
\end{Sol}

\begin{Qn}{Medium-Hard}{Lect 19}{CS-537 Fall 2004 Midterm 1 Q6}{The WhatsItFor: inverted semaphore as a lock}
A WhatsItFor has an initial value and two atomic routines.
\texttt{One()} waits until the value is \emph{less than zero}, then
increments it by one. \texttt{Done()} decrements the value by one, then
wakes one waiting thread (if any).

Show how to use a WhatsItFor to build a lock around a critical section.
Specify the initial value.
\end{Qn}
\begin{Sol}
The WhatsItFor is a semaphore with the sense of the value flipped:
\texttt{One} is ``wait until a deficit exists, then consume one unit of
deficit''; \texttt{Done} is ``create one unit of deficit and wake a
waiter''.

Init the WhatsItFor to $-1$ (already in deficit, so the first
\texttt{One} succeeds immediately). Acquire is \texttt{One()}; release
is \texttt{Done()}.

Trace. Init $-1$. T1 \texttt{One}: value $< 0$, increments to $0$,
enters CS. T2 \texttt{One}: value is $0$, not $< 0$, sleeps. T1
\texttt{Done}: decrements to $-1$, wakes T2. T2 increments to $0$ and
enters. Mutual exclusion holds; the lock is free iff the value is
$-1$.

(Init $0$ does \emph{not} work: the first \texttt{One} waits forever
because the value is never made negative. Pairing \texttt{Done} then
\texttt{One} for acquire with init $0$ also fails to serialise two
threads.)
\end{Sol}

\begin{Qn}{Medium-Hard}{Lect 19}{Mandar Mitra, ISI Kolkata, 2010 End-sem Q5(b)(i)}{Bounded buffer with a \emph{single} semaphore}
A producer and a consumer share an $n$-slot circular buffer and a counter
(\texttt{in}, \texttt{out}, \texttt{count}):
\begin{lstlisting}[numbers=none]
producer                         consumer
p = produce();                   while (count == 0);   /* empty */
while (count == n);              p = buffer[out];
buffer[in] = p;                  out = (out+1) % n;
in = (in+1) % n;                 count--;
count++;                         consume(p);
\end{lstlisting}
They must be synchronised using a \emph{single} semaphore $S$. Insert
\texttt{wait(S)} and \texttt{signal(S)} so the processes are properly
synchronised. (Do \emph{not} use \texttt{semget}/\texttt{semctl}/
\texttt{semop}.)
\end{Qn}
\begin{Sol}
$S$ is a mutex, init 1. The original busy-waits on \texttt{count} are
racy (and deadlock if they spin \emph{holding} $S$). Release $S$ before
waiting for a slot / item, then re-check under $S$:
\begin{lstlisting}[numbers=none]
producer:                            consumer:
  p = produce();                       wait(S);
  wait(S);                             while (count == 0) {
  while (count == n) {                     signal(S);
      signal(S);                           wait(S);
      wait(S);                         }
  }                                    p = buffer[out];
  buffer[in] = p;                      out = (out+1) % n;
  in = (in+1) % n;                     count--;
  count++;                             signal(S);
  signal(S);                           consume(p);
\end{lstlisting}
\texttt{count}, \texttt{in}, \texttt{out}, and the slot they name are
touched only with $S$ held, so lost updates are impossible. A full
producer drops $S$ before retrying so a consumer can run (and vice versa):
no deadlock. This recheck loop works for both one and multiple producers/consumers,
though it busy-spins while full/empty (which is why the three-semaphore
solution of Lect~19 with \texttt{empty}/\texttt{full} is preferred).
(Note: for strictly \emph{one} producer and \emph{one} consumer where
\texttt{in} is private to the producer and \texttt{out} is private to the
consumer, the minimal Mandar fix is simply wrapping \texttt{count++} and
\texttt{count--} in \texttt{wait(S)}/\texttt{signal(S)}.)
\end{Sol}

\begin{Qn}{Medium-Hard}{Lect 19}{Mandar Mitra, ISI Kolkata, 2015 End-sem Q1; 2024 End-sem Q2}{Blocking binary wait/signal, then with TAS / swap}
A binary semaphore's integer is only 0 or 1.

\begin{enumerate}
  \item[(a)] Write pseudo-code for the \emph{blocking} \texttt{wait} and
        \texttt{signal} on a binary semaphore (not the counting version).
  \item[(b)] Implement those operations using an atomic
        \texttt{swap(int *x, int *y)} (2015) or
        \texttt{TestAndSet} (2024). You may introduce a spinlock that
        protects the semaphore's internals, plus a wait queue.
\end{enumerate}
\end{Qn}
\begin{Sol}
(a)~Value $\in \{0,1\}$; increment/decrement operators therefore do not
apply. Sleep a waiter rather than spinning.
\begin{lstlisting}[numbers=none]
wait(S):                             signal(S):
  if (S.value == 1)                    if (S.queue is nonempty)
      S.value = 0;                         wakeup(dequeue(S.queue));
  else {                               else
      enqueue(self, S.queue);              S.value = 1;
      sleep();
  }
\end{lstlisting}
Invariant: $S.value = 1$ iff nobody is in the CS and the queue is
empty. A \texttt{signal} with waiters hands the resource directly to
the woken thread (value stays 0) --- the lock-handoff style of Lect 17
queue locks.

(b)~Protect the bodies with a TAS/swap spinlock \texttt{S.guard}
(init 0 $=$ free). Use a \emph{local} temporary \texttt{t = 1} for
\texttt{swap} (never a shared global, which \texttt{swap} would overwrite):
\begin{lstlisting}[numbers=none]
/* Swap version (2015): */
void lock_guard(S)   { int t = 1; do { swap(&S.guard, &t); } while (t == 1); }
/* TestAndSet version (2024): */
/* void lock_guard(S) { while (TestAndSet(&S.guard) == 1) ; } */
void unlock_guard(S) { S.guard = 0; }   /* plain store; or local int z=0; swap(&S.guard, &z) */

wait(S):
  lock_guard(S);
  if (S.value == 1) { S.value = 0; unlock_guard(S); }
  else { enqueue(self); setpark(); unlock_guard(S); park(); }

signal(S):
  lock_guard(S);
  if (queue nonempty) { unpark(dequeue()); unlock_guard(S); }
  else { S.value = 1; unlock_guard(S); }
\end{lstlisting}
The guard is released \emph{before} sleeping, matching Lecture~17. Calling
\texttt{setpark()} before \texttt{unlock\_guard} (or assuming
\texttt{wakeup}/\texttt{unpark} remembers a pending wakeup if it runs
before \texttt{park}/\texttt{sleep}) closes the lost-wakeup window between
dropping \texttt{S.guard} and going to sleep. (Note: a simple TAS busy-wait
lock uses \texttt{0 = free, 1 = held} and spins, which is \emph{not} what
the blocking binary semaphore in part~(a) asked for.)
\end{Sol}

\begin{Qn}{Medium-Hard}{Lect 19}{CS-537 Spring 2001 Midterm, Part II Q1}{Binary semaphore from a Hoare (then Mesa) monitor}
You have a language that provides only monitors. Implement a binary
semaphore as a monitor class \texttt{BinarySemaphore} with \texttt{P()}
and \texttt{V()}. You have a condition class with \texttt{wait()},
\texttt{signal()}, \texttt{broadcast()}. Use \emph{Hoare} semantics first.

\begin{enumerate}
  \item[(a)] Variables and initial values.
  \item[(b)] Pseudo-code for \texttt{P()} and \texttt{V()}.
  \item[(c)] What must change under Mesa semantics? What goes wrong if you
        do not change it?
\end{enumerate}
\end{Qn}
\begin{Sol}
(a)~An integer \texttt{value = 1} (free) and a condition variable
\texttt{c} (the waiter queue). (A boolean \texttt{busy = 0} is
equivalent.)

(b)~Hoare: the signalled waiter runs \emph{immediately}, before the
signaller continues, and the waiter is the first to re-occupy the
monitor. An \texttt{if} is therefore sufficient:
\begin{lstlisting}[numbers=none]
monitor BinarySemaphore {
    int value = 1;
    condition c;
    void P() { if (value == 0) c.wait();  value = 0; }
    void V() { value = 1;  c.signal(); }
}
\end{lstlisting}
If a waiter exists, Hoare \texttt{V} hands it the monitor with
\texttt{value} already 1; the waiter then sets \texttt{value = 0} before
anyone else can sneak in. If nobody waits, \texttt{value} stays 1.

(c)~Mesa: \texttt{signal} is only a hint; the signaller keeps the
monitor and the waiter later reacquires it, possibly after another
thread has consumed the resource. Replace \texttt{if} by
\texttt{while (value == 0) c.wait();}. If you leave the \texttt{if}, a woken waiter skips the re-check even if a
third thread entered the monitor in the meantime and consumed
\texttt{value} (setting it back to 0), so two threads end up inside the
critical section simultaneously. \texttt{broadcast} is unnecessary for a
binary semaphore (at most one waiter needs to run per \texttt{V}).
\end{Sol}

\begin{Qn}{Medium-Hard}{Lect 19}{Mandar Mitra, ISI Kolkata, 2014 End-sem Q1(a)}{Bounded waiting: busy-wait vs blocking mutex}
$n$ processes run
\texttt{\{ ... wait(mutex); critical\_section(); signal(mutex); ... \}}
with \texttt{mutex} a semaphore init 1. Explain how this does or does
not satisfy \emph{bounded waiting} if the semaphore is implemented with
(i)~busy-waiting, (ii)~blocking. Refer to the implementation of
\texttt{wait}/\texttt{signal}.
\end{Qn}
\begin{Sol}
Bounded waiting: there exists a bound $k$ such that after a process $P$
has requested entry, at most $k$ other entries occur before $P$ enters.

(i)~\emph{Busy-wait.}
\texttt{wait}: \texttt{while (S <= 0) ; S--;}\quad
\texttt{signal}: \texttt{S++;}\quad
(atomic). There is no waiter queue. Who next observes $S=1$ is entirely
up to the scheduler / memory arbiter. A process that has already
started spinning can be skipped by an arbitrary number of luckier
processes that see $S=1$ first. No bound $k$ exists; bounded waiting
fails.

(ii)~\emph{Blocking.}
\texttt{wait}: \texttt{S--; if (S < 0) \{ enqueue(self); sleep(); \}}\quad
\texttt{signal}: \texttt{S++; if (S <= 0) wakeup(dequeue());}.
Waiters form an explicit queue. If that queue is FIFO (the usual OS
implementation, and the Lect 19 picture ``negative value $=$ number of
waiters''), a process that has slept is behind at most the waiters
already queued --- at most $n-1$ entries --- before its wakeup. Bounded
waiting holds (with $k = n-1$). (A LIFO or priority sleep queue would
break it; the argument needs the FIFO policy.)
\end{Sol}

\subsection{Lecture 20: Concurrency Bugs \& Deadlock}

\begin{Qn}{Medium-Hard}{Lect 20}{CS-537 Fall 2013 Midterm, Q16}{Four threads, locks $A,B,C,D$ ``in some order''}
``I wrote some code and it keeps deadlocking. How frustrating! Here is how
the code works. Thread~1 grabs Locks $A$ and $B$ (in some order); Thread~2
grabs Locks $B$ and $C$ (in some order); Thread~3 grabs Locks $C$ and $D$
(in some order); Thread~4 grabs Locks $D$ and $A$ (in some order). Why in
the world does this code deadlock? Somebody help me; I just find
concurrency hard.''
\end{Qn}
\begin{Sol}
Official CS-537 Fall 2013 answer:

Assume Thread~1 grabs $A$ but then is waiting for $B$; Thread~2 grabs $B$
and is waiting for $C$; Thread~3 grabs $C$ but is waiting for $D$;
Thread~4 grabs $D$ but is waiting for $A$. Thus a cycle and a deadlock.

``Solve in many possible ways. A good way is to order the lock
acquisitions: always grab $A$ before $B$ before $C$ before $D$. Doing so
would prevent deadlock.''

That global order breaks circular wait (Coffman condition~4). Even though
each thread still takes a pair, the wait-for graph can only have edges
toward a later letter, so it is acyclic. (Thread~4 must take $A$ then
$D$, never $D$ then $A$, which is exactly the edge that closed the cycle.)

Other legal preventions (the official key's ``many possible ways''):
\texttt{trylock} on the second lock of the pair, dropping the first on
failure (break no-preemption; livelock unless you back off); acquire both
locks of the pair atomically under a meta-lock (break hold-and-wait);
lock-free / CAS structures (break mutual exclusion). Encapsulation makes
the meta-lock and the global order harder if the pairs are hidden inside
library calls --- which is why the ``in some order'' wording is so
dangerous.
\end{Sol}

\begin{Qn}{Medium-Hard}{Lect 20}{CS-537 Spring 2002 Midterm, Q8}{Four Coffman conditions, and the problems of lock preemption}
For deadlock to occur, four conditions must hold: mutual exclusion,
hold-and-wait, no preemption, and circular wait. If any one condition does
not hold, no deadlock can occur. Assume we want to allow ``preemption'',
and thus get out of deadlocks; in other words, if a deadlock is detected,
we will forcibly take a lock away from a thread; by repeatedly doing this,
we will eventually undo the deadlock. What new problems are introduced by
this preemptive approach?
\end{Qn}
\begin{Sol}
The four Coffman conditions (all necessary) are exactly Lect~20:
\begin{enumerate}
  \item \textbf{Mutual exclusion} --- the resource is not shareable.
  \item \textbf{Hold-and-wait} --- a thread holds at least one resource
        while waiting for another.
  \item \textbf{No preemption} --- a held resource is released only
        voluntarily.
  \item \textbf{Circular wait} --- a cycle in the wait-for graph.
\end{enumerate}
Breaking (3) by \emph{forcibly stealing} a lock from a thread that is
already in its critical section introduces a new, worse problem:
\emph{the stolen lock's invariant is gone}. The victim may have been
mid-update (half-linked list, half-transferred money). The thief then
enters the same critical section on corrupt state. You now need rollback
/ transactions just to make preemption safe --- far more mechanism than a
lock.

Further problems: the victim does not know it lost the lock and will
later \texttt{unlock} a lock it no longer owns (or double-unlock);
choosing \emph{which} lock to steal is policy; repeated theft can starve
a thread.

Lect~20's practical attack on (3) is therefore \emph{not} forcible theft.
It is \texttt{trylock()}: the \emph{would-be waiter} refuses to sit in
hold-and-wait. On failure it drops what it already holds and retries
(with exponential back-off, to avoid livelock). The holder of the
contended lock is never preempted; it finishes its CS and releases
voluntarily. That is prevention, not recovery.
\end{Sol}

\begin{Qn}{Medium-Hard}{Lect 20}{Mandar Mitra, ISI Kolkata, 2011 End-sem Q2(a); 2012 End-sem Q1(b)}{Prevention by breaking hold-and-wait, and by breaking circular wait}
\begin{enumerate}
  \item[(a)] [2011 End Q2(a)] Briefly describe a deadlock prevention
        approach which ensures that the hold-and-wait condition is never
        fulfilled in a system. Clearly explain the disadvantages of such
        an approach.
  \item[(b)] [2012 End Q1(b)] Briefly describe and justify a deadlock
        prevention approach which ensures that the circular-wait
        condition is never fulfilled in a system.
\end{enumerate}
\end{Qn}
\begin{Sol}
(a)~Break hold-and-wait. Two equivalent statements of the same protocol:
\begin{itemize}
  \item A process must request \emph{all} resources it will ever need
        in a single atomic act, before it holds any of them. If any
        request cannot be granted, it holds nothing and waits.
  \item Lect~20 code form: a process-wide \emph{meta-lock}. To acquire
        a set of locks, take the meta-lock, take every lock in the set,
        drop the meta-lock.
\end{itemize}
Disadvantages (Mandar + Lect~20): (i)~you must know the whole set in
advance --- encapsulation hides locks inside library calls
(\texttt{set\_intersection}, \texttt{vector\_add}); (ii)~utilization
collapses (resources sit reserved for a process that is not using them
yet); (iii)~the meta-lock serializes \emph{all} multi-lock acquires, so
the system degenerates toward one big lock; (iv)~a process that
periodically needs a surprising extra resource must drop everything and
start over (starvation / convoying).

(b)~Break circular wait with a \emph{total order} on resource types
(lock IDs, or pointer addresses). A process may request $R_j$ only if
it currently holds no $R_i$ with $i \ge j$ (equivalently: always acquire
in increasing ID / increasing address). Justification: every wait-for
edge goes from a held lower-ID resource to a wanted higher-ID resource,
so the wait-for graph is a DAG and cannot contain a cycle. This is the
same idea as the $A < B < C < D$ order in CS-537 Fall 2013 Q16 and the
address-order fix of \texttt{vector\_add}.
\end{Sol}

\clearpage

\section{Lectures 21--24 (incl.\ 22A): Persistence --- I/O Devices, HDDs, RAID \& File Systems}

\subsection{Lectures 21, 22 \& 22A: I/O Devices, Hard Disk Drives \& RAID}

\begin{Qn}{Hard}{Lect 22A}{CS-537 Fa13 Final, Q2}{Worst-case RAID-0/1/4 time with real disk parameters}
A RAID has \textbf{4} disk drives. Each drive has peak transfer
$100$\,MB/s, worst-case seek $9$\,ms, and $15{,}000$\,RPM. Chunk size is
$4$\,KB.
\begin{enumerate}[(a)]
\item RAID-0 (striping). Given four $4$\,KB reads, pick the locations
      (\texttt{[disk, offset]} pairs) to get worst-case performance.
\item How long would it take to complete these four reads?
\item What is the best-case time to complete four reads on this RAID-0?
      State any assumptions.
\item Now a mirrored (RAID-1) system. What is the worst-case time a
      single $4$\,KB write takes?
\item How long will sixteen $4$\,KB writes take on this 4-disk RAID-1,
      assuming the writes are spread randomly across the RAID?
\item RAID-4 (one of the four disks used for parity). When a single-block
      write takes place, how many I/Os take place on the RAID in the
      worst case?
\item How long will such a ``small write'' take, assuming the disk
      parameters above and a worst-case seek?
\end{enumerate}
\end{Qn}
\begin{Sol}
Full rotation $= 60/15{,}000 = 4$\,ms. Transfer of $4$\,KB at
$100$\,MB/s $= 0.04$\,ms. Average seek $\approx 3$\,ms, average rotation
$= 2$\,ms.

(a) Place all four reads on the \textbf{same} disk at distinct sectors
on opposite ends of the stroke (avoiding track-cache hits), e.g.\
\texttt{[0,0]}, \texttt{[0, last]}, \texttt{[0,1]}, \texttt{[0, last-1]}.
Striping cannot help: they serialize on one arm.

(b) $4\times(9+4+0.04) = \mathbf{52.16}$\,ms.

(c) One $4$\,KB read on each of the four disks, heads already on-track
and on-sector: just transfer, in parallel, $\mathbf{0.04}$\,ms.
(If a positioning delay is required, the four disks still proceed in
parallel: $9+4+0.04 = 13.04$\,ms.)

(d) A RAID-1 write must update \textbf{both} copies (consistent-update).
The two writes issue in parallel, so worst-case latency is still one
positioning: $9+4+0.04 = \mathbf{13.04}$\,ms.

(e) Four disks $\Rightarrow$ two mirrored pairs. Random writes spread
evenly: $16/2 = 8$ writes per pair, each pair acting as one logical disk
for writes. Assuming standard average seek $\approx \tfrac13\times 9 = 3$\,ms
and average rotation $2$\,ms,
$8\times(3+2+0.04) = \mathbf{40.32}$\,ms
(under worst-case positioning instead: $8\times 13.04 = 104.32$\,ms; note
also that a mirrored write waits for the slower of the two arms).
(Throughput formula: random-write bandwidth is $(N/2)\,R_{\mathrm{disk}}$, here $2R_{\mathrm{disk}}$.)

(f) Subtractive small write: read old data, read old parity, write new
data, write new parity --- \textbf{4 I/Os} (2 reads then 2 writes).
Additive reconstruction of parity (read the other two data chunks of
the stripe) is also 4 I/Os on this 4-disk array. Worst case $= 4$.

(g) The two reads proceed in parallel (worst seek $+$ rotation $+$
transfer $= 9+4+0.04 = 13.04$\,ms). To rewrite the \emph{same} sector,
the platter rotates one full revolution back to the end of that sector
(waiting $4-0.04$\,ms then transferring $0.04$\,ms, i.e.\ exactly
$4$\,ms total):
\[
(9+4+0.04) + 4 = \mathbf{17.04}\,\mathrm{ms}.
\]
\end{Sol}

\begin{Qn}{Hard}{Lect 22}{CS-537 Fa11 Final, Q3}{FIFO versus SATF on a 1-track and a 3-track disk}
A highly simplified disk: rotational delay $R$, transfer free.
\begin{enumerate}[(a)]
\item Single track, FIFO. Approximate worst-case time for three requests
      to different blocks.
\item Same disk, SATF scheduler. Worst-case time for three requests.
\item Now three tracks. Seek between adjacent tracks costs $S$; seeking
      across two tracks costs $2S$. FIFO, worst-case time for three
      requests.
\item Same three-track disk, SATF. Worst-case time for three requests.
\end{enumerate}
\end{Qn}
\begin{Sol}
Official key (CS-537 Fa11 Final answers, p.~6):

(a) $\boldsymbol{\sim 3R}$. FIFO can be forced to miss each requested
sector just after it passes, paying nearly a full rotation per request.

(b) $\boldsymbol{\sim R}$. SATF services the three requests in
rotational order; one revolution collects all of them.

(c) A single request's worst case is a two-track seek plus a full
rotation: $2S+R$. FIFO can be made to repeat that independently three
times (outer $\to$ inner $\to$ outer): $\boldsymbol{6S+3R}$.

(d) SATF can collect the three requests in roughly one sweep across the
disk (outer to inner, or vice versa): $\boldsymbol{\sim 2S+3R}$
(might still wait $\sim R$ at each stop).
\end{Sol}

\begin{Qn}{Hard}{Lect 22A}{CS-537 Sp11 Final, Q4}{Unveiling the RAID box from bread/bwrite traces}
All systems have 4 disks. Inputs are \texttt{read(blk)}/\texttt{write(blk)};
outputs are low-level \texttt{bread}/\texttt{bwrite} to a specific disk.
Possible answers: RAID-0, RAID-1, RAID-4, RAID-5, or more than one of
these.
\begin{enumerate}[(a)]
\item
\begin{lstlisting}[numbers=none]
write(0) -> [ bwrite(disk=0, block=0), bwrite(disk=1, block=0) ] in parallel
write(1) -> [ bwrite(disk=2, block=0), bwrite(disk=3, block=0) ] in parallel
write(2) -> [ bwrite(disk=0, block=1), bwrite(disk=1, block=1) ] in parallel
write(3) -> [ bwrite(disk=2, block=1), bwrite(disk=3, block=1) ] in parallel
\end{lstlisting}
\item \texttt{read(0) -> [ bread(disk=0, block=0) ]}
\item
\begin{lstlisting}[numbers=none]
write(0) -> [ bread(disk=0, block=0), bread(disk=3, block=0) ] in parallel,
            then [ bwrite(disk=0, block=0), bwrite(disk=3, block=0) ]
write(3) -> [ bread(disk=0, block=1), bread(disk=3, block=1) ] in parallel,
            then [ bwrite(disk=0, block=1), bwrite(disk=3, block=1) ]
\end{lstlisting}
\item
\begin{lstlisting}[numbers=none]
read(0) -> [ bread(disk=0, block=0) ]
read(1) -> [ bread(disk=1, block=0) ]
read(2) -> [ bread(disk=0, block=0), bread(disk=1, block=0),
             bread(disk=3, block=0) ] in parallel
read(3) -> [ bread(disk=0, block=1) ]
read(4) -> [ bread(disk=1, block=1) ]
read(5) -> [ bread(disk=0, block=1), bread(disk=1, block=1),
             bread(disk=3, block=1) ] in parallel
\end{lstlisting}
\end{enumerate}
\end{Qn}
\begin{Sol}
(a) \textbf{RAID-1} (two 2-way mirrored pairs, striped across the
pairs). Each logical write updates both copies of one pair; consecutive
logical blocks alternate pairs. Not RAID-4/5 (no RMW, no parity disk)
and not RAID-0 (two writes per logical request).

(b) A single data-disk read is possible for RAID-0, RAID-1 (read one
mirror), RAID-4, and RAID-5. \textbf{Cannot distinguish}: any of
RAID-0/1/4/5.

(c) \textbf{RAID-4}, subtractive small-write. Data of logical 0 lives
on disk 0 offset 0; parity of that stripe lives on disk 3 (the dedicated
parity disk). Logical 3 is the next stripe on the same data disk
(offset 1), still RMW-ing disk 3. RAID-5 would have rotated the parity
off disk 3 for the second stripe.

(d) \textbf{RAID-4, degraded}: disk 2 has failed. Logical layout is
3 data $+$ 1 parity (disk 3). Logical 2 and 5 live on disk 2, so they
are reconstructed by XOR of the other data disks and the parity disk.
(RAID-5 would have placed logical 5 on a different disk because parity
rotates; the reconstruction pattern here is RAID-4.)
\end{Sol}

\begin{Qn}{Hard}{Lect 22A}{CS-537 Sp09 Final, Q9}{Software RAID-1 read/write and RAID-4 reconstruct-read}
A software RAID exports \texttt{RAIDRead()}/\texttt{RAIDWrite()} on top
of \texttt{fliskRead(flisk, buf, offset)} / \texttt{fliskWrite(...)},
each returning $0$ on success and $-1$ on failure. Each flisk has
$1$ million $4$\,KB blocks; every RAID request is one block.
\begin{enumerate}[(a)]
\item \textbf{RAID-1} (two flisks). Write
      \texttt{RAIDWrite(char *buffer, int offset)} so that it returns
      $0$ iff the data was committed to \textbf{both} disks.
\item Write \texttt{RAIDRead()} for RAID-1 (success if the data can be
      obtained from either copy).
\item Now \textbf{RAID-4} with 3 flisks (two data, one parity). Write
      \texttt{RAIDRead()} that still succeeds when a \emph{data} disk
      fails, reconstructing via XOR. An \texttt{xor()} helper is
      available.
\end{enumerate}
\end{Qn}
\begin{Sol}
(a) Both mirrors must be updated (consistent-update problem; hardware
RAID would log this in NVRAM). Return success only if both writes
succeed:
\begin{lstlisting}[numbers=none]
int RAIDWrite(char *buffer, int offset) {
    int r0 = fliskWrite(0, buffer, offset);
    int r1 = fliskWrite(1, buffer, offset);
    return (r0 == 0 && r1 == 0) ? 0 : -1;
}
\end{lstlisting}

(b) Try one copy, fall back to the other:
\begin{lstlisting}[numbers=none]
int RAIDRead(char *buffer, int offset) {
    if (fliskRead(0, buffer, offset) == 0) return 0;
    if (fliskRead(1, buffer, offset) == 0) return 0;
    return -1;
}
\end{lstlisting}

(c) Mapping with 2 data disks: $\mathrm{disk}=\texttt{offset}\bmod 2$,
$\mathrm{off}=\lfloor\texttt{offset}/2\rfloor$, parity on flisk 2 at
the same \texttt{off}. On a data-disk failure, $D = D_{\mathrm{other}}
\oplus P$:
\begin{lstlisting}[numbers=none]
int RAIDRead(char *buffer, int offset) {
    int disk = offset % 2, other = 1 - disk, off = offset / 2;
    if (fliskRead(disk, buffer, off) == 0) return 0;
    char otherbuf[4096], parity[4096];
    if (fliskRead(other, otherbuf, off) != 0) return -1;
    if (fliskRead(2,     parity,   off) != 0) return -1;
    xor(buffer, otherbuf, parity);  /* buffer = other XOR parity */
    return 0;
}
\end{lstlisting}
\end{Sol}

\begin{Qn}{Hard}{Lect 22A}{CS-537 Fa08 Final, Q9}{RAID-4 I/O flows and RAID-5 placement for two small writes}
RAID-4 with a single parity disk:
\begin{lstlisting}[numbers=none]
 0   1   2   3   P0
 4   5   6   7   P1
 8   9  10  11   P2
12  13  14  15   P3
\end{lstlisting}
Notation: \texttt{(r0 r1)} means those reads issue in parallel;
\texttt{r0, r1} means they are sequential. Write the I/O flow.
\begin{enumerate}[(a)]
\item Read blocks 0, 5, 10, and 15 as fast as possible.
\item Read blocks 100, 109, 126, and 85 as fast as possible.
\item Read blocks 0, 4, 8, and 12.
\item Write block 5.
\item Write blocks 5 and 10.
\item Switch to RAID-5 (rotating parity). Arrange data and parity so
      that the writes to 5 and 10 go as fast as possible. Draw the
      layout.
\end{enumerate}
\end{Qn}
\begin{Sol}
Four data disks: $\mathrm{disk}=B \bmod 4$,
$\mathrm{offset}=\lfloor B/4\rfloor$, parity always on disk 4.

(a) $0\to\mathrm{d0}$, $5\to\mathrm{d1}$, $10\to\mathrm{d2}$,
$15\to\mathrm{d3}$: four different disks.
\textbf{\texttt{(r0 r5 r10 r15)}}.

(b) $100\to\mathrm{d0}$, $109\to\mathrm{d1}$, $126\to\mathrm{d2}$,
$85\to\mathrm{d1}$. 109 and 85 collide on disk 1:
\textbf{\texttt{(r100 r109 r126), r85}}
(or \texttt{(r100 r85 r126), r109}).

(c) All four live on disk 0 (same column). Must serialize:
\textbf{\texttt{r0, r4, r8, r12}}.

(d) Subtractive small write of 5 (stripe with P1):
\textbf{\texttt{(r5 rP1), (w5 wP1)}}.

(e) 5 uses P1 and 10 uses P2, but \textbf{both parities sit on the same
dedicated parity disk}, so the two small writes cannot proceed in
parallel:
\textbf{\texttt{(r5 rP1), (w5 wP1), (r10 rP2), (w10 wP2)}}.
This is the RAID-4 parity-disk bottleneck.

(f) In RAID-5, rotate parity across stripes so that all four blocks
involved in the two small writes --- \texttt{5}, \texttt{P1},
\texttt{10}, and \texttt{P2} --- reside on \textbf{four distinct disks}.
(Note: in standard left-asymmetric rotation \texttt{P1} and \texttt{10}
both land on d3, and in left-symmetric rotation \texttt{5} and
\texttt{10} both land on d0!) One conflict-free RAID-5 placement
puts \texttt{5} on d1, \texttt{P1} on d4, \texttt{10} on d3, and
\texttt{P2} on d0:
\begin{lstlisting}[numbers=none]
d0   d1   d2   d3   d4
 0    1    2    P0   3
 4    5    6    7    P1     /* 5 on d1, P1 on d4 */
P2    8    9   10   11     /* P2 on d0, 10 on d3 */
\end{lstlisting}
Now \texttt{(r5 rP1 r10 rP2), (w5 wP1 w10 wP2)} runs both small writes
completely in parallel across four distinct disks. RAID-5 does
\emph{not} remove the 4-I/O small-write cost, it only removes the
dedicated parity-disk bottleneck.
\end{Sol}

\subsection{Lectures 23--24: File System API \& Implementation}

\begin{Qn}{Hard}{Lect 23--24}{CS-537 Fa11 Final, Q5}{Open/read/close times under a 12$+$1 inode}
Simple disk model: each block I/O takes $D$ time units. Layout is VSFS-like.
All data and metadata begin on disk. Every inode lives in its own block, and
every directory is one block.
\begin{enumerate}[(a)]
\item How long does it take to \texttt{open} the file \texttt{/a/b/c/d.txt}?
\item After opening, the file is read in its entirety. It contains 1036
      blocks. The inode has 12 direct pointers and 1 indirect pointer.
      Disk addresses are 4 bytes; disk blocks are 4\,KB. After opening,
      how long does it take to read the entire file?
\item A new inode has only one pointer: a double-indirect pointer (to a
      double-indirect block of 1024 indirect blocks, each of 1024 data
      pointers). After opening, how long do 50 random reads in a very
      large file take?
\item How long does it take to \texttt{close} a file, approximately
      (disk accesses dominate)?
\end{enumerate}
\end{Qn}
\begin{Sol}
Pointers per block $= 4096/4 = 1024$.
\begin{enumerate}[(a)]
\item Path walk of \texttt{/a/b/c/d.txt}: five inodes (root, \texttt{a},
      \texttt{b}, \texttt{c}, \texttt{d.txt}) and four directory data
      blocks. Each is a separate disk block, so
      \[
        9D.
      \]
\item $12 + 1024 = 1036$: the file uses every direct pointer and a
      \emph{full} single-indirect block. The inode is already in memory
      from \texttt{open}. Remaining I/O: 1 indirect block $+$ 1036 data
      blocks, so
      \[
        1037D.
      \]
\item The double-indirect pointer is in the in-memory inode. Each random
      read, with no extra caching of index blocks, walks
      double-indirect $\to$ indirect $\to$ data:
      \[
        50 \times 3D = 150D.
      \]
      If the double-indirect block remains cached after the first read,
      the bound drops to $D + 50\times 2D = 101D$ (still $3D$ on the
      first sample).
\item \texttt{close} only frees the in-memory fd / open-file object.
      \textbf{0} disk I/O, so time $\approx 0$.
\end{enumerate}
\end{Sol}

\begin{Qn}{Hard}{Lect 24}{CS-537 Sp10 Final, P1}{Unix inode max size, then probe it with \texttt{lseek}}
A typical Unix inode has 12 direct pointers, one indirect pointer, and one
double-indirect pointer. Block size 4\,KB; disk addresses 32 bits.
\begin{enumerate}[(a)]
\item What is the maximum file size? Show the arithmetic.
\item Write the most efficient program that discovers the maximum file
      size on the file system you are actually running, using only
      \texttt{open}, \texttt{read}, \texttt{write}, \texttt{close}, and
      \texttt{lseek}.
\end{enumerate}
\end{Qn}
\begin{Sol}
\begin{enumerate}[(a)]
\item Pointers per block $= 4096/4 = 1024$. Data blocks addressable:
      \[
        12 + 1024 + 1024^{2} = 1\,049\,612.
      \]
      \[
        1\,049\,612 \times 4096 = 4\,299\,210\,752\ \text{bytes}
        \approx 4.004\,\mathrm{GiB}.
      \]
      (Note: if the OS also imposed a 32-bit \emph{unsigned} offset/size
      field, $2^{32}-1 = 4\,294\,967\,295$ would be the \emph{tighter}
      bound, and a signed 32-bit \texttt{off\_t} would cap it at
      $2^{31}-1 = 2\,147\,483\,647$; since no offset width is specified
      here, the TOC figure above is the intended answer.)
      Same template with 8\,KB blocks (CS-537 Sp02 Final, Q6(a)):
      $N=2048$, max $= (12+2048+2048^{2})\times 8192 = 34\,376\,613\,888$
      bytes.
\item Use exponential search followed by binary search with sparse
      \texttt{lseek} $+$ 1-byte \texttt{write} (writing a hole allocates
      only 1 data block plus at most the index blocks along its path,
      not the entire file; \texttt{lseek} itself does no disk I/O):
\begin{lstlisting}[numbers=none]
int fd = open("tmp", O_CREAT | O_RDWR | O_TRUNC, 0600);
off_t lo = 0, hi = 1;
char b = 0;
while (lseek(fd, hi, SEEK_SET) == hi && write(fd, &b, 1) == 1) {
    lo = hi; hi *= 2;              /* exponential probe */
}
while (lo + 1 < hi) {
    off_t mid = lo + (hi - lo) / 2;
    if (lseek(fd, mid, SEEK_SET) == mid && write(fd, &b, 1) == 1)
        lo = mid;
    else
        hi = mid;
}
close(fd);                         /* max file size = lo + 1 bytes */
\end{lstlisting}
\end{enumerate}
\end{Sol}

\begin{Qn}{Hard}{Lect 24}{CS-537 Sp16 Final, Q2}{directfs, extentfs, indirectfs, compactfs (linked)}
File-system blocks are 4\,KB throughout.
\begin{enumerate}[(a)]
\item \texttt{directfs}: inode has 10 direct pointers (32-bit). Max file
      size?
\item \texttt{extentfs}: inode holds exactly one extent
      (base pointer $+$ 8-bit length in blocks). Max file size?
\item \texttt{indirectfs}: inode has 1 direct, 1 indirect, and 1
      double-indirect pointer (4-byte addresses). Max file size
      (equation OK)?
\item \texttt{compactfs}: inode stores only a pointer to the first
      block; each block holds 4\,KB of user data plus a \texttt{next}
      field (linked list). Draw an inode and a 12\,KB file.
\item Max file size for \texttt{compactfs}?
\end{enumerate}
\end{Qn}
\begin{Sol}
Official (OCR) answers: (a) $10\times 4\,\mathrm{KB}=40\,\mathrm{KB}$;
(b) $2^{8}=256$ blocks $\Rightarrow 256\times 4\,\mathrm{KB}=1\,\mathrm{MB}$;
(c) $(1+1024+1024^{2})\times 4\,\mathrm{KB}$.
\begin{enumerate}[(a)]
\item $10\times 4096 = 40\,\mathrm{KB}$.
\item Length field 8 bits $\Rightarrow$ up to $2^{8}-1 = 255$ blocks
      ($1020\,\mathrm{KB}$) for an unsigned $0..255$ encoding, or
      $256$ blocks ($1\,\mathrm{MB}$) if 1-biased ($\text{length}-1$
      stored, as on the CS-537 key). (Extents give fast sequential/random
      arithmetic, but a short length field caps file size; storing a few
      extents per inode is the standard fix.)
\item $N=1024$, so
      \[
        \bigl(1 + 1024 + 1024^{2}\bigr)\times 4096
        = 4\,299\,165\,696\ \text{bytes}.
      \]
\item Inode $\to$ block$_0$ $\to$ block$_1$ $\to$ block$_2$ $\to$
      \texttt{NULL}. Three data blocks, each 4\,KB, chained by
      \texttt{next}. This is linked allocation (FAT is the same idea
      with the \texttt{next} pointers moved into a table).
\item Unbounded except by disk size (or by the width of
      \texttt{next}: a 32-bit block number yields at most $2^{32}$
      blocks). Random access is $O(\#\text{blocks})$ I/Os --- the
      linked-allocation penalty.
\end{enumerate}
\end{Sol}

\begin{Qn}{Hard}{Lect 24}{ISI Mandar End-sem 2012 Q5(c) [cf.\ 2017E Q4(d)]}{\texttt{lseek} then write/read: count disk I/Os, including holes}
SVR2-like FS, 2\,KB blocks. Process $P$ has opened \texttt{abc.txt}
(size 300 bytes) \texttt{O\_RDWR}. Ignore caching. Count disk blocks
read and written for:
\begin{enumerate}[(i)]
\item seek to offset 1000, write 10 bytes;
\item seek to offset 22\,000, write 10 bytes;
\item seek to offset 9000, read 10 bytes.
\end{enumerate}
\end{Qn}
\begin{Sol}
Logical block $= \lfloor \mathrm{offset}/2048\rfloor$. Directs cover
blocks $0..9$; block 10 is the first SI block. Original file occupies
only logical block 0. \texttt{lseek} itself is memory-only.
Assume parts (i)--(iii) execute sequentially on the same file (so
offset 9000 in (iii) lies inside a hole created by (ii); if run
independently on the 300-byte file, offset 9000 is past EOF and returns
0 bytes with 0 data reads). Also assume the inode is read from disk per
operation (if the open file's in-core inode is already cached, subtract
1 inode read from each part). Count inode $+$ file index/data (superblock
free-list updates sit in memory):
\begin{enumerate}[(i)]
\item Offset 1000 is still block 0. Read inode, read block 0, write
      block 0, write inode (size $1010$). \textbf{$2\mathrm{R}+2\mathrm{W}$.}
\item Offset 22\,000: block $10$, in-block offset $1520$. Allocate a
      single-indirect block and data block 10; blocks $1..9$ remain
      holes. Read inode; write SI, write data 10, write inode (size
      $22\,010$). \textbf{$1\mathrm{R}+3\mathrm{W}$.}
\item Offset 9000: block 4, a hole. A hole read returns zeros
      \emph{without} allocating or reading a data block. Read inode
      only. \textbf{$1\mathrm{R}+0\mathrm{W}$.}
\end{enumerate}
\textbf{Cited clone:} 2017 End Q4(d), 4096\,B blocks, sequential read of
size 276\,520: $\lceil 276520/4096\rceil=68$ data blocks ($10$ direct
$+$ $58$ via the SI), so plus one SI $\Rightarrow$ \textbf{69} file
blocks (70 if the inode is not already resident).
\end{Sol}

\begin{Qn}{Hard}{Lect 24}{ISI Mandar End-sem 2010 Q2(a)}{Inode-number to disk-block offset}
SVR2 file system: disk block size 2048 bytes, inode size 128 bytes, disk
block numbers 4 bytes. The on-disk inode list starts at block 3.
Calculate the exact disk location (block number, byte offset) of inode
number 549.
\end{Qn}
\begin{Sol}
Inodes per block $= 2048/128 = 16$. Unix/SVR2 inode numbers start at 1
(inode 0 is the null dirent), so inode 549 is the 548th inode in the
table (0-based index 548):
\[
  \biggl\lfloor\frac{548}{16}\biggr\rfloor = 34,\qquad
  548 \bmod 16 = 4.
\]
\[
  \text{block } = 3+34 = 37,\qquad
  \text{byte offset } = 4\times 128 = 512.
\]
(If an answer key treated numbers as 0-based, the offset would be
$5\times 128=640$ in the same block; the 1-based convention matches
2014 End Q4(a).)
\end{Sol}

\begin{Qn}{Hard}{Lect 24}{ISI Mandar End-sem 2024 Q5}{Max ordinary files, all in the root}
SVR2-like FS: 512\,B blocks, 128\,B inodes, 64\,B directory entries,
disk $=2^{20}$ blocks plus boot and superblock. $k$ blocks store the
inode list; the remaining $2^{20}-k$ are data. $N$ is the number of
ordinary files (not directories). To maximise $N$, put every file in
the root (no subdirectories).
\begin{enumerate}[(a)]
\item Maximum possible $N$.
\item The corresponding $k$.
\end{enumerate}
\end{Qn}
\begin{Sol}
Empty files need no data blocks. Ignoring indirect-block overhead for the
root directory itself, the data region holds only the root directory
entries. Inodes per block $=512/128=4$; dirents per block
$=512/64=8$. Need $N+1$ inodes (files $+$ root) and $N+2$ dirents
(\texttt{.} , \texttt{..} , plus $N$ names):
\[
  k = \Bigl\lceil\frac{N+1}{4}\Bigr\rceil,\qquad
  d = \Bigl\lceil\frac{N+2}{8}\Bigr\rceil,\qquad
  k+d \le 2^{20}.
\]
The largest $N$ satisfying the inequality is
\[
  N = 2\,796\,199,\qquad k = 699\,050
\]
(check: $699\,050\times 4 = 2\,796\,200 \ge N+1$,
$d=\lceil 2\,796\,201/8\rceil=349\,526$,
$k+d=1\,048\,576=2^{20}$). $N=2\,796\,200$ already overflows.
(Note: if the root directory's own SVR2 indirect blocks --- 1 SI $+$ 1
DI $+$ 128 IBs $+$ 1 TI $+$ 21 DIBs $+$ 2595 IBs $= 2\,747$ index blocks
--- are also subtracted from the data region, the exact maximum drops to
$N = 2\,788\,875$, $k = 697\,219$, since $697\,219 + 348\,610 + 2\,747 = 2^{20}$.)
\end{Sol}

\clearpage

\section{Lectures 1--7 (Pre-Minor): CPU Virtualization (Processes, Traps \& Scheduling)}

\begin{Qn}{Hard}{Lect 3}{Mandar ISI 2022 Mid Q1}{\texttt{fork} loop: break on child, break on parent, or no break}
Consider a process $P$ corresponding to the following program.
\begin{lstlisting}[numbers=none]
int main(int ac, char **av) {
    pid_t p;
    for (int i = 0; i < N; i++) {
        p = fork();
        if (p < 0) exit(1);
        if (p == 0) break;    /* LINE 6 */
    }
    sleep(100);
    return 0;
}
\end{lstlisting}
(a)~Suppose the descendants of $P$ are $P_1,P_2,\ldots,P_n$. Calculate $n$, and draw the initial process hierarchy for $P,P_1,P_2,\ldots$.

(b)~Redo part (a) if the condition on line 6 is changed to \texttt{p > 0}.

(c)~Calculate $n$ if the statement on line 6 is deleted or commented out.
\end{Qn}
\begin{Sol}
Assume every \texttt{fork} succeeds.

(a)~The child takes \texttt{p == 0} and \texttt{break}s, so it never forks. The original parent stays in the loop and forks once per iteration. Result: $n=N$ children, all of $P$. Star:
\begin{center}
$P$ with $N$ children $P_1,\ldots,P_N$ (siblings; none has children).
\end{center}

(b)~Now the \emph{parent} breaks after its first successful fork. The child (\texttt{p == 0}) does \emph{not} break and continues the remaining iterations, itself becoming a parent, and so on. Result: a chain of length $N$,
\[
P \to P_1 \to P_2 \to \cdots \to P_N,
\]
so still $n=N$ descendants, but a path rather than a star.

(c)~Every process that reaches iteration $i$ forks, and \emph{both} parent and child continue. The process tree is a full binary tree of height $N$; the total number of processes including $P$ is $2^N$, so
\[
n = 2^N - 1.
\]
\end{Sol}

\begin{Qn}{Hard}{Lect 3}{CS-537 Spring 2018 Mid Problems VII--VIII}{\texttt{fork}/\texttt{execv}/\texttt{wait}: success vs.\ any-call-may-fail}
Here is code. Assume the program \texttt{/bin/true}, when it runs, never prints anything and just returns 0 in all cases.
\begin{lstlisting}[numbers=none]
int main(int argc, char *argv[]) {
  int rc = fork();
  if (rc == 0) {
    char *my_argv[] = { "/bin/true", NULL };
    execv(my_argv[0], my_argv);
    printf("1");
  } else if (rc > 0) {
    wait(NULL);
    printf("2");
  } else {
    printf("3");
  }
  return 0;
}
\end{lstlisting}
\texttt{printf()} prints immediately (no buffering). Mark A for possible, B for not possible.

\textbf{Problem VII.} Assuming \emph{all} system calls succeed, which outputs are possible?
31.\ \texttt{123}\quad 32.\ \texttt{12}\quad 33.\ \texttt{2}\quad 34.\ \texttt{23}\quad 35.\ \texttt{3}

\textbf{Problem VIII.} Assuming \emph{any} of the system calls above might fail (returning an error), which outputs are possible?
36.\ \texttt{123}\quad 37.\ \texttt{12}\quad 38.\ \texttt{2}\quad 39.\ \texttt{23}\quad 40.\ \texttt{3}
\end{Qn}
\begin{Sol}
\textbf{VII (all succeed).} The child (\texttt{rc==0}) \texttt{execv}s \texttt{/bin/true} and therefore never prints \texttt{1}. The parent (\texttt{rc>0}) \texttt{wait}s, then prints \texttt{2}. The \texttt{else} branch is not taken. Only \texttt{2}.
\[
31\ \mathrm{B},\quad 32\ \mathrm{B},\quad 33\ \mathrm{A},\quad 34\ \mathrm{B},\quad 35\ \mathrm{B}.
\]

\textbf{VIII (any call may fail).}
\begin{itemize}
  \item All succeed: still just \texttt{2}.
  \item \texttt{fork} succeeds, \texttt{execv} fails: child returns from \texttt{execv} and prints \texttt{1}; parent \texttt{wait}s then prints \texttt{2}. Output \texttt{12} (if \texttt{wait} also failed immediately, \texttt{21} could occur too, though \texttt{21} is not among the five choices).
  \item \texttt{fork} fails (\texttt{rc<0}): only \texttt{3}.
  \item \texttt{2} and \texttt{3} together would require \texttt{fork} to both succeed and fail. Impossible. Three-digit \texttt{123} likewise impossible (\texttt{1} needs failed \texttt{execv} after a successful \texttt{fork}, which cannot also take the \texttt{rc<0} branch).
\end{itemize}
\[
36\ \mathrm{B},\quad 37\ \mathrm{A},\quad 38\ \mathrm{A},\quad 39\ \mathrm{B},\quad 40\ \mathrm{A}.
\]
(Official key, CS-537 Sp18 Mid answers, Problems VII--VIII.)
\end{Sol}

\begin{Qn}{Hard}{Lect 4}{Mandar ISI 2011 Mid Q1(e); repeated Mandar ISI 2016 Mid Q1(a)}{Faithful restore of a saved context is wrong}
Consider the following sequence of steps used by the kernel to effect a context-switch from process $P$ to $P'$.
\begin{enumerate}
  \item Run the scheduling algorithm. Let $P'$ be the selected process.
  \item If $P == P'$, goto 5.
  \item Save context of $P$.
  \item Restore context of $P'$.
  \item Return to user mode.
\end{enumerate}
(a)~[both papers] Now consider a context switch from process $P''$ to $P$. What problem arises if, in step 4, the context of $P$ is restored faithfully using exactly the values saved in step 3 during the switch from $P$ to $P'$? Your answer should clearly explain what happens at the individual machine-instruction level.

(b)~[2011 Mid only] How should the above sequence of steps be changed to circumvent this problem?
\end{Qn}
\begin{Sol}
(a)~Step 3 saves $P$'s registers \emph{including the PC}, and that PC points at the next kernel instruction after the save --- i.e.\ at step 4, ``restore context of $P'$''. Later, when the kernel switches $P''\to P$ and restores that saved context \emph{faithfully}, $P$ resumes in kernel mode at step 4 of the \emph{old} switch: it executes ``restore context of $P'$'' and then ``return to user mode''. So $P$ loads $P'$'s registers and returns-from-trap into $P'$'s user mode, or otherwise corrupts the switch. In short: the saved PC is in the middle of the context-switch sequence, not at a safe resume point on $P$'s kernel stack (the original trap frame).

(b)~Two classic fixes:
\begin{itemize}
  \item \textbf{Bach / \texttt{setjmp}-style return value (direct fix to steps 3--5):} have \texttt{save\_context(P)} return \texttt{0} when saving and \texttt{1} when $P$ is later restored (much like \texttt{setjmp}/\texttt{longjmp}). Then step~3 checks \texttt{if (save\_context(P) == 1) goto 5;}, so when $P$ is resumed it skips step~4 and goes straight to returning to user mode.
  \item \textbf{Kernel-stack switch (OSTEP / xv6 \texttt{swtch}):} push a return address pointing \emph{after} the switch onto $P$'s kernel stack, save callee-saved registers, switch the kernel stack pointer to $P'$'s kernel stack, pop $P'$'s registers, and \texttt{ret} into $P'$'s post-switch return-to-user-mode path.
\end{itemize}
\end{Sol}

\begin{Qn}{Hard}{Lect 5--6}{Mandar ISI 2016 Mid Q2}{Choose bursts so TAT inequalities hold for FCFS, SJF, RR}
Consider 3 processes, $P_0$, $P_1$ and $P_2$, that are created in that order at almost the same time. Assume that each of these processes has exactly 1 CPU burst (of duration $d_0$, $d_1$ and $d_2$ resp.), and no I/O bursts. Construct scenarios by choosing appropriate values of $d_0$, $d_1$ and $d_2$ ($\le 10$\,ms), and any other parameter if necessary, such that
\begin{itemize}
  \item[(a)] $T_F = T_R = T_S$;
  \item[(b)] $T_F > T_R > T_S$;
  \item[(c)] $T_R > T_F > T_S$;
  \item[(d)] $T_R > T_F = T_S$;
\end{itemize}
where $T_F$, $T_S$ and $T_R$ are the average turnaround times for the 3 processes if they are scheduled using First-Come-First-Served, Shortest-Job-First, and Round-Robin respectively. Justify your answer in each case.
\end{Qn}
\begin{Sol}
All three arrive together, so SJF is non-preemptive and runs the bursts in increasing order. FCFS runs $P_0,P_1,P_2$. RR needs a quantum $q$; if $q$ is at least the longest burst, RR degenerates to FCFS. Turnaround of equal-length jobs is \emph{worse} under RR than under FCFS (OSTEP).

(a)~$T_F=T_R=T_S$. Take $d_0=d_1=d_2=5$ and $q\ge 5$. Then FCFS $=$ SJF (all lengths equal) and RR $=$ FCFS (one slice each, run to completion). Average TAT $= (5+10+15)/3 = 10$ under all three.

(b)~$T_F > T_R > T_S$. Convoy: a long job first, plus a small RR quantum so short jobs finish early. Take $d_0=10$, $d_1=d_2=1$, $q=1$.
\begin{itemize}
  \item FCFS: completions $10,11,12$, TAT avg $33/3=11$.
  \item SJF: order $P_1,P_2,P_0$; completions $1,2,12$, TAT avg $15/3=5$.
  \item RR: $P_0(1)$, $P_1$ done at 2, $P_2$ done at 3, $P_0$ done at 12. TAT avg $17/3 \approx 5.67$.
\end{itemize}
So $11 > 5.67 > 5$.

(c)~$T_R > T_F > T_S$. Nearly equal jobs (RR stretches everyone) but not already in SJF order. Take $d_0=5$, $d_1=4$, $d_2=3$, $q=1$.
\begin{itemize}
  \item FCFS: completions $5,9,12$, avg $26/3 \approx 8.67$.
  \item SJF: order $P_2,P_1,P_0$; completions $3,7,12$, avg $22/3 \approx 7.33$.
  \item RR: three full rounds through all three (9 time units) finish $P_2$ at $t=9$; then $P_0,P_1,P_0$ finish $P_1$ at 11 and $P_0$ at 12. TAT avg $32/3 \approx 10.67$.
\end{itemize}
So $10.67 > 8.67 > 7.33$.

(d)~$T_R > T_F = T_S$. FCFS already in SJF order, and RR still stretches. Take $d_0=3$, $d_1=4$, $d_2=5$, $q=1$. FCFS $=$ SJF: completions $3,7,12$, avg $22/3 \approx 7.33$. RR: $P_0$ done at 7, $P_1$ at 10, $P_2$ at 12, avg $29/3 \approx 9.67 > 7.33$.
\end{Sol}

\begin{Qn}{Hard}{Lect 7}{CS-537 Fall 2009 Mid Q3}{Reverse-engineer lost MLFQ parameters from a trace}
In the latest LostOS, there is an MLFQ scheduler. Unfortunately, all of its parameters have been lost! Your job is to examine a scheduling trace and determine those critical parameters.

Two CPU-bound (no I/O) jobs, A and B, run. After 500 milliseconds the behaviour repeats indefinitely (until the jobs are done). A close-up of the first $\sim 220$\,ms shows three distinct slice lengths being used, the shortest first (about 10\,ms), then about 20\,ms, then a long slice of about 55\,ms, with A and B taking turns; the whole pattern (high-priority short slices, then medium, then long RR, then a return to short slices) restarts at 500\,ms.
\begin{itemize}
  \item[(a)] How many queues do you think there are in this MLFQ scheduler?
  \item[(b)] How long is the time slice at the top-most (high priority) queue?
  \item[(c)] How long is the time slice at the bottom-most (low priority) queue?
  \item[(d)] How often do processes get moved back to the topmost queue?
  \item[(e)] Why does the scheduling policy MLFQ move processes to higher priority levels (i.e., the topmost queue) sometimes?
\end{itemize}
\end{Qn}
\begin{Sol}
Official key (CS-537 Fa09 Mid answers, Q3):
\begin{itemize}
  \item[(a)] \textbf{Three} queues --- three distinct time slices in the close-up (10\,ms, 20\,ms, 55\,ms).
  \item[(b)] Top-most slice $= \mathbf{10}$\,ms (the first slice you get).
  \item[(c)] Bottom-most slice $= \mathbf{55}$\,ms (the longest slice shown).
  \item[(d)] Boost period $S \approx \mathbf{500}$\,ms (the trace repeats from the high-priority slices at 500\,ms).
  \item[(e)] Most important: \textbf{avoid starvation}. Also: \textbf{reassess} a job that may now be interactive (behaviour change).
\end{itemize}
This is Lect~7 in pictures: short slices at high priority (approximate SJF / good response), long slices at low priority (amortise CS for CPU hogs), periodic boost $S$.
\end{Sol}

\clearpage

\section{Lectures 8--12 (Pre-Minor): Early Memory Virtualization (B\&B, Seg., Free Space \& Paging)}

\begin{Qn}{Hard}{Lect 10}{CS-537 Fall 2004 Final, Q1}{Broken segment-table translations (explicit 2-bit SEG\_MASK)}
A segmentation MMU uses the following segment table. The segment number
is taken from the \textbf{top two bits} of each $32$-bit virtual
address; the remaining bits are the offset. Bounds is a size.

\begin{center}
\begin{tabular}{c|c|c|c}
Segment & Base & Bounds & Protection \\
\hline
0 & \texttt{0x1000} & \texttt{0x100} & Read \\
1 & \texttt{0x2000} & \texttt{0x200} & Read/Write \\
2 & \texttt{0x5000} & \texttt{0x500} & Read/Write \\
\end{tabular}
\end{center}

Each of the following observed translations is wrong. For each, state
what the access \emph{should} have done.
\begin{enumerate}
  \item[(a)] Load \texttt{0x00000010} $\to$ segmentation violation
  \item[(b)] Load \texttt{0x40000300} $\to$ loaded from PA
        \texttt{0x00002300}
  \item[(c)] Load \texttt{0x80000300} $\to$ loaded from PA
        \texttt{0x00002300}
  \item[(d)] Store \texttt{0x00000050} $\to$ stored to PA
        \texttt{0x00001050}
  \item[(e)] Load \texttt{0xC0000010} $\to$ loaded from PA
        \texttt{0x00000010}
\end{enumerate}
\end{Qn}
\begin{Sol}
$\mathrm{SEG}=(\mathrm{VA}\gg 30)$ and
$\mathrm{offset}=\mathrm{VA}\ \&\ \texttt{0x3FFFFFFF}$. Valid iff
$\mathrm{offset}<\mathrm{bounds}[\mathrm{SEG}]$ and the access matches
the protection bits; then
$\mathrm{PA}=\mathrm{base}[\mathrm{SEG}]+\mathrm{offset}$.

(a) $\mathrm{SEG}=00$, offset \texttt{0x10}$=16<$\texttt{0x100},
protection Read: a load is legal. Should have succeeded at
PA $\texttt{0x1000}+\texttt{0x10}=\texttt{0x1010}$. The fault is
wrong.

(b) \texttt{0x40000300} has top bits $01$ (SEG 1), offset
\texttt{0x300}$=768>$\texttt{0x200}$=512$. Should have been a
\textbf{segmentation fault}. The MMU ignored the bounds and computed
$\texttt{0x2000}+\texttt{0x300}=\texttt{0x2300}$.

(c) \texttt{0x80000300} has top bits $10$ (SEG 2), offset
\texttt{0x300}$<$\texttt{0x500}. Should have loaded from
$\texttt{0x5000}+\texttt{0x300}=\texttt{0x5300}$. The MMU used
segment 1's base instead.

(d) $\mathrm{SEG}=0$, offset \texttt{0x50}$<$\texttt{0x100}, but
segment 0 is \textbf{Read}-only. A store must raise a
\textbf{protection fault}. The MMU ignored R/W/X.

(e) Top bits $11$ select SEG 3, which does not exist. Should have
faulted. The MMU treated the offset as a physical address.
\end{Sol}

\begin{Qn}{Hard}{Lect 11}{CS-537 Fall 2004 Mid2, Q2}{First-fit vs.\ best-fit on a four-chunk free list}
A free list consists of the following free chunks, in this order from
the head of the list: $10$ bytes, $20$ bytes, $40$ bytes, and
$10$ bytes. Then come the allocation requests: allocate $10$, allocate
$15$, allocate $10$, allocate $35$.
\begin{enumerate}
  \item Using first-fit allocation, will all requests succeed?
  \item Using best-fit allocation, will all requests succeed?
  \item In general, which is better, first fit or best fit?
\end{enumerate}
\end{Qn}
\begin{Sol}
Assume a request of $n$ bytes is satisfied from the chosen chunk and
any remainder replaces that chunk in place (no header overhead, as the
exam stated only sizes).

\textbf{First fit.}
\begin{itemize}
  \item alloc $10$: takes the first $10$. List: $20,40,10$.
  \item alloc $15$: first large enough is $20$, remainder $5$.
        List: $5,40,10$.
  \item alloc $10$: $5$ is too small, so take $40$, remainder $30$.
        List: $5,30,10$.
  \item alloc $35$: no chunk $\ge 35$. \textbf{Fails.}
\end{itemize}

\textbf{Best fit} (smallest chunk $\ge$ request).
\begin{itemize}
  \item alloc $10$: two exact $10$-byte chunks; take the first.
        List: $20,40,10$.
  \item alloc $15$: smallest $\ge 15$ is $20$, remainder $5$.
        List: $5,40,10$.
  \item alloc $10$: exact $10$. List: $5,40$.
  \item alloc $35$: take $40$, remainder $5$. List: $5,5$.
        \textbf{All four succeed.}
\end{itemize}

\textbf{In general} neither dominates. Best fit tends to leave many
tiny leftover fragments (external fragmentation) and scans the whole
list. First fit is faster (stops at the first hit) but can strip large
chunks early, as it did to the $40$-byte block above, and then fail a
later large request. Worst fit (always the largest chunk) tries to
keep leftovers usable, but also scans the whole list and still
fragments. Next fit is first fit that resumes at the last cursor.
\end{Sol}

\begin{Qn}{Hard}{Lect 12}{CS-537 Spring 2011 Mid, Problem 5}{Which physical bytes are touched to fetch and execute one load?}
A linear page table with $1$-byte PTEs. Address space $128$ bytes,
$32$-byte pages, physical memory $128$ bytes. The page-table base
register is physical address $16$. Contents:

\begin{center}
\begin{tabular}{c|c}
VPN & PFN \\
\hline
0 & $1$ \\
1 & not valid \\
2 & $3$ \\
3 & not valid \\
\end{tabular}
\end{center}

The instruction
\begin{lstlisting}[numbers=none]
10: LOAD 70, R1
\end{lstlisting}
resides at virtual address $10$ and loads a single byte from virtual
address $70$ into \texttt{R1}. There is no TLB. On a diagram of
physical memory:
\begin{enumerate}
  \item[(a)] Box each valid virtual page (and label them).
  \item[(b)] Box the page table (and label it).
  \item[(c)] Circle every physical address referenced during execution
        of the instruction, including instruction fetch and the data
        access.
\end{enumerate}
\end{Qn}
\begin{Sol}
Four virtual pages ($128/32$), so a $2$-bit VPN and a $5$-bit offset.
The linear table lives at PA $16$ and occupies $4$ bytes
(PA $16..19$).

(a) VPN $0$ is in frame $1$ (PA $32..63$). VPN $2$ is in frame $3$
(PA $96..127$). VPN $1$ and $3$ are not present in physical memory.

(b) The page table is the $4$ bytes at PA $16..19$.

(c) No TLB $\Rightarrow$ every VA costs a PTE read, then the data
(assuming a 1-byte instruction encoding at VA $10$):
\begin{itemize}
  \item Fetch at VA $10$: VPN $0$, offset $10$. PTE is at
        $\mathrm{PTBR}+\mathrm{VPN}=16+0=16$. Instruction byte is at
        $\mathrm{PFN}\times 32+\mathrm{offset}=32+10=42$.
  \item Data VA $70$: VPN $2$ ($70=2\times 32+6$), offset $6$. PTE is
        at $16+2=18$. Data byte is at $3\times 32+6=102$.
\end{itemize}
The four circled physical addresses are therefore \textbf{16, 42, 18,
102} --- two page-table references plus the instruction byte plus the
data byte. This is the Lect 12 ``paging costs two memory accesses''
tax, before any TLB (Lect 13) is introduced.
\end{Sol}

\clearpage

\part{Additional Practice, MCQs \& Conceptual Drills (Extra Practice)}
% ===========================================================================

\section{Lectures 13--15: Virtual Memory --- TLB, Advanced Paging \& Swapping}

\subsection{Lectures 13--14: Paging with TLB \& Advanced Paging}

\subsubsection*{Additional Challenging \& Analytical Problems}

\begin{Qn}{Medium-Hard}{Lect 14}{CS-537 Fa21 Mid, Q25--26}{Base$+$bounds over a page table (segmented paging)}
To shrink page tables, ZOS combines base/bounds with paging. The address
space is still chopped into pages. A base register points at the page
table; a bounds register holds the size of that table (max valid VPN,
plus one).
\begin{enumerate}[(a)]
\item This approach (choose one):
\begin{enumerate}[(A)]
\item Enables fast translation with only two extra memory references to
      fetch a PTE.
\item Enables a compact page table, if you use the virtual address space
      carefully.
\item Is always smaller than a linear page table.
\item Supports a sparse address space while still also minimizing page
      table memory usage.
\item None of the above.
\end{enumerate}
\item 32-bit address space, 1\,KB pages, 4-byte PTEs, 100 processes.
      Each process uses only \textbf{one} virtual page. What is the
      \textbf{worst-case} total size of all of these page tables?
\begin{enumerate}[(A)]
\item 16\,MB
\item 160\,MB
\item 1600\,MB
\item 16\,GB
\item None of the above.
\end{enumerate}
\end{enumerate}
\end{Qn}
\begin{Sol}
This is Lecture 14's \textbf{segmented page table}: one (smaller) linear
table per segment, hardware tracking $(\mathrm{base},\mathrm{bounds})$ of
that table. VA layout is $\mathrm{seg}\vert\mathrm{page}\vert\mathrm{offset}$;
here there is a single segment whose bounds is ``how many VPNs exist at
the low end of the space''.
\begin{enumerate}[(a)]
\item \textbf{(B).} If the process only uses the low VPNs, bounds can be
      tiny (one PTE if only VPN 0 is valid). It is \emph{not} always
      smaller than a full linear table (use the last VPN and bounds must
      cover the whole space). It does \emph{not} handle a sparse space
      (code at 0, stack at $2^{32}-1$) without many empty PTEs in the
      hole --- that is what multi-level tables fix. One extra memory
      reference still fetches the PTE (not two).
\item \textbf{(C) 1600\,MB.} VPN bits $= 32-10 = 22$, so a full table
      is $2^{22}$ PTEs $\times 4\,\mathrm{B} = 16\,\mathrm{MB}$. Worst
      case: the single valid page is the \emph{last} VPN, so bounds must
      be $2^{22}$ and each of 100 processes still needs a full 16\,MB
      table. Total $100\times 16\,\mathrm{MB} = 1600\,\mathrm{MB}$.
      (Best case, VPN 0 only: $100\times 4\,\mathrm{B} = 400\,\mathrm{B}$.)
\end{enumerate}
\end{Sol}

\begin{Qn}{Medium-Hard}{Lect 13}{CS-537 Sp16 Mid, Q14}{Hit/miss strings for sequential and strided array walks}
Trace TLB behaviour. Produce a string of hits (\texttt{H}) and misses
(\texttt{m}); e.g.\ \texttt{mHmH} means miss, hit, miss, hit. Ignore
instruction references. The array is page-aligned. Page size is 32 bytes;
each array element is a 4-byte integer.
\begin{enumerate}[(a)]
\item 1-entry TLB. Access contiguous 4-byte integers in a large array,
      starting at index 0. Hit/miss pattern?
\item 2-entry TLB with LRU replacement. Same contiguous scan. Pattern?
\item 1-entry TLB. Access \emph{every other} 4-byte integer (indices
      $0,2,4,\ldots$). Pattern?
\item 16-entry TLB with FIFO replacement. Repeatedly scan all 4-byte
      integers in a small contiguous array of 24 integers. Pattern on
      the \emph{fourth} run through the loop?
\end{enumerate}
\end{Qn}
\begin{Sol}
A 32-byte page holds $32/4 = 8$ consecutive integers, so a dense scan
touches 8 elements per page.
\begin{enumerate}[(a)]
\item One live page at a time. The first integer of each page misses;
      the next seven hit:
      \[
      \texttt{mHHHHHHH}\ \textrm{repeated}.
      \]
      Miss rate $1/8 = 12.5\%$.
\item Sequential access never needs two pages at once, so the second
      TLB slot is unused. Same pattern:
      $\texttt{mHHHHHHH}$ repeated.
\item Every-other-integer hits 4 elements per page (indices
      $0,2,4,6$ then the next page). Pattern:
      \[
      \texttt{mHHH}\ \textrm{repeated}.
      \]
\item 24 integers occupy $24\times 4 = 96$ bytes $= 3$ pages. A 16-entry
      TLB holds all three translations after the first pass. The fourth
      run is therefore all hits: $\texttt{H}$ repeated 24 times.
\end{enumerate}
\end{Sol}

\begin{Qn}{Medium-Hard}{Lect 13}{CS-537 Fa10 Mid, Q4}{Software-managed TLB entry: VPN, PFN, ASID}
A MIPS-style software-managed TLB entry contains (among other fields) a
VPN, a PFN, a valid bit, a dirty bit, and an ASID. On a miss the hardware
traps; the OS installs a translation and returns-from-trap; the hardware
retries the instruction.

For each of VPN, PFN, and ASID: (i)~what is the field for?
(ii)~what goes wrong if a bit in that field is corrupted just before the
OS writes the entry into the TLB?
\end{Qn}
\begin{Sol}
\begin{enumerate}[(a)]
\item \textbf{VPN.} Tag used to match the faulting virtual page.
      A flipped VPN-bit installs the \emph{right} PFN under the
      \emph{wrong} tag: later accesses to the true VPN miss (or hit a
      different row), while a different VPN may hit and be sent to this
      PFN --- a protection/aliasing bug, possibly a write into another
      process's page if ASIDs also match.
\item \textbf{PFN.} The physical frame of a matching VPN.
      A flipped PFN-bit sends a correct VPN to the wrong frame: reads
      return foreign data; writes corrupt another page (or a non-RAM
      physical address). Valid+dirty bits would not save you; the
      mapping itself is wrong.
\item \textbf{ASID.} Address-space identifier (typically $\ll$ a full
      PID, e.g.\ 8 bits $\Rightarrow$ 256 concurrent tagged spaces).
      Match is $(\mathrm{VPN},\mathrm{ASID})$. A flipped ASID either
      (1)~hides the translation from its owner (extra misses / faults)
      or (2)~lets a \emph{different} process hit this PFN when it uses
      the same VPN --- a cross-process leak. Because ASIDs are recycled,
      the OS must flush rows whose ASID is being reused.
\end{enumerate}
Hardware-managed TLBs (x86, ARM) would walk the page table via CR3/PTBR
on a miss and never trap if the PTE is valid; software-managed TLBs
(MIPS, some SPARC) \emph{always} trap, which is why a buggy OS fill is
immediately visible, and why inverted/hashed tables are a natural fit
(Lecture 14).
\end{Sol}

\begin{Qn}{Medium}{Lect 13}{CS-537 Fa08 Final, Q10(a)--(c)}{Control flow: TLB hit, miss-present, miss-not-present}
A user process runs on an OS with a \textbf{software-managed} TLB and a
linear page table in physical memory. Numbered events:
\begin{enumerate}[0.]
\item hardware checks the TLB for the VPN$\to$PPN translation
\item hardware issues a load to a physical address
\item OS code at the start of the TLB-miss handler runs
\item OS code at the end of the TLB-miss handler runs (return from trap)
\item OS code that accesses the page table to look up a translation
\item OS code that updates the TLB with a new mapping
\item OS code that updates the page table with a new VPN$\to$PPN mapping
\item a disk request is initiated
\item a disk request completes
\end{enumerate}
Write the flow for: (a)~a TLB hit; (b)~a TLB miss to a valid page that
is present in memory; (c)~a TLB miss to a valid page that is not present
(must be fetched from disk).
\end{Qn}
\begin{Sol}
\begin{enumerate}[(a)]
\item Hit: $0,1$. The MMU never leaves hardware; no trap, no PT walk.
\item Miss, present: $0,2,4,5,3$, then retry $0,1$.
      The miss handler reads the PTE (valid and present), fills the TLB,
      returns; the retry is now a hit and the data load proceeds.
\item Miss, not present: $0,2,4,6,7,8,5,3$, then retry $0,1$.
      The handler finds a valid-but-not-present PTE, allocates a frame,
      starts disk I/O, blocks; after completion it writes the PPN into
      the PTE, fills the TLB, and returns. (If a dirty victim must be
      written back, another $7,8$ pair occurs first.) The process is
      blocked for milliseconds; a scheduler may run someone else.
\end{enumerate}
On a \textbf{hardware-managed} TLB, step 2 would not be an OS trap when
the PTE is present: the MMU itself walks the page table via CR3. The OS
is invoked only for true page faults (not-present or invalid).
\end{Sol}

\begin{Qn}{Medium-Hard}{Lect 14}{CS-537 Fa04 Mid2, Q12}{Reverse translation: which VPN maps to physical address $P$?}
In a system with a typical linear page table, you are given a physical
address $P$.
\begin{enumerate}[(a)]
\item How would you find out which virtual address(es) are mapped to $P$?
\item What data structures might you add to speed this up?
\end{enumerate}
\end{Qn}
\begin{Sol}
\begin{enumerate}[(a)]
\item A linear table is $\mathrm{VPN}\mapsto\mathrm{PPN}$. Going the other
      way requires a scan of \emph{every} PTE of \emph{every} process,
      comparing the stored PPN to $P$'s frame and, on a match, joining
      the VPN with $P$'s offset. Shared pages produce several VAs.
      Cost is $O((\textrm{\#processes})\times(\textrm{\#VPNs}))$.
\item An \textbf{inverted page table} (or a hash from PPN to a list of
      $(\mathrm{VPN},\mathrm{ASID})$) makes the reverse map $O(1)$ or
      $O(\textrm{\#sharers})$. A frame table maintained by the OS
      (coremap) is the same idea. This reverse map is also what a
      software TLB-miss handler uses when the forward structure is
      inverted.
\end{enumerate}
\end{Sol}

\begin{Qn}{Medium-Hard}{Lect 14}{CS-537 Fa04 Mid2, Part II Q1(a)}{Combined segmentation and paging}
Describe how a combined segmentation-and-paging (hybrid) approach works,
and how it improves on strict paging and on strict segmentation.
\end{Qn}
\begin{Sol}
VA is split $\mathrm{seg}\# \mid \mathrm{page}\# \mid \mathrm{offset}$.
Hardware holds, per segment, a \textbf{base} (physical address of that
segment's page table) and a \textbf{bounds} (number of PTEs / pages in
the segment). Walk:
\begin{enumerate}
\item Use $\mathrm{seg}\#$ to select $(\mathrm{base},\mathrm{bounds})$.
\item If $\mathrm{page}\# \ge \mathrm{bounds}$, segmentation fault.
\item $\mathrm{PTE} = \mathrm{Mem}[\mathrm{base} + \mathrm{page}\#]$.
\item $\mathrm{PA} = (\mathrm{PFN}\ll \mathrm{offset\text{-}bits})
      \mid \mathrm{offset}$.
\end{enumerate}
Versus \textbf{strict paging}: each segment's table only has as many PTEs
as that segment has pages, so the huge hole between heap and stack is not
stored (the linear-PT space problem). Versus \textbf{strict segmentation}:
pages inside a segment need not be contiguous, so external fragmentation
of segment placement goes away; sharing and protection are still per
segment (and per page). Quiz-scale example from lecture: 2-bit seg, 18-bit
page, 12-bit offset, 4-byte PTE $\Rightarrow$ one segment's table is
$2^{18}\times 4\,\mathrm{B} = 1\,\mathrm{MB}$ contiguous --- already
better than a 4\,MB linear table, but a single segment table is still
contiguous, which is why we then page the page table (multi-level).
\end{Sol}

\begin{Qn}{Medium-Hard}{Lect 14}{CS-537 Sp16 Final, Q9(a)--(e)}{Three-level table for two valid pages; large-page trade-off}
\begin{enumerate}[(a)]
\item 32-bit VAs, 4\,KB pages. How many PTEs in a linear page table?
\item Now a \textbf{three-level} table: one root directory, mid-level
      directories, then leaves. Only the very first and very last
      virtual pages are valid. How many pages does this page table use?
\item \texttt{int *p = 0; printf("\%d\textbackslash n", *p);} --- what
      happens, given (b)?
\item Replace 4\,KB pages by 4\,MB pages. Two major benefits of large
      pages?
\item Two major disadvantages?
\end{enumerate}
\end{Qn}
\begin{Sol}
\begin{enumerate}[(a)]
\item $2^{20} = 1\,048\,576$ PTEs (20-bit VPN).
\item Five pages: 1 root $+$ 2 mid-level pages (one for the low end of
      the space, one for the high end) $+$ 2 leaf pages (one PTE each).
      A two-level table would use 3 pages for the same sparsity (1
      directory $+$ 2 leaves); the extra two levels exist so that each
      \emph{directory} itself still fits in a page when the VA grows.
\item Virtual address 0 is valid by assumption, so \texttt{*p} does
      \emph{not} segfault: it prints the four bytes at the start of
      whatever physical page was mapped at VPN 0 (often the first
      bytes of the code or a mapped zero page, depending on the
      process image). (If VPN 0 were \emph{invalid}, this would be a
      null-pointer fault.)
\item Smaller page tables (VPN shrinks by 10 bits when the page grows
      from 4\,KB to 4\,MB: $2^{10}$ PTEs, 4\,KB linear table instead of
      4\,MB); much larger \textbf{TLB reach}
      ($N_{\mathrm{TLB}}\times 4\,\mathrm{MB}$ instead of
      $N_{\mathrm{TLB}}\times 4\,\mathrm{KB}$), hence fewer compulsory
      TLB misses.
\item Internal fragmentation (a 1-byte mapping wastes almost 4\,MB);
      slower/heavier paging I/O; harder physical allocation (need
      contiguous aligned 4\,MB frames); worse cache colouring. Large
      pages do \emph{not} by themselves shrink a multi-level table
      below the ``one page per used level'' floor when the space is
      sparse --- they help when the \emph{used} region is huge and
      contiguous.
\end{enumerate}
\end{Sol}

\begin{Qn}{Medium}{Lect 14}{CS-537 Fa13 Mid, Q11}{Chopping a tiny linear table into page-sized leaves}
64-byte address space, 4-byte pages, 1-byte PTEs. A linear table thus
has $64/4=16$ PTEs (16 bytes). McFlub draws an 8-entry page directory
and pairs of PTEs (PTE0--1, 2--3, \ldots), marking directory slots
0, 4, 7 valid.

If the page size is 4 bytes, how \emph{must} the page table be chopped,
and what does the page directory look like?
\end{Qn}
\begin{Sol}
A page holds 4 bytes $= 4$ PTEs, so leaves are groups of \textbf{four},
not two. There are $16/4=4$ possible leaves, so the directory has
\textbf{4} PDEs (and itself fits in a page). McFlub's 8-entry directory
would be right for 2-byte pages, not 4-byte pages.

Given McFlub's valid PTEs (PTE0 valid; PTE10 valid; PTE14 and PTE15
valid), the four-PTE groups are:
\begin{center}
\begin{tabular}{cl}
leaf 0 (PTE0--3)  & valid (contains PTE0) \\
leaf 1 (PTE4--7)  & empty \\
leaf 2 (PTE8--11) & valid (contains PTE10) \\
leaf 3 (PTE12--15)& valid (contains PTE14, PTE15) \\
\end{tabular}
\end{center}
So PDEs: valid, invalid, valid, valid. (PDE valid $\ne$ PTE valid: a
valid PDE only promises that the inner page exists.)
\end{Sol}

\begin{Qn}{Medium-Hard}{Lect 14}{ISI Mandar 2017 End-sem, Q3(b)(ii)}{Maximum page-table size on i386 (4\,KB pages)}
The i386 page size is 4\,KB. Compute the total size of the page tables
(the conventional 2-level IA-32 scheme: 10-bit directory, 10-bit table,
12-bit offset, 4-byte entries). Assume a fully populated address space.
\end{Qn}
\begin{Sol}
One page directory: $2^{10}$ PDEs $\times 4\,\mathrm{B} = 4\,\mathrm{KB}$.
Fully populated, every PDE is valid, so there are $2^{10}$ leaf page
tables, each $4\,\mathrm{KB}$, totalling $4\,\mathrm{MB}$. Grand total
$4\,\mathrm{KB}+4\,\mathrm{MB}$. (This is exactly the linear-table size
of 4\,MB plus one extra directory page --- the full-space tax of
multi-level paging.)
\end{Sol}

\begin{Qn}{Medium-Hard}{Lect 13--14}{CS-537 Fa22 Mid, Q3}{Large pages, TLB reach, and a PDE large-page bit}
xOS/x537 supports 4\,KB pages and 4\,MB pages on a 32-bit virtual
address space. PTEs and PDEs are 4 bytes. The top-level page directory
has 1024 entries (4\,KB). A PDE is
\[
[\,\textrm{valid}\ \vert\ \textrm{large-page bit}\ \vert\ \textrm{PFN}\,].
\]
If valid and large-page $=0$, the PFN names a 4\,KB leaf page-table
page. If valid and large-page $=1$, the PFN names a 4\,MB data page
directly.
\begin{enumerate}[(a)]
\item With only 4\,KB pages, how many PTEs would a linear table have?
\item With only 4\,MB pages, how many PTEs would a linear table have?
\item Main problem with using (only) large pages?
\item True statement about large pages and the TLB?
\item First and last PDEs valid, all others invalid, both large-page
      bits set. How much memory can the process legally access?
\item Same two PDEs valid, but both large-page bits clear, and every
      PTE in those two leaves is valid. Legal memory?
\item Which bits of a 32-bit VA index the page directory?
\item If the VA grew to 33 bits, which of the following is true?
      (A)~page size would have to grow; (B)~the page directory would
      have to be bigger; (C)~you could not use large pages;
      (D)~you could only use large pages; (E)~none.
\end{enumerate}
For (a)--(d) and (g)--(h) the options are of the form ``roughly
$10^k$'' / qualitative; for (e)--(f) they are 0/1/2/4/8\,MB (and
``none''). Give the correct choice and the derivation.
\end{Qn}
\begin{Sol}
\begin{enumerate}[(a)]
\item $2^{20}\approx 1$ million PTEs. Answer: roughly $1{,}000{,}000$.
\item Offset is now 22 bits ($4\,\mathrm{MB}=2^{22}$), VPN $=10$ bits,
      $2^{10}=1024$ PTEs. Answer: roughly $1{,}000$.
\item Internal fragmentation: a sparsely used 4\,MB page wastes almost
      the whole frame. (They do \emph{not} make PT access slower, nor
      do they consume more \emph{page-table} memory.)
\item Large pages \textbf{reduce TLB miss rates} by increasing TLB
      reach: 64 entries $\times 4\,\mathrm{KB} = 256\,\mathrm{KB}$ vs.\
      64 $\times 4\,\mathrm{MB} = 256\,\mathrm{MB}$. They do not make
      the TLB physically larger.
\item Two valid large-page PDEs $\Rightarrow$ two 4\,MB mappings:
      \textbf{8\,MB}.
\item Each leaf has 1024 PTEs $\times 4\,\mathrm{KB} = 4\,\mathrm{MB}$,
      two leaves $\Rightarrow$ \textbf{8\,MB} again (same coverage,
      different granularity).
\item 1024-entry directory $\Rightarrow$ \textbf{top 10 bits}.
      Split: 10 (PD) $+$ 10 (PT) $+$ 12 (offset), or 10 $+$ 22 when the
      large-page bit is set.
\item \textbf{(B)}: one extra VA bit doubles the number of PDEs if we
      keep 4\,KB directory pages and 4-byte entries, so the directory
      (or the number of directory pages) grows. This is why 64-bit
      machines add levels rather than a 16-PB linear table.
\end{enumerate}
\end{Sol}

\begin{Qn}{Medium-Hard}{Lect 14}{CS-537 Fa22 Mid, Q5}{Tiny two-level table: linear size, sparsity, PDE decode}
Same 15-bit / 32-byte / 1-byte PTE parameters as the paging homework
(1024 virtual pages, 128 physical pages, 5+5+5 split).
\begin{enumerate}[(a)]
\item A \textbf{linear} page table for this space uses how many pages?
      (A)~3 (B)~32 (C)~64 (D)~1024 (E)~none
\item A multi-level table for an address space with just \textbf{one}
      valid page uses how many pages? (A)~1 (B)~2 (C)~32 (D)~1024 (E)~none
\item Which is true of multi-level tables?
      (A)~they are never larger than linear;
      (B)~they can be bigger than linear if the space is fully allocated;
      (C)~they can only be two levels deep;
      (D)~they cannot contain a page directory;
      (E)~none.
\item Page directory (page 81), 32 bytes, high bit $=$ valid:
\begin{lstlisting}[numbers=none]
84 fb 8b b7 9e b9 fa ca 90 b4 f0 86 e3 74 99 a4
a2 7f aa 83 a3 e7 7f bb eb ad e2 7f ...
\end{lstlisting}
      How many \textbf{invalid} directory entries are there?
      (count high-bit-zero bytes; the printed dump shows three
      \texttt{0x7f} bytes among the first 29 plus any remaining
      \texttt{0x7f}s --- use that valid $= \mathrm{byte}\ge\texttt{0x80}$.)
\item First PDE is \texttt{0x84}. Which physical frame (decimal) does it
      name?
\item Physical page 2 (a leaf) contains two bytes other than
      \texttt{0x7f}: \texttt{0xa0} at index 14 and \texttt{0xc6} at
      index 21. How many valid PTEs?
\item Which directory entry refers to page 2?
      (A)~\texttt{0x84} (B)~\texttt{0x8b} (C)~\texttt{0x82} (D)~\texttt{0x7f}
\item Physical page 93 holds data. A physical address \texttt{0x48cb}
      refers to a byte on that page; which byte value is returned?
      Bytes of page 93 begin
      \texttt{15 08 08 00 11 15 07 02 12 12 0c 03 1e 1a 13 11} \ldots
\end{enumerate}
\end{Qn}
\begin{Sol}
\begin{enumerate}[(a)]
\item 1024 PTEs $\times 1\,\mathrm{B} = 1024\,\mathrm{B} = 32$ pages.
      \textbf{(B)}.
\item 1 directory page $+$ 1 leaf $= 2$ pages. \textbf{(B)}.
\item \textbf{(B)}: fully populated, multi-level $=$ linear leaves $+$
      directory pages (strictly larger). Levels can be 3, 4, 5, \ldots
\item A byte is invalid iff its high bit is 0, i.e.\ value
      $< \texttt{0x80}$. In the official key the directory contains
      \textbf{three} invalid entries (the three \texttt{0x7f} bytes
      visible in the dump; all other printed bytes are $\ge\texttt{0x80}$).
      Answer: 3.
\item \texttt{0x84} $= 1000\,0100_2$: valid, PFN $=0000100_2=4$.
      Answer: \textbf{4}.
\item Two bytes with high bit 1 (\texttt{0xa0}, \texttt{0xc6}); the rest
      are \texttt{0x7f} (invalid). Answer: \textbf{2}.
\item We need valid $+$ PFN $=2$: $1\,0000010_2=\texttt{0x82}$.
      \textbf{(C)}.
\item \texttt{0x48cb}: PFN of page 93 is $93=0x5d$; more directly, the
      offset is the low 5 bits of \texttt{0xcb} $= 01011_2=11$. Byte 11
      of the dump is \texttt{0x0c}? Wait: indices
      0:\texttt{15}, 1:\texttt{08}, \ldots, 11:\texttt{0c}. The official
      key for this offering returns \texttt{0x03}, which is index 11
      of a slightly different dump layout in which \texttt{0x03} sits
      at offset 11. Using the numbers printed in the question stem
      above, offset 11 is \texttt{0x0c}. (If the stem's second group
      starts at index 8:
      8:\texttt{12}, 9:\texttt{12}, 10:\texttt{0c}, 11:\texttt{03} ---
      yes, wrapping after 8 bytes of the first row of 8, then
      \texttt{12 12 0c 03}: index 11 $=\texttt{0x03}$.)
      Answer: \texttt{0x03}.
\end{enumerate}
\end{Sol}

\begin{Qn}{Medium-Hard}{Lect 13}{CS-537 Fa22 Mid, Q6}{Linear PTE decode, LRU TLB traces, software-managed TLB}
16-bit VAs, 4\,KB pages (so a 4-bit VPN and 12-bit offset), 128\,KB of
physical memory. Linear PTEs (hex), high bit $=$ valid, low bits $=$ PFN:
\begin{center}
\begin{tabular}{cc@{\hspace{1.5em}}cc}
VPN & PTE    & VPN & PTE \\
\hline
0 & \texttt{0x8018} & 8 & \texttt{0x0000} \\
1 & \texttt{0x0000} & 9 & \texttt{0x801e} \\
2 & \texttt{0x0000} & a & \texttt{0x801c} \\
3 & \texttt{0x800c} & b & \texttt{0x0000} \\
4 & \texttt{0x8009} & c & \texttt{0x800f} \\
5 & \texttt{0x0000} & d & \texttt{0x0000} \\
6 & \texttt{0x801d} & e & \texttt{0x0000} \\
7 & \texttt{0x8013} & f & \texttt{0x8008} \\
\end{tabular}
\end{center}
\begin{enumerate}[(a)]
\item VA \texttt{0xce16} translates to what physical address?
\item Physical address \texttt{0x18398} was accessed. What VA was issued?
\item 4-entry LRU TLB, initially empty. Stream
      \texttt{0xce16, 0xce17, 0xce18, 0xce19}. How many TLB hits?
\item Same TLB, stream
      \texttt{0x0000, 0x3001, 0x4002, 0x6003, 0x0001, 0x3002, 0x4003, 0x6004}.
      How many hits?
\item TLB entry format $[\mathrm{valid}\mid\mathrm{VPN}\mid\mathrm{PFN}]$.
      Which statement is \textbf{false}?
      (A)~the TLB can never be empty;
      (B)~the TLB should be flushed on a process switch;
      (C)~the TLB valid bit is equivalent to the PTE valid bit;
      (D)~the TLB lives in the MMU;
      (E)~none (all true).
\item Four TLB rows
      $[1,3,\texttt{0xc}]$, $[1,0,\texttt{0x18}]$,
      $[1,4,\texttt{0x9}]$, $[1,6,\texttt{0x1d}]$.
      Which row is inconsistent with the page table?
\item TLB rows
      $[0,7,\texttt{0x13}]$, $[1,9,\texttt{0x1e}]$,
      $[1,3,\texttt{0xc}]$, $[0,\texttt{0xf},\texttt{0x08}]$.
      Stream \texttt{0x9000, 0x9001, 0x9002, 0x9003}. Hits?
\item Best true statement about a \textbf{software-managed} TLB:
      (A)~allows more flexible page-table structures;
      (B)~makes miss lookup faster;
      (C)~makes the MMU more complex;
      (D)~adds bits to each TLB entry;
      (E)~all false.
\end{enumerate}
\end{Qn}
\begin{Sol}
VPN $=$ top 4 hex-bits of a 16-bit VA; offset $=$ low 12 bits.
\begin{enumerate}[(a)]
\item \texttt{0xce16}: VPN \texttt{0xc}, offset \texttt{0xe16}.
      PTE \texttt{0x800f} is valid, PFN \texttt{0x0f}.
      PA $= \texttt{0x0fe16}$. (If that is not among the listed
      options, the matching choice is ``none'' / the computed PA.)
      Checking options in the exam: they listed
      \texttt{0xce16}, \texttt{0x18e16}, \texttt{0x0fe16}-style, fault.
      Computed: \texttt{0x0fe16}.
\item PFN \texttt{0x18} is VPN 0's mapping (\texttt{0x8018}). Offset
      \texttt{0x398}, so VA $\texttt{0x0398}$.
\item All four VAs are VPN \texttt{0xc}: \texttt{mHHH} $\Rightarrow$
      \textbf{3 hits}.
\item VPNs $0,3,4,6,0,3,4,6$. Four distinct pages fill the 4-entry TLB
      on the first four accesses (all misses); the next four all hit:
      \textbf{4 hits}.
\item \textbf{(A) and (C) are false}; the exam asked for a single
      false statement and accepted (A) or (C). The TLB \emph{can} be
      empty (all valid bits 0 after a flush). The TLB valid bit means
      ``this row is a cached translation'', not ``this VPN is a valid
      mapping in the page table'' --- a cached row should only exist
      for valid PTEs, but the bits are not the same bit. (B) is true
      without ASIDs; (D) is true.
\item Row 2 claims VPN 0 $\mapsto$ PFN \texttt{0x18}, matching
      \texttt{0x8018}. Row 1: VPN 3 $\mapsto$ \texttt{0xc}, matches
      \texttt{0x800c}. Row 3: VPN 4 $\mapsto$ \texttt{0x9}, matches
      \texttt{0x8009}. Row 4: VPN 6 $\mapsto$ \texttt{0x1d}, matches
      \texttt{0x801d}. All four match --- answer \textbf{(E)} if the
      printed row-2 PFN is \texttt{0x18}; if the exam printed row 2 as
      PFN \texttt{0x08} (VPN \texttt{0xf}'s frame), that row is the
      bad one. The stem above uses \texttt{0x18}, consistent with the
      page table, so (E).
\item All four VAs are VPN \texttt{0x9}, which is cached and valid in
      row 2: \textbf{4 hits}.
\item \textbf{(A)}. The OS miss handler may walk a linear, multi-level,
      hashed, or inverted table; the MMU only has to trap and retry.
      Hardware walkers must understand the table format, so they are
      \emph{less} flexible and \emph{more} complex. Software miss
      handling is typically \emph{slower}, not faster.
\end{enumerate}
\end{Sol}

\begin{Qn}{Medium-Hard}{Lect 13}{CS-537 Fa21 Mid, Q29}{Unique VPNs instead of ASID (no TLB flush)}
Hardware has \textbf{no} ASID/PID field in the TLB. Instead of flushing
on a context switch, ZOS guarantees that each process uses a set of VPNs
disjoint from every other process. Which of the following is \textbf{not}
true of this scheme?
\begin{enumerate}[(A)]
\item Allows fully flexible use of the virtual address space by each process.
\item Allows faster context switches than a flush-everything approach.
\item Allows fully flexible use of physical memory.
\item Allows sharing of code pages between processes.
\item None of the above (all are true).
\end{enumerate}
\end{Qn}
\begin{Sol}
\textbf{(A).} Each process is barred from using VPNs that another process
already uses, so its virtual address space is \emph{not} fully flexible
(no process may start its stack at the same canonical high VPN, etc.).
(B)~is true: no flush, so previously cached translations of the
switched-to process (and of others) survive. (C)~is true: physical
frames are still allocated freely; only the \emph{virtual} numbers are
partitioned. (D)~is true: two processes can map the same PFN under
\emph{different} VPNs, or, if they share a VPN, they could not both
cache it --- but code sharing is possible via distinct VPNs (less
convenient) or by accepting that a shared VPN would be a shared TLB
row pointing at a shared PFN, which this design forbids for
\emph{private} pages but can still represent a shared PFN at two VPNs.
The exam key: A is the statement that is \emph{not} true.
\end{Sol}

\subsubsection*{Conceptual Drills, MCQs \& Further Practice}

\begin{Qn}{Easy-Medium}{Lect 13}{CS-537 Sp18 Mid, Problem XVII}{Does a 4-entry LRU TLB help this address trace?}
Page size 64 bytes; TLB holds 4 entries; LRU replacement. Each item is a
virtual-address trace. Fill \textbf{A} if the TLB can speed up the
program (i.e., at least one hit occurs after the compulsory misses),
\textbf{B} if it cannot (every access misses).
\begin{enumerate}[(81)]
\item $0, 100, 200, 1, 101, 201, \ldots$ (repeats this pattern)
\item $0, 100, 200, 300, 0, 100, 200, 300, \ldots$
\item $0, 1000, 2000, 3000, 4000, 0, 1000, 2000, 3000, 4000, \ldots$
\item $0, 1, 2, 3, 4, 5, 6, 0, 1, 2, 3, 4, 5, 6, \ldots$
\item $300, 200, 100, 0, 300, 200, 100, 0, \ldots$
\end{enumerate}
\end{Qn}
\begin{Sol}
VPN $= \lfloor \mathrm{VA}/64\rfloor$. A 4-entry LRU TLB hits iff the
number of distinct pages in the repeating working set is $\le 4$.
\begin{enumerate}[(81)]
\item Pages $\lfloor 0/64\rfloor=0$, $\lfloor 100/64\rfloor=1$,
      $\lfloor 200/64\rfloor=3$ --- three pages. After three compulsory
      misses, the rest of every cycle hits. \textbf{A.}
\item Pages 0, 1, 3, 4 --- four pages, all cached after the first
      cycle. \textbf{A.}
\item Pages 0, 15, 31, 46, 62 --- \textbf{five} pages in a cycle of
      length 5. LRU on 4 slots with a cyclic working set of 5 misses
      on \emph{every} access (the lecture's thrash example). \textbf{B.}
\item Addresses $0..6$ are all on page 0. One compulsory miss, then
      all hits. \textbf{A.}
\item Same four pages as (82), reverse order. \textbf{A.}
\end{enumerate}
\end{Sol}

\begin{Qn}{Easy-Medium}{Lect 13}{CS-537 Fa21 Mid, Q21--23}{What a TLB is; reverse translation; strided hit rate}
\begin{enumerate}[(21)]
\item Best description of a TLB (choose one):
\begin{enumerate}[(A)]
\item A memory to store page-table entries.
\item A collection of base/bounds pairs to help segmentation work quickly.
\item An OS component to speed up translation.
\item A hardware feature that ensures fair use of memory.
\item An address-translation cache.
\end{enumerate}
\item A 4-entry TLB (all valid, decimal):
      $\mathrm{VPN}\,0\to\mathrm{PFN}\,1$,
      $1\to 100$, $2\to 101$, $3\to 102$.
      14-bit virtual addresses, 4\,KB pages. What virtual address
      accesses physical address 50 (decimal)?
\begin{enumerate}[(A)]
\item \texttt{0x0032}
\item \texttt{0x1023}
\item \texttt{0x3012}
\item \texttt{0x3132}
\item None of the above.
\end{enumerate}
\item 4\,KB pages; the process accesses every 128th byte on a series of
      pages. TLB starts empty; ignore code. Hit rate?
\begin{enumerate}[(A)]
\item $\approx 99\%$
\item $\approx 88\%$
\item $\approx 50\%$
\item $\approx 17\%$
\item None of the above.
\end{enumerate}
\end{enumerate}
\end{Qn}
\begin{Sol}
\begin{enumerate}[(21)]
\item \textbf{(E).} It is hardware (not an OS component), it caches
      VPN$\to$PPN translations (not arbitrary data, not base/bounds),
      and it has nothing to do with fairness.
\item Offset of PA 50 is 50; that requires \textbf{PFN 0}. No TLB row
      has PFN 0, so no generated VA can produce PA 50. \textbf{(E).}
      (14-bit VA, 12-bit offset $\Rightarrow$ only 4 VPNs exist, and
      they map to frames 1, 100, 101, 102.)
\item Intended reading: bytes $0,128,256,\ldots$ --- $4096/128=32$
      touches per page, first misses, next 31 hit, hit rate
      $31/32\approx 97\%$, which is none of (A)--(D). \textbf{(E).}
      Alternative reading (byte 128 of every page): all compulsory
      misses, hit rate $0\%$, also (E).
\end{enumerate}
\end{Sol}

\begin{Qn}{Easy-Medium}{Lect 13}{CS-537 Sp02 Mid, Q12}{Circle every true statement about a TLB}
A translation lookaside buffer is generally used to (circle \emph{all}
that apply):
\begin{enumerate}[(A)]
\item translate virtual page numbers into physical page numbers
\item translate physical page numbers into virtual page numbers
\item make segmentation have the benefits of a pure paging approach
\item translate the addresses generated by loads
\item translate the addresses generated by stores
\item translate the addresses generated by instruction fetches
\item remove the need for a full-sized page table
\item make translations happen quickly
\end{enumerate}
\end{Qn}
\begin{Sol}
\textbf{(A), (D), (E), (F), (H).}

The TLB is a cache of VPN$\to$PPN translations used on \emph{every}
kind of virtual reference (loads, stores, fetches). It makes those
translations fast when they hit. It does \emph{not} invert the map
(that is an inverted page table), does not magically give segmentation
paging's benefits, and does \emph{not} shrink or eliminate the page
table --- the PT is still the backing store for misses. (Lecture: TLBs
hide the \emph{time} cost of paging, not the \emph{space} cost.)
\end{Sol}

\begin{Qn}{Easy-Medium}{Lect 13}{CS-537 Fa04 Final, Q7}{Proper uses of a TLB (Broken OS)}
Circle every entry that represents a proper use of a TLB:
\begin{enumerate}[(A)]
\item To cache frequently accessed data
\item To cache frequently used address translations
\item To speed up the time to load data from memory
\item To enforce protection across pages
\item To help determine whether a page is resident in memory or not
\end{enumerate}
\end{Qn}
\begin{Sol}
\textbf{(B), (C), (D).}

(A)~is a \emph{data} cache (L1/L2), not a TLB. (B)~is the definition.
(C)~follows: a hit skips the page-table walk, so the subsequent data
load starts sooner. (D)~protection bits (R/W/X) live in the TLB row so
the MMU can deny a write to a read-only page without walking the PT.
(E)~residency is the PTE \textbf{present} bit (Lecture 15); a TLB hit
implies the page is present (else the OS would not have filled the
row), but the TLB is not the structure used to \emph{discover}
residency on a miss.
\end{Sol}

\begin{Qn}{Easy-Medium}{Lect 13}{CS-537 Fa11 Mid, Q15}{Compulsory misses on a 1024-int zeroing loop}
\begin{lstlisting}[numbers=none]
int i;
int p[1024];
for (i = 0; i < 1024; i++)
    p[i] = 0;
\end{lstlisting}
Assume a 1\,KB page size and 4-byte \texttt{int}s, page-aligned
\texttt{p}, and ignore instruction fetches and the loop index.
How many TLB hits and misses occur the \emph{first} time this code runs?
\end{Qn}
\begin{Sol}
The array is $1024\times 4 = 4096$ bytes $= 4$ pages of 1\,KB. Sequential
zeroing touches each page 256 times. Four compulsory misses, then 1020
hits (256 hits on page 0 after its miss, etc.).
Miss rate $4/1024 = 0.39\%$, the same arithmetic as the lecture's
$2/2048 = 0.1\%$ array-walk example.
\end{Sol}

\begin{Qn}{Medium}{Lect 13}{CS-537 Fa04 Mid2, Q7}{Two TLB rows with the same PPN, or the same VPN}
\begin{enumerate}[(a)]
\item When can a TLB contain two entries with the same physical address?
\item When can it contain two entries with the same virtual address?
\item Do either of these cases need extra hardware support?
\end{enumerate}
\end{Qn}
\begin{Sol}
\begin{enumerate}[(a)]
\item \textbf{Sharing}: two VPNs (possibly in two address spaces, or
      two aliases in one) map to the same PFN --- shared library, two
      mappings of a file, copy-on-write before the write. Both rows
      may be cached. Also a process that maps the same frame at two
      VAs.
\item \textbf{Same VPN, different ASID/PID}: process 0 and process 1
      both use VPN 0, mapped to different (or even the same) frames.
      Without ASIDs this must not happen: a flush on context switch
      prevents two live rows with the same VPN. With ASIDs it is the
      normal case (see Fa08 Mid P3, VPN \texttt{00000} for PID 00 and
      01).
\item Same-PPN sharing needs no extra TLB hardware (two tags, one
      PPN). Same-VPN coexistence \textbf{requires an ASID (or PID)
      tag}, plus a match rule $(\mathrm{VPN},\mathrm{ASID})$. Global
      bits (the MIPS $G$ bit) are an extra option so a row matches any
      ASID (kernel mappings).
\end{enumerate}
\end{Sol}

\begin{Qn}{Medium}{Lect 13}{CS-537 Sp12 Mid, Q5}{Timer interrupt driven by TLB misses}
The usual timer interrupt fires every $T$ milliseconds. In a
subtly-different system the CPU is interrupted after every $N$ TLB
misses instead. How does this change the timer and its usefulness?
\end{Qn}
\begin{Sol}
A timer's job is to guarantee that the OS \emph{eventually} regains the
CPU, even from a tight user loop. A TLB-miss counter fails at that:
a process whose working set fits in the TLB (a tight arithmetic loop
on a handful of pages) generates \textbf{zero} misses and can run
forever --- the OS loses control. Conversely, a process with a large
strided working set is interrupted constantly, so it is punished for
being memory-unfriendly. The mechanism also requires a live paging
system (no TLB misses on a machine without paging). Compared to a
real timer it is therefore \textbf{risky}, not equally effective.
(Same idea: CS-537 Fa21 Mid Q28, answer ``risky''.)
\end{Sol}

\begin{Qn}{Medium}{Lect 13}{CS-537 Sp12 Mid, Q6}{A TLB with no valid bit}
A TLB often has a valid bit per row. In this variant there is no such
bit. Implications?
\end{Qn}
\begin{Sol}
The valid bit means ``ignore this row''. Without it:
\begin{itemize}
\item You cannot \textbf{flush} by clearing bits; you must overwrite
      every row with a translation that will never match (e.g.\ a
      reserved VPN, or an ASID that is never allocated), or the MMU
      must treat ``VPN $=0$ and PPN $=0$'' as empty by convention.
\item At boot / after a context-switch flush the TLB may contain
      garbage that spuriously matches a VPN and returns a wild PFN ---
      a serious protection bug.
\item Hardware needs some other empty-marker (a dedicated ASID, a
      match-never pin, or a forced rewrite of all $N$ rows).
\end{itemize}
The PTE valid bit is a different bit: it says whether the \emph{virtual
page} is mapped. The TLB valid bit says whether the \emph{row} holds a
cached translation.
\end{Sol}

\begin{Qn}{Medium}{Lect 13}{CS-537 Sp12 Mid, Q7}{Lazy TLB fill: install only on the third miss}
A software-managed TLB does not install a translation on the first miss
of a VPN. It increments a counter in the PTE; only when the counter
reaches 3 does it fill the TLB. How does this lazy update change
behaviour?
\end{Qn}
\begin{Sol}
Each distinct VPN now suffers \textbf{three} miss traps before it is
cached (and the counter occupies extra PTE bits). Cold-start and
one-shot touches become more expensive; a page that is touched twice
and never again never occupies a TLB slot, which can help a small TLB
if the workload has many ephemeral pages (a crude filter for
compulsory-only pages). Once a page is ``hot'' ($\ge 3$ touches) the
steady-state hit rate matches an ordinary TLB. Replacement of TLB rows
is unchanged. Hardware-managed walkers cannot easily do this: they
fill on the first miss by construction.
\end{Sol}

\begin{Qn}{Easy-Medium}{Lect 13}{CS-537 Sp09 Final, Q2}{What happens on a software-managed TLB miss?}
A system uses 1\,MB pages, linear page tables, and a software-managed
TLB. Describe what happens on a TLB miss. What structures (hardware
and software) are accessed?
\end{Qn}
\begin{Sol}
The MMU traps (no hardware PT walk). The OS miss handler, running in
kernel mode on the kernel stack, reads PTBR/CR3 from the PCB/hardware,
indexes the linear table at
$\mathrm{PTBR}+\mathrm{VPN}\times\mathrm{sizeof(PTE)}$, and:
\begin{itemize}
\item if the PTE is invalid: deliver SIGSEGV / kill the process;
\item if valid but not present: page-fault path (disk, Lecture 15);
\item if valid and present: choose a TLB victim (LRU or random), write
      $(\mathrm{VPN},\mathrm{PPN},\mathrm{prot}[,\mathrm{ASID}])$ into
      that row, return-from-trap.
\end{itemize}
The hardware then retries the instruction, hits, and issues the
physical load/store. Structures touched: trap hardware (IDT / exception
vector), kernel stack, the linear page table in DRAM, the TLB SRAM, and
(on a fault) the swap device. Large 1\,MB pages make the linear table
small (a 32-bit space has only $2^{12}$ VPNs) and make TLB reach huge,
which is the only reason a linear table plus a tiny TLB is even
plausible here.
\end{Sol}

\begin{Qn}{Easy-Medium}{Lect 13}{CS-537 Fa21 Mid, Q28}{Interrupt the CPU after every $N$ TLB misses}
ZOS interrupts the CPU after every $N$ TLB misses instead of using a
periodic timer. As compared to a timer, this approach
\begin{enumerate}[(A)]
\item is equally effective
\item is faster to program
\item is risky
\item requires less memory
\item requires virtual memory support
\end{enumerate}
\end{Qn}
\begin{Sol}
\textbf{(C)} (and (E) is also true, but the best comparison-to-timer
answer on the key is ``risky''). A loop that never misses the TLB never
yields. See Sp12 Mid Q5 in this section.
\end{Sol}

\begin{Qn}{Easy-Medium}{Lect 13}{CS-537 Fa21 Mid, Q30}{Why large pages?}
Regular pages are 1\,KB; large pages are 1\,MB. When a contiguous
aligned region is in use, the OS prefers one large page. Why is that a
good idea?
\begin{enumerate}[(A)]
\item They reduce system complexity.
\item They speed up trap handling.
\item They can increase TLB hit rates.
\item They make physical memory allocation easier.
\item None of the above.
\end{enumerate}
\end{Qn}
\begin{Sol}
\textbf{(C).} One TLB row now covers 1\,MB instead of 1\,KB, so TLB
reach grows by $1024\times$ and miss rates drop for sequential / large
working sets. Complexity and physical allocation both get \emph{harder}
(two page sizes, alignment, fragmentation); trap handling is essentially
unchanged.
\end{Sol}

\begin{Qn}{Easy-Medium}{Lect 13}{CS-537 Sp09 Final, Q3}{Justifying large pages plus linear page tables}
The kilOS team uses 1\,MB pages and linear page tables on machines that
ship with large amounts of memory. Justify the decision.
\end{Qn}
\begin{Sol}
Linear-table size is $(\#\mathrm{VPNs})\times(\mathrm{PTE\ size})$.
Raising the page size from 4\,KB to 1\,MB cuts the VPN by 8 bits, so
the table shrinks by $256\times$ (a 32-bit, 4-byte-PTE table goes from
4\,MB to 16\,KB). With plentiful DRAM, 16\,KB per process is cheap, and
the table is small enough to keep contiguous. TLB reach grows by the
same $256\times$, so a 64-entry TLB covers 64\,MB instead of 256\,KB.
The costs --- internal fragmentation, 1\,MB disk transfers on a fault
--- are acceptable when memory is large and the working set is not
full of tiny objects. This is the lecture's ``why large pages exist''
paragraph, applied to linear tables.
\end{Sol}

\begin{Qn}{Easy-Medium}{Lect 14}{CS-537 Sp18 Mid, Problem XXX}{Five true/false claims about multi-level page tables}
Fill \textbf{A} if true, \textbf{B} if false.
\begin{enumerate}[(146)]
\item A multi-level page table may use more pages than a linear page table.
\item It is easier to allocate pages of the page table in a multi-level
      table (as compared to a linear page table).
\item Multi-level page-table lookups take longer than linear page-table
      lookups.
\item With larger virtual address spaces, usually more levels are used.
\item TLBs are useful in making multi-level page tables even smaller.
\end{enumerate}
\end{Qn}
\begin{Sol}
\begin{enumerate}[(146)]
\item \textbf{A (true).} If every virtual page is valid, you still store
      every PTE \emph{plus} the directory page(s). Fully populated
      multi-level is strictly larger than linear.
\item \textbf{A (true).} Linear tables must be one contiguous
      allocation of (e.g.) 4\,MB. Multi-level tables are page-sized
      chunks, allocated only where VPNs are used, from the ordinary
      free-frame list.
\item \textbf{A (true).} A $k$-level walk is $k$ memory references
      before the data; linear is one. (TLB hits hide this; they do not
      change the miss-path length.)
\item \textbf{A (true).} Each level's table is sized to fit in a page.
      More VPN bits $\Rightarrow$ more levels (IA-32: 2; x86-64 48-bit:
      4; a full 64-bit 4\,KB/4\,B design would need $\sim 6$).
\item \textbf{B (false).} TLBs make translation \emph{faster}; they do
      not store fewer PTEs. Space reduction is the job of inverted /
      segmented / multi-level tables.
\end{enumerate}
\end{Sol}

\begin{Qn}{Easy-Medium}{Lect 13--14}{CS-537 Fa09 Mid, Q6(a)--(g)}{T/F: why multi-level tables and TLBs exist}
Circle True or False. (VMM items (h)--(j) of the original exam are
omitted as out of syllabus.)
\begin{enumerate}[(a)]
\item The main reason to have a multi-level page table is to speed up
      address translation.
\item The main reason to have a hardware TLB is to speed up address
      translation.
\item Using a multi-level page table increases TLB hit time.
\item As the amount of addressable physical memory grows, the size of
      each PTE must grow as well.
\item As the size of the virtual address space grows, the amount of
      space occupied by a linear page table also grows.
\item In a 2-level multi-level page table, each page-directory entry
      should have a valid bit.
\item In a 2-level multi-level page table, the size of the page
      directory increases as the virtual address space gets larger.
\end{enumerate}
\end{Qn}
\begin{Sol}
\begin{enumerate}[(a)]
\item \textbf{False.} The main reason is to \textbf{save memory} (and
      to avoid a contiguous allocation) by omitting inner tables for
      unused VPN ranges. Multi-level walks are \emph{slower}.
\item \textbf{True.} The TLB caches VPN$\to$PPN so the common case is
      one array access, not a PT walk.
\item \textbf{False.} TLB hit time is a hardware CAM/SRAM lookup; it
      does not depend on how the backing page table is shaped.
\item \textbf{True} (generally): the PFN field must address more
      frames, so the PTE widens (unless it already had spare bits).
\item \textbf{True.} $\#\mathrm{PTEs}=2^{\mathrm{VA\ bits}-\lg(\mathrm{page})}$.
\item \textbf{True.} An invalid PDE means ``no inner table allocated
      for this VPN range'' --- that is the space-saving mechanism.
      (PTE valid means ``this virtual page is mapped'', a different
      statement.)
\item \textbf{True.} More VPN bits beyond those used to index a
      page-sized leaf must be decoded by more PDEs (or by another
      level).
\end{enumerate}
\end{Sol}

\begin{Qn}{Easy-Medium}{Lect 14}{CS-537 Fa13 Final, warmup}{Worst-case page-table references on a two-level TLB miss}
What is the worst-case number of memory accesses required to do a
\textbf{page-table lookup} on a TLB miss on a system with a two-level
multi-level page table?
\end{Qn}
\begin{Sol}
\textbf{Two}: one for the page directory (PDE) and one for the leaf PTE.
(The subsequent data access is a third memory reference but is not part
of the \emph{page-table lookup}.) Three-level: 3; four-level: 4. A TLB
hit: 0 page-table accesses.
\end{Sol}

\begin{Qn}{Medium}{Lect 14}{CS-537 Fa11 Mid, Q16}{Space saved by a sparse two-level directory}
A two-level table: the page directory is a single page of 32 PDEs
(valid bit $+$ PFN of a leaf). Each leaf is one page of PTEs. One
process's directory (32 bytes, high bit $=$ valid) contains 8 valid
PDEs and 24 invalid PDEs. Page size is 32 bytes.

How much space does this multi-level table save, in bytes, compared to
a typical linear page table for the same parameters?
\end{Qn}
\begin{Sol}
Linear table: 32 PDEs $\times$ (one page of PTEs each) would be 32 leaf
pages if stored densely, i.e.\ $32\times 32\,\mathrm{B}=1024\,\mathrm{B}$
of PTEs (plus no directory). Multi-level: 1 directory page $+$ 8 leaves
$= 9\times 32 = 288\,\mathrm{B}$. Savings:
$32-8=24$ leaf pages not allocated $= 24\times 32 = 768\,\mathrm{B}$.
(If one also charges the linear table no directory and the multi-level
table the extra 32-byte directory, the net saving is
$1024-(256+32)=736\,\mathrm{B}$. The exam's expected accounting is the
24 omitted leaf pages, 768 bytes.)
\end{Sol}

\begin{Qn}{Medium}{Lect 13--14}{CS-537 Fa04 Mid2, Part II Q2}{Two page sizes: costs, TLB reach, and PT changes}
A system supports a small page size (4\,KB) and a big one (1\,MB).
\begin{enumerate}[(a)]
\item If we have only a single page size and it is quite big, what
      hurts? If it is quite small, what hurts?
\item Using big pages can improve performance. Which hardware resources
      are better utilized?
\item To support multiple page sizes, what must change in a simple
      linear page table, and how would you implement it?
\item The OS first hands out small pages, then periodically tries to
      promote a contiguous aligned batch into one big page. What must
      be true, and what structures are updated?
\end{enumerate}
\end{Qn}
\begin{Sol}
\begin{enumerate}[(a)]
\item Big: internal fragmentation; expensive faults (move 1\,MB
      to/from disk); harder allocation of contiguous aligned frames.
      Small: huge page tables; tiny TLB reach $\Rightarrow$ more
      compulsory TLB misses; more PTEs to scan for replacement.
\item The \textbf{TLB} (reach $= N\times$ page size) and, on a fault,
      disk sequential bandwidth (one large transfer vs.\ many small
      ones). Caches may also see fewer compulsory misses if the
      large page is used densely.
\item Each PTE must record the page size (a ``large'' bit) \emph{or}
      the table must be split into a 4\,KB table and a 1\,MB table
      (two PTBRs). A linear 4\,KB table cannot index 1\,MB pages
      without a different VPN/offset split; typical hardware uses
      the PDE large-page bit of Fa22 Mid Q3: a directory entry may
      point at a 1\,MB frame instead of at a 4\,KB leaf.
\item The $256$ consecutive 4\,KB VPNs that form a 1\,MB aligned
      region must all be valid, map to $256$ consecutive aligned
      physical frames, and have compatible protection. Then: allocate
      (or already hold) the 1\,MB frame, rewrite the PDE with the
      large-page bit, free the 256 leaf PTEs / small frames, and
      \textbf{flush} (or demote) any TLB rows for those small VPNs,
      installing one large-page TLB entry instead.
\end{enumerate}
\end{Sol}

\begin{Qn}{Easy-Medium}{Lect 13}{CS-537 Fa13 Mid, Q6}{OS involvement in a hardware-managed paging system}
McFlub: ``In a hardware-managed paging system, the OS has to get
involved sometimes. For example, when a process starts running, the OS
has to set the page-table base register (PTBR) to point to the page
table of the process. Actually, that's about it for OS involvement with
paging.''

What did McFlub miss?
\end{Qn}
\begin{Sol}
PTBR (CR3 on Intel) is necessary but nowhere near sufficient. The OS
must also:
\begin{itemize}
\item allocate and initialize a page table on process creation, and
      free it on exit;
\item \textbf{switch PTBR on every context switch} (and, without ASIDs,
      flush the TLB);
\item grow/shrink tables on \texttt{sbrk}/\texttt{mmap}, stack
      expansion, and \texttt{munmap};
\item handle page faults (present bit, swap) and protection faults;
\item set and harvest referenced/dirty bits for replacement;
\item shoot down or update TLB rows when a mapping changes (even on a
      hardware-walked TLB, a PTE store is not enough until the old
      translation is invalidated).
\end{itemize}
What the OS does \emph{not} do on a hardware-managed TLB: walk the
page table on a miss to a \emph{present} page --- the MMU does that via
CR3. On a software-managed TLB that walk \emph{is} OS code.
\end{Sol}

\subsection{Lecture 15: Swapping \& Page Replacement Policies}

\subsubsection*{Additional Challenging \& Analytical Problems}

\begin{Qn}{Hard}{Lect 15}{CS-537 Fa22 Mid, Q25--32}{2Q: a FIFO filter in front of an LRU queue}
xOS 2Q: two queues sharing a 4-page cache. A brand-new page enters the
\textbf{bottom} queue at the first-in end (bottom is FIFO). A hit on a
bottom-queue page \textbf{promotes} it to the MRU end of the
\textbf{top} queue (top is LRU). A page evicted from the top is demoted
to the first-in end of the bottom. Empty frames of a queue are filled
before that queue evicts.
\begin{enumerate}
  \item[25.] 2 pages top + 2 bottom. Trace $1,2,3,4,1,2,3,4$. Misses?
        (a) 0 (b) 2 (c) 4 (d) 8 (e) none of the above
  \item[26.] Same trace, plain LRU with 4 pages. Misses?
        (a) 0 (b) 4 (c) 6 (d) 8 (e) none of the above
  \item[27.] 2+2 2Q, trace $1,2,1,2,3,4,5,6,1,2,1,2,1,2$. Misses?
        (a) 0 (b) 1 (c) 4 (d) 8 (e) none of the above
  \item[28.] Same trace, plain LRU-4. Misses?
        (a) 0 (b) 1 (c) 4 (d) 8 (e) none of the above
  \item[29.] Trace $1,2,3,1,2,3,4,5,6,7,1,2,3,1,2,3$. Which split of 4
        pages, top (LRU) / bottom (FIFO), maximises hits?
        (a) 4/0 (b) 3/1 (c) 2/2 (d) 1/3 (e) 0/4
  \item[30.] Both queues FIFO, 2+2, trace $1,2,3,4,1,2,3,4$. Misses?
        (a) 0 (b) 1 (c) 2 (d) 4 (e) none of the above
  \item[31.] FIFO/FIFO 2+2, trace $1,2,1,2,3,4,5,6,5,6,1,2$. Misses?
        (a) 2 (b) 4 (c) 6 (d) 8 (e) none of the above
  \item[32.] Compared with a single queue, 2Q is always better / always
        worse / identical / simpler / none of the above?
\end{enumerate}
\end{Qn}
\begin{Sol}
Notation: top $[\,]_{\mathrm{LRU}}$ (LRU left, MRU right), bottom
$[\,]_{\mathrm{FIFO}}$ (oldest left).

25. $1,2$ fill bottom; $3$ evicts 1; $4$ evicts 2; the repeat $1,2,3,4$
never finds a page still in bottom, so nothing is ever promoted.
\textbf{8 misses (d)}. (A looping working set of size 4 is the FIFO
filter's worst case.)

26. LRU-4: four compulsory misses, then four hits. \textbf{4 (b)}.

27. $1*,2*$ bottom$=[1,2]$; $1$ promotes $\to$ top$=[1]$, bottom$=[2]$;
$2$ promotes $\to$ top$=[1,2]$, bottom empty. Then $3*,4*,5*,6*$ only
touch bottom (evicting 3 then 4). Then $1,2,1,2,1,2$ all hit in top.
Misses: $1,2,3,4,5,6$ --- \textbf{6}, which is not 0, 1, 4 or 8
$\to$ \textbf{(e)}.

28. LRU-4 on the same string: $1*,2*$, hits on the repeats, then
$3*,4*,5*$ (evict 1), $6*$ (evict 2), then $1*,2*$, then hits.
\textbf{8 (d)}. So 2Q (6) beats LRU (8) on a working set that is
referenced, then blown out by a scan, then reused: the FIFO filter
absorbs the scan.

29. Direct count of hits (16 references):
\begin{center}
\begin{tabular}{lcc}
\hline
split & misses & hits \\
\hline
4 top / 0 bottom (pure LRU) & 10 & 6 \\
3 / 1 & 16 & 0 \\
2 / 2 & 16 & 0 \\
1 / 3 &  9 & 7 \\
0 / 4 (pure FIFO) & 10 & 6 \\
\hline
\end{tabular}
\end{center}
Most hits: \textbf{(d) 1 in top, 3 in bottom} (7 hits). The 3-page FIFO
filter holds the looping triple $\{1,2,3\}$ through the $4,5,6,7$ scan
long enough that the 1-page LRU can pin the promoted page.

30. No promotions (same as 25): \textbf{8 misses $\to$ (e)}.

31. $1*,2*$, promote 1 then 2 (top$=[1,2]$), $3*,4*,5*,6*$ (bottom
ends $[5,6]$), then $5,6$ promote (evicting 1 then 2 from top into
bottom), then $1,2$ promote back. Misses: $1,2,3,4,5,6$ ---
\textbf{6 (c)}.

32. 2Q is more complex and is better or worse depending on the
workload (25 vs.\ 26 vs.\ 27 vs.\ 28). Not always better, worse,
identical, or simpler. \textbf{(e) none of the above}.
\end{Sol}

\begin{Qn}{Medium-Hard}{Lect 15}{CS-537 Fa04 Mid2, Q6}{TLB hit, TLB miss, and page-fault control flow}
A system uses pure paging and a \textbf{software-managed} hardware TLB
(the OS, not the MMU, installs TLB entries). Describe, step by step:
\begin{enumerate}
  \item[(a)] what happens on a TLB hit;
  \item[(b)] what happens on a TLB miss;
  \item[(c)] what happens on a page fault.
\end{enumerate}
\end{Qn}
\begin{Sol}
(a) \textbf{TLB hit.} Hardware matches the VPN (and ASID, if any), checks
protection bits, forms $\mathrm{PA}=(\mathrm{PFN}\ll\mathrm{offset})+\mathrm{offset}$,
and completes the load/store. The OS is not entered.

(b) \textbf{TLB miss (page present).} Hardware traps. The OS walks the page
table. If valid$\,=1$ and present$\,=1$, it installs
$(\mathrm{VPN}\to\mathrm{PFN},\mathrm{prot})$ in the TLB, returns from the
trap, and the hardware \textbf{retries} the instruction --- which now TLB-hits.
If valid$\,=0$, this is a protection/segmentation fault, not a page fault.

(c) \textbf{TLB miss + page fault} (valid$\,=1$, present$\,=0$). The OS
must run disk I/O, so it always handles this path.
\begin{enumerate}
  \item Read the disk-block number out of the PTE.
  \item If no free frame, choose a victim. If the victim is
        \textbf{dirty}, write it back to swap first; a clean victim is
        discarded.
  \item Issue the read of the missing page into the (now free) frame.
        The faulting process is \textbf{blocked}; the scheduler runs
        another ready process (disk I/O is millisecond-scale).
  \item On the disk interrupt: set PFN, present$\,=1$, clear dirty,
        (usually) set the reference bit; invalidate any stale TLB entry
        for the victim; optionally fill the TLB for the new page.
  \item Make the process ready. When it is rescheduled, the instruction
        is \textbf{restarted} (this is why the hardware must deliver a
        precise interrupt).
\end{enumerate}
The three outcomes in one line: TLB hit (hardware only);
TLB miss + present (OS fills TLB, no disk); TLB miss + not-present
(OS page-fault handler + block + restart).
\end{Sol}

\begin{Qn}{Medium-Hard}{Lect 15}{CS-537 Fa04 Mid2, Q8}{Clock / second-chance: mechanism, hardware, and relation to LRU}
\begin{enumerate}
  \item[(a)] Describe how clock works. What hardware support is needed?
        What software structures must be kept?
  \item[(b)] Can clock ever behave exactly like perfect LRU?
\end{enumerate}
\end{Qn}
\begin{Sol}
(a) Frames are arranged as a circular list with a clock hand. Each PTE
(or frame-table entry) has a hardware \textbf{use/reference} bit.
\begin{itemize}
  \item On every access the \textbf{hardware} sets the page's use bit
        to 1. The OS is not involved.
  \item On a fault, if a free frame exists, use it. Otherwise the OS
        walks from the hand: if use$\,=1$, the OS \textbf{clears} it to
        0 and advances (the page gets a second chance); if use$\,=0$,
        that frame is the victim. The new page is installed with
        use$\,=1$ and the hand advances past it.
\end{itemize}
Software: a circular frame list (or an index into the frame table) and
the current hand position. Hardware: one use bit per PTE, set
automatically on a hit (the same bit the MMU already uses for
accessed/dirty accounting).

If every use bit is 1, the hand sweeps a full circle, clears them all,
and evicts the frame it started at --- clock \textbf{degenerates to
FIFO}.

(b) \textbf{Only in special cases}, not in general. Clock has no total
order on recency, only a 1-bit age. Counterexamples:
\begin{itemize}
  \item All-bits-1 degeneration to FIFO, which is not LRU.
  \item A page referenced just after the hand passed (use already
        cleared) looks as cold as a page that has been unused for a
        full rotation; LRU would distinguish them.
\end{itemize}
Clock \emph{equals} LRU on some short traces (e.g.\ when at most one
use bit is set, so the unique 0-bit page is exactly the LRU page), but
it is an approximation, not an implementation, of LRU.
\end{Sol}

\begin{Qn}{Medium-Hard}{Lect 15}{CS-537 Fa10 Mid, Q5}{Why approximate LRU, how reference bits work, and bit-flip damage}
Iron OS approximates LRU with reference bits. Page-table entries also
carry a present bit.
\begin{enumerate}
  \item[(a)] Why approximate LRU rather than implement it directly?
  \item[(b)] How does a reference-bit scheme work?
  \item[(c)] Some reference bits flip by accident. What goes wrong? Does
        the system remain correct?
  \item[(d)] How do present bits work, and why are they needed?
  \item[(e)] Some present bits flip by accident. What goes wrong? Does
        the system remain correct?
\end{enumerate}
\end{Qn}
\begin{Sol}
(a) Perfect LRU needs either a list update on every reference (too much
work in the hit path) or a scan of every frame's timestamp at
replacement (too much work in the fault path). LRU itself is already
only a guess at OPT. Finding \emph{an old} page is enough; clock does
that with one hardware bit and a cheap sweep.

(b) Hardware sets the bit on access. The OS periodically (or during the
clock sweep) clears bits. A victim is a page whose bit is still 0 ---
not referenced since the last clear. Clock, second-chance, and the
additional-reference-bits / aging register are all instances.

(c) $1\to 0$: a hot page looks cold and may be evicted, causing extra
faults. $0\to 1$: a cold page looks hot and is skipped, so a hotter
page may be evicted instead. Replacement is \textbf{suboptimal but
semantically correct}: the bit is only a hint. No wrong translation is
installed.

(d) Present$\,=1$: PFN is a frame. Present$\,=0$: the field is a swap
address; an access is a page fault. Needed because a valid page may
live on disk; without the bit the MMU would treat a disk block number
as a frame number.

(e) This is \textbf{not} correctable by ``just extra faults''.
Present $1\to 0$: a resident page looks swapped; the fault handler
reads ``from disk'' using the PFN as a block number --- garbage, or a
panic. Present $0\to 1$: a swapped page looks resident; the disk
address is used as a PFN and the process reads/writes some other
process's frame. Either flip can corrupt memory or crash the kernel.
\end{Sol}

\begin{Qn}{Medium-Hard}{Lect 15}{CS-537 Sp01 Final, Q3(a)--(c)}{LRU intuition, the cyclic worst case, and a policy that beats it}
\begin{enumerate}
  \item[(a)] Describe LRU and the intuition for using it.
  \item[(b)] Give a worst-case scenario for LRU.
  \item[(c)] What replacement policy works well on that scenario?
\end{enumerate}
(Part (c) of the same MEMS-themed exam asked for a hybrid that switches
policies; the in-syllabus core is (a)--(c). The looping-scan worst case
also appears as CS-537 Fa11 Mid, Q18.)
\end{Qn}
\begin{Sol}
(a) On a miss, evict the page whose last reference is the oldest. Past
recency is used as a predictor of future recency (temporal locality):
a page used recently is likely to be used again soon, so it is a bad
victim; a page unused for a long time is a good victim. LRU is the
standard online stand-in for OPT (Belady's MIN), which evicts the page
used furthest in the \emph{future}.

(b) A cyclic scan of $N+1$ pages with only $N$ frames:
$1,2,\ldots,N,N+1,1,2,\ldots$. Just before each new page the LRU page
is exactly the one that will be needed next. Hit rate is 0 after the
cold start. (The same pattern also makes FIFO fault on every access;
OPT, which evicts the page just used because it is the furthest in the
future, also faults on every access here.)

(c) \textbf{MRU} (evict the most recently used page) keeps the $N-1$
oldest pages resident and takes almost $N-1$ hits per cycle.
\textbf{Random} also beats LRU in expectation, because it sometimes
evicts a page that is \emph{not} the next one. A practical hybrid
(the exam's (d)--(e)) watches the miss rate: if it spikes on a looping
scan, switch from LRU/clock to MRU or random until the miss rate
drops.
\end{Sol}

\begin{Qn}{Medium-Hard}{Lect 15}{CS-537 Fa04 Final, Q15}{Broken LRU: which traces a 2-frame cache may not produce}
A Broken-LRU (BLRU) 2-page cache produced the following
hit/miss labels. For each sequence say whether the labels are possible
under true LRU (``LRU'') or not (``Not LRU''), and name the first
violation.
\begin{enumerate}
  \item[(a)] A miss, B miss, C miss, B miss
  \item[(b)] A miss, B miss, A miss
  \item[(c)] A miss, B miss, C miss, B hit, C hit
  \item[(d)] A miss, B miss, A hit, A hit, B hit, C miss, B miss
  \item[(e)] A miss, B miss, C miss, A miss, B miss, C miss
\end{enumerate}
\end{Qn}
\begin{Sol}
True LRU, 2 frames, LRU on the left, MRU on the right.

(a) $A*$ $\{A\}$, $B*$ $\{A,B\}$, $C*$ evict $A$ $\{B,C\}$, then $B$
is resident so it must be a \textbf{hit}. The reported miss is
impossible. \textbf{Not LRU}.

(b) After $A*,B*$ the cache is $\{A,B\}$, so the next $A$ is a hit.
Reported miss: \textbf{Not LRU}.

(c) $A*,B*,C*$ $\to\{B,C\}$, then $B$ hit $\{C,B\}$, $C$ hit $\{B,C\}$.
All labels match. \textbf{LRU}.

(d) $A*,B*$ $\{A,B\}$; $A$ hit $\{B,A\}$; $A$ hit; $B$ hit $\{A,B\}$;
$C*$ evict $A$ $\{B,C\}$; $B$ is MRU so the next $B$ is a
\textbf{hit}. Reported miss: \textbf{Not LRU}.

(e) Cyclic scan of 3 pages with 2 frames: LRU (and FIFO) miss on every
access after the first two. \textbf{LRU} (this is the LRU worst case).
\end{Sol}

\begin{Qn}{Medium-Hard}{Lect 15}{CS-537 Fa04 Mid2, Part II Q1(c)--(d)}{Should LRU change when RAM is split into a fast tier and a slow tier?}
A machine has two kinds of physical memory behind one physical address
space: a fast primary region on the memory bus and a slower secondary
region on the I/O bus (still orders of magnitude faster than disk). The
OS can see the boundary.
\begin{enumerate}
  \item[(c)] You decide to use LRU for the page cache. Describe LRU.
  \item[(d)] Should the basic LRU algorithm change on this hardware?
\end{enumerate}
\end{Qn}
\begin{Sol}
(c) Resident pages are totally ordered by last-reference time. A hit
moves the page to MRU. A miss with a full cache evicts the LRU page
(writeback first if dirty). Intuition: recency $\approx$ future use.

(d) \textbf{Yes.} One LRU list treats every frame as equal, but a
primary-frame victim costs much more to re-fill than a secondary-frame
victim. A cost-aware policy:
\begin{itemize}
  \item Prefer evicting from the \textbf{slow} tier (or evicting a
        \textbf{clean} slow page) before evicting a hot primary page.
  \item Use two clocks / two LRU lists, and \textbf{promote} a page
        from secondary to primary on reuse (and demote a cold primary
        page to secondary instead of to disk).
  \item Disk remains the last resort; secondary RAM is an intermediate
        victim cache.
\end{itemize}
Unmodified LRU is still correct, just not sensitive to the two-level
miss penalty. (The same idea is why clock with a dirty bit prefers
$(0,0)$ then $(0,1)$ before evicting a recently used page.)
\end{Sol}

\subsubsection*{Conceptual Drills, MCQs \& Further Practice}

\begin{Qn}{Easy-Medium}{Lect 15}{CS-537 Sp18 Mid, Q86--90}{True/false: OPT, LRU, FIFO, Random, and Belady}
Which statements about page-replacement policies are true? Mark A for
true, B for false.
\begin{enumerate}
  \item[86.] The LRU policy always outperforms the FIFO policy.
  \item[87.] The OPT (optimal) policy always performs at least as well as
        LRU.
  \item[88.] A bigger cache's hit percentage is always $\ge$ a smaller
        cache's hit percentage, if they use the same replacement policy.
  \item[89.] A bigger cache's hit percentage is always $\ge$ a smaller
        cache's hit percentage, if they use LRU.
  \item[90.] Random replacement is always worse than LRU.
\end{enumerate}
\end{Qn}
\begin{Sol}
86. \textbf{B (false).} LRU has a well-known looping worst case (scan
$N+1$ pages with $N$ frames) on which FIFO is no better, and there exist
strings where FIFO records strictly fewer misses than LRU. ``Always
outperforms'' is too strong.

87. \textbf{A (true).} OPT (Belady's MIN) minimises the number of misses
over the future reference string, so no online policy --- LRU included
--- can beat it on hit rate. (Ties are allowed; ``at least as well''
holds.)

88. \textbf{B (false).} Counterexample: Belady's anomaly for FIFO
(CS-537 Fa04 Mid2 Q10 / Mandar 2012 Q3: 9 misses with 3 frames, 10 with
4). The claim would require the stack property.

89. \textbf{A (true).} LRU \emph{is} a stack algorithm: the $k$-frame
resident set is always the $k$ most recently used pages, a subset of
the $(k+1)$-frame set. Every $k$-frame hit is a $(k+1)$-frame hit, so
the hit rate is monotone in the number of frames (it may stay equal).

90. \textbf{B (false).} Random can beat LRU on the cyclic-scan workload
(it sometimes evicts a page that is \emph{not} the next one). Luck of
the draw, but not ``always worse''.
\end{Sol}

\begin{Qn}{Easy-Medium}{Lect 15}{CS-537 Sp18 Mid, Q91--95}{LRU hits and misses in a 4-page memory}
A memory holds 4 pages and uses LRU. The first four references are
$6,7,7,9$. Of the next five references $7,9,0,4,9$, which hit? Mark A
for a hit, B for a miss.
\begin{enumerate}
  \item[91.] 7 \quad 92.\ 9 \quad 93.\ 0 \quad 94.\ 4 \quad 95.\ 9
\end{enumerate}
\end{Qn}
\begin{Sol}
After priming (LRU left, MRU right):
$6*$ $[6]$, $7*$ $[6,7]$, $7$ hit $[6,7]$, $9*$ $[6,7,9]$.
Then
\begin{center}
\begin{tabular}{ccl}
\hline
ref & result & cache after \\
\hline
7 & hit  & $[6,9,7]$ \\
9 & hit  & $[6,7,9]$ \\
0 & miss & $[6,7,9,0]$ \\
4 & miss & $[7,9,0,4]$ \quad (evict 6) \\
9 & hit  & $[7,0,4,9]$ \\
\hline
\end{tabular}
\end{center}
91 A, 92 A, 93 B, 94 B, 95 A.
(Official key, CS-537 Sp18 Mid answers, Problem XIX.)
\end{Sol}

\begin{Qn}{Easy-Medium}{Lect 15}{CS-537 Sp16 Mid, Q5}{FIFO vs LRU traces, small and huge caches}
Virtual-page trace (memory empty at the start, all pages begin on swap):
\[
0,\;1,\;4,\;0,\;1,\;3,\;0,\;1,\;4,\;1.
\]
For each access write H (resident) or M (must be fetched from swap):
\begin{enumerate}
  \item[(a)] FIFO, 3 frames;
  \item[(b)] FIFO, 5 frames;
  \item[(c)] LRU, 3 frames;
  \item[(d)] LRU, 1000 frames.
\end{enumerate}
\end{Qn}
\begin{Sol}
(a) FIFO-3: M M M H H M M M M H \quad (7 misses).
After $0,1,4$ the cache is $[0,1,4]$; $3$ evicts 0; $0$ evicts 1; $1$
evicts 4; $4$ evicts 3; last $1$ hits.

(b) FIFO-5: only four distinct pages $\{0,1,3,4\}$, five frames, so
never a capacity miss: M M M H H M H H H H \quad (4 misses).

(c) LRU-3: M M M H H M H H M H \quad (5 misses).
$0,1,4$ fill; hits on $0,1$; $3$ evicts 4; hits on $0,1$; $4$ evicts 3;
hit on $1$.

(d) LRU-1000: same 4 compulsory misses as (b): M M M H H M H H H H.

The 3-vs-5 comparison for FIFO is \emph{not} Belady (misses drop);
Belady needs a string whose FIFO-4 resident set is not a superset of
the FIFO-3 set.
\end{Sol}

\begin{Qn}{Easy-Medium}{Lect 15}{CS-537 Fa11 Mid, Q17--19}{Grade Ime's LRU answers}
Fall 2011 is a grading exam: Ime A.\ Stoodant answered the three LRU
questions below. For each, decide whether Ime is correct; if not, give
the right answer.
\begin{enumerate}
  \item[17.] ``LRU is considered a good page-replacement policy. Describe
        what a good policy should do, and what LRU is.''
        \emph{Ime:} ``A good policy minimises the miss rate by keeping
        popular pages in memory. LRU does this by keeping the
        least-recently-used pages in memory.''
  \item[18.] ``LRU also has its downsides. Describe a case where LRU
        behaves really poorly.''
        \emph{Ime:} ``Trick question, LRU never does poorly.''
  \item[19.] ``LRU is hard to implement perfectly. Instead most systems
        use reference bits to approximate LRU. How are reference bits
        used?''
        \emph{Ime:} ``The reference bit in the PTE is set when the page
        is accessed. To evict, the OS scans and finds a page whose
        reference bit is 0.''
\end{enumerate}
\end{Qn}
\begin{Sol}
17. \textbf{Wrong} (half-right). A good policy does minimise misses by
keeping pages that will be used soon. LRU predicts that from the
\emph{past}: it keeps the \textbf{most} recently used pages and
\textbf{evicts} the least recently used page. Ime swapped the
quantifier.

18. \textbf{Wrong.} Cyclic scan of $N+1$ pages with $N$ frames:
$1,2,\ldots,N+1,1,2,\ldots$. Just before each access the LRU page is
the next one needed; hit rate drops to 0. (Same answer as CS-537 Sp01
Final Q3(b).)

19. \textbf{Mostly right, incomplete.} Hardware sets the bit on a
reference; a 0-bit page is a candidate victim. Missing pieces that a
full answer names: the OS must \textbf{clear} bits (otherwise they all
stick at 1), and the scan is organised as a \textbf{clock} so that a
1-bit page is given a second chance rather than being skipped forever.
Without clearing, the scheme degenerates to ``never evict anyone''.
\end{Sol}

\begin{Qn}{Medium}{Lect 15}{CS-537 Sp12 Mid, Q9}{A 32-bit reference counter instead of a reference bit}
A per-page reference bit is sometimes used to track active pages. A
subtle change turns it into a per-page 32-bit counter that is
incremented on every reference. What is the impact on page replacement?
Can (or should) the OS policy change to use it?
\end{Qn}
\begin{Sol}
The OS now has a \textbf{frequency} (and, if the counter is never
reset, a lifetime-use) signal, not just a 1-bit recency flag.

What becomes possible:
\begin{itemize}
  \item \textbf{LFU} (evict the smallest counter) or an aging variant
        (periodically right-shift every counter, which is the
        additional-reference-bits / NFU scheme).
  \item A closer approximation to perfect LRU if the counter is a
        timestamp rather than a tally (then replacement scans for the
        min timestamp --- the Fa09 LostOS design).
\end{itemize}
Cost: incrementing 32 bits on every hit is heavier than setting one
bit, and finding the min counter still requires a scan (or a heap).
The OS \emph{should} change the policy --- using the extra bits as a
clock use-bit would throw information away --- but it should
\textbf{age} the counters. An un-aged LFU never forgets a once-hot page
and will not evict it even after the working set has moved on.

(This is still an approximation of OPT; it does not make the
replacement decision free.)
\end{Sol}

\begin{Qn}{Medium}{Lect 15}{ISI Mandar 2014 End-sem, Q2}{PTE extra bits, Belady's anomaly, and additional-reference-bits LRU}
32-bit addresses, 2\,KB pages, 4-byte PTEs.
\begin{enumerate}
  \item[(a)] (i) How many bits of each PTE remain for extra information
        besides the frame number? (ii) List and briefly explain three
        kinds of such information.
  \item[(b)] What is Belady's anomaly?
  \item[(c)] Explain the additional-reference-bits approximation to LRU.
\end{enumerate}
\end{Qn}
\begin{Sol}
(a)(i) Offset $=11$ bits ($2\,\mathrm{KB}=2^{11}$). A 32-bit physical
address then has a 21-bit PFN. A 4-byte PTE has $32-21=11$ bits left
for flags. (If physical memory is smaller the PFN shrinks and more
flag bits appear; 11 is the number when PA width equals VA width.)

(a)(ii) Any three of:
\begin{itemize}
  \item \textbf{Valid:} VPN is in the address space (else segfault).
  \item \textbf{Present:} page is in RAM vs.\ on swap (else page fault).
  \item \textbf{Dirty/modified:} must write back on eviction.
  \item \textbf{Reference/accessed:} clock / second-chance / ARB input.
  \item \textbf{Protection} R/W/X.
\end{itemize}

(b) For some policies (FIFO is the textbook example) \emph{increasing}
the number of allocated frames can \emph{increase} the number of page
faults on a fixed reference string. Stack algorithms (LRU, OPT) cannot
do this. Classic string: $1,2,3,4,1,2,5,1,2,3,4,5$ yields 9 FIFO
misses with 3 frames and 10 with 4.

(c) Each PTE holds an $n$-bit shift register (often 8). Every $T$
milliseconds a timer interrupt: for every resident page, shift the
register right one bit and insert the hardware reference bit at the
MSB, then clear the hardware bit. The $n$-bit word is a decaying
history of use; the page with the \textbf{smallest} word is the
approximate LRU victim. $n=1$ with a continuous scan is clock;
larger $n$ is strictly more accurate and strictly more work
(the Nth-chance / aging idea).
\end{Sol}

\begin{Qn}{Medium}{Lect 15}{ISI Mandar 2012 End-sem, Q3(e)}{Minimum frames a demand-paged process must be given}
How can you determine the minimum number of frames that must be
allocated to a process in a demand-paging system?
\end{Qn}
\begin{Sol}
The bound is set by the \textbf{ISA}, not by the replacement policy.
A single instruction may touch several pages at once (instruction
fetch, possibly spanning a page boundary, plus each memory operand,
plus an auto-incremented stack slot, \ldots). If the process is given
fewer frames than that maximum, the restarted instruction can fault
again on a different page, evict the page it just brought in, and
\textbf{livelock}. So allocate at least
\[
  m = \max_{\text{instructions }I}(\#\text{ of distinct pages }I\text{ may reference}).
\]
Typical numbers: a two-page instruction plus two operand pages $\to$
$m=4$; an auto-increment addressing mode can add one more. Replacement
among the remaining frames (FIFO, LRU, clock, \ldots) is then free to
run.
\end{Sol}

\begin{Qn}{Medium}{Lect 15}{ISI Mandar 2024 End-sem, Q4(d); cf.\ Q4(c)}{Compulsory faults vs.\ extra frames and the replacement algorithm}
A 32-bit machine, 4\,KB pages, \texttt{int}s of 4 bytes,
\texttt{rows=10}, \texttt{cols=1024}. Fragment II allocates each of
the 10 rows by a separate \texttt{malloc}. The row-start offsets
modulo 4096 are
\[
1920,1936,1952,1968,1984,2000,2016,2032,2048,2064
\]
(an arithmetic progression of 16 bytes, the usual 16-byte malloc
header). The process then sums every matrix element in row-major
order. Part (c) of the same question asks how many \textbf{distinct}
pages that summation references; part (d) asks:
\begin{enumerate}
  \item[(d)] How many of those references fault? How does the number
        change if the process is given more frames, or if the
        replacement algorithm changes?
\end{enumerate}
You may assume demand paging, no prepaging, and that none of the
matrix pages is resident at the start of the summation.
\end{Qn}
\begin{Sol}
Each row is 4096 bytes starting at a non-zero offset, so it straddles
\textbf{two} pages. Starts are 4112 bytes apart ($4096+16$), so row
$i$ occupies physical pages $P+i$ and $P+i+1$. Ten rows therefore
touch pages $P,P+1,\ldots,P+10$: \textbf{11 distinct data pages}.
(The pointer array \texttt{A[10]} is a 12th page if pointer
dereferences are counted; the argument below is unchanged.)

(d) Under demand paging with an empty cache, every distinct page
faults \textbf{exactly once} the first time it is touched --- 11
compulsory misses --- \emph{provided} the process has at least 11
frames (or, more weakly, at least as many frames as the number of
pages that are live together during the scan). Extra frames beyond
that cannot reduce the count: there is nothing left to prefetch in
this model.

If the process has \textbf{fewer} than 11 frames, the scan's working
set does not fit. Then:
\begin{itemize}
  \item additional \textbf{capacity} faults appear;
  \item their number \textbf{does} depend on the replacement algorithm
        (LRU/clock will keep the two pages of the current row and
        perhaps the previous row; FIFO may evict a page of the current
        row and re-fault it; OPT is the lower bound);
  \item Belady's anomaly is possible under FIFO if one compared two
        small frame counts.
\end{itemize}
Giving more frames, or switching LRU $\leftrightarrow$ FIFO, cannot
create \emph{fewer} than 11 faults here, and cannot create
\emph{more} than 11 once $k\ge 11$.
\end{Sol}

\begin{Qn}{Medium}{Lect 15}{CS-537 Sp11 Mid, P6}{Classify each reference: TLB hit, TLB miss/page hit, or page fault}
TLB has 4 entries, RAM holds 8 pages, both start empty, LRU replacement
in both. Write H (TLB hit), M (TLB miss, page present), or PF (page
fault) for each reference. (Trace reconstructed from the OCR'd exam
pages; the three patterns are the standard ones used with these cache
sizes.)
\begin{enumerate}
  \item[(a)] $0,1,2,3,0,1,2,3$
  \item[(b)] $0,1,2,3,4,1,2$
  \item[(c)] RAM reduced to 3 pages: $1,2,0,1,3,0,3,1,2$
\end{enumerate}
\end{Qn}
\begin{Sol}
A TLB hit requires the VPN to sit in the 4-entry LRU TLB. A page fault
occurs only if the VPN is not among the RAM pages; the OS then fills
RAM (and typically the TLB) and retries. A resident page whose TLB
entry was evicted is M.

(a) $0,1,2,3$ all PF (cold TLB and cold RAM). The four translations
now occupy the whole TLB and four of eight frames, so the repeat is
HHHH.
\[
\mathrm{PF},\;\mathrm{PF},\;\mathrm{PF},\;\mathrm{PF},\;
\mathrm{H},\;\mathrm{H},\;\mathrm{H},\;\mathrm{H}.
\]

(b) Four PFs, then $4$ is a PF that evicts VPN 0 from the LRU TLB
(RAM still has room). $1$ and $2$ remain in the TLB: H, H.
\[
\mathrm{PF},\;\mathrm{PF},\;\mathrm{PF},\;\mathrm{PF},\;
\mathrm{PF},\;\mathrm{H},\;\mathrm{H}.
\]

(c) 3-frame LRU RAM, 4-entry LRU TLB (victim pages are invalidated in
the TLB).
\begin{center}
\small
\begin{tabular}{c|ccccccccc}
ref  & 1 & 2 & 0 & 1 & 3 & 0 & 3 & 1 & 2 \\
\hline
RAM  & 1 & 12& 120& 201& 013& 130& 103& 031& 312 \\
     &   &   &    &    & evict 2 & & & & evict 0 \\
out  & PF& PF& PF & H  & PF & H  & H  & H  & PF \\
\end{tabular}
\end{center}
(LRU order: least-recent left, most-recent right.) Page 2 is evicted on
the fault to 3 and is still gone when it is reused, so the last access
is a page fault. Sequence:
$\mathrm{PF},\mathrm{PF},\mathrm{PF},\mathrm{H},\mathrm{PF},\mathrm{H},\mathrm{H},\mathrm{H},\mathrm{PF}$.
\end{Sol}

\begin{Qn}{Medium}{Lect 15}{CS-537 Fa10 Mid, Q4(e)}{The dirty bit: purpose, and what a flipped dirty bit does}
A software-managed TLB entry contains a dirty bit D, among other
fields. (i) What is the dirty bit for? (ii) What happens if it is
flipped just before the OS installs the TLB entry?
\end{Qn}
\begin{Sol}
(i) Hardware sets D when the page is written. Replacement uses it to
decide whether the victim must be \textbf{written back to swap}
(D$\,=1$) or can be discarded (D$\,=0$, the swap copy is still
current). Skipping writebacks of clean pages is the entire reason the
bit exists; it is also the reason enhanced clock classifies
$(\mathrm{use},\mathrm{dirty})$ pairs as
$(0,0)<(0,1)<(1,0)<(1,1)$.

(ii) $0\to 1$: the page looks dirty. Evicting it causes an
\textbf{unnecessary writeback} (slower, still correct). $1\to 0$: the
page looks clean. Evicting it \textbf{drops the only copy of the
stores} that happened since it was paged in --- silent data loss, or
a stale version coming back on the next fault. Unlike a flipped
reference bit (which only disturbs the replacement hint), a flipped
dirty bit can violate correctness.
\end{Sol}

\clearpage

\section{Lectures 16--20: Concurrency --- Threads, Locks, CVs, Semaphores \& Deadlock}

\subsection{Lectures 16--17: Concurrency, Threads \& Locks}

\subsubsection*{Additional Challenging \& Analytical Problems}

\begin{Qn}{Medium-Hard}{Lect 17}{CS-537 Fall 2004 Mid1 Long Q2}{Hash table: why lock, one lock, then many locks}
A multi-threaded hash table is an array of linked lists. Insert hashes, then inserts into that list:

\begin{lstlisting}[numbers=none]
void hash_insert(int key, object_t *obj) {
    int whichList = key % HASH_SIZE;
    list_insert(hashLists[whichList], key, obj);
}
\end{lstlisting}

The list insert (details omitted) is the usual
\texttt{tmp->next = head; head = tmp;}

(a)~Why must \texttt{list\_insert} be synchronized?

(b)~Add a single lock into \texttt{hash\_insert} to fix the problem.

(c)~What is the biggest performance problem with this single-lock solution?

(d)~Add more locks to improve potential concurrency.
\end{Qn}
\begin{Sol}
(a)~Two threads inserting into the \emph{same} bucket can both read the same \texttt{head} and the later store orphans a node (exactly the Lecture~17 linked-list race).

(b)~One global mutex around the call:
\begin{lstlisting}[numbers=none]
mutex_lock(&table_lock);
list_insert(hashLists[whichList], key, obj);
mutex_unlock(&table_lock);
\end{lstlisting}

(c)~The lock serializes \emph{all} inserts, even into disjoint buckets. Unrelated keys cannot proceed in parallel; the whole table becomes a single critical section.

(d)~Fine-grained / per-bucket locks: an array \texttt{mutex\_t bucket[HASH\_SIZE]}, and lock only \texttt{bucket[whichList]}. Inserts into different buckets then run concurrently. (Hand-over-hand per-node locking is a further step, with deadlock risk if lock order is not fixed --- Lecture~17 flags that.)
\end{Sol}

\begin{Qn}{Medium-Hard}{Lect 17}{CS-537 Fall 2004 Mid1 Q4}{Strict-alternation lock: ME without progress?}
\begin{lstlisting}[numbers=none]
while (turn == (1 - threadID))
    ; /* spin */
balance = balance + amount; /* critical section */
turn = 1 - threadID;
\end{lstlisting}

\texttt{turn} is initialized to 0, and only two threads run. Does this provide mutual exclusion? Explain.
\end{Qn}
\begin{Sol}
Yes, ME holds. \texttt{turn} is a single shared token. A thread enters only when \texttt{turn} equals its id, then flips \texttt{turn} to the other id on exit. At most one thread sees the matching value at a time.

It is \emph{not} a satisfactory CSP solution: it enforces strict alternation, so it fails the \textbf{progress} (deadlock-freedom / no unnecessary waiting) requirement. If thread~0 never wants the CS, thread~1 spins forever even though the CS is free. It also fails bounded wait for a thread that is skipped because the other never runs. This is the ``broken flag / take-turns'' family Lecture~17 contrasts with Peterson.
\end{Sol}

\begin{Qn}{Hard}{Lect 17}{CS-537 Spring 2002 Mid Long Q2}{Synchronize \texttt{balance += amount} with LL/SC and with fetch-and-add}
\texttt{update(amount)} must atomically do \texttt{balance = balance + amount}. Amount arrives in \texttt{R1}. Unsynchronized code:

\begin{lstlisting}[numbers=none]
load <balance>, R2
add  R1, R2, R3
store <balance>, R3
\end{lstlisting}

(a)~Implement \texttt{update} with load-linked / store-conditional only. \texttt{ll} loads; \texttt{sc} stores iff the address was not written since the matching \texttt{ll}, and sets the source register to 1 on success and 0 on failure.

(b)~Any limitations of (a)?

(c)~Implement \texttt{update} with atomic fetch-and-add (\texttt{fetchadd <addr>, RS} adds RS into \texttt{*addr}).

(d)~Which primitive is more appropriate here?
\end{Qn}
\begin{Sol}
(a)~Retry until the store-conditional succeeds:

\begin{lstlisting}[numbers=none]
top:
  ll    <balance>, R2      /* linked load of current balance */
  add   R1, R2, R3         /* R3 = balance + amount */
  sc    <balance>, R3      /* try to commit; R3 <- 1/0 */
  if (R3 == 0) goto top    /* lost the race: retry */
\end{lstlisting}

(If \texttt{sc} overwrites the register that holds the new balance, copy \texttt{R3} to a temp first, or use a dedicated success register as the ISA allows.)

(b)~Under contention the loop can starve (no fairness). A delayed thread may livelock if others keep writing \texttt{balance}. LL/SC also fails if \emph{any} write hits the reservation granule (false sharing), even a write of the same value. It is a lock-free add, not a wait-free add.

(c)~One instruction:
\begin{lstlisting}[numbers=none]
fetchadd <balance>, R1     /* atomically balance += amount */
\end{lstlisting}

(d)~Fetch-and-add is the better fit: the update \emph{is} an add. LL/SC implements a general read-modify-write via a retry loop; fetch-and-add does it in hardware with no lost-update window and no extra branching.
\end{Sol}

\begin{Qn}{Medium-Hard}{Lect 17}{CS-537 Spring 2009 Final Q3}{TTAS lock, and unlocking with TAS}
\begin{lstlisting}[numbers=none]
void lock(lock_t *lock) {
  do {
    while (lock->flag) ;           /* unprotected check */
  } while (TestAndSet(&lock->flag, 1));
}
void unlock(lock_t *lock) { lock->flag = 0; }
\end{lstlisting}

(a)~Does this lock work? (b)~When is it better / worse than a pure TAS spinlock?

They then change unlock to \texttt{TestAndSet(\&lock->flag, 0)}.

(c)~Does that unlock work? (d)~Better, worse, or the same?
\end{Qn}
\begin{Sol}
(a)~Yes. Same TTAS argument as 2012 Mid Q14.

(b)~Better when many CPUs spin: ordinary loads vs.\ repeated exclusive atomics. Worse (or no win) on a uniprocessor, or when the lock is uncontended (you still do a TAS at the end). Also slightly more work in the uncontended path (load then TAS).

(c)~Yes. TAS to 0 writes 0 regardless of the old value, so the lock becomes free. ME is preserved.

(d)~Worse. Unlock does not need an atomic RMW; a plain store of 0 suffices and is cheaper (no lock prefix / no exclusive line from the unlocker beyond the write). TAS-on-unlock also returns the old value, which is ignored.
\end{Sol}

\begin{Qn}{Medium-Hard}{Lect 17}{CS-537 Fall 2021 Midterm 2 Q25--Q26}{\texttt{FREE} and \texttt{HELD} polarity}
\begin{lstlisting}[numbers=none]
void init(lock_t *lock) { lock->flag = FREE; }
void acquire(lock_t *lock) {
     while (xchg(&lock->flag, HELD) == HELD)
         ;
}
void release(lock_t *lock) { lock->flag = 1; }
\end{lstlisting}

Q25. What should \texttt{FREE} be so this is a mutex? (a)~0 (b)~1 (c)~2 (d)~0 or 2 (e)~0 or 1

Q26. Same code: what should \texttt{HELD} be? (a)~0 (b)~1 (c)~2 (d)~0 or 2 (e)~0 or 1
\end{Qn}
\begin{Sol}
Release writes $1$, so the free value \emph{is} 1. Acquire TASes in \texttt{HELD} and spins while the old value equals \texttt{HELD}. Therefore \texttt{HELD} must be \emph{unequal} to 1, i.e.\ 0 or 2.

Q25: (b)~1. \quad Q26: (d)~0 or 2.
\end{Sol}

\begin{Qn}{Hard}{Lect 17}{Mandar ISI 2016 Mid Q4(b)}{\texttt{Swap} solution for the $n$-process CSP, including bounded waiting}
A processor provides \texttt{Swap(x,y)} that atomically interchanges two variables. Implement an $n$-process CSP solution that satisfies mutual exclusion, progress, \emph{and} bounded waiting.
\end{Qn}
\begin{Sol}
This is the textbook waiting-array lock (Galvin), with \texttt{Swap} in place of TAS. \texttt{lock==true} means held. A waiter sets \texttt{waiting[i]} and Swap-s a local \texttt{key} with \texttt{lock} until either it receives \texttt{key==false} (lock was free) or another exiting thread clears \texttt{waiting[i]}. Exit hands the lock to the next waiter in cyclic order, which gives bounded waiting ($\le n-1$ overtaking).

\begin{lstlisting}[numbers=none]
shared bool lock = false;
shared bool waiting[n];          /* all false */
/* local: bool key; int j; */

/* entry */
waiting[i] = true;
key = true;
while (waiting[i] && key)
    Swap(&lock, &key);           /* atomic exchange of lock and key */
waiting[i] = false;

critical_section();

/* exit: pass to the next waiter, else release */
j = (i + 1) % n;
while ((j != i) && !waiting[j])
    j = (j + 1) % n;
if (j == i)
    lock = false;                /* nobody waiting */
else
    waiting[j] = false;          /* wake j; it skips Swap and enters */
\end{lstlisting}

ME: only a thread that observes \texttt{key==false} or a directed wake can enter; \texttt{lock} is a single token. Progress: if the CS is free and someone wants in, either \texttt{Swap} succeeds or a waiter is handed \texttt{waiting[j]=false}. Bounded wait: the cyclic scan starts at $i+1$, so a waiting $P_k$ is skipped at most $n-1$ times.
\end{Sol}

\begin{Qn}{Medium-Hard}{Lect 17}{Mandar ISI 2022 Mid Q4(a)}{TAS in a kernel whose CPU has no TAS instruction}
Implement an atomic \texttt{TestAndSet()} system call on a processor that has no atomic test-and-set instruction.
\end{Qn}
\begin{Sol}
On a \textbf{uniprocessor}, disable interrupts around the load/store so the sequence cannot be interleaved with another process (or with an interrupt handler that might also take the lock):

\begin{lstlisting}[numbers=none]
int sys_TestAndSet(int *addr) {
    int old;
    disable_interrupts();     /* clear IF / equivalent */
    old = *addr;
    *addr = 1;
    enable_interrupts();
    return old;
}
\end{lstlisting}

This is exactly Lecture~17's ``disable interrupts'' lock, used \emph{inside the kernel} and only briefly. It does \emph{not} work on a multiprocessor (another CPU can still interleave). There one needs a real atomic, or a per-CPU interrupt disable \emph{plus} a hardware TAS/CAS on a shared flag. User-level code must not be allowed to disable interrupts (malicious delay of the timer; missed I/O) --- that is why this is a \emph{system call} / kernel primitive, not a library routine.
\end{Sol}

\begin{Qn}{Medium-Hard}{Lect 17}{Mandar ISI 2018 Mid Q5}{Smallest TAS critical section for a bounded buffer}
Write producer/consumer (bounded buffer) synchronized with atomic test-and-set only, with the critical section as small as possible. Initialize all variables.
\end{Qn}
\begin{Sol}
TAS gives a mutex, not a CV: waiters \emph{spin} on the buffer-full / buffer-empty conditions, but they must \emph{drop the lock} while spinning, otherwise the other side can never fill/empty the buffer (progress failure). Smallest CS: lock only while reading/writing the shared indices and the slot.

\begin{lstlisting}[numbers=none]
#define N 10
int buf[N], in = 0, out = 0, count = 0;
int lock = 0;
void tas_lock()   { while (TestAndSet(&lock)) ; }
void tas_unlock() { lock = 0; }

void producer(int item) {
    int done = 0;
    while (!done) {
        tas_lock();
        if (count < N) {          /* CS: check + insert + bookkeeping */
            buf[in] = item;
            in = (in + 1) % N;
            count++;
            done = 1;
        }
        tas_unlock();             /* drop lock before retrying if full */
    }
}

void consumer(int *item) {
    int done = 0;
    while (!done) {
        tas_lock();
        if (count > 0) {
            *item = buf[out];
            out = (out + 1) % N;
            count--;
            done = 1;
        }
        tas_unlock();
    }
}
\end{lstlisting}

Locking across the whole wait-for-space loop would deadlock the system: the consumer could never acquire TAS to make space. (Condition variables, owned by Lecture~18, would \texttt{park} instead of spinning.)
\end{Sol}

\begin{Qn}{Hard}{Lect 16}{CS-537 Spring 2001 Mid Long Q2}{Race to the Finish}
Shared \texttt{int z = 0}. Many threads call

\begin{lstlisting}[numbers=none]
void update(int x, int y) { z += x + y; }
\end{lstlisting}

with \texttt{x} in R1 and \texttt{y} in R2. Instructions: \texttt{load [addr], Rdest}, \texttt{add Rdest, Rsrc1, Rsrc2}, \texttt{store Rsrc, [addr]}.

(a)~Write assembly for \texttt{update}.

(b)~What is a race condition here?

(c)~Mark interrupt points that allow a race.

(d)~Now \texttt{z = x + y} (not additive). Two threads call \texttt{update(3,4)} and \texttt{update(10,20)}. Possible outcomes? Does a lock around the assignment change that?
\end{Qn}
\begin{Sol}
(a)
\begin{lstlisting}[numbers=none]
load  [z], R3
add   R4, R1, R2      /* R4 = x + y */
add   R5, R3, R4      /* R5 = z + (x+y) */
store R5, [z]
\end{lstlisting}

(b)~A race: concurrent updates whose loads/stores interleave so that one add is lost. Final \texttt{z} is not the sum of all \texttt{(x+y)} contributions.

(c)~Any preemption after the \texttt{load [z]} and before the matching \texttt{store} --- the other thread's full update is then overwritten. Preemption \emph{before} the load, or \emph{after} the store, is safe for this snippet.

(d)~Without a lock the possible values of \texttt{z} are $7$ or $30$ (last writer wins; there is no ``add into old z'', so 37 cannot appear). A lock serializes the two assignments but the last writer still wins: still $\{7,30\}$. The lock removes non-determinism of a \emph{torn} assignment, not the fact that both writes target the same cell with absolute stores. Contrast with (a)--(c), where a lock \emph{would} make \texttt{z} uniquely equal to the sum of all contributions.
\end{Sol}

\begin{Qn}{Hard}{Lect 16}{CS-537 Fall 2021 Midterm 2 Q1--Q2, Q9--Q10}{Check-then-act and racy multiply}
\textbf{Q1.} Global \texttt{int i = 0}. Two threads each run \texttt{if (i == 0) i++;}. Possible finals of \texttt{i}?
(a)~0 (b)~1 (c)~2 (d)~1 or 2 (e)~none

\textbf{Q2.} Same, but the increment is wrapped in \texttt{mutex\_acquire}/\texttt{release} \emph{inside} the \texttt{if}. Possible finals?
(same options)

\textbf{Q9.} Global \texttt{count = 100}. Two threads each run \texttt{count = count * 2;} as load, mul-in-register, store. Possible finals?
(a)~100 (b)~200 (c)~400 (d)~200 or 400 (e)~none

\textbf{Q10.} Same multiply, now inside a mutex. Possible finals?
\end{Qn}
\begin{Sol}
\textbf{Q1.} Two plausible compilations. If both threads can pass the test on $i=0$ and then each do a load/inc/store of $i$, finals are 1 (one wins the store, or they serialize) or 2 (both entered and both stores committed). $i=0$ is impossible (at least one increment stores a positive value). Official key accepted (b) or (d) depending on the assumed assembly; (d) is the conservative Lecture~16 answer.

Trace for 2:
\begin{center}
\begin{tabular}{c|l|l|l}
Step & T1 & T2 & $i$ \\
\hline
1 & load $i\to r=0$; $0==0$, enter & & 0 \\
2 & & load $i\to r=0$; enter & 0 \\
3 & $r{+}{+}{=}1$; store $i=1$ & & 1 \\
4 & & $r{+}{+}{=}1$; store $i=1$ & 1 \\
\end{tabular}
\end{center}
That particular store-of-1 twice yields 1. For 2, each must reload $i$ inside the then-branch (or compile \texttt{i++} against a fresh load). If the then-branch is \texttt{load; inc; store} after a \emph{separate} test load, both can test on 0, then serialize two increments: $0\to1\to2$. So (d).

\textbf{Q2.} The lock is \emph{after} the test. Both can still pass \texttt{i==0}, then serialize \texttt{i++}. Same set as Q1. Official: ``the lock doesn't help''. (d).

\textbf{Q9.} Serialize: $100\to200\to400$. Both load 100, both store 200: lost update. (d)~200 or 400.

\textbf{Q10.} Mutex makes the multiply atomic wrt the other multiply: $100\to200\to400$ only. (c)~400.
\end{Sol}

\begin{Qn}{Medium-Hard}{Lect 16}{CS-537 Fall 2021 Midterm 2 Q5}{Double-checked list initialization}
Global \texttt{List\_t *L = NULL}. Two threads run

\begin{lstlisting}[numbers=none]
if (L == NULL) {
    L = malloc(sizeof(List_t));
    List_Init(&L); /* sets L->head = NULL */
}
\end{lstlisting}

\texttt{malloc} never fails. Possible outcomes:
(a)~crash (b)~memory may be leaked (c)~list properly initialized once and works
(d)~both (b) and (c) (e)~each of (a), (b), and (c)
\end{Qn}
\begin{Sol}
If one thread runs the whole \texttt{if} first, the other sees non-\texttt{NULL} $L$: (c). If both pass \texttt{L==NULL}, both \texttt{malloc}, both \texttt{List\_Init}; the later store to $L$ wins and the other block is leaked: (b). Crash is out (malloc succeeds; no \texttt{NULL} deref required). Answer (d).

This is the same check-then-act pattern as Q1--Q2; a lock must cover the test \emph{and} the publish of $L$.
\end{Sol}

\begin{Qn}{Medium-Hard}{Lect 16--17}{CS-537 Spring 2002 Mid Q6}{Double-checked locking to open a pipe}
Original (correct, coarse):

\begin{lstlisting}[numbers=none]
void MakeSurePipeIsOpen() {
    mutex_lock(&lock);
    if (pipe == -1)
        pipe = open("/tmp/fifo", O_WRONLY);
    mutex_unlock(&lock);
}
\end{lstlisting}

``Clever'' rewrite:

\begin{lstlisting}[numbers=none]
void MakeSurePipeIsOpen() {
    if (pipe == -1) {
        mutex_lock(&lock);
        if (pipe == -1)
            pipe = open("/tmp/fifo", O_WRONLY);
        mutex_unlock(&lock);
    }
}
\end{lstlisting}

Does it still work? If so, what is the advantage? If not, why?
\end{Qn}
\begin{Sol}
It is the classic double-checked locking pattern. Under sequential consistency for aligned integer \texttt{pipe}, it \emph{does} work: every thread that sees \texttt{pipe != -1} observes a completed \texttt{open}; every thread that sees $-1$ takes the lock and rechecks. Advantage: the common path (pipe already open) does not acquire the mutex.

Caveats the exam expects: (i)~without a memory barrier / \texttt{volatile}, a weakly ordered CPU may let a thread see a non-$-1$ \texttt{pipe} before the \texttt{open} side effects are visible; (ii)~two threads can both call \texttt{open} if they both pass the outer test before either publishes --- the inner lock prevents that, so (ii) is actually fine. Net: accepted as correct with a performance win on the uncontended already-open path, provided the store to \texttt{pipe} is visible. Lecture~16's ``malicious scheduler'' still cannot break the inner locked recheck.
\end{Sol}

\begin{Qn}{Hard}{Lect 16--17}{CS-537 Fall 2004 Final Q3}{Five broken \texttt{update} snippets}
\texttt{update} should add \texttt{amount} to global \texttt{balance}. For each, say Race, Works, or Broken (not a race):

(a)
\begin{lstlisting}[numbers=none]
int tmp = amount;
lock(); tmp = tmp + balance; unlock();
balance = tmp;
\end{lstlisting}
(b)
\begin{lstlisting}[numbers=none]
int tmp = amount;
lock(); tmp = tmp + amount; balance = tmp; unlock();
\end{lstlisting}
(c)
\begin{lstlisting}[numbers=none]
lock(); balance = balance + amount; lock();
\end{lstlisting}
(d)
\begin{lstlisting}[numbers=none]
int *tmp = &balance;
lock(); *tmp = *tmp + amount; unlock();
\end{lstlisting}
(e)
\begin{lstlisting}[numbers=none]
int tmp = balance;
lock(); tmp = tmp + amount; unlock();
\end{lstlisting}
\end{Qn}
\begin{Sol}
(a)~\textbf{Race.} The store to \texttt{balance} is outside the lock. Two threads can each compute \texttt{tmp = amount + old\_balance} and the later store loses an add.

(b)~\textbf{Broken} (not a race on \texttt{balance}'s atomicity). Inside the lock it sets \texttt{balance = 2*amount}, ignoring the old balance. Deterministically wrong.

(c)~\textbf{Broken.} Second \texttt{lock()} deadlocks the caller (non-recursive mutex). No matching unlock.

(d)~\textbf{Works.} \texttt{tmp} aliases \texttt{balance}; the add happens inside the lock.

(e)~\textbf{Race} (and also drops the write). Load of \texttt{balance} is outside the lock; computed \texttt{tmp} is never stored back. Even with a store, the load/store would need to sit together inside the lock.
\end{Sol}

\begin{Qn}{Medium-Hard}{Lect 16--17}{CS-537 Fall 2004 Mid1 Q3}{The rogue thread that refuses to take the lock}
Properly synchronized:

\begin{lstlisting}[numbers=none]
void update(int amount) {
    mutex_lock(lock);
    balance = balance + amount;
    mutex_unlock(lock);
}
\end{lstlisting}

A rogue thread instead runs \texttt{balance = balance + amount;} with no lock. Can we prevent threads from updating shared data in this rogue fashion?
\end{Qn}
\begin{Sol}
Not in C / pthreads. A lock is a \emph{convention}: every accessor must cooperate. The compiler will still emit ordinary loads and stores to \texttt{balance} for the rogue. Hardware does not tag the cell as ``only accessible while \texttt{lock} is held''. Isolation of the variable into a module whose only exported functions take the lock is a software discipline, not an OS guarantee. (Language-level monitors / private object fields are the usual way to make the rogue not compile; that is outside Lecture~16--17.)
\end{Sol}

\begin{Qn}{Medium-Hard}{Lect 17}{CS-537 Fall 2021 Midterm 2 Q29}{Where to put the list lock}
\begin{lstlisting}[numbers=none]
void List_Insert(list_t *L, int key) {
    node_t *new = malloc(sizeof(node_t)); /* L1 */
    new->key = key;                       /* L2 */
    new->next = L->head;                  /* L3 */
    L->head = new;                        /* L4 */
}
\end{lstlisting}

\texttt{malloc} is thread-safe. Unlock is after L4. Best place for \texttt{mutex\_lock}?
(a)~before L1 (b)~before L2 (c)~before L3 (d)~before L4 (e)~none
\end{Qn}
\begin{Sol}
(c)~right before L3. L1--L2 touch only the new node (private until published). L3--L4 is the atomic swing of the shared head. Locking earlier reduces concurrency; locking only L4 lets two threads both read the same head at L3 and lose a node.
\end{Sol}

\begin{Qn}{Medium-Hard}{Lect 16}{CS-537 Fall 2021 Midterm 2 Q14}{Child prints \texttt{*p} before \texttt{p} is set}
\begin{lstlisting}[numbers=none]
int *p = NULL;
void child(void *arg) { printf("%d\n", *p); }
int main() {
    thread_t p1;
    int x = 3;
    thread_create(&p1, child, NULL);
    p = &x;
    return 0;
}
\end{lstlisting}

What can this print? (a)~3 \quad (b)~prints nothing and exits \quad (c)~crashes \quad (d)~3, or nothing, or crash \quad (e)~none
\end{Qn}
\begin{Sol}
(d). Three schedules:
\begin{itemize}
\item Child runs while \texttt{p} is still \texttt{NULL}: crash.
\item Parent sets \texttt{p=\&x}, then child runs, then parent is still alive: prints 3.
\item Parent sets \texttt{p}, then \texttt{return 0} exits the process before the child runs (or while it is running): nothing, or a crash if \texttt{x}'s stack died with \texttt{main}. The official key also includes ``prints nothing'' from parent exiting first.
\end{itemize}
No join $\Rightarrow$ lifetime of \texttt{x} and of the process are not synchronized with the child. Answer (d).
\end{Sol}

\begin{Qn}{Hard}{Lect 16--17}{CS-537 Spring 2011 Final Problem 2(a)--(d)}{Coarse lock vs.\ fine lock vs.\ a bare flag}
Initial: \texttt{int g = 10; int *f = NULL;} Mutex \texttt{m}.

(a)~T1: \texttt{lock; f=\&g; /* other stuff */; printf("\%d", *f); unlock;}
T2: \texttt{lock; f=NULL; unlock;}
Possible outputs?

(b)~T1 splits the lock: \texttt{lock; f=\&g; unlock; /* other stuff */; lock; printf(*f); unlock;}
(T2 unchanged.) Possible outputs?

(c)~Do locks around critical sections prevent ``bad things'' in (b)?

(d)~T1: \texttt{pending=1; while(pending) printf("hello\\n");} \quad T2: \texttt{pending=0;}
Possible outputs?
\end{Qn}
\begin{Sol}
(a)~T2 cannot run between \texttt{f=\&g} and the \texttt{printf}, because T1 holds \texttt{m} across both. T2 either runs entirely first (\texttt{f=NULL}, then T1 sets \texttt{f=\&g} and prints 10) or entirely after T1's CS (T1 prints 10, then T2 nulls \texttt{f}). Only output: \texttt{10}. No crash.

(b)~T2 can now sneak in during ``other stuff'':
\begin{center}
\begin{tabular}{c|l|l}
Step & T1 & T2 \\
\hline
1 & lock; \texttt{f=\&g}; unlock & \\
2 & & lock; \texttt{f=NULL}; unlock \\
3 & lock; \texttt{printf(*f)} $\to$ deref NULL & \\
\end{tabular}
\end{center}
Outputs: \texttt{10}, or a crash. (If T2 runs first then T1's first CS republishes \texttt{\&g} before the printf, 10.)

(c)~No. Locks prevent races on the \emph{variables they cover for the duration they are held}. Splitting the CS re-opens a TOCTOU window on \texttt{f}: the check-like publish and the use are no longer atomic. Fine-grained locking is not automatically safe.

(d)~Any finite number of \texttt{hello} lines (including zero), or infinitely many.
\begin{itemize}
\item T2 runs first, then T1 sets \texttt{pending=1} and loops forever.
\item T1 sets \texttt{pending=1}, prints $k\ge 0$ times, T2 stores 0, T1's next while-test fails.
\item T1 tests \texttt{pending} before either store (if \texttt{pending} started 0, zero hellos then T1 sets 1 and maybe loops).
\end{itemize}
This is an ordering problem; a lock on each assignment does not force ``hello at least once before T2''. (Lecture~18 CVs would; do not use them here.)
\end{Sol}

\begin{Qn}{Medium-Hard}{Lect 16}{Mandar ISI 2017 End Q2(c)}{\texttt{pthread\_create(\&tid, NULL, main, NULL)} plus an infinite print}
\begin{lstlisting}[numbers=none]
int main(int ac, char *av[]) {
    static unsigned int i = 0;
    pthread_t tid;
    pthread_create(&tid, NULL, main, NULL);
    while (1) {
        printf("Hello %u\n", i);
        i++;
    }
    return 0;
}
\end{lstlisting}

Explain what the program does. (It compiles and runs.)
\end{Qn}
\begin{Sol}
\texttt{pthread\_create} starts a \emph{second} thread at \texttt{main}. That thread again creates a third thread at \texttt{main}, and so on: threads are created exponentially (each new thread clones another \texttt{main}) until \texttt{pthread\_create} fails or the machine is exhausted.

Every thread that returns from \texttt{create} (success or not) falls into \texttt{while(1)} and prints \texttt{Hello} with the \emph{shared} static \texttt{i}. Prints interleave; \texttt{i++} races, so the printed sequence is not $0,1,2,\ldots$ and values may repeat or skip. The process never reaches \texttt{return 0} unless a thread's \texttt{create} path is somehow cut off; typically it runs until resources run out.
\end{Sol}

\begin{Qn}{Hard}{Lect 16--17}{Mandar ISI 2024 End Q3(a)--(c)}{Three threads on \texttt{strtok}}
\begin{lstlisting}[numbers=none]
char str[] = "Example string, containing words, separated by spaces, and commas.";
char *cp = strtok(str, " ,");
do
    puts(cp);
while (NULL != (cp = strtok(NULL, " ,")));
\end{lstlisting}

(a)~Create 3 threads in \texttt{main} so each runs the do-while.

(b)~What goes wrong in the output? Give examples.

(c)~Fix with \texttt{pthread\_mutex} so each thread prints \emph{all} tokens of \texttt{str}. You may not use \texttt{strtok\_r}.
\end{Qn}
\begin{Sol}
(a)~Sketch:
\begin{lstlisting}[numbers=none]
void *worker(void *arg) {
    char *cp = strtok(str, " ,");
    do puts(cp);
    while ((cp = strtok(NULL, " ,")) != NULL);
    return NULL;
}
/* in main: */
pthread_t t[3];
for (int k = 0; k < 3; k++) pthread_create(&t[k], NULL, worker, NULL);
for (int k = 0; k < 3; k++) pthread_join(t[k], NULL);
\end{lstlisting}

(b)~\texttt{strtok} keeps a \emph{single static} saved pointer and \emph{mutates} \texttt{str} by writing \texttt{NUL}s. Three threads share both. Races:
\begin{itemize}
\item Two threads call \texttt{strtok(str,...)} and both reset the saved pointer; tokens are skipped or repeated.
\item One thread's \texttt{strtok(NULL,...)} uses another thread's saved offset: a thread prints from the middle, or sees \texttt{NULL} immediately.
\item Concurrent writes of \texttt{'\textbackslash0'} tear the string; \texttt{puts} may print fragments, empty lines, or crash on a stale pointer.
\end{itemize}
Example: T1 prints \texttt{Example}, T2 resets and also prints \texttt{Example}, T3 then sees \texttt{NULL} and prints nothing; remaining words vanish.

(c)~One mutex around the \emph{entire} tokenize-and-print loop so a thread walks the whole string without interleaving. Because \texttt{strtok} destroys \texttt{str}, each thread needs its \emph{own} copy:

\begin{lstlisting}[numbers=none]
pthread_mutex_t m = PTHREAD_MUTEX_INITIALIZER;
char str[] = "Example string, containing words, ...";

void *worker(void *arg) {
    char buf[256];
    pthread_mutex_lock(&m);
    strcpy(buf, str);                 /* snapshot under lock if str is shared */
    char *cp = strtok(buf, " ,");     /* strtok only this thread's buf */
    do puts(cp);
    while ((cp = strtok(NULL, " ,")) != NULL);
    pthread_mutex_unlock(&m);
    return NULL;
}
\end{lstlisting}

If the mutex covers the whole worker, even a shared \texttt{strtok} on a per-thread \texttt{buf} is serialized (three full identical listings). That matches ``each individual thread correctly prints all the tokens''. Without the private \texttt{buf}, serializing \texttt{strtok} on the one \texttt{str} makes the second and third threads see an already-consumed string (\texttt{strtok} is not restartable). Private buffers are required; the mutex then only protects \texttt{puts} ordering if desired. The intended full-credit solution is: mutex + per-thread copy of \texttt{str} + \texttt{strtok} on the copy.
\end{Sol}

\begin{Qn}{Hard}{Lect 17}{CS-537 Spring 2016 Mid Q10}{Which spinlock timelines are possible?}
One lock. In the traces, \texttt{c}~=~compute, \texttt{S}~=~spin, \texttt{A}~=~acquire, \texttt{R}~=~release. For each, say what happened, or why the trace is impossible.

(1)
\begin{lstlisting}[numbers=none]
T1: ccccAccccRccccc
T2: cccccccAcccccRccc
\end{lstlisting}

(2)
\begin{lstlisting}[numbers=none]
T1: ccccAccccccccccRcccc
T2: cccccSSSSSSSSSSSSSSS
\end{lstlisting}

(3)~Both traces show \texttt{A} overlapping (both appear to hold the lock).

(4)
\begin{lstlisting}[numbers=none]
T1: ccccAcc
T2: cccSSSSSSSSSSSSS... (forever)
\end{lstlisting}

(5)
\begin{lstlisting}[numbers=none]
T1: AR AR AR AR
T2: AR AR AR AR
\end{lstlisting}

(6)~Both spin forever, nobody ever \texttt{A}s.
\end{Qn}
\begin{Sol}
(1)~\textbf{Possible, good case.} T1 acquires, critical section, releases; then T2 acquires. No overlap of hold intervals. Uncontended (or well-ordered) spinlock on 1 or $N$ CPUs.

(2)~\textbf{Possible on 1 CPU, bad.} T1 holds the lock and is \emph{preempted} (the long \texttt{c} after \texttt{A} is T1 running other code \emph{or} descheduled while still holding). T2 runs and spins the rest of the trace; it cannot acquire until T1's \texttt{R}. Classic uniprocessor spinlock disaster.

(3)~\textbf{Impossible} for a correct lock: two threads cannot both have acquired without an intervening release. If the lock is the \emph{broken flag} lock, it \emph{would} be possible --- but the problem says we are using a spinlock to protect a CS, i.e.\ a working one. Call it BAD.

(4)~\textbf{Possible} if T1 acquired and then was delayed/killed/deadlocked without \texttt{R}, while T2 spins forever. Also possible if T1 is preempted holding the lock and the scheduler is unfair (T1 never runs again). Matches Lecture~17 ``basic spinlocks are unfair''.

(5)~\textbf{Possible on $N$ CPUs}: rapid uncontended acquire/release by both (interleaved across processors). Not ``spin forever''.

(6)~\textbf{Impossible} for a correct lock starting free: the first TAS/xchg succeeds. Possible only if the lock was initialized held and nobody releases, or both wait on a broken protocol. Official OCR notes: ``no one was able to get the lock --- bad''.
\end{Sol}

\begin{Qn}{Medium-Hard}{Lect 17}{CS-537 Spring 2009 Final Q2}{Why spinlocks hurt on 1 CPU and help on $N$ CPUs}
kilOS places each thread on a dedicated application processor (no time sharing).

(1)~On a normal system, why do simple spinlocks sometimes perform poorly on a single processor?

(2)~On kilOS, why would spinlocks perform well between threads on different application processors?

(3)~Write TAS \texttt{lock()} on \texttt{lock->flag} (init 0). Skip unlock.

(4)~Even when they perform ``well'', what negative properties remain?
\end{Qn}
\begin{Sol}
(1)~If the holder is preempted, waiters spin for the rest of the timeslice doing no useful work. Wait cost $\propto T$, not $\propto$ CS length.

(2)~The holder is \emph{running on another CPU} and will release soon (short CS). Spinning is then cheaper than a context switch / park. Lecture~17: spinlocks are for many CPUs + short hold.

(3)
\begin{lstlisting}[numbers=none]
void lock(lock_t *lock) {
    while (TestAndSet(&lock->flag, 1) == 1)
        ;
}
\end{lstlisting}

(4)~Unfair (no FIFO); can starve; cache-line ping-pong under contention; no waiting queue, so the scheduler cannot prefer the lock holder; does not sleep, so it wastes cycles if the hold is long. Ticket locks fix fairness; park/unpark queue locks fix the waste.
\end{Sol}

\begin{Qn}{Medium-Hard}{Lect 17}{CS-537 Spring 2002 Final Q3}{Turning interrupts off as a lock}
(a)~What is a critical section, and why do concurrent programs need them?

(b)~Why is exposing \texttt{TurnOffInterrupts()} / \texttt{TurnOnInterrupts()} to user-level programs dangerous?

(c)~A new primitive \texttt{TurnOffInterrupt(float amount)} disables interrupts for at most \texttt{amount} seconds, then the hardware re-enables them. How does this fix (b)?

(d)~Is the matching \texttt{TurnOnInterrupts()} still useful?

(e)~What about modern systems makes a timed interrupt-off mechanism hard to use?
\end{Qn}
\begin{Sol}
(a)~A CS is a code sequence that must execute as if atomic wrt other threads' CS on the same shared state. Needed so invariants (list structure, counter value) are not broken by interleaving.

(b)~A user thread can disable the timer and I/O interrupts and monopolize the CPU (or deadlock the kernel's ability to regain control). Lecture~17: disable-interrupts is uniprocessor-only and for \emph{brief kernel} use.

(c)~The hardware bound restores the timer, so a malicious/buggy user cannot disable interrupts forever. The OS gets control after \texttt{amount}.

(d)~Yes: if the CS finishes early, re-enable so the rest of the timeslice can take interrupts (I/O, preemption of \emph{other} threads). Otherwise you delay unrelated interrupt handling for the full \texttt{amount}.

(e)~Multiprocessors: disabling interrupts on \emph{this} CPU does not stop another CPU. Also: choosing \texttt{amount} is hard (too small $\Rightarrow$ interrupt in the CS; too large $\Rightarrow$ bad I/O latency). Virtualization / nested interrupt controllers make ``off'' not mean what user code thinks.
\end{Sol}

\subsubsection*{Conceptual Drills, MCQs \& Further Practice}

\begin{Qn}{Easy-Medium}{Lect 16}{Mandar ISI 2011 End Q1(a); repeated Mandar ISI 2016 End Q1(a) and 2022 End Q5(a)}{Shared versus private: process context vs.\ thread context}
(a)~[2011 End] List the most important constituents of a process's context that must be maintained on a \emph{per-thread} basis.

(b)~[2016 End] For a multi-threaded application, is each of the following shared by all threads or not: (i)~global variables, (ii)~local variables, (iii)~open files?

(c)~[2022 End] List 3 per-thread items and 3 per-process items.
\end{Qn}
\begin{Sol}
Lecture~16 split:

\textbf{Private per thread:} thread id (TID), register file (including IP/PC and SP), stack (the stack \emph{lives} in the shared address space but each thread has its own stack pointer and frames). Optionally: thread-local storage, scheduling priority / state of that thread.

\textbf{Shared by all threads of the process:} PID, address space (page table / PTBR, code, heap, globals), open file descriptors, cwd, uid/gid.

Hence:
\begin{itemize}
\item (a)~TID, registers (PC, SP, GPRs), stack pointer / private stack. Not: page table, FD table.
\item (b)~(i)~globals: shared (heap/data). (ii)~locals: private (on that thread's stack), unless the programmer takes their address and publishes it. (iii)~open files: shared (same FD table).
\item (c)~Per-thread: TID, registers, stack. Per-process: PID, address space / heap / code, open FDs.
\end{itemize}

\textbf{Excluded from this topic:} user-level vs.\ kernel-level thread-library trivia (async I/O for ULTs; \texttt{fork} copying all threads vs.\ one) --- those appear on the same Mandar questions as (b)/(c) of 2011/2017/2022 and are out of Lect~16--17.
\end{Sol}

\begin{Qn}{Easy-Medium}{Lect 16}{Mandar ISI 2011 End Q1(c)}{Why thread-local storage?}
Using an example, explain why thread-local storage (TLS) is necessary.
\end{Qn}
\begin{Sol}
Library routines that hide a static/global buffer are unsafe with threads. Classic: \texttt{strtok} (see 2024 End Q3 in Part~I) or \texttt{errno}. Two threads calling \texttt{strtok} concurrently share one saved offset; tokens are lost or mixed. TLS gives each thread its own copy of that hidden state (\texttt{errno}, \texttt{strtok}'s saved pointer, a per-thread RNG seed) without forcing every API to take an extra context pointer. Stack locals are already per-thread; TLS is for data that must persist across calls but not be shared.
\end{Sol}

\begin{Qn}{Easy-Medium}{Lect 16}{CS-537 Fall 2021 Midterm 2 Q20--Q23}{How many threads, and is \texttt{balance} racy?}
\begin{lstlisting}[numbers=none]
void mythread(void *arg) { return; }
thread_t p1, p2;
thread_create(&p1, mythread, NULL);
thread_create(&p2, mythread, NULL);
thread_join(p1);
thread_join(p2);
\end{lstlisting}

Q20. Maximum number of threads that can exist at one moment?
(a)~1 (b)~2 (c)~3 (d)~4 (e)~none

Q21. Each child does \texttt{int result=0; result = result+200; printf("result \%d", result);}. Printed value?
(a)~0 (b)~200 (c)~400 (d)~indeterminate / race (e)~other constant

Q22. Global \texttt{balance=0}. Each child does \texttt{balance = balance+200}. Create p1, \emph{join} p1, create p2, join p2. Final \texttt{balance}?
(a)~0 (b)~200 (c)~400 (d)~race (e)~other

Q23. Same update, but both creates happen before either join. Final \texttt{balance}?
\end{Qn}
\begin{Sol}
Q20. Parent plus two children: 3. (c).

Q21. \texttt{result} is a stack local. Each thread prints 200 (twice). (b).

Q22. Joins serialize the two updates: $0\to200\to400$. No overlap. (c).

Q23. Both children run together; \texttt{balance = balance+200} is a racy load/add/store. Final 200 or 400. (d).
\end{Sol}

\begin{Qn}{Easy-Medium}{Lect 16}{CS-537 Spring 2018 Mid Problem III}{Two workers decrement a counter}
\begin{lstlisting}[numbers=none]
volatile int counter = 1000;
void *worker(void *arg) { counter--; return NULL; }
int main() {
  pthread_t p1, p2;
  pthread_create(&p1, NULL, worker, NULL);
  pthread_create(&p2, NULL, worker, NULL);
  pthread_join(p1, NULL); pthread_join(p2, NULL);
  printf("%d\n", counter);
}
\end{lstlisting}

Which outputs are possible? A~=~possible, B~=~not.
11.~0 \quad 12.~1000 \quad 13.~999 \quad 14.~998 \quad 15.~1002
\end{Qn}
\begin{Sol}
\texttt{counter--} is load, sub, store (Lecture~16's three x86 ops). Serialized: $1000\to999\to998$. Both load 1000, both store 999: lost update, print 999. Cannot print 0, 1000 (both workers ran), or 1002 (nothing adds). Official: 11~B, 12~B, 13~A, 14~A, 15~B.
\end{Sol}

\begin{Qn}{Easy-Medium}{Lect 16}{CS-537 Spring 2016 Final Q8}{\texttt{pthread\_create} / \texttt{pthread\_join} with integer arguments}
\begin{lstlisting}[numbers=none]
void *func(void *arg) {
    int arg_int = (int) arg;
    int x = /* compute from arg_int */;
    return (void *) x;
}
\end{lstlisting}

(a)~Call \texttt{func} directly with argument 100; put the return value in \texttt{int foo}. Casts required.

(b)~Create exactly one thread that calls \texttt{func} with 100.

(c)~\texttt{pthread\_join} that thread and print the integer return value.

(d)~Create 100 threads, pass arguments $1..100$, join all, print the sum of the returned integers.
\end{Qn}
\begin{Sol}
\begin{lstlisting}[numbers=none]
/* (a) */
int foo = (int) func((void *) 100);

/* (b) */
pthread_t p;
pthread_create(&p, NULL, func, (void *) 100);

/* (c) */
void *ret;
pthread_join(p, &ret);
printf("%d\n", (int) ret);

/* (d) */
pthread_t th[100];
int sum = 0;
for (int i = 1; i <= 100; i++)
    pthread_create(&th[i-1], NULL, func, (void *) i);
for (int i = 0; i < 100; i++) {
    void *r;
    pthread_join(th[i], &r);
    sum += (int) r;
}
printf("%d\n", sum);
\end{lstlisting}

Passing small integers through \texttt{void *} is the exam convention; a heap/struct argument is safer in real code (and must not be a loop stack slot --- cf.\ 2018 Mid XXV in Part~I).
\end{Sol}

\begin{Qn}{Easy-Medium}{Lect 17}{Mandar ISI 2015 Mid Q5}{TAS spinlock vs.\ the three CSP properties}
\begin{lstlisting}[numbers=none]
shared char lock = 0;
while (TestAndSet(&lock)) ;  critical_section();  lock = 0;
\end{lstlisting}

For each of the three desirable CSP properties: (i)~state the property; (ii)~does this solution have it?
\end{Qn}
\begin{Sol}
\begin{enumerate}
\item \textbf{Mutual exclusion:} at most one process in the CS. Yes: TAS is atomic; only the thread that sees old value 0 enters, and it leaves \texttt{lock=1} until the store of 0 on exit.
\item \textbf{Progress} (deadlock-freedom): if the CS is free and someone wants in, someone enters. Yes: a free lock is 0, so the next TAS returns 0.
\item \textbf{Bounded waiting} (starvation-freedom / fairness): there is a bound on how many times others may enter after a process has requested the CS. \textbf{No.} TAS spinlocks are unfair (Lecture~17): a thread can lose the TAS race forever. Ticket locks (fetch-and-add) restore FIFO / bounded wait.
\end{enumerate}
\end{Sol}

\begin{Qn}{Easy-Medium}{Lect 17}{CS-537 Fall 2011 Mid Q20; also CS-537 Fall 2004 Mid1 Q5}{Implement a spinlock with atomic exchange}
The atomic exchange (test-and-set) instruction returns the old value of a memory location while setting it to a new value, atomically. Implement a spinlock (\texttt{init}, \texttt{lock}/\texttt{acquire}, \texttt{unlock}/\texttt{release}).
\end{Qn}
\begin{Sol}
\begin{lstlisting}[numbers=none]
typedef struct { int flag; } lock_t;     /* 0 = free, 1 = held */
void init(lock_t *m)    { m->flag = 0; }
void lock(lock_t *m)    { while (xchg(&m->flag, 1) == 1) ; }
void unlock(lock_t *m)  { m->flag = 0; }  /* or xchg(&m->flag, 0) */
\end{lstlisting}

\texttt{xchg(\&flag,1)} (TAS) atomically does \texttt{old = *addr; *addr = 1; return old}. Spin while that previous value was already 1. Unlock publishes 0 with a plain store. Two threads cannot both see old$=0$ because the read-and-write is one memory operation (ME). It is not FIFO. On 1~CPU the holder must be preemptable or waiters spin a full timeslice (Lecture~17).
\end{Sol}

\begin{Qn}{Easy-Medium}{Lect 17}{CS-537 Spring 2002 Mid Q10}{Disable interrupts: how, why it works, why it is a bad idea}
(a)~How can turning interrupts on/off implement a critical section?

(b)~Why does it work?

(c)~Why is it generally a bad idea?
\end{Qn}
\begin{Sol}
(a)
\begin{lstlisting}[numbers=none]
disable_interrupts();
/* critical section */
enable_interrupts();
\end{lstlisting}

(b)~On a uniprocessor the only way another thread runs is a timer (or other) interrupt that enters the scheduler. Masking interrupts makes the CS atomic wrt other threads on that CPU.

(c)~User code can hang the machine; I/O is delayed; multiprocessors are not covered; cannot disable for long. The kernel may use it briefly around a TAS or around the inner \texttt{guard} of a \texttt{park}/\texttt{unpark} lock (Lecture~17).
\end{Sol}

\begin{Qn}{Medium}{Lect 16}{CS-537 Fall 2021 Midterm 2 Q15}{Join the child that was passed \texttt{\&x}}
\begin{lstlisting}[numbers=none]
void child(void *arg) {
    int *p = (int *) arg;
    printf("%d\n", *p);
}
int main() {
    int x = 3;
    thread_t p1;
    thread_create(&p1, child, &x);
    thread_join(p1);
    return 0;
}
\end{lstlisting}

What will this print?
(a)~3 \quad (b)~nothing, then exit \quad (c)~crash \quad (d)~3, or nothing, or crash \quad (e)~none
\end{Qn}
\begin{Sol}
(a)~3. The parent passes a pointer to its live stack variable \texttt{x} and \emph{joins} before returning, so \texttt{x} remains valid for the whole child. Contrast Q14 in Part~I (no join, \texttt{p} published late).
\end{Sol}

\begin{Qn}{Easy-Medium}{Lect 16}{CS-537 Spring 2002 Final Q3(a)}{Race, critical section, mutual exclusion, atomicity}
What is a critical section, and why are they important to concurrent programs? Define, in one or two sentences each, race condition, mutual exclusion, and atomicity, using \texttt{counter++} as the witness.
\end{Qn}
\begin{Sol}
A \textbf{critical section} is a code sequence that touches shared state and must execute as if atomic wrt other threads' CSs on that state. Concurrent programs need them so invariants (a list's pointer structure, a counter's value) are not broken by interleaving.

\begin{itemize}
\item \textbf{Race condition:} the outcome depends on a non-deterministic interleaving; some interleavings violate the intended invariant (lost update of \texttt{counter}).
\item \textbf{Mutual exclusion:} at most one thread executes a CS for that state at a time.
\item \textbf{Atomicity:} the CS appears to happen all-at-once; intermediate register values are never visible. A lock (or an atomic RMW) restores atomicity of \texttt{counter++}.
\end{itemize}

Non-determinism of whole-CS order is not itself a bug; \emph{partial} interleaving of a CS is. (Parts (b)--(e) of this same 2002 Final question are the timed-interrupt lock in Part~I.)
\end{Sol}

\subsection{Lecture 18: Condition Variables}

\subsubsection*{Conceptual Drills, MCQs \& Further Practice}

\begin{Qn}{Easy-Medium}{Lect 18}{CS-537 Fa11 Mid, Q21}{Define a CV and show a join snippet}
Condition variables are useful in certain types of multi-threaded
programs. Describe what a condition variable is, and show an example of
how you would use one in a little code snippet.
\end{Qn}
\begin{Sol}
A condition variable is an \textbf{explicit queue} of sleeping threads,
with two operations:
\begin{itemize}
\item \texttt{wait(c, m)} --- the caller must hold mutex \texttt{m};
      \texttt{wait} atomically releases \texttt{m} and sleeps on
      \texttt{c}'s queue; before returning it reacquires \texttt{m}.
\item \texttt{signal(c)} --- wake \emph{one} waiter, if any. If the
      queue is empty, Mesa \texttt{signal} is a no-op (not remembered).
\item \texttt{broadcast(c)} --- wake \emph{all} waiters.
\end{itemize}
CVs enforce \textbf{ordering} (``do not proceed until that other thread
has done $X$''). Locks enforce \textbf{mutual exclusion}. You almost
always use both: the mutex protects the state variable that records
whether $X$ has happened; the CV is how a thread sleeps until that
state changes.

Official-key example: parent waiting for a child (join).
\begin{lstlisting}[numbers=none]
int done = 0;
mutex_t m;  cond_t c;

void child() {                       void parent_join() {
    /* child's work */                   Mutex_lock(&m);
    Mutex_lock(&m);                      while (done == 0)
    done = 1;                                Cond_wait(&c, &m);
    Cond_signal(&c);                     Mutex_unlock(&m);
    Mutex_unlock(&m);
}
\end{lstlisting}
A signal with no waiter is lost unless \texttt{done} records the
completion; that is why the state variable is not optional.
\end{Sol}

\begin{Qn}{Medium}{Lect 18}{CS-537 Sp09 Final, Q4}{kilOS condition variables: vs.\ locks, vs.\ the scheduler, vs.\ spinning}
kilOS provides condition variables. Answer each part.
\begin{enumerate}[(a)]
\item In a typical system, how are CVs \emph{different} from locks?
\item In a typical system, how are CVs \emph{similar} to locks? In
      particular, in what cases would these primitives interact with
      the system scheduler?
\item In kilOS, \texttt{wait()} on a CV just \textbf{spins} until the
      condition becomes true, instead of putting the caller to sleep.
      Is this simple implementation OK for kilOS? Why or why not?
\end{enumerate}
\end{Qn}
\begin{Sol}
(a)~A lock provides \textbf{mutual exclusion}: at most one holder.
A CV provides \textbf{ordering}: a thread waits for a predicate on
shared state to become true. A lock has a binary held/free state; a
CV has \emph{no} integer state, only a queue. \texttt{unlock} always
has an effect (the lock is free); Mesa \texttt{signal} with no waiter
is forgotten. You can (and should) hold a lock while signalling; you
must not \texttt{wait} without holding the associated mutex.

(b)~Both may \textbf{block} the caller and therefore call into the
scheduler: \texttt{lock} blocks when the lock is held;
\texttt{cond\_wait} blocks (after releasing the mutex) when the
predicate is false. Both also \textbf{wake} a thread
(\texttt{unlock} may hand the lock to a waiter; \texttt{signal}
moves a CV-waiter to ready). After wakeup the thread becomes
runnable and the scheduler may switch to it. On a uniprocessor a
spinning lock \emph{must} yield or be preemptable, otherwise the
holder never runs; a sleeping CV wait is the blocking analogue.

(c)~\textbf{Not OK} as a general CV, for the same reason a spinlock
is a bad lock on a uniprocessor with a long critical section: the
waiter burns a CPU until the signaller runs. If kilOS is
uniprocessor (or has fewer CPUs than spinning waiters), a spinning
\texttt{wait} can prevent the thread that would make the condition
true from ever being scheduled --- livelock / complete hang. Even
on a multiprocessor it wastes cycles. Sleeping \texttt{wait} is the
point of a CV (Lect 18: spin-on-\texttt{done} is correct but
inefficient; the CV exists to avoid that spin). A spinning wait
is also not a Mesa \texttt{wait}: it does not release the mutex
while ``waiting'', so if the caller spins \emph{holding} the lock
the signaller can never take the lock to update the condition.
\end{Sol}

\begin{Qn}{Easy-Medium}{Lect 18}{CS-537 Fa21 Mid2, Q32}{Two CVs named \texttt{bug} and \texttt{patterns}}
Bug ``fixes'' the previous code:
\begin{lstlisting}[numbers=none]
/* producer */                       /* consumer */
mutex_lock(&m);                      mutex_lock(&m);
while (numfull == max)               while (numfull == 0)
    cond_wait(&bug, &m);                 cond_wait(&patterns, &m);
do_fill(i);                          int tmp = do_get();
cond_signal(&patterns);              cond_signal(&bug);
mutex_unlock(&m);                    mutex_unlock(&m);
\end{lstlisting}
What is wrong from a correctness standpoint?
\begin{enumerate}[(a)]
\item Too much signalling and waiting
\item Signal signals the wrong CV; wait waits on the wrong one
\item Producer can incorrectly wake a producer
\item Consumer can incorrectly wake a consumer
\item None of the above
\end{enumerate}
\end{Qn}
\begin{Sol}
\textbf{(e)}, official key. ``Well done, Bug.''

This is the textbook two-CV solution with odd names:
producers wait on \texttt{bug} (``not full'' / empty slots) and signal
\texttt{patterns} (``not empty''); consumers wait on \texttt{patterns}
and signal \texttt{bug}. Same-kind wakeup is impossible, \texttt{while}
re-checks, the mutex covers \texttt{numfull}. (a)~is a performance
complaint, not a correctness bug. (b)--(d)~are false: the queues are
the right way around.
\end{Sol}

\begin{Qn}{Easy-Medium}{Lect 18}{CS-537 Sp02 Mid, Q7}{How many CVs for an unbounded buffer?}
You are implementing a producer--consumer shared buffer that producers
use to pass data to consumers, but the buffer is \textbf{unbounded}
(no limit on how big it can get).
\begin{enumerate}[(a)]
\item How many condition variables will you need in order to implement
      this buffer properly, and why?
\item How is this different than a standard bounded-buffer
      implementation?
\end{enumerate}
\end{Qn}
\begin{Sol}
(a)~\textbf{One} CV, plus one mutex and a \texttt{count} (or equivalent
empty-test). Only consumers ever wait (buffer empty). Producers never
wait, so there is no second predicate and no second queue.
\begin{lstlisting}[numbers=none]
mutex_t m;  cond_t not_empty;  int count = 0;  /* plus a list */

void put(item) {
    mutex_lock(&m);
    enqueue(item); count++;
    cond_signal(&not_empty);     /* one consumer is enough */
    mutex_unlock(&m);
}
item get() {
    mutex_lock(&m);
    while (count == 0)
        cond_wait(&not_empty, &m);
    item x = dequeue(); count--;
    mutex_unlock(&m);
    return x;
}
\end{lstlisting}
Same-kind wakeup of extra consumers is harmless: they re-check
\texttt{while (count == 0)} and sleep. No producer is sleeping on
that queue, so the everyone-asleep bug cannot occur.

(b)~A bounded buffer has two wait reasons (empty \emph{and} full) and
two kinds of thread, so Mesa's single-CV bug appears: you need
\texttt{empty} and \texttt{fill} (or one CV plus \texttt{broadcast}).
Unbounded drops the ``full'' predicate. The mutex is still required
(the list/\texttt{count} is shared), and \texttt{while} is still
required (several consumers may be waiting; one \texttt{put} enables
only one \texttt{get}).
\end{Sol}

\subsection{Lecture 19: Semaphores}

\subsubsection*{Additional Challenging \& Analytical Problems}

\begin{Qn}{Medium-Hard}{Lect 19}{Mandar Mitra, ISI Kolkata, 2017 Mid-sem Q3}{Atomic wait/signal with no TAS and no interrupt disable}
A counting semaphore $S$ is accessed only through
\begin{lstlisting}[numbers=none]
wait(S) { while (S <= 0) ;  S--; }     signal(S) { S++; }
\end{lstlisting}
A proper implementation must ensure (i)~two processes never change $S$
concurrently, and (ii)~two processes never concurrently execute the two
statements in the body of \texttt{wait}. Implement \texttt{wait} and
\texttt{signal} \emph{without} TAS, interrupt disabling, or other special
instructions. You may assume load/store/test are atomic, and you may use
a standard software solution of the $n$-process critical-section problem
(identify which).
\end{Qn}
\begin{Sol}
Wrap both bodies in the entry/exit sections of a software CS algorithm
that already guarantees mutual exclusion, progress, and (if you need
bounded waiting) bounded waiting. For $n=2$, Peterson; for general $n$,
Lamport's bakery (or Eisenberg--McGuire). The two semaphore operations
share one CS, so they never interleave on $S$.

\begin{lstlisting}[numbers=none]
/* Bakery: ticket[n], choosing[n]; process id i */
wait(S):                             signal(S):
  enter_bakery(i);                     enter_bakery(i);
  while (S <= 0) {                     S++;
      exit_bakery(i);                  exit_bakery(i);
      enter_bakery(i);               }
  }
  S--;
  exit_bakery(i);
\end{lstlisting}
The inner \texttt{while} must \emph{drop} the bakery CS; otherwise a
zero $S$ would spin holding the CS and deadlock with every
\texttt{signal}. That is the same lock-ordering lesson as
\texttt{wait(mutex)} before \texttt{wait(empty)}. Requirements (i) and
(ii) hold because $S$\texttt{--} and $S$\texttt{++} and the test
\texttt{S <= 0} all occur inside \texttt{enter}/\texttt{exit}. (This is
a \emph{busy-wait} semaphore. A blocking implementation also needs a
sleep queue, which cannot be built from loads and stores alone.)
\end{Sol}

\begin{Qn}{Medium-Hard}{Lect 19}{Mandar Mitra, ISI Kolkata, 2022 End-sem Q1}{Blocking wait/signal inside the bakery}
The blocking counting-semaphore operations are
\begin{lstlisting}[numbers=none]
void wait(semaphore S) {            void signal(semaphore S) {
    S.value--;                          S.value++;
    if (S.value < 0) {                  if (S.value <= 0) {
        add process to S.L;                 remove process P from S.L;
        sleep();                            wakeup(P);
    }                                   }
}                                   }
\end{lstlisting}
A proper implementation ensures that two processes never change
\texttt{S.value} or \texttt{S.L} concurrently, and that lines 3--5 of
\texttt{wait} are effectively atomic. Given
\texttt{enter\_bakery(id)} / \texttt{exit\_bakery(id)}, implement
\texttt{wait} and \texttt{signal}.
\end{Qn}
\begin{Sol}
The bakery lock is the mutex around the semaphore internals. As in the
previous problem, a process that will sleep must leave the bakery
\emph{before} calling \texttt{sleep}, but it must be on \texttt{S.L}
before it drops the lock (otherwise a concurrent \texttt{signal} can
miss it: the lost-wakeup of Lect 18).
\begin{lstlisting}[numbers=none]
void wait(semaphore S) {
    enter_bakery(id);
    S.value--;
    if (S.value < 0) {
        add self to S.L;
        exit_bakery(id);     /* still on S.L */
        sleep();
        return;
    }
    exit_bakery(id);
}
void signal(semaphore S) {
    enter_bakery(id);
    S.value++;
    if (S.value <= 0) {
        remove process P from S.L;
        wakeup(P);
    }
    exit_bakery(id);
}
\end{lstlisting}
Lines 3--5 of wait (decrement, test, enqueue) all run inside one bakery
critical section, so they are atomic wrt other \texttt{wait}/
\texttt{signal}s. Negative \texttt{S.value} equals $|S.L|$, matching
Lect 19.
\end{Sol}

\begin{Qn}{Medium}{Lect 19}{CS-537 Spring 2001 Midterm Q7}{Initial values of \texttt{mutex} and \texttt{signal}}
Two threads, two semaphores \texttt{mutex} and \texttt{signal}. $x$ is
always $\ge 0$, and these are the only two threads.
\begin{lstlisting}[numbers=none]
Thread 1                            Thread 2
mutex.P();                          mutex.P();
if (x > 0)                          x++;
    signal.V();                     signal.V();
mutex.V();                          mutex.V();
signal.P();
\end{lstlisting}
\texttt{mutex} should provide mutual exclusion; \texttt{signal} should
let thread 2 activate thread 1 when $x > 0$. Give both initial values.
\end{Qn}
\begin{Sol}
\texttt{mutex = 1} (binary lock: the first \texttt{P} must succeed).
\texttt{signal = 0} (thread 1 must block at \texttt{signal.P()} until
some \texttt{V} has happened).

If $x$ is already $> 0$, thread 1 posts \texttt{signal} itself and then
\texttt{P}s it, so it does not rely on thread 2 --- correct, because the
activation condition already holds. If $x=0$ and thread 1 runs first, it
skips the \texttt{V} and blocks on \texttt{P}; thread 2 increments $x$
and posts, waking thread 1. If thread 2 runs first, it posts
\texttt{signal} (remembered); thread 1 then either posts a second time
(if it also sees $x>0$) and consumes one, or just consumes the remembered
post. Extra posts are harmless because nobody else waits on
\texttt{signal}.

\texttt{signal = 1} would let thread 1 pass \texttt{signal.P()} even when
$x=0$ and thread 2 has not run. \texttt{mutex = 0} deadlocks both
threads immediately.
\end{Sol}

\subsubsection*{Conceptual Drills, MCQs \& Further Practice}

\begin{Qn}{Easy-Medium}{Lect 19}{CS-537 Spring 2018 Midterm, Problem XXVII}{T/F: semaphore semantics}
A semaphore is a useful synchronisation primitive. For each statement fill
\textbf{A} if true, \textbf{B} if not true.

\begin{enumerate}
  \item[131.] Each semaphore has an integer value.
  \item[132.] If a semaphore is initialised to 1, it can be used as a lock.
  \item[133.] Semaphores can be initialised to values higher than 1.
  \item[134.] A single lock and condition variable can be used in tandem to
        implement a semaphore.
  \item[135.] Calling \texttt{sem\_post()} may block, depending on the
        current value of the semaphore.
\end{enumerate}
\end{Qn}
\begin{Sol}
Official CS-537 Spring 2018 answers:

131.~\textbf{A (True).} By definition a semaphore carries an integer
value (Lect 19: that is the whole difference from a CV).

132.~\textbf{A (True).} This is a binary semaphore / mutex; first
\texttt{wait} takes it from 1 to 0, a second waiter sleeps.

133.~\textbf{A (True).} Counting semaphores: $N$ resource units
(\texttt{emptyBuffer} init $N$, a connection pool of size 4, \ldots).

134.~\textbf{A (True).} The Zemaphore: integer + lock + CV (Lect 19
equivalence).

135.~\textbf{B (False).} Only \texttt{sem\_wait} blocks.
\texttt{sem\_post} increments, wakes at most one waiter, and returns.
\end{Sol}

\begin{Qn}{Easy-Medium}{Lect 19}{CS-537 Spring 2018 Midterm, Problem XXVIII}{T/F: producer/consumer initial values}
Classic semaphore producer/consumer:
\begin{lstlisting}[numbers=none]
void *producer(void *arg) {
    for (i = 0; i < num; i++) {
        sem_wait(&empty); sem_wait(&mutex);
        put(i);
        sem_post(&mutex); sem_post(&full);
    }
}
void *consumer(void *arg) {
    while (!done) {
        sem_wait(&full); sem_wait(&mutex);
        int tmp = get(i);
        sem_post(&mutex); sem_post(&empty);
    }
}
\end{lstlisting}
A $=$ true, B $=$ not true.

\begin{enumerate}
  \item[136.] The semaphore \texttt{full} must be initialised to 0.
  \item[137.] The semaphore \texttt{full} must be initialised to 1.
  \item[138.] The semaphore \texttt{empty} must be initialised to 1.
  \item[139.] The semaphore \texttt{empty} \emph{can} be initialised to 1.
  \item[140.] The semaphore \texttt{mutex} must be initialised to 1.
\end{enumerate}
\end{Qn}
\begin{Sol}
Official CS-537 Spring 2018 answers:

136.~\textbf{A (True).} \texttt{full} counts full slots; the buffer
starts empty, so 0. A consumer must wait until a producer posts.

137.~\textbf{B (False).} Directly contradicts 136.

138.~\textbf{B (False).} \texttt{empty} counts empty slots and must be
initialised to the buffer size. That is 1 only for a 1-slot buffer, not
in general.

139.~\textbf{A (True).} For a size-1 buffer, \texttt{empty = 1} is
exactly right (Lect 19 single-slot solution).

140.~\textbf{A (True).} \texttt{mutex} is a lock, so init 1.
\end{Sol}

\begin{Qn}{Easy-Medium}{Lect 19}{CS-537 Fall 2021 Midterm 2 Q6}{Join with one child: what initial value?}
Bug uses a semaphore to wait for a child thread. Textbook / Lect 19
\texttt{sem\_wait} decrements, then sleeps if the value is negative.
\begin{lstlisting}[numbers=none]
sem_t s;
void child(void *arg) { /* ... */ sem_post(&s); }
int main(void) {
    thread_t p1;
    sem_init(XXX);
    thread_create(&p1, child);
    sem_wait(&s);
    return 0;
}
\end{lstlisting}
What should $s$ be initialised to?
\begin{enumerate}
  \item[(a)] 0
  \item[(b)] 1
  \item[(c)] 2
  \item[(d)] 3
  \item[(e)] None of the above
\end{enumerate}
\end{Qn}
\begin{Sol}
Official CS-537 Fall 2021 Midterm 2 answer: \textbf{(a) 0}.

Init 0: if the parent waits first, value goes $0\to -1$ and the parent
sleeps; the child's post takes it to 0 and wakes the parent. If the
child posts first, value goes $0\to 1$ (the post is remembered); the
parent's wait decrements to 0 and returns immediately. Init 1 would let
the parent pass \texttt{sem\_wait} before the child had run --- the
join would not join.
\end{Sol}

\begin{Qn}{Medium}{Lect 19}{CS-537 Fall 2021 Midterm 2 Q7}{Join with three children: what initial value?}
Same code as the previous question, except the parent creates three
children and then calls \texttt{sem\_wait(\&s)} \emph{once}.
\begin{lstlisting}[numbers=none]
sem_init(XXX);
thread_create(&p1, child);
thread_create(&p2, child);
thread_create(&p3, child);
sem_wait(&s);
\end{lstlisting}
What should $s$ be initialised to?
\begin{enumerate}
  \item[(a)] 0
  \item[(b)] 1
  \item[(c)] 2
  \item[(d)] 3
  \item[(e)] None of the above
\end{enumerate}
\end{Qn}
\begin{Sol}
Official CS-537 Fall 2021 Midterm 2 answer: \textbf{(e) None of the above}.

A single \texttt{sem\_wait} on a 0-init semaphore returns after
\emph{one} post, so the parent would proceed after the first child,
not the third. Under the exam-room reading in which the parent waits
until the value becomes positive and children only post, the right
initial value is $-2$ ($1-N$ for $N=3$): three posts take $-2 \to -1
\to 0 \to 1$, and only then does the parent proceed. $-2$ is not among
(a)--(d).

(The Lect 19 / OSTEP join is ``one child, init 0''. For $N$ children
the clean code is init 0 and a loop \texttt{for (i=0;i<N;i++)
sem\_wait(\&s);}, which \emph{is} among the spirit of (a) but is not
what this stem asked.)
\end{Sol}

\begin{Qn}{Medium}{Lect 19}{CS-537 Fall 2021 Midterm 2 Q8}{A lock and a semaphore protecting the same counter}
\begin{lstlisting}[numbers=none]
int count = 0;
mutex_t m;
sem_t s;            /* initialised to 1 */
/* Thread 1 */                 /* Thread 2 */
mutex_acquire(&m);             sem_wait(&s);
count++;                       count++;
mutex_release(&m);             sem_post(&s);
\end{lstlisting}
The final value of \texttt{count} is:
\begin{enumerate}
  \item[(a)] 0
  \item[(b)] 1
  \item[(c)] 2
  \item[(d)] 1 or 2
  \item[(e)] None of the above
\end{enumerate}
\end{Qn}
\begin{Sol}
Official answer: \textbf{(d) 1 or 2}.

The mutex and the semaphore are \emph{different} objects. Thread 1's
critical section is not exclusive of thread 2's. They can run one after
the other (\texttt{count = 2}) or both load 0, both increment, both
store 1 (\texttt{count = 1}). Using a lock and a semaphore to ``protect
the same critical section'' does not prevent the data race.
\end{Sol}

\begin{Qn}{Medium}{Lect 19}{CS-537 Fall 2021 Midterm 2 Q11--Q13}{Semaphore as a lock: init 1, 0, and 2}
In each variant two threads run
\begin{lstlisting}[numbers=none]
sem_wait(&s);
count--;
sem_post(&s);
\end{lstlisting}
with Lect 19 / textbook \texttt{sem\_wait} (decrement; sleep if
negative). \texttt{count} starts at 0.

\begin{enumerate}
  \item[(a)] \textbf{Q11.} $s$ is initialised to 1. Final \texttt{count}?
  \item[(b)] \textbf{Q12.} $s$ is initialised to 0. Final \texttt{count}?
  \item[(c)] \textbf{Q13.} $s$ is initialised to 2. Final \texttt{count}?
\end{enumerate}
Options for each: (a)~0\quad (b)~$-1$\quad (c)~$-2$\quad
(d)~$-1$ or $-2$\quad (e)~none of the above.
\end{Qn}
\begin{Sol}
Official CS-537 Fall 2021 Midterm 2 answers (textbook definition):

(a)~\textbf{Q11: (c) $-2$.} Init 1 is a lock. Each thread's decrement
runs exclusively, so \texttt{count} goes $0\to -1\to -2$.

(b)~\textbf{Q12: (a) 0.} Init 0: the first \texttt{wait} takes the
value to $-1$ and sleeps; the second takes it to $-2$ and sleeps. Nobody
ever \texttt{post}s, so both threads are stuck and \texttt{count}
stays 0. (This is the ``lock initialised to 0'' deadlock of Lect 19.)

(c)~\textbf{Q13: (d) $-1$ or $-2$.} Init 2 lets both threads into the
critical section at once. They may serialise (\texttt{count = -2}) or
both load 0, both store $-1$ (\texttt{count = -1}).
\end{Sol}

\begin{Qn}{Easy-Medium}{Lect 19}{CS-537 Fall 2021 Midterm 2 Q27--Q28}{Reader/writer lock: \texttt{LOCK\_INIT} and \texttt{WRITELOCK\_INIT}}
The Lect 19 / OSTEP reader/writer lock:
\begin{lstlisting}[numbers=none]
typedef struct _rwlock_t {
    sem_t lock; sem_t writelock; int readers;
} rwlock_t;
void rwlock_init(rwlock_t *rw) {
    rw->readers = 0;
    sem_init(&rw->lock, LOCK_INIT);
    sem_init(&rw->writelock, WRITELOCK_INIT);
}
void rwlock_acquire_readlock(rwlock_t *rw) {
    sem_wait(&rw->lock);
    rw->readers++;
    if (rw->readers == 1) sem_wait(&rw->writelock);
    sem_post(&rw->lock);
}
void rwlock_release_readlock(rwlock_t *rw) {
    sem_wait(&rw->lock);
    rw->readers--;
    if (rw->readers == 0) sem_post(&rw->writelock);
    sem_post(&rw->lock);
}
void rwlock_acquire_writelock(rwlock_t *rw) {
    sem_wait(&rw->writelock);
}
\end{lstlisting}
\begin{enumerate}
  \item[(a)] \textbf{Q27.} Value of \texttt{LOCK\_INIT}?
  \item[(b)] \textbf{Q28.} Value of \texttt{WRITELOCK\_INIT}?
\end{enumerate}
Options: (a)~0\quad (b)~1\quad (c)~2\quad (d)~$N$ (number of readers you
wish to allow)\quad (e)~none of the above.
\end{Qn}
\begin{Sol}
Official CS-537 Fall 2021 Midterm 2 answers (textbook / Lect 19):

(a)~\textbf{Q27: (b) 1.} \texttt{lock} protects the \texttt{readers}
counter, i.e.\ it is used as a mutex, so init 1.

(b)~\textbf{Q28: (b) 1.} \texttt{writelock} lets in one writer \emph{or}
the whole set of readers (the first reader takes it). Init 1. Init 0
would deadlock the first reader and the first writer; init $N$ would
let $N$ writers in at once.
\end{Sol}

\begin{Qn}{Easy-Medium}{Lect 19}{Mandar Mitra, ISI Kolkata, 2018 End-sem Q1(a)}{Busy-wait binary wait and signal}
The integer value of a binary semaphore can only be 0 or 1. Write
pseudo-code for the \emph{busy-waiting} \texttt{wait} and
\texttt{signal} on a binary semaphore $S$. Increment (\texttt{++}) and
decrement (\texttt{--}) operators do \emph{not} apply. (Atomicity of
the bodies is not required in this part.)
\end{Qn}
\begin{Sol}
\begin{lstlisting}[numbers=none]
wait(S) {                            signal(S) {
    while (S.value == 0) ;               S.value = 1;
    S.value = 0;                     }
}
\end{lstlisting}
\texttt{wait} spins until the lock is free, then takes it by writing 0.
\texttt{signal} simply writes 1; it never blocks (Lect 19:
\texttt{sem\_post} never waits). Using \texttt{S++}/\texttt{S--} would
let the value leave $\{0,1\}$.
\end{Sol}

\subsection{Lecture 20: Concurrency Bugs \& Deadlock}

\subsubsection*{Additional Challenging \& Analytical Problems}

\begin{Qn}{Hard}{Lect 20}{CS-537 Fall 2008 Final, Q4}{Deadlock flows: two threads, inverted mutex pair}
Assume two threads may be arbitrarily interleaved. They repeatedly run:
\begin{lstlisting}[numbers=none]
Thread 1                          Thread 2
1. lock(&mutex1);                 a. lock(&mutex2);
2. lock(&mutex2);                 b. lock(&mutex1);
3. // critical section            c. // critical section
4. unlock(&mutex2);               d. unlock(&mutex1);
5. unlock(&mutex1);               e. unlock(&mutex2);
\end{lstlisting}
A \emph{flow} is a sequence of these labels (as in the rest of that exam).
Write:
\begin{enumerate}
  \item[(a)] a flow that leads to deadlock;
  \item[(b)] a deadlock flow with the \emph{fewest} possible context
        switches;
  \item[(c)] a deadlock flow with the \emph{most} possible context switches
        that does not re-use a flow element more than once, with Thread~1
        starting just before ``1'' and Thread~2 just before ``a'';
  \item[(d)] a flow that has one thread in the critical section and the
        other waiting on a lock, and \emph{no} deadlock, again starting
        Thread~1 just before ``1'' and Thread~2 just before ``a'';
  \item[(e)] a flow that has one thread holding \emph{both} locks and the
        other holding just one.
\end{enumerate}
\end{Qn}
\begin{Sol}
Classic lock-order inversion on $\{\texttt{mutex1},\texttt{mutex2}\}$.

(a)~Any flow in which each thread takes its first lock before either takes
its second. One such:
\[
1,\; a,\; 2,\; b.
\]
After $1$, Thread~1 holds \texttt{mutex1}; after $a$, Thread~2 holds
\texttt{mutex2}; $2$ and $b$ then both wait. Cycle, deadlock.

(b)~Fewest context switches: $1,\; a$ (one switch, Thread~1$\to$Thread~2).
Each thread now holds one lock of the inverted pair; the next acquire on
either side blocks. (The expanded form $1,a,2,b$ is the same deadlock
with the blocking acquires written out.)

(c)~Longest deadlock flow with unique labels, Thread~1 first: $1, a, 2, b$.
You cannot include $3$--$5$ or $c$--$e$: those sit \emph{after} a thread
already holds both locks, which is exactly the non-deadlock case. Three
switches (alternate threads until both are stuck).

(d)~Thread~1 takes both locks and enters the CS; Thread~2 then waits for
\texttt{mutex2}:
\[
1,\; 2,\; 3,\; a.
\]
Thread~1 is in the critical section holding both mutexes; Thread~2 is
blocked on $a$. No cycle: Thread~1 will run $4,5$ and release. (The
shorter $1,2,a$ is the same idea before the CS comment.)

(e)~Empty / impossible. If Thread~1 has executed $1$ and $2$ it holds
\emph{both} mutexes, so Thread~2 holds none. If Thread~2 has executed $a$
it holds \texttt{mutex2}, so Thread~1 cannot acquire both. The two-lock
set cannot be partitioned into ``both'' plus ``one''.
\end{Sol}

\subsubsection*{Conceptual Drills, MCQs \& Further Practice}

\begin{Qn}{Easy-Medium}{Lect 20}{CS-537 Spring 2001 Midterm, Q9}{Which of these does \emph{not} guarantee deadlock is avoided?}
Which of the following will NOT guarantee that deadlock is avoided? Please
circle all that apply.
\begin{enumerate}
  \item[(a)] Acquire all resources (locks) all at once, atomically
  \item[(b)] Use locks sparingly
  \item[(c)] Acquire resources (locks) in a fixed order
  \item[(d)] Be willing to release a held lock if another lock you want is
        held, and then try the whole thing over again
\end{enumerate}
\end{Qn}
\begin{Sol}
\textbf{(b)} only.

(a)~Breaks hold-and-wait (Lect~20 meta-lock / ``acquire the whole set
atomically''). Deadlock is impossible; the cost is utilization and
encapsulation.

(c)~Breaks circular wait (global lock order). Deadlock is impossible.

(d)~Breaks no-preemption. This is \texttt{trylock}: on failure, drop what
you hold and retry. Deadlock is impossible. (Livelock is possible, which
is why Lect~20 adds exponential back-off --- but livelock is not
deadlock.)

(b)~``Use locks sparingly'' is advice, not a protocol. Two locks used
``sparingly'' in opposite orders still deadlock (the
\texttt{vector\_add} pattern). It does not \emph{guarantee} anything.

No other option fails the guarantee, so the only circle is (b).
\end{Sol}

\begin{Qn}{Medium}{Lect 20}{CS-537 Spring 2018 Midterm, Problem XXIX (Q141--145)}{T/F: deadlock possible under which schedules?}
One way to avoid deadlock is to schedule threads carefully. Assume the
following characteristics of threads $T_1$, $T_2$, and $T_3$:
\begin{itemize}
  \item $T_1$ (at some point) acquires and releases locks $L_1$, $L_2$
  \item $T_2$ (at some point) acquires and releases locks $L_1$, $L_3$
  \item $T_3$ (at some point) acquires and releases locks $L_3$, $L_1$,
        and $L_4$
\end{itemize}
For which schedules below is deadlock possible?
To answer: fill in \textbf{A} if deadlock is possible, \textbf{B} for not
possible.
\begin{enumerate}
  \item[141.] $T_1$ runs to completion, then $T_2$ to completion, then
        $T_3$ runs
  \item[142.] $T_1$ and $T_2$ run concurrently to completion, then $T_3$
        runs
  \item[143.] $T_1$, $T_2$, and $T_3$ run concurrently
  \item[144.] $T_3$ runs to completion, then $T_1$ and $T_2$ run
        concurrently
  \item[145.] $T_1$ and $T_3$ run concurrently to completion, then $T_2$
        runs
\end{enumerate}
\end{Qn}
\begin{Sol}
Official CS-537 Spring 2018 answers. Key (from the answer sheet):
``Do threads that CAN deadlock ever run at the same time? If so, they
might deadlock. If not, it isn't. From the above, only $T_2$ and $T_3$
can deadlock, because they each grab two locks ($L_1$ and $L_3$) in some
order.''

$T_1$ shares only $L_1$ with $T_2$ and only $L_1$ with $T_3$ --- one
shared lock cannot close a cycle. $T_3$'s extra $L_4$ is irrelevant.
$T_2$ and $T_3$ share \emph{two} locks with unspecified order, so
lock-order inversion is possible \emph{exactly when those two run
concurrently}.

\begin{itemize}
  \item 141.~\textbf{B (not possible).} $T_2$ and $T_3$ never overlap.
  \item 142.~\textbf{B (not possible).} $T_2$ and $T_3$ never overlap.
  \item 143.~\textbf{A (possible).} $T_2$ and $T_3$ run at the same time.
  \item 144.~\textbf{B (not possible).} $T_2$ and $T_3$ never overlap.
  \item 145.~\textbf{B (not possible).} $T_2$ and $T_3$ never overlap.
\end{itemize}

Scheduling that never co-schedules a deadlock-capable pair is a
(heavy-handed) prevention trick; the cheap fix remains a global lock
order on $\{L_1,L_3\}$.
\end{Sol}

\begin{Qn}{Easy-Medium}{Lect 20}{Mandar Mitra, ISI Kolkata, 2012 End-sem Q1(a)}{What a RAG can tell you about deadlock}
Given a resource-allocation graph for a set of processes and resources,
what possible conclusions can be drawn about the processes with respect
to deadlock?
\end{Qn}
\begin{Sol}
Three standard conclusions (Coffman circular wait vs.\ actual deadlock):
\begin{enumerate}
  \item If the RAG has \textbf{no cycle}, the system is \textbf{not}
        deadlocked. Circular wait is necessary; without it, deadlock is
        impossible.
  \item If the RAG has a cycle \textbf{and every resource type on that
        cycle has a single instance}, the processes on the cycle
        \textbf{are} deadlocked. With one instance, a cycle is also
        sufficient.
  \item If the RAG has a cycle \textbf{and some resource type on the
        cycle has multiple instances}, deadlock \textbf{may or may not}
        exist. A free instance, or a holder who is not itself waiting,
        can still satisfy a request and break the wait. (This is why
        Mandar 2010 End Q4 asks you to draw both a deadlocking cycle and
        a harmless cycle when each resource has two instances.)
\end{enumerate}
This is a statement about Coffman conditions and instance counts. It is
not Banker's algorithm (no Max / Need / safe-sequence calculation) and
not a periodic detection/recovery procedure.
\end{Sol}

\clearpage

\section{Lectures 21--24 (incl.\ 22A): Persistence --- I/O Devices, HDDs, RAID \& File Systems}

\subsection{Lectures 21, 22 \& 22A: I/O Devices, Hard Disk Drives \& RAID}

\subsubsection*{Additional Challenging \& Analytical Problems}

\begin{Qn}{Hard}{Lect 22}{CS-537 Sp11 Final, Q3}{Black-box disk scheduler identification}
The disk has one surface, 100 tracks of 100 sectors: track $0$ holds
sectors $0\ldots 99$, track $1$ holds $100\ldots 199$, and so on. For
each model, the requests are issued together; you observe the completion
order. Identify the scheduler and a plausible initial head position.
\begin{enumerate}[(a)]
\item Requests \texttt{0, 1, 2, 3, 4, 5}. Completed
      \texttt{0, 1, 2, 3, 4, 5}.
\item Requests \texttt{0, 500, 200, 400, 300, 100}. Completed
      \texttt{0, 100, 200, 300, 400, 500}.
\item Requests \texttt{0, 500, 200, 400, 300, 100}. Completed
      \texttt{200, 300, 400, 500, 100, 0}.
\item Requests \texttt{0, 50, 110, 600}. Completed
      \texttt{0, 110, 50, 600}.
\item Could disks A--D actually be using the \emph{same} scheduler?
\end{enumerate}
\end{Qn}
\begin{Sol}
(a) Sequential in issue order on track 0. Consistent with
\textbf{FIFO}, and also with SSTF/SATF if the head started at sector
$0$ (or just before it). Initial state: track 0, near sector 0.

(b) Strictly increasing LBA / track. \textbf{SCAN/C-SCAN} (elevator)
sweeping outward from track 0, or SSTF from the outer rim. Initial:
track 0.

(c) Starts at track 2 (sector 200), sweeps inward/up through 300, 400,
500, then the remaining 100 and 0. \textbf{SCAN} (or C-SCAN, then a
return) with initial head at track 2. SSTF from 200 would also pick
300 next (nearest seek).

(d) After 0, the scheduler skips 50 (same track) and takes 110 (track 1)
before returning to 50. That is \textbf{SATF/SPTF}: sector 50 has just
passed, so a one-track seek to 110 is cheaper than waiting a full
rotation. Initial: track 0, just after sector 0.

(e) \textbf{Yes}: a single SATF/SPTF scheduler produces all four traces
for suitable initial (track, sector) and rotation phase. FIFO cannot
explain (b)--(d); pure SSTF (ignore rotation) cannot explain skipping
50 in (d).
\end{Sol}

\begin{Qn}{Hard}{Lect 22}{CS-537 Sp16 Final, Q7(a)--(e)}{SSTF on a rotating ``matrix disk''}
A matrix disk has 3 rows $\times$ 8 columns. Columns shift left every
$R$ time units (wrap-around). The head sits over column 0 of the
current row and reads/writes that block in $R$ time units. Seeking to
another row costs $S$ per row of distance. Starting position (head over
block 0):
\begin{lstlisting}[numbers=none]
[0]  1   2   3   4   5   6   7
 8   9  10  11  12  13  14  15
16  17  18  19  20  21  22  23
\end{lstlisting}
After 2 shifts with $R=1$ the head would be over block 2 of row 0.
\begin{enumerate}[(a)]
\item $R=1$, $S=1$. Minimum time to read block 21?
\item $R=2$, $S=1$. Time to read blocks 6 and 18, if 6 is read before 18?
\item $R=2$, $S=1$. Time if 18 is read before 6?
\item $R=1$. For what values of $S$ should the scheduler read 18 before 6?
\item Fill in SSTF's seek-time function on this 8-column disk:
\begin{lstlisting}[numbers=none]
int how_long(block, curr_row, S) {
    dest_row = ________;
    seek_dist = ________;
    seek_time = ________;
    return seek_time;
}
\end{lstlisting}
\end{enumerate}
\end{Qn}
\begin{Sol}
Official key (CS-537 Sp16 Final answers, p.~12).

(a) Block 21 is row $2$, column $5$. Seek two rows ($2S=2$); during
those 2 time units the columns shift twice, so 21 is then 3 columns
from the head. Wait 3, read 1:
$2~\mathrm{seek}+3~\mathrm{wait}+1~\mathrm{read}=\mathbf{6}$.

(b) 6 then 18:
$12~\mathrm{wait}+2~\mathrm{read}+2~\mathrm{seek}+4~\mathrm{wait}+2~\mathrm{read}
=\mathbf{22}$.

(c) 18 then 6:
$2~\mathrm{seek}+2~\mathrm{wait}+2~\mathrm{read}+2~\mathrm{seek}+4~\mathrm{wait}+2~\mathrm{read}
=\mathbf{14}$.

(d) Prefer 18-before-6 when that path is no slower:
$\boldsymbol{S\le 1}$.

(e) Destination row is $\lfloor\mathrm{block}/8\rfloor$; seek distance
is the absolute row difference; time is distance times $S$
(work-conserving SSTF, rotation ignored --- that is why SATF/SPTF can
beat it):
\begin{lstlisting}[numbers=none]
dest_row  = block / 8;
seek_dist = abs(dest_row - curr_row);
seek_time = seek_dist * S;
\end{lstlisting}
\end{Sol}

\begin{Qn}{Medium-Hard}{Lect 22}{CS-537 Sp12 Final, Q2(a)--(d)}{32 random versus sequential 4\,KB I/Os on a 15k-RPM disk}
The disk has average seek $4$\,ms, $15{,}000$\,RPM, and transfer rate
$100$\,MB/s.
\begin{enumerate}[(a)]
\item How long does it take to perform 32 reads, of size $4$\,KB, to
      random locations?
\item How long will it take to perform 32 writes, of size $4$\,KB, to
      random locations?
\item How long does it take to read 32 consecutive $4$\,KB chunks?
\item How long does it take to write 32 consecutive $4$\,KB chunks?
\end{enumerate}
\end{Qn}
\begin{Sol}
Full rotation $= 4$\,ms, average rotation $= 2$\,ms.
Transfer of $4$\,KB: $4/10^{5} = 0.04$\,ms.
A single random $4$\,KB I/O: $4+2+0.04 = 6.04$\,ms.

(a) $32\times 6.04 = \mathbf{193.28}$\,ms
($\approx 32\times 6 = 192$\,ms if transfer is neglected).
Random throughput $\approx 32\times 4\,\mathrm{KB}/193.28\,\mathrm{ms}
\approx 0.66$\,MB/s, versus sequential $100$\,MB/s.

(b) With no drive write cache (write-through), a random write pays the
same positioning as a random read: $\mathbf{193.28}$\,ms.
(Write-back / immediate reporting would hide this, at the cost of
durability.)

(c) One positioning, then a $128$\,KB sequential transfer:
$4+2+128/100 = \mathbf{7.28}$\,ms.

(d) Same as (c) for write-through sequential writes:
$\mathbf{7.28}$\,ms.
\end{Sol}

\begin{Qn}{Medium-Hard}{Lect 22}{CS-537 Fa04 Mid3, Q1}{SSTF, two-queue SSTF, track buffer, and mirrored SSTF}
\begin{enumerate}[(a)]
\item Describe shortest-seek-time-first (SSTF) and its primary problem.
\item Two-queue SSTF: the first $N$ requests go into queue 1 and are
      served SSTF; arrivals during that period go into queue 2 and wait
      until queue 1 is empty; then the roles swap. Compare to SSTF.
\item Disks have a small internal \textbf{track buffer}. Assume you
      know which blocks will hit that cache. How would you modify SSTF
      to use this knowledge?
\item In a mirrored RAID, a read may be served from either copy. How
      would you change SSTF on top of two disks with identical contents
      (for reads)?
\end{enumerate}
\end{Qn}
\begin{Sol}
(a) SSTF always picks the pending request whose cylinder is closest to
the current arm position (the OS approximation of SPTF, ignoring
rotation). Primary problem: \textbf{starvation} --- a steady stream of
requests near the arm can delay a far-away request indefinitely.

(b) Two-queue SSTF is a \textbf{batch / anti-starvation} variant: a
request that arrives while the other queue is being served cannot cut
in. Far-away requests in the active batch still complete before new
nearby arrivals. It is fairer than SSTF, at some cost in average seek
(it is closer in spirit to SCAN, and is work-conserving inside a batch
but not across batches --- related to anticipatory/idle-wait ideas).

(c) Treat a track-buffer hit as \textbf{zero positioning cost} (the
request is already ``at'' the head). Prefer cache hits before any real
seek; among misses, keep SSTF. This is SPTF with a $0$-cost option.

(d) For each logical read, compute the seek to \emph{both} replicas and
pick the replica with the \textbf{shorter seek} (then apply SSTF across
the resulting physical requests). Two heads give twice the chance of a
short seek; this is why RAID-1 random-read bandwidth is $NR$, not
$N/2\cdot R$.
\end{Sol}

\begin{Qn}{Medium-Hard}{Lect 22A}{CS-537 Fa04 Mid3, Q2}{RAID-1 and RAID-4 reconstruction onto a hot spare}
Each disk has $D$ blocks; reading or writing one block takes $T$\,ms.
CPU/XOR cost is negligible. A hot spare replaces a failed disk.
\begin{enumerate}[(a)]
\item RAID-1. One disk fails. How many blocks must be read and written
      to fully reconstruct?
\item How long does RAID-1 reconstruction take (best implementation)?
\item RAID-4. One disk fails (not the parity disk). How many blocks
      must be read and written?
\item How long does RAID-4 reconstruction take (best implementation)?
\end{enumerate}
\end{Qn}
\begin{Sol}
(a) Copy the surviving mirror: \textbf{$D$ reads and $D$ writes}.

(b) Stream sequentially from the good disk onto the spare. The read and
the write of successive blocks overlap, so time $\boldsymbol{\approx DT}$
(one sequential pass), not $2DT$.

(c) Let the array have $N$ disks ($N-1$ data $+$ 1 parity). Every
surviving disk must be read in full and XOR-ed; the spare is written in
full: \textbf{$(N-1)D$ reads and $D$ writes}.

(d) Read all $N-1$ survivors \textbf{in parallel} (sequential scan of
each), XOR on the fly, write the spare in parallel with those reads.
Limited by one sequential pass: $\boldsymbol{\approx DT}$.
(Rebuild does \emph{not} need a small-write RMW; it is a full-stripe
sequential reconstruction.)
\end{Sol}

\begin{Qn}{Medium-Hard}{Lect 22A}{CS-537 Sp16 Final, Q6(a)--(d)}{Even parity, RAID-4 map, small write, RAID-5 rotation}
\begin{enumerate}[(a)]
\item RAID-4/5 use even parity. Compute the parity bit:
      (i) $0\,0\,1\,0$\quad
      (ii) $1\,0\,1\,0$\quad
      (iii) $0\,1\,0\,1$\quad
      (iv) $1\,1\,1\,1$.
\item Draw a RAID-4 array, chunk size $4$\,KB, 4 data disks $+$ 1 parity
      disk.
\item In a 5-disk RAID-4, a single-block (small) write: which blocks
      are read and written, and in what order? Use the I/O-minimal
      (subtractive) path.
\item RAID-5 rotates parity. What is the primary reason? Does it solve
      the small-write problem?
\end{enumerate}
\end{Qn}
\begin{Sol}
Official key (CS-537 Sp16 Final answers, p.~10). Even parity means the
XOR of data bits and $P$ is $0$, so $P$ is the XOR of the data bits.

(a) (i) $0\oplus 0\oplus 1\oplus 0 = \mathbf{1}$;
(ii) $1\oplus 0\oplus 1\oplus 0 = \mathbf{0}$;
(iii) $0\oplus 1\oplus 0\oplus 1 = \mathbf{0}$;
(iv) $1\oplus 1\oplus 1\oplus 1 = \mathbf{0}$.

(b)
\begin{lstlisting}[numbers=none]
  D0     D1     D2     D3     Dparity     (4KB chunks)
   0      1      2      3      P_0123
   4      5      6      7      P_4567
   8      9     10     11      P_8,9,10,11
\end{lstlisting}

(c) 2 reads, then 2 writes: read \textbf{target data} and the
\textbf{corresponding parity} (in parallel); XOR old data, new data,
and old parity to obtain new parity; write \textbf{target} and
\textbf{new parity} (in parallel).

(d) Primary reason: avoid the \textbf{parity-disk bottleneck}. It does
\textbf{not} solve the small-write problem --- RAID-5 still does the
same 4-I/O RMW, just on rotating disks.
\end{Sol}

\begin{Qn}{Medium-Hard}{Lect 22}{CS-537 Sp18 Final, Q17--19}{FIFO / SSTF / SATF last-block on the 3-track simulator}
The drive has three tracks and rotates counter-clockwise. The arm starts
over block 6 on the outer track. Sector layout (outer $\to$ inner):
\begin{lstlisting}[numbers=none]
outer :  6  7  8  9 10 11  0  1  2  3  4  5     (then wrap)
middle: 18 19 20 21 22 23 12 13 14 15 16 17
inner : 30 31 32 33 34 35 24 25 26 27 28 29
\end{lstlisting}
Rotation carries increasing outer-track numbers under the head
($6\to 7\to 8\to\cdots\to 5$). Request stream:
\texttt{21, 17, 11, 7, 13}.
\begin{enumerate}
\item[17.] FIFO. Last block read?
      (a) 21 (b) 17 (c) 11 (d) 7 (e) 13
\item[18.] SSTF, \textbf{very fast} (essentially infinite) seek. Last
      block read?
\item[19.] SATF, very fast seek. Last block read?
\end{enumerate}
\end{Qn}
\begin{Sol}
Official key: 17\,e, 18\,a, 19\,b. (Shingled-disk Q11--16 and Q20 of
the same exam are out of syllabus and omitted.)

17. \textbf{(e) 13.} FIFO preserves issue order; 13 is last.

18. \textbf{(a) 21.} Seek is free, so SSTF is by track distance. Start
on the outer track: service outer requests 11 and 7 first, then the
middle-track trio $\{13,17,21\}$. On a track, service in rotational
order ($13$, then $17$, then $21$). Last $= 21$.

19. \textbf{(b) 17.} SATF/SPTF with free seek is nearest-sector-first
in rotation. From 6 the rotational order of the five requests is
$7$, $11$, $13$, $21$, then $17$ (17 sits with 5, just before a full
revolution back to 6). Last $= 17$.
\end{Sol}

\begin{Qn}{Medium-Hard}{Lect 22}{CS-537 Fa21 Final, Q5--8}{FIFO / SSTF / SATF on the disk simulator (requests 13, 4, 35, 34, 5)}
Same 3-track geometry and rotation as the previous question (head starts
over block 6; rotate $6\to 7\to\cdots$). Requests: \texttt{13, 4, 35, 34, 5}.
\begin{enumerate}
\item[5.] FIFO. Which request is serviced last?
      (a) 13 (b) 4 (c) 35 (d) 34 (e) 5
\item[6.] SSTF, infinitely fast seek (on a track, read in the fastest
      available order). Last request?
\item[7.] SATF, infinitely fast seek. Last request?
\item[8.] With infinitely fast seek, order the schedulers from fastest
      to slowest on this five-request set:
      (a) SSTF, SATF, FIFO\quad
      (b) SSTF, FIFO, SATF\quad
      (c) SATF, SSTF, FIFO\quad
      (d) SATF, FIFO, SSTF\quad
      (e) FIFO, SATF, SSTF
\end{enumerate}
\end{Qn}
\begin{Sol}
Official arrows (CS-537 Fa21 Final solutions.pdf): 5\,e, 6\,c, 7\,e, 8\,c.

Tracks: outer $\{4,5\}$, middle $\{13\}$, inner $\{34,35\}$.

5. \textbf{(e) 5.} FIFO last in the issue list.

6. \textbf{(c) 35.} SSTF stays on the outer track first ($4,5$), then
middle ($13$), then inner ($34$ then $35$ in rotational order). Last
$= 35$.

7. \textbf{(e) 5.} SATF follows rotation from 6: $34/35$ (near the top)
come first, then $13$, then $4$, and $5$ is just before 6 --- a nearly
full rotation. Last $= 5$.

8. \textbf{(c) SATF, SSTF, FIFO.} Seek is free, so SATF (rotation-aware)
beats SSTF (track-then-rotation); both beat FIFO, which ignores geometry
entirely. SATF is the disk's SPTF.
\end{Sol}

\subsubsection*{Conceptual Drills, MCQs \& Further Practice}

\begin{Qn}{Easy-Medium}{Lect 22}{CS-537 Fa21 Final, Q1--4, Q9--10}{Hard-drive vocabulary (Word Thief fill-in)}
Each question has exactly one answer.
\begin{enumerate}
\item[1.] ``The drive consists of a large number of sectors (512-byte
      blocks), each of which can be read or \underline{\hspace{1.2cm}}.''
      (a) deleted\quad (b) written\quad (c) read from\quad
      (d) erased\quad (e) truncated
\item[2.] ``We start with a \underline{\hspace{1.2cm}}, a circular hard
      surface on which data is stored persistently by inducing magnetic
      changes to it.''
      (a) platter\quad (b) surface\quad (c) track\quad
      (d) sector\quad (e) block
\item[3.] ``Data is encoded on each surface in concentric circles of
      sectors; we call one such concentric circle a
      \underline{\hspace{1.2cm}}.''
      (a) platter\quad (b) surface\quad (c) track\quad
      (d) sector\quad (e) block
\item[4.] ``Outer tracks tend to have more sectors than inner tracks
      \ldots These disks are often referred to as
      \underline{\hspace{1.2cm}} disk drives.''
      (a) wicked smart\quad (b) extended\quad (c) regular\quad
      (d) multi-zoned\quad (e) shingled
\item[9.] ``The three basic components of disk I/O time are
      \underline{\hspace{1.2cm}}, \underline{\hspace{1.2cm}}, and
      \underline{\hspace{1.2cm}}.''
      (a) transition, position, constitution\quad
      (b) move head, wait for head to stop moving, transfer\quad
      (c) position, transfer, react\quad
      (d) shake, rattle, roll\quad
      (e) seek, rotate, transfer
\item[10.] ``A disk scheduling algorithm that avoids starvation is
      called \underline{\hspace{1.2cm}}.''
      (a) PASS\quad (b) CRANK\quad (c) SCAN\quad
      (d) CHECK\quad (e) FAIRPLAY
\end{enumerate}
\end{Qn}
\begin{Sol}
Official arrows (CS-537 Fa21 Final solutions.pdf):
1\,b, 2\,a, 3\,c, 4\,d, 9\,e, 10\,c.

1. \textbf{(b) written.} Sectors are the addressable read/write unit
(512\,B or 4\,KB). They are not ``erased'' like flash.

2. \textbf{(a) platter.} A spindle of platters; each platter has two
surfaces.

3. \textbf{(c) track.} A cylinder is the same-numbered track on every
surface (the arm need not seek to switch heads).

4. \textbf{(d) multi-zoned.} Zone-bit recording (ZBR): more sectors on
outer tracks because of the larger circumference. (e) is SMR, which is
out of this syllabus.)

9. \textbf{(e) seek, rotate, transfer.} Average seek $\approx
\tfrac13$ max seek; RPM $\to$ rotation time $=60/\mathrm{RPM}$.

10. \textbf{(c) SCAN} (the elevator). SSTF/SPTF are greedy and can
starve; SCAN/C-SCAN bound the wait by sweeping. (Work-conserving SCAN
vs.\ anticipatory/idle-wait is a further refinement, not asked here.)
\end{Sol}

\begin{Qn}{Easy-Medium}{Lect 22A}{CS-537 Fa21 Final, Q11--20}{RAID fill-in: capacity, mirroring, small write, throughput}
Each question has exactly one answer. $N$ disks; a single disk delivers
sequential bandwidth $S$\,MB/s and random bandwidth $R$\,MB/s, and
holds $B$ bytes.
\begin{enumerate}
\item[11.] RAIDs improve performance and capacity. Finally, RAIDs can
      improve \underline{\hspace{1.2cm}}.
      (a) size\quad (b) the odds\quad (c) reliability\quad
      (d) latency\quad (e) distribution
\item[12.] A 2-way mirrored system can tolerate 1 disk failure for
      certain, and up to \underline{\hspace{1.2cm}} failures depending
      on which disks fail.
      (a) 2\quad (b) $N/4$\quad (c) $N/2$\quad (d) $N/3$\quad (e) $N$
\item[13.] Maximum bandwidth during sequential writing to a 2-way
      mirrored array is \underline{\hspace{1.2cm}}.
      (a) $S$\quad (b) $2S$\quad (c) $NS$\quad
      (d) $NS/2$\quad (e) $NS^{2}$
\item[14.] RAID-4, 2-bit blocks:
      \texttt{Block0=00}, \texttt{Block1=10}, \texttt{Block2=11},
      \texttt{Block3=10}, \texttt{Parity=??}
      (a) 00\quad (b) 01\quad (c) 10\quad (d) 11\quad (e) none
\item[15.] RAID-4 useful capacity is \underline{\hspace{1.2cm}}.
      (a) $N B$\quad (b) $N$\quad (c) $B$\quad
      (d) $N(B-1)$\quad (e) $(N-1)B$
\item[16.] Random-read performance of RAID-5 as compared to RAID-4 is
      \underline{\hspace{1.2cm}}.
      (a) a little better\quad (b) a little worse\quad
      (c) a lot better\quad (d) a lot worse\quad
      (e) about the same
\item[17.] RAID level(s) \underline{\hspace{1.2cm}} encounter(s) the
      ``small write'' problem.
      (a) 0\quad (b) 1\quad (c) 4\quad (d) 5\quad (e) 4 and 5
\item[18.] A single write to RAID-4 requires \underline{\hspace{0.8cm}}
      read(s) and then \underline{\hspace{0.8cm}} write(s) (subtractive
      parity).
      (a) 2, 1\quad (b) 1, 2\quad (c) 2, 2\quad (d) 1, 1\quad
      (e) none
\item[19.] RAID-5 with $N$ disks under a random-write workload delivers
      \underline{\hspace{1.2cm}} MB/s.
      (a) $NR$\quad (b) $NR/2$\quad (c) $NR/4$\quad
      (d) $R/2$\quad (e) $R$
\item[20.] If you strictly want performance and do not care about
      reliability, \underline{\hspace{1.2cm}} is obviously best.
      (a) rebooting\quad (b) a parity-based approach\quad
      (c) mirroring\quad (d) thinking\quad (e) striping
\end{enumerate}
\end{Qn}
\begin{Sol}
Official arrows (CS-537 Fa21 Final solutions.pdf):
11\,c, 12\,c, 13\,d, 14\,d, 15\,e, 16\,a, 17\,e, 18\,c, 19\,c, 20\,e.

11. \textbf{(c) reliability.} The third RAID goal (with capacity and
bandwidth). Latency of a single request is still one-disk latency $D$.

12. \textbf{(c) $N/2$.} 2-way mirroring: one failure is always
tolerated; if no two failures hit the same pair, up to $N/2$ disks may
fail.

13. \textbf{(d) $NS/2$.} Sequential write must update both copies, so
only $N/2$ independent sequential streams.

14. \textbf{(d) 11.} Bitwise even parity:
bit 0: $0\oplus 1\oplus 1\oplus 1 = 1$;
bit 1: $0\oplus 0\oplus 1\oplus 0 = 1$.

15. \textbf{(e) $(N-1)B$.} One disk of parity. Same useful capacity as
RAID-5; RAID-1 is $NB/2$; RAID-0 is $NB$.

16. \textbf{(a) a little better.} RAID-4 random read uses $N-1$ data
disks; RAID-5 uses all $N$ (parity is rotated, so every disk holds
data). One extra disk: a little, not a lot.

17. \textbf{(e) 4 and 5.} Both use parity and therefore RMW for a
partial-stripe write. RAID-0/1 do not.

18. \textbf{(c) 2, 2.} Subtractive: read old data $+$ old parity, then
write new data $+$ new parity.

19. \textbf{(c) $NR/4$.} Each random write generates 4 physical I/Os
spread across $N$ disks. RAID-4 would be $R/2$ (parity bottleneck).

20. \textbf{(e) striping} (RAID-0): capacity $NC$, throughput $NS$ /
$NR$, but 0 failures tolerated. JBOD is similar but not a single
logical disk.
\end{Sol}

\begin{Qn}{Medium}{Lect 22A}{CS-537 Fa04 Final, Q14}{Broken RAID (BRAID): find the bad parity / mirror bits}
BRAID sometimes fills in redundant bits incorrectly. Circle the wrong
parity/redundancy bits, or state that there is no problem. Even parity.
Each row is one stripe. For RAID-1, disks $(0,1)$ are a pair and disks
$(2,3)$ are a pair.

(a) BRAID Level 4:
\begin{lstlisting}[numbers=none]
Disk0  Disk1  Disk2  Disk3(=P)
  0      0      0      0
  0      1      1      1
  0      1      0      1
  1      1      1      1
\end{lstlisting}

(b) BRAID Level 5 (parity rotates; even parity of each whole stripe):
\begin{lstlisting}[numbers=none]
Disk0  Disk1  Disk2  Disk3
  0      0      0      0
  0      1      0      1
  0      0      0      1
  0      1      1      1
\end{lstlisting}

(c) BRAID Level 1 (mirroring):
\begin{lstlisting}[numbers=none]
Disk0  Disk1  Disk2  Disk3
  0      0      0      0
  0      0      1      1
  1      1      0      1
  1      1      1      1
\end{lstlisting}

(d) BRAID Level 0:
\begin{lstlisting}[numbers=none]
Disk0  Disk1  Disk2  Disk3
  0      0      0      0
  0      1      1      1
  0      1      0      1
  1      1      1      1
\end{lstlisting}
\end{Qn}
\begin{Sol}
(a) $P = D_0 \oplus D_1 \oplus D_2$.
Row 1: $0\oplus 0\oplus 0 = 0$ (stored 0) OK.
Row 2: $0\oplus 1\oplus 1 = 0$ but stored $1$ --- \textbf{wrong}.
Row 3: $0\oplus 1\oplus 0 = 1$ (stored 1) OK.
Row 4: $1\oplus 1\oplus 1 = 1$ (stored 1) OK.
Circle Disk 3, row 2.

(b) Even parity over all four bits of a RAID-5 stripe (the parity disk
rotates, but the XOR of the row must still be 0).
Row 1: $0$ OK; row 2: $0\oplus 1\oplus 0\oplus 1 = 0$ OK;
row 3: $0\oplus 0\oplus 0\oplus 1 = 1$ --- \textbf{fails};
row 4: $0\oplus 1\oplus 1\oplus 1 = 1$ --- \textbf{fails}.
Circle rows 3 and 4 (any one bit in each row may be viewed as the bad
parity).

(c) Pair $(0,1)$ matches in every row. Pair $(2,3)$:
rows 1, 2, 4 match; row 3 has $0$ vs.\ $1$ --- \textbf{Disk 3, row 3}
is the broken mirror.

(d) \textbf{No problem.} RAID-0 has no redundant bits to get wrong
(and no reliability).
\end{Sol}

\begin{Qn}{Medium-Hard}{Lect 22A}{CS-537 Sp02 Final, Q1}{The small-write problem, RAID-5, and delayed parity}
\begin{enumerate}[(a)]
\item RAID-4 uses parity. What is the ``small write'' problem, and why
      is it bad?
\item RAID-5 is a slight variant of RAID-4. How does RAID-5 work?
\item In what way does RAID-5 address some of the problems of RAID-4?
\item One idea is to \textbf{delay} the update of the parity block
      (update it only after some time). How could this improve
      performance?
\item What problems are introduced by delaying updates to the parity
      block?
\end{enumerate}
\end{Qn}
\begin{Sol}
(a) A write of one data chunk is not a full stripe, so parity cannot be
recomputed from the new stripe alone. Subtractive RMW reads old data
and old parity, then writes both --- \textbf{4 I/Os for 1 logical
write}. Worse, every small write hits the \textbf{dedicated parity
disk}, so random-write bandwidth is $\sim R/2$ independent of $N$.

(b) RAID-5 is RAID-4 with \textbf{rotating (left-symmetric) parity}:
the parity chunk of stripe $i$ lives on disk $(N-1-i \bmod N)$ rather
than always on disk $N-1$. Capacity, sequential bandwidth, and
single-failure tolerance are unchanged: $(N-1)C$, $(N-1)S$, 1 disk.

(c) Random reads can use all $N$ disks ($NR$ vs.\ RAID-4's $(N-1)R$).
Random writes spread the RMW across disks, giving $N/4\cdot R$ rather
than $R/2$. RAID-5 is strictly better than RAID-4 on random I/O; it
does \textbf{not} eliminate the 4-I/O small write.

(d) Delayed / batched parity: several data writes in the same stripe
can share one parity update (or a full-stripe write can skip RMW
entirely). Fewer parity I/Os, better sequential-write behaviour.

(e) \textbf{Consistent-update}: a crash between the data write and the
delayed parity write leaves a stripe whose XOR no longer matches $P$.
Reconstruction from that parity would be wrong. Hardware RAID stores
the intent in \textbf{NVRAM}; without it, delayed parity is unsafe.
\end{Sol}

\subsection{Lectures 23--24: File System API \& Implementation}

\subsubsection*{Additional Challenging \& Analytical Problems}

\begin{Qn}{Hard}{Lect 24}{CS-537 Sp11 Final, Q5}{Black-box VSFS: recover the operation sequence}
INITIAL STATE, \texttt{state(i)}:
\begin{lstlisting}[numbers=none]
Inode Bitmap : 10000000
Inode Table  : [size=1,ptr=0,type=d] [] [] [] [] [] [] []
Data Bitmap  : 10000000
Data         : [("." 0), (".." 0)] [] [] [] [] [] [] []
\end{lstlisting}
Eight inodes and eight data blocks; files are at most one block; types
\texttt{d} (directory) and \texttt{r} (regular); root is inode 0. Each
part starts from the previous part's FINAL STATE. Name the operation(s).
\begin{enumerate}[(a)]
\item To
\begin{lstlisting}[numbers=none]
Inode Bitmap : 11000000
Inode Table  : [size=1,ptr=0,type=d] [size=0,ptr=-,type=r] [] ...
Data Bitmap  : 10000000
Data         : [("." 0),(".." 0),("f" 1)] [] ...
\end{lstlisting}
\item Then inode 1 becomes \texttt{[size=1,ptr=7,type=r]}, data bitmap
      \texttt{10000001}, data block 7 holds \texttt{SOMEDATA}.
\item Then the root directory also contains \texttt{("b" 1)}.
\item Then the root directory is
      \texttt{("." 0),(".." 0),("b" 1)} (the \texttt{f} dirent is gone).
\item Then data block 7 holds \texttt{OTHERDATA}.
\item Then the file system returns to \texttt{state(i)}.
\end{enumerate}
\end{Qn}
\begin{Sol}
\begin{enumerate}[(a)]
\item \texttt{creat("/f")} (or \texttt{open("/f", O\_CREAT)}): allocate
      inode 1 as an empty regular file, add dirent \texttt{(f,1)}. No
      data block ($ptr=-$). Root nlink stays 2.
\item \texttt{write} of one block to \texttt{/f}: allocate data block 7,
      set size $=1$ and $ptr=7$.
\item \texttt{link("/f","/b")}: a second name for inode 1. Inode bitmap
      and data bitmap unchanged; the inode's link count is not drawn in
      this notation, but both dirents share inode 1.
\item \texttt{unlink("/f")}: remove the \texttt{f} dirent. Inode 1 remains
      (still named \texttt{/b}) with its data.
\item \texttt{write} overwriting \texttt{/b}'s existing block (same
      pointer 7; bitmap unchanged).
\item \texttt{unlink("/b")}: last link; free inode 1 and data block 7;
      root directory back to \texttt{.} / \texttt{..}.
\end{enumerate}
\end{Sol}

\begin{Qn}{Medium-Hard}{Lect 23}{ISI Mandar End-sem 2015 Q3(c--d) [cf.\ 2011E Q4(b), 2017E Q4(b)]}{Directory block after mkdir/touch, and why nlink $=5$}
\begin{lstlisting}[numbers=none]
$ mkdir myNewDirectory
$ mkdir myNewDirectory/a myNewDirectory/b myNewDirectory/c
$ touch  myNewDirectory/d myNewDirectory/e
\end{lstlisting}
\begin{enumerate}[(c)]
\item Draw the data block of \texttt{myNewDirectory}. What changes if
      \texttt{myNewDirectory/d} is deleted?
\item Link count of an inode $n$ $=$ number of dirents (over the whole
      FS) whose inumber field is $n$. Why is the link count of
      \texttt{myNewDirectory} equal to 5 after the commands above?
\end{enumerate}
\end{Qn}
\begin{Sol}
\begin{enumerate}[(c)]
\item Directory $=$ list of \texttt{(name, inumber)}:
\begin{lstlisting}[numbers=none]
.     -> inumber of myNewDirectory
..    -> inumber of the parent
a     -> inumber of subdirectory a
b     -> inumber of subdirectory b
c     -> inumber of subdirectory c
d     -> inumber of file d
e     -> inumber of file e
\end{lstlisting}
      Deleting \texttt{d} does \emph{not} compact the block: the slot
      is marked free (inumber 0 / \texttt{IFREE}); the name may linger
      until the slot is reused. (2011 End Q4(b) is the same idea after
      compiling then deleting \texttt{a.out} in \texttt{os/}.)
\item Five dirents name that directory inode:
      (1)~the parent's entry \texttt{myNewDirectory};
      (2)~its own \texttt{.};
      (3--5)~\texttt{..} inside each of \texttt{a}, \texttt{b},
      \texttt{c}. Regular files \texttt{d} and \texttt{e} do \emph{not}
      bump a directory's nlink. A new subdirectory raises the parent's
      nlink by 1 because of \texttt{..}.
\end{enumerate}
\textbf{Cited clone:} 2017 End Q4(b), after
\texttt{ls > out1; df > out2; cp; rm out1; mv out2 df-output;}
\texttt{ln -s ls-output file-list; ln df-output used-and-free-space},
the mounted-root dirent set is
\texttt{.}, \texttt{..}, \texttt{ls-output} (regular),
\texttt{df-output} (regular), \texttt{file-list} (symlink),
\texttt{used-and-free-space} (hard link to \texttt{df-output}).
\texttt{out1} is gone; \texttt{out2} was renamed.
\end{Sol}

\begin{Qn}{Medium-Hard}{Lect 24}{CS-537 Sp12 Final, Q3(a)}{Which VSFS blocks are read for two root files?}
VSFS-like layout: one superblock, one data-bitmap block (DB), one inode-bitmap
block (IB), inode blocks (I), data blocks (D); all 4\,KB. Each inode is
(unrealistically) 4\,KB, i.e.\ one inode per block. Superblock is already in
memory; once a block is read it stays cached; no write traffic. The root
directory contains only the two files \texttt{foo} and \texttt{bar}, each
4\,KB. In reading both files in their entirety, which blocks does the file
system read from disk?
\end{Qn}
\begin{Sol}
Bitmaps are not needed on a pure read. Because each inode occupies a whole
block, root / \texttt{foo} / \texttt{bar} cannot share an inode block.
The walk is:
\begin{enumerate}
\item root inode $I_{\mathrm{root}}$;
\item root directory data $D_{\mathrm{root}}$ (lookup \texttt{foo} and
      \texttt{bar}; one block, then cached);
\item \texttt{foo}'s inode $I_{\mathrm{foo}}$ and data $D_{\mathrm{foo}}$;
\item \texttt{bar}'s inode $I_{\mathrm{bar}}$ and data $D_{\mathrm{bar}}$.
\end{enumerate}
\textbf{Six blocks:}
$I_{\mathrm{root}}$, $D_{\mathrm{root}}$, $I_{\mathrm{foo}}$,
$D_{\mathrm{foo}}$, $I_{\mathrm{bar}}$, $D_{\mathrm{bar}}$.
(Parts (b)--(e) of this exam ask HDD/flash/FFS timings and are out of
scope here.)
\end{Sol}

\begin{Qn}{Medium-Hard}{Lect 24}{CS-537 Fa21 Final, Q32--33, 36--37, 39--40}{VSFS size, open/append I/O, and simulator ops}
\begin{enumerate}[(32)]
\item An inode has 12 direct pointers and a single indirect pointer.
      Each pointer names a 4\,KB block; disk addresses are 4 bytes.
      The file can grow to
      \begin{enumerate}[(a)]
      \item 4096\,KB \quad (b) 4100\,KB \quad (c) 4104\,KB \quad
      (d) 4044\,KB \quad (e) 4144\,KB
      \end{enumerate}
\item Twelve directs plus a single \emph{and} a double indirect; 4\,KB
      blocks, 4-byte addresses. Approximate max file size:
      \begin{enumerate}[(a)]
      \item $\sim$4\,KB \quad (b) $\sim$1\,MB \quad (c) $\sim$4\,MB \quad
      (d) $\sim$1\,GB \quad (e) $\sim$4\,GB
      \end{enumerate}
\item Open \texttt{/foo/bar}, read it, and close it. The file system will
      read \underline{\hspace{1.2em}} inodes.
      \begin{enumerate}[(a)]
      \item 0 \quad (b) 1 \quad (c) 2 \quad (d) 3 \quad (e) 4
      \end{enumerate}
\item Once \texttt{/foo/bar} is opened, a process appends a data block
      three times. With no write buffering, each append writes
      \begin{enumerate}[(a)]
      \item the file's inode, the data bitmap, and the data block itself
      \item the directory data block, the data block, and the superblock
      \item the inode bitmap, the directory data block, and the inode
      \item the journal, the data block, and the directory
      \item the department, the chair, the entire university
      \end{enumerate}
\item Empty VSFS $\to$ final state with inode 1
      \texttt{[f a:-1 r:1]} and root dirent \texttt{(m,1)}. The operation?
      \begin{enumerate}[(a)]
      \item \texttt{creat("/m")}
      \item \texttt{mkdir("/m")}
      \item \texttt{unlink("/m")}
      \item append a block to the root directory
      \item append a block to the root inode
      \end{enumerate}
\item Then inode 1 has \texttt{r:2} and the root also contains
      \texttt{(o,1)}. The operation?
      \begin{enumerate}[(a)]
      \item \texttt{mkdir("/o")}
      \item \texttt{unlink("/m")}
      \item \texttt{read("/m")}
      \item \texttt{link("/m","/o")}
      \item \texttt{creat("/o")}
      \end{enumerate}
\end{enumerate}
\end{Qn}
\begin{Sol}
\begin{enumerate}[(32)]
\item Pointers per indirect block $=4096/4=1024$.
      \[
        (12+1024)\times 4\,\mathrm{KB} = 4144\,\mathrm{KB}.
      \]
      \textbf{(e)}
\item Double-indirect dominates:
      $1024^{2}\times 4\,\mathrm{KB} = 4\,\mathrm{GB}$, plus $4\,\mathrm{MB}$
      from the single and $48\,\mathrm{KB}$ from the directs
      $\approx 4\,\mathrm{GB}$. \textbf{(e)}
\item Path inodes: root, \texttt{foo}, \texttt{bar}. \texttt{read} uses
      the already-open inode; \texttt{close} is memory-only. \textbf{(d)~3}
\item Each unbuffered append of a small file: allocate (data bitmap),
      write data, update inode. Directory is unchanged. \textbf{(a)}
\item Empty regular file named \texttt{m}, no data block
      (\texttt{a:-1}), nlink 1. \textbf{(a)} \texttt{creat("/m")}
\item Same inode, nlink $1\to 2$, extra dirent. \textbf{(d)}
      \texttt{link("/m","/o")}
\end{enumerate}
\end{Sol}

\begin{Qn}{Hard}{Lect 24}{ISI Mandar End-sem 2017 Q4(c)}{Offset ranges of each TOC entry}
Newly created SVR2 filesystem, 4096-byte blocks, 4-byte block numbers,
standard 13-slot TOC (10 direct, 1 SI, 1 DI, 1 TI). Calculate the range
of file offsets covered by each TOC entry.
\end{Qn}
\begin{Sol}
$N=1024$. Each data block covers 4096 bytes.
\begin{center}
\begin{tabular}{ll}
\hline
TOC slot & offsets (inclusive) \\
\hline
direct $[0]$ & $0$ -- $4095$ \\
direct $[1]$ & $4096$ -- $8191$ \\
direct $[k]$ & $4096k$ -- $4096(k+1)-1$, \ $k=0..9$ \\
direct $[9]$ & $36\,864$ -- $40\,959$ \\
single indirect & $40\,960$ -- $4\,235\,263$ \\
double indirect & $4\,235\,264$ -- $4\,299\,202\,559$ \\
triple indirect & $4\,299\,202\,560$ -- $4\,402\,345\,713\,663$ \\
\hline
\end{tabular}
\end{center}
Single: $10N$ blocks starting at byte $10\times 4096$.
Double: $N^{2}$ further blocks. Triple: $N^{3}$ further blocks.
\end{Sol}

\begin{Qn}{Medium-Hard}{Lect 23}{ISI Mandar End-sem 2010 Q3(a) [cf.\ 2018E Q3(c--d)]}{When do inode-cache and GFT refcounts exceed 1?}
Describe the situations in which the reference count becomes greater
than one for entries in (i)~the inode cache and (ii)~the global file
table.
\end{Qn}
\begin{Sol}
\begin{enumerate}[(i)]
\item In-core inode \texttt{i\_count} $>1$: several independent
      reasons to pin the inode --- two processes have it open (even via
      different GFT entries / different \texttt{open}s), a process's
      cwd or root is that inode, it is a mount point, or a hard-link
      lookup is in flight. Multiple names (hard links) do \emph{not}
      by themselves raise the in-core count; they raise the on-disk
      \texttt{nlink}.
\item GFT / open-file refcount $>1$: \texttt{dup}/\texttt{dup2}, or
      \texttt{fork} (child inherits the parent's fds). Both share the
      \emph{offset}. (2018 End Q3(d) is the same question under the
      name ``GFT indegree $>1$''. 2018 End Q3(c): four GFT fields ---
      current offset, pointer to in-core inode, refcount, and open
      mode/flags.)
\end{enumerate}
\end{Sol}

\begin{Qn}{Medium}{Lect 23}{ISI Mandar End-sem 2011 Q4(a)}{Is the directory graph a tree?}
The directory hierarchy of an SVR2-like filesystem is a graph: nodes are
files/directories, and an edge points from a directory to each name it
contains. Is this graph always acyclic? Is it necessarily a tree?
\end{Qn}
\begin{Sol}
Not a tree, and not always acyclic.
\begin{itemize}
\item Hard links to \emph{regular files} give two parents, so the
      structure is a DAG, not a tree.
\item \texttt{.} and \texttt{..} are cycles (self-loop; back-edge to
      parent). Unix therefore usually \emph{forbids} hard links to
      directories, otherwise a true cycle of directories could be
      formed and \texttt{nlink} / deletion become ill-defined.
\item Symlinks can add further cycles in the \emph{pathname} graph
      (Fa13 Final Q3(h)) without changing the inode DAG.
\end{itemize}
\end{Sol}

\begin{Qn}{Medium-Hard}{Lect 23--24}{CS-537 Fa04 Mid3, Q3}{Inode-free FS vs Unix: path I/O and hard links}
A non-standard design, IFFS, stores most of a file's metadata in the
directory entry instead of in an inode.
\begin{enumerate}[(a)]
\item Which of the following are usually in a standard Unix inode?
      Circle all that apply:
      (i) direct pointers to data blocks;
      (ii) the name of the file;
      (iii) timestamps (accessed, updated, \ldots);
      (iv) the inode number of the parent directory;
      (v) the current file-pointer position.
\item In a standard Unix FS, how many disk reads to read a single block
      of \texttt{/this/path/is/toolong}? Nothing is cached.
\item Same question for IFFS.
\item IFFS does not support hard links. Why?
\end{enumerate}
\end{Qn}
\begin{Sol}
\begin{enumerate}[(a)]
\item \textbf{(i)} and \textbf{(iii)} only. The name lives in the
      directory; the parent is \texttt{..} in the directory (and is
      not stored in a regular-file inode); the file offset lives in
      the \emph{open-file} table, not the inode.
\item Path \texttt{/this/path/is/toolong}: 5 inodes (root $+$ 4
      components) $+$ 4 directory data blocks $+$ 1 file data block
      $= \mathbf{10}$ reads.
\item IFFS folds inode contents into the dirent, so the walk is 4
      directory-data reads $+$ 1 file-data read $= \mathbf{5}$
      (root's location is well-known; no separate inode-table reads).
\item A hard link is a second name for the \emph{same} metadata. With
      metadata copied into every dirent, two names would be two
      independent copies: a write through one name would not update
      size/pointers in the other. The inode exists precisely so that
      \texttt{nlink} names share one copy of the pointers.
\end{enumerate}
\end{Sol}

\begin{Qn}{Medium}{Lect 24}{CS-537 Fa04 Final, Q10}{Linked allocation: which claims are true?}
The Broken File System uses linked allocation: each block of the file
tells you where the next block is. Circle those of the following which
are true.
\begin{enumerate}[(a)]
\item Sequentially overwriting all of the blocks of a file takes about
      the same time as overwriting the last block.
\item Each read of a random block of a file requires only about one I/O.
\item Computing the size of a file takes about as much time as reading
      the entire file.
\item Reading an entire file sequentially takes about the same time as
      reading just the last block of the file.
\item The size of the inode is smaller than in a more typical Unix
      inode.
\end{enumerate}
\end{Qn}
\begin{Sol}
\begin{enumerate}[(a)]
\item \textbf{False.} Overwriting the last block still walks the chain
      (same reads) but does only one write; overwriting all does a
      write per block.
\item \textbf{False.} Random access is $O(\text{block index})$ I/Os ---
      the classic linked-allocation failure. (FAT moves the chain into
      a cacheable table so random access becomes a memory walk $+$ one
      data I/O.)
\item \textbf{True} if the inode does not store a size and you must
      walk to a nil terminator. (If size is cached in the inode, this
      is false; the question's spirit is the walk.)
\item \textbf{True} for the \emph{read} traffic: getting to the last
      block already reads every predecessor to follow \texttt{next}.
\item \textbf{True:} the inode need only store the first-block pointer
      (and maybe a size), not a multi-level TOC.
\end{enumerate}
So (c), (d), (e) are the true claims; (a), (b) are false.
\end{Sol}

\begin{Qn}{Medium-Hard}{Lect 24}{ISI Mandar End-sem 2012 Q5(a--b) [cf.\ 2018E Q3(a)]}{How many inode-list vs data blocks?}
Create an SVR2-like filesystem on a 512\,MB disk, 2\,KB block size.
Average file size is about 16\,KB; an inode occupies 128 bytes.
\begin{enumerate}[(a)]
\item How many blocks for the inode list, and how many remain as data
      blocks? (Reasonable approximations OK.)
\item What goes wrong if that average-size estimate is (i)~an
      overestimate (files shorter than 16\,KB) or (ii)~an underestimate?
      Which is more serious?
\end{enumerate}
\end{Qn}
\begin{Sol}
\begin{enumerate}[(a)]
\item Total blocks $= 512\,\mathrm{MB}/2\,\mathrm{KB} = 2^{18} = 262\,144$.
      Inodes per block $= 2048/128 = 16$. Want roughly one inode per
      file and $16\,\mathrm{KB}/2\,\mathrm{KB}=8$ data blocks per file,
      so if $I$ inode-list blocks and $D$ data blocks,
      \[
        16I \approx D/8 \quad\Rightarrow\quad D \approx 128I,
      \]
      \[
        I + 128I \approx 262\,144 \quad\Rightarrow\quad
        I \approx 2032,\quad D \approx 260\,112
      \]
      ($\approx 32\,512$ inodes). Super/boot/bitmap blocks are a
      rounding error at this scale.
\item (i) Overestimate of file size $\Rightarrow$ too few inodes: data
      blocks remain free but \texttt{creat} fails (\texttt{ENOSPC} on
      inodes). (ii) Underestimate $\Rightarrow$ too many inodes: the
      inode table wastes space that could have been data, but files
      can still be created until data runs out. \textbf{(i) is worse:}
      stranded free data is unusable without reformatting. (2018 End
      Q3(a) is the converse calculation:  $2^{10}$ inode-list blocks
      and $2^{18}$ data blocks of 4\,KB, 128\,B inodes $\Rightarrow$
      $32\,768$ inodes. At avg 16\,KB $=4$ blocks, $32\,768\times 4$
      data blocks used and $131\,072$ data blocks stranded; at avg
      16\,MB only 64 files fit and $32\,704$ inodes are unused. Zero
      waste when avg $= 8$ blocks $=32\,\mathrm{KB}$.)
\end{enumerate}
\end{Sol}

\begin{Qn}{Medium-Hard}{Lect 24}{ISI Mandar End-sem 2017 Q4(a)}{Free-block-list size on a huge volume}
Newly created SVR2 filesystem, 4096-byte blocks, $2^{32}$ data blocks,
4-byte block numbers. How many disk blocks are required initially to
store the list of free disk blocks? State assumptions.
\end{Qn}
\begin{Sol}
SVR2 stores free data-block numbers in a linked list of ordinary data
blocks (plus a small cache in the superblock). Each list block holds
$4096/4=1024$ 4-byte words: 1 word is a pointer to the next list
block and 1023 are free-block numbers. Let $L$ be the number of list
blocks. Those $L$ blocks are themselves taken from the $2^{32}$ data
blocks, and they name the remaining free blocks:
\[
  L + L\cdot 1023 = 2^{32} \quad\Rightarrow\quad L\cdot 1024 = 2^{32}
  \quad\Rightarrow\quad L = 2^{22} = 4\,194\,304.
\]
Assumption: the entire data region starts free, and the superblock
cache is negligible compared with $2^{32}$. (VSFS would instead keep
a data \emph{bitmap} of $2^{32}$ bits $= 2^{29}$ bytes
$=2^{17}$ 4\,KB blocks --- more compact, which is why the lecture
prefers bitmaps.)
\end{Sol}

\begin{Qn}{Medium}{Lect 23}{ISI Mandar End-sem 2015 Q3(b) [cf.\ 2016E Q3(c)]}{\texttt{/} and \texttt{/home} both show inode 2}
\texttt{stat} reports that both \texttt{/} and \texttt{/home} have
inode number 2. How is that possible?
\end{Qn}
\begin{Sol}
Inode numbers are unique only \emph{within one filesystem}.
\texttt{/home} is a \emph{mount point}: a second filesystem is mounted
there, and \texttt{stat} on \texttt{/home} returns the root inode of
the mounted volume. Unix roots are traditionally inode 2, so both
roots report 2. (The underlying directory that was the mount point
has a different inode, visible with \texttt{stat} on the unmounted
volume or via the mount table.) 2016 End Q3(c): \texttt{umount}
refuses if any in-core inode of that volume has a positive refcount
(open file, cwd, mount, \ldots) --- the kernel's ``filesystem busy''
check.
\end{Sol}

\subsubsection*{Conceptual Drills, MCQs \& Further Practice}

\begin{Qn}{Easy-Medium}{Lect 23}{CS-537 Sp18 Final, Q21--30}{File descriptors, \texttt{open}/\texttt{unlink}/\texttt{lseek}/\texttt{fsync}/\texttt{rename}}
\begin{enumerate}[(21)]
\item A file descriptor is:
\begin{enumerate}[(a)]
\item a system-wide object used to access files
\item a per-process integer used to access files
\item readily forged
\item returned to a process via the \texttt{close()} system call
\item hard to understand
\end{enumerate}
\item Adding the \texttt{O\_TRUNC} flag to the \texttt{open()} call will
\begin{enumerate}[(a)]
\item put the data in a locked ``trunk'' file
\item cause the open call to fail (usually)
\item imply you must add \texttt{O\_TRUNC} to \texttt{close()} as well
\item create the file (if it does not exist)
\item truncate the file to size $=0$
\end{enumerate}
\item The \texttt{unlink()} system call is to the program \texttt{rm} as
      THIS is to the \texttt{rmdir} program:
\begin{enumerate}[(a)]
\item \texttt{unlink()}
\item \texttt{delete()}
\item \texttt{rmdir()}
\item \texttt{link()}
\item \texttt{fork()}
\end{enumerate}
\item The \texttt{lseek()} call is used to
\begin{enumerate}[(a)]
\item reposition the disk head
\item do a long disk seek, immediately
\item change the current file offset
\item force changes to disk
\item close files after a layoff period
\end{enumerate}
\item The following information is NOT available within a typical inode:
\begin{enumerate}[(a)]
\item owner
\item size (bytes)
\item blocks allocated
\item file name
\item last access time
\end{enumerate}
\item The \texttt{read()} call is to the \texttt{cat} program as the
      BLANK call is to the \texttt{ls} program:
\begin{enumerate}[(a)]
\item \texttt{read()}
\item \texttt{readdir()}
\item \texttt{stat()}
\item \texttt{fstat()}
\item \texttt{umount()}
\end{enumerate}
\item Which type of links do most UNIX file systems support?
\begin{enumerate}[(a)]
\item hyperlinks
\item soft links
\item forked links
\item sausage links
\item unlinks
\end{enumerate}
\item A file is NOT
\begin{enumerate}[(a)]
\item a container for data
\item a byte array that can be read or written
\item something with a low-level name
\item easily deleted
\item something that can be referred to via a high-level name, thanks
      to directories
\end{enumerate}
\item To write data to a file and then force the contents to disk, we
      should call:
\begin{enumerate}[(a)]
\item \texttt{write()}
\item \texttt{write()} then \texttt{fsync()}
\item \texttt{write()} then \texttt{fopen()}
\item \texttt{write()} then \texttt{falloc()}
\item \texttt{write()} then \texttt{fit()}
\end{enumerate}
\item To atomically replace the contents of a file \texttt{foo}:
      \texttt{open()} a temporary file, \texttt{write()} the new data,
      force it to disk (as in Q29), and finally:
\begin{enumerate}[(a)]
\item \texttt{close()} the file
\item \texttt{rename()} the temporary file to the desired name
\item \texttt{link()} the file to another name
\item \texttt{unlink()} the original \texttt{foo}
\item none of the above
\end{enumerate}
\end{enumerate}
\end{Qn}
\begin{Sol}
Official key: 21\,b, 22\,e, 23\,c, 24\,c, 25\,d, 26\,b, 27\,b, 28\,d,
29\,b, 30\,b.
\begin{enumerate}[(21)]
\item An fd is a small \emph{per-process} integer indexing that
      process's fd table. The system-wide object is the open-file
      (GFT) entry. \textbf{(b)}
\item \texttt{O\_TRUNC} sets the size to 0 if the file exists.
      (\texttt{O\_CREAT} is what creates; they are often combined.)
      \textbf{(e)}
\item \texttt{rm} $\leftrightarrow$ \texttt{unlink}; \texttt{rmdir}
      $\leftrightarrow$ \texttt{rmdir} (empty directories only).
      \textbf{(c)}
\item \texttt{lseek} changes the GFT offset (\texttt{SEEK\_SET} /
      \texttt{CUR} / \texttt{END}); it does \emph{no} disk I/O.
      \textbf{(c)}
\item The file \emph{name} lives in a directory entry, not in the
      inode. That is why hard links are possible. \textbf{(d)}
\item \texttt{ls} enumerates directory entries via
      \texttt{readdir} (which uses \texttt{getdents} internally).
      \textbf{(b)}
\item Among the given options, soft (symbolic) links. Unix also
      supports hard links; that option is not listed. \textbf{(b)}
\item ``Easily deleted'' is the claim that is not always true
      (nlink, open fds, permissions, mount). \textbf{(d)}
\item Writes may sit in the buffer cache; \texttt{fsync} forces them
      (and, for a new name, one must sometimes \texttt{fsync} the
      directory too). \textbf{(b)}
\item \texttt{rename} is atomic: a crash leaves either the old
      contents or the new ones, never a mix. This is the standard
      atomic-update pattern (\texttt{mv} uses it). \textbf{(b)}
\end{enumerate}
\end{Sol}

\begin{Qn}{Easy-Medium}{Lect 24}{CS-537 Sp18 Final, Q39--40}{Why bitmaps; what VSFS does not provide}
\begin{enumerate}[(39)]
\item Bitmaps are useful as allocation structures because of all of
      the following EXCEPT:
\begin{enumerate}[(a)]
\item they are compact
\item they are human readable
\item they allow a quick lookup of free space
\item they readily find consecutive free blocks
\item updates to them do not add any disk traffic
\end{enumerate}
\item VSFS has all of the following features EXCEPT:
\begin{enumerate}[(a)]
\item regular files
\item directories
\item hard links
\item simple allocation structures
\item fast crash consistency
\end{enumerate}
\end{enumerate}
\end{Qn}
\begin{Sol}
Official key: 39\,b \emph{or} e; 40\,e.
\begin{enumerate}[(39)]
\item Bitmaps are compact ($1$ bit/block), support $O(1)$-word scans
      for a free bit, and consecutive $0$-bits are a free extent.
      They are not particularly human-readable \textbf{(b)}, and
      flipping a bit \emph{does} generate a disk write of that bitmap
      block \textbf{(e)}. Either exception is accepted.
\item VSFS is the ``very simple'' on-disk layout (superblock, inode
      bitmap, data bitmap, inode table, data region). It has files,
      directories, hard links, and bitmaps. Crash consistency is
      left as a future question (journaling / FSCK are out of this
      course). \textbf{(e)}
\end{enumerate}
\end{Sol}

\begin{Qn}{Easy-Medium}{Lect 24}{CS-537 Fa21 Final, Q31, 34--35, 38}{Inode, directory pairs, bitmaps, write buffering}
\begin{enumerate}[(31)]
\item ``The \underline{\hspace{1.5em}} is the generic name used in many
      file systems for the structure that holds a file's metadata
      (length, permissions, location of constituent blocks).''
\begin{enumerate}[(a)]
\item superblock
\item inode
\item data block
\item directory
\item journal
\end{enumerate}
\item ``In VSFS (as in many file systems), a directory is basically a
      list of (\underline{\hspace{1em}}, \underline{\hspace{1em}}) pairs.''
\begin{enumerate}[(a)]
\item directory name, file attribute
\item file name, inode number
\item file name, parent location
\item inode number, file type
\item inode type, file directory
\end{enumerate}
\item ``Free-space management is important for all file systems. In
      VSFS we have two simple \underline{\hspace{1.5em}} for this task.''
\begin{enumerate}[(a)]
\item free lists
\item management teams
\item hamburgers
\item bitmaps
\item inodes
\end{enumerate}
\item ``Write buffering has a number of performance benefits. They
      are \underline{\hspace{1em}}, \underline{\hspace{1em}}, and
      \underline{\hspace{1em}}.''
\begin{enumerate}[(a)]
\item skipping, tricking, and faking out
\item batching, scheduling, and avoiding writes altogether
\item batching, scheduling, and smoothing writes out
\item batching, avoiding writes, and spacing writes out over time
\item anticipating, logging, and batch amortizing
\end{enumerate}
\end{enumerate}
\end{Qn}
\begin{Sol}
\begin{enumerate}[(31)]
\item \textbf{(b) inode.} The superblock is FS-wide parameters; a
      directory is just a file of \texttt{(name, inumber)} pairs; a
      journal is crash-consistency machinery (out of syllabus).
\item \textbf{(b) file name, inode number.} That is the lecture
      definition of a directory, including \texttt{.} and \texttt{..}.
\item \textbf{(d) bitmaps} --- one for inodes, one for data blocks.
      (SVR2's free-block \emph{list} is the older alternative; VSFS
      uses bitmaps because a list ages poorly.)
\item \textbf{(b)}. OSTEP: delay $\Rightarrow$ (1)~batch several
      updates into fewer I/Os, (2)~schedule the resulting I/Os,
      (3)~avoid some writes entirely (create-then-delete of a temp
      file). Apps that need durability call \texttt{fsync}.
\end{enumerate}
\end{Sol}

\begin{Qn}{Medium}{Lect 23--24}{CS-537 Sp09 Final, Q7}{How many reads to open \texttt{/1/2/3/4.txt}, then a sibling?}
Memory caching plays a big role in any file system. Directories are one
block; the target files are one block.
\begin{enumerate}[1.]
\item Reading \texttt{/1/2/3/4.txt} (exactly one file-system block of
      data) on a typical FFS-like / VSFS file system: how many reads,
      including data and metadata, assuming nothing is cached?
\item No block is ever read from disk more than once. After reading
      \texttt{/1/2/3/4.txt} we read \texttt{/1/2/3/5.txt}. How many
      \emph{extra} reads take place?
\item In a system with a vast amount of memory, does any disk I/O
      ever take place? In what circumstances?
\end{enumerate}
\end{Qn}
\begin{Sol}
\begin{enumerate}[1.]
\item Five inodes (\texttt{/}, \texttt{1}, \texttt{2}, \texttt{3},
      \texttt{4.txt}) $+$ four directory data blocks $+$ one file
      data block $= \mathbf{10}$ reads.
\item Prefix \texttt{/1/2/3} is cached (inodes and directory blocks).
      Remaining: inode of \texttt{5.txt} $+$ its data block
      $= \mathbf{2}$ extra reads.
\item Yes: the \emph{first} access to a never-before-seen block;
      writes that miss the cache (or that \texttt{fsync} forces);
      cold start after reboot; and any workload larger than the
      cache (capacity misses). Reads of a working set that fits
      never touch the disk again.
\end{enumerate}
(The FFS layout-policy parts of this exam are out of syllabus; the
path-walk / cache counting is not.)
\end{Sol}

\begin{Qn}{Medium}{Lect 24}{CS-537 Sp02 Final, Q2(a--c)}{Buffer-cache hit, miss, and LRU thrash}
\texttt{search(file)} opens a file and reads it one byte at a time
counting \texttt{'a'}s. Disk: 10\,ms per 4\,KB block. Memory: 50\,$\mu$s
per 4\,KB. Ignore metadata I/O; file I/O dominates.
\begin{enumerate}[(a)]
\item 100 files, each 4\,KB, none cached. How long is the first run?
\item After that run everything is in the buffer cache. How long is a
      second run of the same 100 files?
\item Cache holds 999 4\,KB blocks (LRU). The program is run twice,
      each time over the \emph{same} 1000 4\,KB files. How long is
      the \emph{second} run?
\end{enumerate}
\end{Qn}
\begin{Sol}
\begin{enumerate}[(a)]
\item $100\times 10\,\mathrm{ms} = 1\,\mathrm{s}$.
\item $100\times 50\,\mu\mathrm{s} = 5\,\mathrm{ms}$
      ($200\times$ faster --- the point of the buffer cache).
\item Working set $=1000$ pages, cache $=999$, cyclic reuse $\Rightarrow$
      LRU misses on \emph{every} file (the lecture's thrash example).
      Second run: $1000\times 10\,\mathrm{ms} = 10\,\mathrm{s}$.
      (Parts (d)--(e) of this exam invent an informing API and are
      omitted.)
\end{enumerate}
\end{Sol}

\begin{Qn}{Medium}{Lect 23}{ISI Mandar End-sem 2010 Q3(b) [exact repeat: 2016E Q3(b)]}{\texttt{creat}: which permissions, existing vs new?}
When a file is ``created'' with \texttt{creat}, what permissions are
checked, and why, if (i)~the file already exists; (ii)~the file does
not exist?
\end{Qn}
\begin{Sol}
\texttt{creat(path, mode)} is \texttt{open(path, O\_WRONLY|O\_CREAT|O\_TRUNC, mode)}.
\begin{enumerate}[(i)]
\item File exists: search permission (\texttt{x}) on every directory
      in the path, plus \emph{write} permission on the file itself
      (it will be truncated and overwritten). The \texttt{mode}
      argument is ignored. The dirent is unchanged, so the
      directory's write bit is not required.
\item File does not exist: search on the path, plus \emph{write and
      execute} on the parent directory (a new dirent is added). The
      new inode's permission bits come from \texttt{mode \& \~{}umask}.
\end{enumerate}
\end{Sol}

\begin{Qn}{Easy-Medium}{Lect 23}{ISI Mandar Mid-sem 2014 Q1(c)ii}{Multiple in-core copies of an inode?}
T/F: Multiple copies of the inode corresponding to \texttt{a.out} may
exist simultaneously in the inode table.
\end{Qn}
\begin{Sol}
\textbf{False.} The in-core inode cache keeps \emph{one} copy per
\texttt{(device, inumber)}. Many processes may run the same
\texttt{a.out} (their fd tables / GFT entries all point at that one
in-core inode). A second copy would split \texttt{nlink}, size, and
block pointers. (Several GFT entries for the same inode \emph{are}
allowed --- that is a second \texttt{open}, each with its own offset.)
\end{Sol}

\begin{Qn}{Medium}{Lect 23}{ISI Mandar End-sem 2018 Q3(c--d) [cf.\ 2010E Q3(a)]}{Four GFT fields; GFT indegree $>1$}
\begin{enumerate}[(c)]
\item List the four fields of each global-file-table entry.
\item Describe the situations that make the indegree of a GFT entry
      greater than 1.
\end{enumerate}
\end{Qn}
\begin{Sol}
\begin{enumerate}[(c)]
\item (1)~current file offset; (2)~pointer to the in-core inode;
      (3)~reference count (how many fd-table slots point here);
      (4)~open mode / flags (\texttt{O\_RDONLY} / \texttt{O\_APPEND} /
      \ldots). This is xv6 \texttt{struct file} / the lecture's
      open-file object.
\item Indegree $>$1 means two fd-table slots share the entry:
      \texttt{dup}/\texttt{dup2}, and \texttt{fork} (the child
      inherits the parent's fds). Shared offset is the observable
      consequence (ProZilla, 2015 End Q4). A second \texttt{open}
      of the same path creates a \emph{new} GFT entry (indegree 1,
      independent offset).
\end{enumerate}
\end{Sol}

\begin{Qn}{Medium}{Lect 23}{ISI Mandar End-sem 2014 Q4(a,d)}{Why inodes/blocks start at 1; what \texttt{ls -l} shows}
\begin{enumerate}[(a)]
\item Why are (i)~inodes and (ii)~disk blocks numbered from 1, not 0?
\item (i)~What timestamps are stored in an SVR2 inode?
      (ii)~Design an experiment to decide which of them \texttt{ls -l}
      prints.
\end{enumerate}
\end{Qn}
\begin{Sol}
\begin{enumerate}[(a)]
\item (i)~Dirent inumber $0$ is the \emph{free/empty slot} sentinel
      (and a null pointer in TOC slots). Numbering live inodes from 1
      (Unix root is traditionally inode 2) keeps 0 available as
      ``none''.
      (ii)~Block number 0 is likewise a null data pointer in the
      inode TOC (``this slot is unused / a hole''). The boot block
      may physically sit at disk address 0, but the FS address space
      of \emph{data} blocks starts at 1.
\item (i)~Three timestamps: \texttt{atime} (last access),
      \texttt{mtime} (last content modification), \texttt{ctime}
      (last inode-status change: chmod, chown, link, \ldots).
      (ii)~Create a file. \texttt{ls -l} vs \texttt{ls -lu} vs
      \texttt{ls -lc}:
      \texttt{chmod} / \texttt{ln} bumps \texttt{ctime} only;
      \texttt{echo x >> f} bumps \texttt{mtime} (and usually
      \texttt{ctime}); \texttt{cat f} bumps \texttt{atime}.
      \texttt{ls -l} tracks \texttt{mtime}.
\end{enumerate}
\end{Sol}

\begin{Qn}{Medium}{Lect 23}{ISI Mandar End-sem 2016 Q3(c)}{How \texttt{umount} knows the FS is busy}
In the \texttt{umount} system call, how does the kernel determine
whether any file/directory on the mounted filesystem is currently
in use?
\end{Qn}
\begin{Sol}
The kernel walks the in-core inode cache (and the mount/cwd/root
bookkeeping) for inodes whose \texttt{i\_dev} is the victim
filesystem and whose in-core refcount is positive: open files
(GFT $\to$ inode), a process whose cwd or root is on that volume,
or a nested mount. If any such pin exists, \texttt{umount} fails
with ``busy'' / \texttt{EBUSY}. (A dirty buffer of that volume
must also be flushed; that is durability, not the busy check.)
\end{Sol}

\begin{Qn}{Medium}{Lect 24}{ISI Mandar End-sem 2022 Q6(c)}{100-byte inodes: pack, pad, or a third option?}
If the inode size were 100 bytes instead of 128, which would you
choose for storing the inode list, and why?
\begin{itemize}
\item store inodes contiguously as usual;
\item allocate 128 bytes of disk space to each inode, but only 100
      bytes hold useful information.
\end{itemize}
Suggest a third alternative that is better than both.
\end{Qn}
\begin{Sol}
With 1024-byte blocks (the paper's parameters), $1024/100=10$
remainder 24, while $1024/128=8$ exactly.
\begin{itemize}
\item Contiguous 100-byte packing: 10 inodes/block, only 24 bytes
      wasted, but inode $n$'s address is
      $\mathrm{start}+n\times 100$, which is \emph{not}
      block-aligned. An inode may span two disk blocks (2 I/Os
      to read one inode) whenever $n\times 100 \bmod 1024 > 924$.
\item Pad to 128: 8 inodes/block, every inode in one aligned
      block, but $28$ bytes wasted per inode ($22\%$ of the table).
\item \textbf{Better third option:} pack the largest $k$ such that
      $k\times 100 \le$ block size and \emph{no inode spans a
      block} --- here $k=10$, with the leftover 24 bytes as
      padding \emph{at the end of the block} (not per inode).
      Alternative: pick an inode size that \emph{divides} the
      block size (100 is a poor choice; 128 is what SVR2 used
      for this reason). The lecture's 256\,B inodes in 4\,KB
      blocks ($16$/block) follow the same rule.
\end{itemize}
\end{Sol}

\begin{Qn}{Easy-Medium}{Lect 23}{ISI Mandar End-sem 2010 Q3(c)}{How the shell implements \texttt{a.out < in > out}}
Clearly explain the steps by which a typical shell handles
\texttt{a.out < input.txt > output.txt}.
\end{Qn}
\begin{Sol}
This is path traversal $+$ \texttt{dup2} on the open-file table:
\begin{enumerate}
\item Shell \texttt{fork}s. The child inherits fd 0/1/2 and the
      parent's GFT entries for the terminal.
\item Child \texttt{open("input.txt", O\_RDONLY)} $\to$ some fd
      $k$, then \texttt{dup2(k, 0)} and \texttt{close(k)}: fd 0
      now names that file's open-file object (offset 0).
\item Child \texttt{open("output.txt", O\_WRONLY|O\_CREAT|O\_TRUNC, 0666)}
      $\to$ $\ell$, \texttt{dup2($\ell$, 1)}, \texttt{close($\ell$)}.
\item Child \texttt{exec}s \texttt{a.out}. The new image keeps the
      fd table: \texttt{read(0,...)} consumes \texttt{input.txt},
      \texttt{write(1,...)} appends to \texttt{output.txt} via the
      shared GFT offset.
\item Parent \texttt{wait}s. \texttt{dup2} is why redirection
      composes with \texttt{fork}: it retargets a well-known fd
      without the child having to know the pathname.
\end{enumerate}
\end{Sol}

\begin{Qn}{Medium}{Lect 24}{ISI Mandar End-sem 2014 Q4(c)}{Why a free-block list but no free-inode list?}
Briefly explain how the list of free data blocks is maintained in
an SVR2 filesystem. Why is an analogous structure not needed for
free inodes?
\end{Qn}
\begin{Sol}
Free \emph{data} blocks are scattered over the disk. SVR2 chains
them in a linked list of blocks of free-block numbers, with a
small cache of those numbers in the superblock (see 2017 End
Q4(a)). A free list ages (becomes random) over time; VSFS
replaces it with a data bitmap.

Free \emph{inodes} already live in a compact, dense table. The
superblock caches a handful of known-free inumbers; when the
cache empties the kernel linearly scans the inode list (or an
inode bitmap, in VSFS) looking for \texttt{nlink}$=0$. There is
no need to thread free inodes into a disk-resident list --- the
table itself is the search structure. (VSFS just uses an inode
bitmap, which is the lecture's solution to both problems.)
\end{Sol}

\clearpage

\section{Lectures 1--7 (Pre-Minor): CPU Virtualization (Processes, Traps \& Scheduling)}

\subsubsection*{Additional Challenging \& Analytical Problems}

\begin{Qn}{Hard}{Lect 4}{CS-537 Spring 2016 Mid Q4}{Limited Direct Execution: OS, hardware, or user?}
This question is about the Limited Direct Execution protocol. Some of these steps are performed by the OS; some are handled by hardware (HW); some are in the user program itself (USER). Mark which actor performs each step in this example of a process being created, running, issuing a system call, and exiting.
\begin{center}
\begin{tabular}{ll}
Create entry for process list & OS / HW / USER \\
Allocate memory for program & OS / HW / USER \\
Load program into memory & OS / HW / USER \\
Setup user stack with argv & OS / HW / USER \\
Fill kernel stack with reg/PC & OS / HW / USER \\
execute return-from-trap instruction & OS / HW / USER \\
restore regs from kernel stack & OS / HW / USER \\
switch to user mode & OS / HW / USER \\
set PC to \texttt{main()} & OS / HW / USER \\
Start running in \texttt{main()} & OS / HW / USER \\
Call a system call & OS / HW / USER \\
execute trap instruction & OS / HW / USER \\
save regs to kernel stack & OS / HW / USER \\
switch to kernel mode & OS / HW / USER \\
set PC to OS trap handler & OS / HW / USER \\
Handle trap & OS / HW / USER \\
Do work of syscall & OS / HW / USER \\
execute return-from-trap instruction & OS / HW / USER \\
restore regs from kernel stack & OS / HW / USER \\
switch to user mode & OS / HW / USER \\
set PC to instruction after earlier trap & OS / HW / USER \\
Call \texttt{exit()} system call & OS / HW / USER \\
\end{tabular}
\end{center}
\end{Qn}
\begin{Sol}
Official key (CS-537 Sp16 Mid Q4), matching the OSTEP LDE protocol:
\begin{center}
\begin{tabular}{ll}
Create entry for process list & \textbf{OS} \\
Allocate memory for program & \textbf{OS} \\
Load program into memory & \textbf{OS} \\
Setup user stack with argv & \textbf{OS} \\
Fill kernel stack with reg/PC & \textbf{OS} \\
execute return-from-trap instruction & \textbf{OS} \\
restore regs from kernel stack & \textbf{HW} \\
switch to user mode & \textbf{HW} \\
set PC to \texttt{main()} & \textbf{HW} \\
Start running in \texttt{main()} & \textbf{USER} \\
Call a system call & \textbf{USER} \\
execute trap instruction & \textbf{USER} \\
save regs to kernel stack & \textbf{HW} \\
switch to kernel mode & \textbf{HW} \\
set PC to OS trap handler & \textbf{HW} \\
Handle trap & \textbf{OS} \\
Do work of syscall & \textbf{OS} \\
execute return-from-trap instruction & \textbf{OS} \\
restore regs from kernel stack & \textbf{HW} \\
switch to user mode & \textbf{HW} \\
set PC to instruction after earlier trap & \textbf{HW} \\
Call \texttt{exit()} system call & \textbf{USER} \\
\end{tabular}
\end{center}
The OS \emph{initiates} return-from-trap (it is a privileged instruction the kernel executes); the hardware then restores registers, drops the mode bit, and sets the PC. The user program executes ordinary instructions and the trap instruction; it never itself switches the mode bit.
\end{Sol}

\begin{Qn}{Hard}{Lect 3}{CS-537 Spring 2016 Mid Q3}{\texttt{fork}/\texttt{exec}/\texttt{wait} call-and-return timeline}
Remember \texttt{fork()}, \texttt{exec()}, and \texttt{wait()}? Trace \emph{when each of these system calls is CALLED, and when they RETURN}. Write down something like ``\texttt{fork} called'' if \texttt{fork()} has been called, and ``\texttt{fork} returned'' if it returns. A particular answer may have more than one action.
\begin{center}
\begin{tabular}{p{0.55\textwidth}p{0.38\textwidth}}
\hline
Event & System call(s) called/returned \\
\hline
Process A is running; it starts to make a child process B
  (a nearly exact copy of itself) & \\
Process A then continues running just after the OS has been
  told to create B & \\
Process B runs for the first time & \\
Process B now overlays its address space with a new program
  and starts running in its \texttt{main()} & \\
Process B runs for a while, then A runs for a while, then B, etc. & \\
Process A now ensures that it will be notified when B exits,
  blocking until that is the case & \\
B now creates C & \\
C starts running & \\
C tries to overlay its address space but fails to do so
  and decides to exit & \\
Process B runs again, then creates D & \\
B and D run, each doing some disk I/O & \\
B finally exits and A runs & \\
\hline
\end{tabular}
\end{center}
\end{Qn}
\begin{Sol}
Official key (CS-537 Sp16 Mid Q3), reading each row as the syscall(s) that the narrative forces:
\begin{center}
\begin{tabular}{p{0.52\textwidth}p{0.42\textwidth}}
\hline
Event & Call / return \\
\hline
A starts to make child B
  & \texttt{fork} called \\
A continues just after the OS has been told to create B
  & \texttt{fork} returned (in A; child's pid) \\
B runs for the first time
  & \texttt{fork} returned (in B; returns 0) \\
B overlays its address space and starts in \texttt{main()}
  & \texttt{exec} called (and does \emph{not} return on success) \\
B runs, then A, then B, \ldots
  & none (pure CPU scheduling / context switch) \\
A will be notified when B exits, blocking until then
  & \texttt{wait} called (A blocks) \\
B now creates C
  & \texttt{fork} called (in B) \\
C starts running
  & \texttt{fork} returned (in C; returns 0) \\
C's overlay fails and C exits
  & \texttt{exec} called, \texttt{exec} returned (failure),
    then \texttt{exit} called \\
B runs again, then creates D
  & \texttt{fork} called/returned in B; \texttt{fork} returned in D \\
B and D each do disk I/O
  & both Blocked on I/O (no \texttt{fork}/\texttt{exec}/\texttt{wait}) \\
B finally exits and A runs
  & \texttt{wait} returned in A (reaps B) \\
\hline
\end{tabular}
\end{center}
Key facts used: successful \texttt{fork} returns \emph{twice} (parent and child); successful \texttt{exec} \emph{never} returns; failed \texttt{exec} returns $-1$ in the same process, which may then \texttt{exit}; \texttt{wait} in the parent returns only after the child has become a zombie.
\end{Sol}

\begin{Qn}{Hard}{Lect 3}{Mandar ISI 2017 Mid Q1(b)}{Fill in a mini-shell: \texttt{fork}/\texttt{exec}/\texttt{wait} with error checks}
The following (pseudo-)code fragment contains the outline of a very basic shell. Fill in the blank portions so that it works as intended. Assume that all arrays are large enough. For full credit, (i)~the syntax for all necessary system calls should be correct; and (ii)~your code should check for errors returned by system calls.
\begin{lstlisting}[numbers=none]
char linebuf[BUF_SIZE];
char **argv[MAX_ARGS];
int argc;
/* Put any declarations here */
while (1) {
      /* write code to print the shell prompt (e.g., "$ ") */
      readline(linebuf); // contains the command entered by the user
      argc = parseline(linebuf); // puts the command (program name) in argv[0],
                                   // and the remaining arguments (if any) in
                                   // argv[1], argv[2] ...
      /* write code to create child process */
      /* execute the command entered by the user in the child */
}
\end{lstlisting}
\end{Qn}
\begin{Sol}
Classic Lect~3 pattern: print prompt; \texttt{fork}; child \texttt{exec}; parent \texttt{wait}; check every return.
\begin{lstlisting}[numbers=none]
char linebuf[BUF_SIZE];
char **argv[MAX_ARGS];
int argc;
pid_t pid;
int status;

while (1) {
    if (write(1, "$ ", 2) < 0) { perror("write"); exit(1); }
    readline(linebuf);
    argc = parseline(linebuf);
    if (argc < 1) continue;

    pid = fork();
    if (pid < 0) { perror("fork"); continue; }

    if (pid == 0) {
        execvp(argv[0], argv);
        perror("exec");   /* reached only on failure */
        exit(1);
    } else {
        if (wait(&status) < 0) perror("wait");
    }
}
\end{lstlisting}
\texttt{execvp} is the usual argv-passing variant (the lecture's \texttt{exec()}). On success it does not return, so the child's \texttt{perror}/\texttt{exit} run only if the executable cannot be loaded. The parent must \texttt{wait} (otherwise the child becomes a zombie). Empty lines are skipped. A production shell would also handle \texttt{cd} in the parent and I/O redirection in the child before \texttt{exec}; those are not required by the printed blanks.
\end{Sol}

\begin{Qn}{Hard}{Lect 4}{Mandar ISI 2017 Mid Q1(a)}{User $\to$ kernel: syscall, exception, interrupt}
List the types of events that cause a process to switch from user mode to kernel mode. For each type of event,
\begin{itemize}
  \item explain in 1--2 lines when the event occurs;
  \item state whether the event is synchronous / asynchronous;
  \item state whether the event is handled by the kernel in user context or system context.
\end{itemize}
\textit{Note: the first part of this question (the three event types) appeared as MTL458 Minor Q2(a).}
\end{Qn}
\begin{Sol}
Three event types (the trap instruction of Lect~4 covers all three; the IDT entry is selected by a number in a register):
\begin{center}
\begin{tabular}{p{0.18\textwidth}p{0.28\textwidth}p{0.16\textwidth}p{0.26\textwidth}}
\hline
Type & When & Sync? & Context \\
\hline
\textbf{System call}
  & User program executes a trap/\texttt{int}/\texttt{syscall} instruction to request an OS service (\texttt{read}, \texttt{fork}, \ldots).
  & Synchronous (the process asked for it)
  & \emph{User context}: kernel runs on the process's kernel stack, on behalf of that process. \\
\textbf{Exception} (fault)
  & Illegal instruction, divide-by-zero, bad memory reference, \ldots, caused by the instruction the process is executing.
  & Synchronous
  & \emph{User context} of the faulting process (the kernel may kill it, or, e.g., handle a page fault and resume). \\
\textbf{Interrupt}
  & Device or timer raises an interrupt line; independent of the instruction the current process is on.
  & Asynchronous
  & \emph{System context} (interrupt handler; not ``the current process asked for this''). After the handler, the OS may return-from-trap into the same process or context-switch. \\
\hline
\end{tabular}
\end{center}
Hardware on every such event: raise privilege, switch to the kernel stack, save user PC/registers on that stack, look up the IDT, jump to the handler. Return-from-trap reverses this and drops to user mode.
\end{Sol}

\begin{Qn}{Hard}{Lect 4}{CS-537 Fall 2013 Mid Q4}{McFlub on traps and interrupts}
Professor McFlub says:
\begin{quote}
``Traps, like system calls, and interrupts, such as from a programmable timer or disk, are handled very similarly in OSes. For example, when each occurs, the hardware saves some state (such as the program counter) and jumps into the OS. At that point, the OS handles the trap or interrupt. When finished, the OS returns, through a normal return instruction, and then retries the trapping or interrupted instruction.''
\end{quote}
Correct McFlub. (If he is completely right, say so.)
\end{Qn}
\begin{Sol}
Official key (CS-537 Fa13 Mid answers, Q4):
\begin{itemize}
  \item Returns from the kernel require a \emph{different} instruction (return-from-trap / \texttt{iret}), which both undoes what the trap did (restores saved state from the kernel stack) \emph{and} changes privileged mode back to user mode. A normal \texttt{ret} does not drop the mode bit.
  \item Retry only happens in some cases (e.g.\ a fault that the OS repaired). For a syscall trap, the return must go to the instruction \emph{after} the trapping instruction, otherwise the syscall would re-execute forever.
\end{itemize}
The first half (hardware saves PC, looks up the IDT, jumps into the OS) is fine.
\end{Sol}

\begin{Qn}{Medium-Hard}{Lect 4}{Mandar ISI 2010 Mid Q1(c)}{Stopping a \texttt{printf} loop vs.\ a tight arithmetic loop}
Consider two processes $P_1$ and $P_2$ corresponding to the following programs:
\begin{center}
\begin{tabular}{cc}
$P_1$ & $P_2$ \\
\begin{lstlisting}[numbers=none]
void main() {
  while (1) {
    printf("Hello!\n");
  }
}
\end{lstlisting}
&
\begin{lstlisting}[numbers=none]
void main() {
  int i = 0;
  while (1) {
    i = i + 1; i = i - 1;
  }
}
\end{lstlisting}
\end{tabular}
\end{center}
For each process, clearly explain what mechanism(s) exist(s) by which the operating system can ensure that the process does not monopolise the CPU.
\end{Qn}
\begin{Sol}
$P_1$: \texttt{printf} eventually issues a \texttt{write} system call. That is a \emph{cooperative} transfer: trap into the kernel, and the kernel need not return-from-trap into the same process (it may context-switch). Even without a timer, $P_1$ keeps entering the kernel.

$P_2$: a pure user-mode infinite loop; it never traps on its own. Cooperative scheduling is not enough (early Macintosh / Xerox Alto: infinite loop $\to$ reboot). The OS needs a \emph{timer interrupt} armed at boot: every few ms the hardware preempts $P_2$, saves state on its kernel stack, and the scheduler may run someone else.

On a modern kernel both processes are also preempted by the timer; $P_1$ has the extra, sufficient cooperative path via the syscall.
\end{Sol}

\begin{Qn}{Hard}{Lect 5--6}{CS-537 Fall 2008 Mid Q2}{RR finish times, and SJF with lying estimates}
Assume three jobs that enter a system and need to be scheduled. The first job that enters is called A, and it needs 10 seconds of CPU time. The second, which arrives just after A, is called B, and it needs 15 seconds of CPU time. The third, C, arrives just after B, and needs 10 seconds of CPU time.
For all questions involving round-robin, assume that there is no cost to context switching. Also assume that if job X arrives just before Y, a round-robin scheduler will schedule X before Y.
\begin{itemize}
  \item[(a)] Assuming a shortest-job-first (SJF) policy, at what time does B finish?
  \item[(b)] Assuming a longest-job-first (LJF) policy, at what time does B finish?
  \item[(c)] Assuming a round-robin policy (time slice 1 second), when does job A finish?
  \item[(d)] Assuming a round-robin policy (time slice 1 second), when does job B finish?
  \item[(e)] Assuming a round-robin policy (unknown time-slice $\le 2$ seconds), when does job B finish?
  \item[(f)] Assuming a round-robin policy (unknown time-slice), for what values of the time-slice will B finish before C?
\end{itemize}
Of course, SJF is unrealistic, because usually the OS doesn't know how long jobs are. In this system, though, the user gives the OS an estimate. If they tell you a job will last $N$ seconds, it might last anywhere between $N-5$ and $N+5$ seconds. But, being a nice, trusting OS, the OS assumes the user is exactly right.
\begin{itemize}
  \item[(g)] Assuming SJF, what estimates (by the user) will lead the OS to make the worst decisions for these jobs in terms of achieving the lowest average response time?
  \item[(h)] In this case, what will the average response time of all jobs be?
  \item[(i)] What would the average response time of all jobs be if you instead had used LJF?
\end{itemize}
\end{Qn}
\begin{Sol}
Official key (CS-537 Fa08 Mid answers, Q2). Arrivals: A then B then C, essentially at $t=0$ for these purposes.

(a)~SJF order A, C, B (A and C both 10, A arrived first; B is 15). B finishes last, at time $10+10+15=\mathbf{35}$.

(b)~LJF: B first, finishes at $\mathbf{15}$.

(c)~RR, $q=1$: the pattern is \texttt{ABC ABC} \ldots for 10 rounds (until A and C's 10 slices are done). A's slices occupy $[0,1],[3,4],\ldots,[27,28]$. A finishes at $\mathbf{28}$.

(d)~B is the longest and therefore last; total work is 35, so B finishes at $\mathbf{35}$.

(e)~Any $q\le 2$ still leaves B last, so B still finishes at $\mathbf{35}$.

(f)~Two windows. If $q\ge 15$, B runs to completion in its first slice (after A has finished its 10) and so finishes before C. If $7.5 \le q < 10$, A, B and C each need two slices, but B's second slice ends before C's second slice (B started its first slice before C). At $q=10$, A finishes in one slice, then B runs 10 of 15, then C runs all 10 and finishes \emph{before} B, so $q=10$ is excluded.

(g)~Worst SJF decisions: any estimates that make the OS run B first, i.e.\ estimate$(A)=$ estimate$(C)=X$, estimate$(B)=Y$ with $X>Y$. Possible because true runtimes are only 5 seconds apart and each estimate may be off by 5.

(h)~Then B's response time is 0, A's is 15, C's is 25. Average $\mathbf{40/3 \approx 13.33}$.

(i)~LJF with those same (wrong) estimates would run C (response 0), then A (10), then B (20). Average $\mathbf{10}$.
\end{Sol}

\begin{Qn}{Hard}{Lect 5--6}{CS-537 Fall 2013 Mid Q2}{McFlub's STCF turnaround arithmetic}
Professor McFlub says:
\begin{quote}
``Shortest Time to Completion First is a great policy, if you care about response time. But today we're going to worry about turnaround time. Let's compute the turnaround times for STCF for these three jobs, which arrive at about the same time: Job A (which needs to run for 10 time units), Job B (needs 20), and Job C (needs 30). Let's assume this happens when the jobs are run:
\texttt{ABBCCC}
The Turnaround for A is thus 0, 10 for B, and 20 for C. Thus, average turnaround for STCF is 10, which is great. See, I told you.''
\end{quote}
Correct McFlub.
\end{Qn}
\begin{Sol}
Official key (CS-537 Fa13 Mid answers, Q2):
\begin{itemize}
  \item STCF is good for \emph{turnaround}, not response time. (Response time is first-scheduled minus arrival; STCF can hold a long job to let short ones finish, which is the opposite of RR's response-time goal.)
  \item The Gantt \texttt{ABBCCC} is the right \emph{shape} (A then B then C) but the lengths are 10, 20, 30, so it is really A then BB then CCC with each letter a block of 10. Completions: $A$ at 10, $B$ at 30, $C$ at 60.
  \item Turnaround (arrivals 0): $10$, $30$, $60$, average $100/3 \approx 33.33$, not 10.
  \item (If you instead computed response times as the correction: $0$, $10$, $30$.)
\end{itemize}
McFlub reported the response times and called them turnaround, then averaged the wrong numbers.
\end{Sol}

\begin{Qn}{Medium-Hard}{Lect 5--6}{CS-537 Fall 2022 Mid Q9--Q11}{SJF problems, STCF preemption, MLFQ vs.\ SJF/STCF}
xOS scheduling (single CPU, context switches take no time).
\begin{enumerate}
  \item[9.] The xOS developers first considered using a Shortest-Job-First scheduler. Which of these problems does a SJF scheduler \emph{not} have?
  \begin{itemize}
    \item[(a)] Long jobs may never run
    \item[(b)] A job's runtime is difficult to predict
    \item[(c)] Average turnaround time may be greater than with other policies even if all jobs arrive at the same time
    \item[(d)] Does not support preemptive scheduling
    \item[(e)] It does have all of these problems
  \end{itemize}
  \item[10.] Consider STCF. Job A arrives at $T=0$ with length $500$\,ms. Job B arrives at $T=100$\,ms with length $100$\,ms. What does the CPU schedule look like?
  \begin{itemize}
    \item[(a)] A runs for $500$\,ms, then B runs for $100$\,ms
    \item[(b)] B runs for $100$\,ms, then A runs for $500$\,ms
    \item[(c)] A runs for $100$\,ms, then B runs for $100$\,ms, then A runs again for $400$\,ms
    \item[(d)] A runs for $100$\,ms, then A and B run in round-robin for $200$\,ms, then A runs for another $300$\,ms
    \item[(e)] A runs for $100$\,ms, then B runs for $100$\,ms, then A runs for $100$\,ms, then B runs for $100$\,ms, then A runs for $300$\,ms
  \end{itemize}
  \item[11.] In what regard is MLFQ better compared to a scheduler like SJF or STCF?
  \begin{itemize}
    \item[(a)] Average response time
    \item[(b)] (other options on the paper include I/O count and CPU utilisation)
    \item[(e)] None of the above
  \end{itemize}
\end{enumerate}
\end{Qn}
\begin{Sol}
9.\ \textbf{(c).} With simultaneous arrivals, SJF (and STCF) minimise average turnaround; that is the theorem. SJF \emph{does} risk starving a long job if short jobs keep arriving (a); it \emph{does} need known runtimes (b); it is defined to be non-preemptive (d). Official key: C.

10.\ \textbf{(c).} At $T=100$, A has 400\,ms left and B has 100\,ms, so STCF preempts A and runs B to completion, then A. Official key: C.

11.\ \textbf{(a).} MLFQ puts new jobs at the highest queue and RR's them with a short slice, so interactive/short jobs get a fast first schedule (good response time) without knowing runtimes in advance. SJF/STCF ignore response time. Official key: A. (Lottery / multi-CPU options on this exam are out of the MTL458 syllabus and are omitted.)
\end{Sol}

\begin{Qn}{Hard}{Lect 7}{CS-537 Spring 2016 Mid Q11}{Circle the real MLFQ rules, then name the rules in each trace}
The MLFQ (multi-level feedback queue) is a scheduling discipline. It consists of a number of rules. First, circle the rules that are actually part of the \emph{final} policy:
\begin{enumerate}
  \item If Priority(A) $>$ Priority(B), A runs (B doesn't).
  \item If Priority(A) $<$ Priority(B), A runs (B doesn't).
  \item If Priority(A) $=$ Priority(B), A \& B run in round-robin fashion.
  \item If Priority(A) $=$ Priority(B), A runs to completion, then B.
  \item When a job enters the system, it is placed at the highest priority (the topmost queue).
  \item When a job enters the system, it is placed in any queue.
  \item When a job enters the system, it is placed in the lowest queue.
  \item Once a job uses up its time slice at a given level, its priority is reduced.
  \item Once a job uses up its time slice at a given level, it moves to the end of the round-robin queue.
  \item Once a job uses up its time slice at a given level, it exits.
  \item After some time period, move each job up to a higher priority queue.
  \item After some time period, move each job down to a lower priority queue.
  \item After some time period, move all the jobs in the system to the top-most (highest-priority) queue.
\end{enumerate}
Now, write down the rule (or rules) that come into play in each of the following example traces of MLFQ behaviour (use the numbers from above). Three queues; Q3 is highest, Q1 is lowest. Note that $*$ marks when A arrives, if the information is relevant.
\begin{itemize}
  \item[(i)] A alone: a short run on Q3, then A occupies Q1 forever (\texttt{AAAAAAAAAAAA}\ldots).
  \item[(ii)] A arrives (short run on Q3) while B is a long job on Q1; A then leaves B to run on Q1.
  \item[(iii)] After A has been demoted, A and B share Q1 as \texttt{AAAABBBBAAAABBBB}\ldots.
  \item[(iv)] Two jobs enter together and RR at every level: Q3 \texttt{AB}, Q2 \texttt{AABB}, Q1 \texttt{AAABBBAAABBB}\ldots.
\end{itemize}
\end{Qn}
\begin{Sol}
Official key (CS-537 Sp16 Mid Q11): the final MLFQ rules that are circled are \textbf{1, 3, 5, 8, 13} --- exactly Lect~7's five rules (priority, RR-within-level, new job at top, demote after using the allotment/slice, periodic boost of \emph{everyone} to the top). (Rule 8 on this paper is the exam's wording of ``once a job uses up its time slice at a given level, its priority is reduced''; the refined Lect~7 Rule~4 is the same idea with allotment accounting against gaming.)

Traces (rule numbers from the list above):
\begin{itemize}
  \item[(i)] New job at the top (5), then repeated demotion until the bottom (8); only one job, so it then runs (1). Key: \textbf{5, 8}.
  \item[(ii)] A enters at the top (5) and therefore outranks B (1). After A is done or demoted, B resumes on Q1. Key: \textbf{5, 1}.
  \item[(iii)] Equal priority, so RR on Q1 (3); A got there by demotion (8). Key: \textbf{8, 3}.
  \item[(iv)] Both enter at the top (5), RR at equal priority (3), demote together after using the slice (8). Key: \textbf{5, 3, 8}.
\end{itemize}
\end{Sol}

\begin{Qn}{Hard}{Lect 7}{CS-537 Fall 2013 Mid Q3}{McFlub's MLFQ rules (random entry, no boost)}
Professor McFlub lists MLFQ rules:
\begin{itemize}
  \item If Priority(A) $>$ Priority(B), A runs (B doesn't).
  \item If Priority(A) $==$ Priority(B), A and B run in round-robin style.
  \item When a job enters the system, it is assigned a \emph{random} priority and put in the queue corresponding to that priority.
  \item Once a job uses up its time allotment at a given level (regardless of how many times it has given up the CPU), its priority is reduced (i.e., it moves down one queue).
\end{itemize}
``That's it! Amazingly, these rules are all you need.'' Correct McFlub.
\end{Qn}
\begin{Sol}
Official key (CS-537 Fa13 Mid answers, Q3):
\begin{itemize}
  \item Jobs enter at \emph{top} priority, not random. Random entry would send some short/interactive jobs straight to a low queue and destroy response time (cf.\ CS-537 Fa21 Mid Q12).
  \item Need a periodic \emph{priority boost} of all jobs to the topmost queue: (i)~to avoid starvation of long-running jobs, and (ii)~to re-evaluate a job that has changed from CPU-bound to interactive.
\end{itemize}
Rules 1, 2 and 4 as stated \emph{are} the final Lect~7 rules 1, 2 and 4 (allotment accounting, which already stops gaming). The missing piece is Rule 5 (boost after period $S$), and Rule 3 is wrong as written.
\end{Sol}

\begin{Qn}{Hard}{Lect 7}{CS-537 Spring 2002 Final Q4(a)}{Game a naive priority scheduler with I/O-reset quanta}
Assume a priority-based scheduler that works as follows. A newly created process enters at the highest priority. If the process uses up its entire time quantum when it runs, it moves to the next lowest priority. If a process initiates an I/O request, it is moved back to the top priority, and its time quantum begins anew. How would you ``game'' this scheduler?
\end{Qn}
\begin{Sol}
Issue a trivial I/O (or yield) just before the quantum expires --- e.g.\ run $99\%$ of the slice, then read one byte from a pipe/file you control. The scheduler treats you as interactive: you are boosted to the top with a full quantum, and you never ``use up'' a quantum, so you never drop. A CPU hog that is honest falls to the bottom. This is exactly the gaming attack that Lect~7's rewritten Rule~4 (charge the \emph{allotment} no matter how many times the job gave up the CPU) is designed to stop.

(Parts (b)--(d) of the original 2002 question use lottery tickets, SJF splitting, and a cache-hit-rate scheduler; lottery is out of the MTL458 syllabus and is omitted here.)
\end{Sol}

\begin{Qn}{Medium-Hard}{Lect 7}{CS-537 Fall 2022 Mid Q14}{MLFQ without a priority boost}
The xOS developers decided not to add boosting to their MLFQ scheduler to reduce code complexity. What problems could this cause?
\begin{itemize}
  \item[(a)] Interactive jobs might not be scheduled at all
  \item[(b)] CPU-heavy (batch) jobs might get starved
  \item[(c)] Batch jobs will stay at the highest priority level
  \item[(d)] There can only be one job at the highest priority level
  \item[(e)] None of the above
\end{itemize}
\end{Qn}
\begin{Sol}
\textbf{(b).} Without Rule~5, a steady stream of new or interactive jobs (always entering --- and staying --- at high priority) prevents long-running CPU-bound jobs that have sunk to the bottom from ever running. Interactive jobs are the ones MLFQ \emph{favours}; they are not the victims. Batch jobs \emph{leave} the highest level after using their allotment, so (c) is false. Official key: B.
\end{Sol}

\begin{Qn}{Medium-Hard}{Lect 2--3}{CS-537 Spring 2016 Mid Q13}{Fill in Running / Ready / Blocked through forks and I/O}
A system with three processes (A, B, C) and a single CPU. Processes can be in one of three states: Running, Ready, Blocked. If the process does not exist (yet), or has exited, put ``---''. Fill in the states of each process:
\begin{center}
\begin{tabular}{p{0.58\textwidth}ccc}
\hline
Event & A & B & C \\
\hline
Process A is loaded into memory and starts executing in \texttt{main()} & & & \\
Process A calls \texttt{fork()} and creates Process B (but A, the parent, keeps running) & & & \\
Process A issues a request to the disk; B starts executing at the return from \texttt{fork()} & & & \\
B calls \texttt{fork()}, creating Process C; B keeps running & & & \\
B's timeslice expires; C runs & & & \\
A's I/O completes (but there are no other changes) & & & \\
C waits for user input. A runs & & & \\
\hline
\end{tabular}
\end{center}
\end{Qn}
\begin{Sol}
Official key (CS-537 Sp16 Mid Q13):
\begin{center}
\begin{tabular}{lccc}
\hline
Event & A & B & C \\
\hline
A loaded, runs in \texttt{main()}
  & Running & --- & --- \\
A \texttt{fork}s B, A keeps running
  & Running & Ready & --- \\
A issues disk I/O; B runs (return from \texttt{fork})
  & Blocked & Running & --- \\
B \texttt{fork}s C; B keeps running
  & Blocked & Running & Ready \\
B's timeslice expires; C runs
  & Blocked & Ready & Running \\
A's I/O completes (no other changes)
  & Ready & Ready & Running \\
C waits for user input; A runs
  & Running & Ready & Blocked \\
\hline
\end{tabular}
\end{center}
Blocked $\to$ Ready on I/O completion, never Blocked $\to$ Running in one step (the scheduler only picks from Ready). C's ``wait for user input'' is a blocking \texttt{read}.
\end{Sol}

\subsubsection*{Conceptual Drills, MCQs \& Further Practice}

\begin{Qn}{Easy-Medium}{Lect 2}{CS-537 Spring 2018 Mid Problem IV}{Which process-state traces are possible?}
Processes exist in a number of different states. We've focused upon a few (Running, Ready, and Blocked) but real systems have slightly more. For example, xv6 also has an Embryo state (used when the process is being created), and a Zombie state (used when the process has exited but its parent hasn't yet called \texttt{wait()} on it).

Assuming you start observing the states of a given process at some point in time (not necessarily from its creation, but perhaps including that), which process states could you possibly observe? Note: once you start observing the process, you will see ALL states it is in, until you stop sampling.

To answer: Fill in A for possible, B for not possible.
\begin{enumerate}
  \item[16.] Running, Running, Running, Ready, Running, Running, Running, Ready
  \item[17.] Embryo, Ready, Ready, Ready, Ready, Ready
  \item[18.] Running, Running, Blocked, Blocked, Blocked, Running
  \item[19.] Running, Running, Blocked, Blocked, Blocked, Ready, Running
  \item[20.] Embryo, Running, Blocked, Running, Zombie, Running
\end{enumerate}
\end{Qn}
\begin{Sol}
Legal transitions: Embryo $\to$ Ready $\to$ Running; Running $\to$ Ready (descheduled), Running $\to$ Blocked (I/O / wait), Running $\to$ Zombie (\texttt{exit}); Blocked $\to$ Ready (I/O done); no transition out of Zombie; never Blocked $\to$ Running in one step; never Embryo $\to$ Running (must be Ready first).
\begin{enumerate}
  \item[16.] \textbf{A.} Run, deschedule, run, deschedule.
  \item[17.] \textbf{A.} Created, then waiting for the CPU.
  \item[18.] \textbf{B.} Blocked must become Ready before Running.
  \item[19.] \textbf{A.} Run, block on I/O, unblock (Ready), run.
  \item[20.] \textbf{B.} Skips Ready after Embryo; also runs after Zombie.
\end{enumerate}
(Official key, CS-537 Sp18 Mid answers, Problem IV.)
\end{Sol}

\begin{Qn}{Easy-Medium}{Lect 2}{CS-537 Spring 2018 Mid Problems I--II}{Two processes printing; a second process cannot reset your globals}
\textbf{Problem I.} A program's \texttt{main} is
\begin{lstlisting}[numbers=none]
int main(int argc, char *argv[]) {
  char *str = argv[1];
  while (1)
    printf("%s", str);
  return 0;
}
\end{lstlisting}
Two processes, both running instances of this program, are currently running (nothing else of relevance is, except perhaps the shell). They were invoked as a parallel command (wish project 2a): \texttt{wish> main a \&\& main b}. Which of the following beginning outputs are possible? A $=$ possible, B $=$ not possible.
1.\ \texttt{abababab} \ldots\quad
2.\ \texttt{aaaaaaaa} \ldots\quad
3.\ \texttt{bbbbbbbb} \ldots\quad
4.\ \texttt{aaaabbbb} \ldots\quad
5.\ \texttt{bbbbaaaa} \ldots

\textbf{Problem II.} \texttt{increment.c} prints \texttt{value++} in a loop starting from 0. While it is running, a separate process \texttt{reset.c} (its own \texttt{int value; value = 0;}) is run once. Which outputs of the increment process are possible?
6.\ \texttt{012345678} \ldots\quad
7.\ \texttt{012301234} \ldots\quad
8.\ \texttt{012345670123} \ldots\quad
9.\ \texttt{01234567891011} \ldots\quad
10.\ \texttt{123456789} \ldots
\end{Qn}
\begin{Sol}
\textbf{I.} Two processes, two address spaces, arbitrarily interleaved stdout. Any mix of \texttt{a} and \texttt{b} is possible, including long runs of one letter (that process just got the CPU). All five \textbf{A}. (Official key, Problem I.)

\textbf{II.} \texttt{reset}'s \texttt{value} is a different process's memory; CPU virtualisation gives each process a private address space (Lect~2). Increment never resets and starts at 0, so it prints \texttt{01234567891011}\ldots (the ``91011'' looking glued is just multi-digit numbers). It cannot skip 0, and cannot jump backwards.
\[
6\ \mathrm{A},\quad 7\ \mathrm{B},\quad 8\ \mathrm{B},\quad 9\ \mathrm{A},\quad 10\ \mathrm{B}.
\]
(Official key, Problem II.)
\end{Sol}

\begin{Qn}{Easy-Medium}{Lect 2}{CS-537 Spring 2002 Mid Q1}{Name the state}
Processes (or threads) can be in one of three states: Running, Ready, or Blocked. For each of the following four examples, write down which state the process (or thread) is in:
\begin{itemize}
  \item[(a)] Waiting in \texttt{DomainRead()} for a message from some other process to arrive.
  \item[(b)] Spin-waiting for a variable \texttt{x} to become non-zero.
  \item[(c)] Having just completed an I/O, waiting to get scheduled again on the CPU.
  \item[(d)] Waiting inside of \texttt{pthread\_cond\_wait()} for some other thread to signal it.
\end{itemize}
\end{Qn}
\begin{Sol}
(a)~Blocked (sleeping on a message / IPC receive).\\
(b)~Running (or Ready if descheduled): a spin-wait \emph{is} using the CPU; it has not blocked.\\
(c)~Ready (I/O completion moves Blocked $\to$ Ready; the scheduler has not yet dispatched it).\\
(d)~Blocked (the thread is off the CPU, waiting on the CV).\\
\textit{Note: (d) is a Lect~18 mechanism, but the question is only asking for the process/thread \emph{state}, which is Lect~2.}
\end{Sol}

\begin{Qn}{Easy-Medium}{Lect 2}{CS-537 Fall 2021 Mid Q1}{BeOS dropped a process state}
BeOS keeps track of a small amount of information about each running process in a data structure called the ``process list''. However, to keep the amount of information in the process list small, Binus reduced the number of states a process could be in. Specifically, in BeOS, there are only two states a process can be in: RUNNING and BLOCKED. What process state found in most OSes was removed from BeOS?
\begin{itemize}
  \item[(a)] STOPPED \quad (b) READY \quad (c) ACTIVE \quad (d) WAITING \quad (e) LOCKED
\end{itemize}
\end{Qn}
\begin{Sol}
\textbf{(b) READY.} The PCB/process-list normally records at least New/Ready/Running/Blocked/Dead (Lect~2). Removing Ready means there is no place to put a runnable process that does not currently have the CPU. Official key: B.
\end{Sol}

\begin{Qn}{Easy-Medium}{Lect 3}{Mandar ISI 2018 Mid Q1}{T/F: many procs per executable; one proc, many executables; PSW}
Are the following statements true or false? Write 1--2 lines to justify.
\begin{itemize}
  \item[(a)] Multiple processes can correspond to the same program / executable file at the same time.
  \item[(b)] The same process can correspond to multiple programs / executable files (not necessarily at the same time).
  \item[(c)] In user mode, the processor status word (PSW) may be modified, but arbitrary write operations on the PSW are not permitted.
\end{itemize}
\end{Qn}
\begin{Sol}
(a)~\textbf{True.} \texttt{fork} (or running the same binary twice) yields two PCBs, two address spaces, one on-disk executable. (Same stem as Mandar 2014 Mid Q1(c)(i).)

(b)~\textbf{True.} \texttt{exec()} replaces the memory image of the \emph{same} PCB/pid with a different executable. That is the whole point of the fork-then-exec pattern.

(c)~\textbf{True} as a description of dual-mode hardware: some bits of the PSW/condition codes are written by ordinary user instructions (flags after an \texttt{add}), but the \emph{mode bit} and other privileged PSW fields cannot be the target of an arbitrary user-mode load/store (otherwise a process would enter kernel mode). A user-mode \texttt{load} into the PSW as a whole is not permitted. (Near-repeat: Mandar 2022 Mid Q5(a).)
\end{Sol}

\begin{Qn}{Easy-Medium}{Lect 3--4}{Mandar ISI 2014 Mid Q1(a)--(c)(i)}{T/F: interrupts, \texttt{fork}, many processes on \texttt{a.out}}
True or false?
\begin{itemize}
  \item[(a)] Interrupts are synchronous events caused by peripheral devices.
  \item[(b)] All processes are created via the \texttt{fork()} system call.
  \item[(c)(i)] Multiple processes may run \texttt{a.out} concurrently.
\end{itemize}
\end{Qn}
\begin{Sol}
(a)~\textbf{False.} Interrupts are \emph{asynchronous} (timer, disk, \ldots). Syscalls and exceptions are the synchronous traps. (Devices cause interrupts, but ``synchronous'' is the wrong word.)

(b)~\textbf{False} on a Unix system if taken literally: \texttt{init} (pid 1) is created by the kernel at boot, not by \texttt{fork}. After that, yes, every other process is forked from a parent. The statement as written is false.

(c)(i)~\textbf{True.} Same as 2018 Mid Q1(a).
\end{Sol}

\begin{Qn}{Easy-Medium}{Lect 3}{CS-537 Spring 2018 Mid Problems V--VI}{\texttt{printf("a"); fork(); printf("b");} with and without \texttt{fork} failure}
\begin{lstlisting}[numbers=none]
int main(int argc, char *argv[]) {
  printf("a");
  fork();
  printf("b");
  return 0;
}
\end{lstlisting}
\texttt{printf} prints immediately (no buffering). A $=$ possible, B $=$ not possible.

\textbf{V.} \texttt{fork()} succeeds. Possible outputs?
21.\ \texttt{ab}\quad 22.\ \texttt{abb}\quad 23.\ \texttt{bab}\quad 24.\ \texttt{bba}\quad 25.\ \texttt{a}

\textbf{VI.} \texttt{fork()} might fail. Possible outputs of the same program?
26.\ \texttt{ab}\quad 27.\ \texttt{abb}\quad 28.\ \texttt{bab}\quad 29.\ \texttt{bba}\quad 30.\ \texttt{a}
\end{Qn}
\begin{Sol}
Unbuffered: parent prints \texttt{a} \emph{before} the fork, so both parent and child then print \texttt{b}. The only success output is \texttt{abb} (the two \texttt{b}'s may be in either process order, but both are \texttt{b}). Nothing can print \texttt{b} before \texttt{a}, and the child is not a copy of a still-buffered \texttt{a} (the problem forbids buffering).
\[
21\ \mathrm{B},\ 22\ \mathrm{A},\ 23\ \mathrm{B},\ 24\ \mathrm{B},\ 25\ \mathrm{B}.
\]
If \texttt{fork} fails, the single process still prints \texttt{a} then \texttt{b}, so \texttt{ab} becomes possible in addition to \texttt{abb}.
\[
26\ \mathrm{A},\ 27\ \mathrm{A},\ 28\ \mathrm{B},\ 29\ \mathrm{B},\ 30\ \mathrm{B}.
\]
(Official key, CS-537 Sp18 Mid answers, Problems V--VI.)
\end{Sol}

\begin{Qn}{Easy-Medium}{Lect 3}{Mandar ISI 2014 Mid Q2(d); also Mandar ISI 2024 Mid Q2(a) and 2016 Mid Q1(b) (fork part)}{What can \texttt{fork()} return?}
Briefly explain what possible value(s) the \texttt{fork()} system call can return.
\end{Qn}
\begin{Sol}
On success, \texttt{fork} returns \emph{twice}: \texttt{0} in the child, and the child's pid ($>0$) in the parent. On failure it returns $-1$ in the caller and no child is created. (2016 Mid Q1(b) also asks about \texttt{execve} --- does not return on success, $-1$ on failure, visible only in the same process --- and \texttt{exit} --- does not return to the caller; the status is collected by the parent's \texttt{wait}.)
\end{Sol}

\begin{Qn}{Easy-Medium}{Lect 3}{Mandar ISI 2024 Mid Q2(b)}{Program that forks 10 children}
Write a program to fork 10 children. Each child process should print its own pid and exit.
\end{Qn}
\begin{Sol}
\begin{lstlisting}[numbers=none]
int main(void) {
    int i;
    pid_t pid;
    for (i = 0; i < 10; i++) {
        pid = fork();
        if (pid < 0) { perror("fork"); return 1; }
        if (pid == 0) {                 /* child */
            printf("%d\n", getpid());
            exit(0);
        }
    }
    while (wait(NULL) > 0) ;            /* reap all 10 */
    return 0;
}
\end{lstlisting}
The child must \texttt{exit} (or \texttt{break}) so it does not itself iterate and create grandchildren. Compare CS-537 Sp01 Mid Q6, where both sides \texttt{while(1);} after the first fork, so the loop never continues and only \emph{two} processes result.
\end{Sol}

\begin{Qn}{Medium}{Lect 3}{CS-537 Spring 2001 Mid Q6}{A \texttt{fork} loop that never loops}
Assume we run the following code snippet. After waiting for a ``long'' time, how many processes will be running on the machine, ignoring all other processes except those involved with this code snippet? You can assume that \texttt{fork()} never fails.
\begin{lstlisting}[numbers=none]
void runMe() {
    for (int i = 0; i < 23; i++) {
        int rc = fork();
        if (rc == 0) {
            while (1) ;
        } else {
            while (1) ;
        }
    }
}
\end{lstlisting}
\end{Qn}
\begin{Sol}
\textbf{2 processes.} After the first \texttt{fork}, \emph{both} the parent and the child enter \texttt{while (1);} and never return to the \texttt{for}-increment. The loop body therefore runs once. (A common trap is to answer $2^{23}$ or $24$.)
\end{Sol}

\begin{Qn}{Medium}{Lect 3}{CS-537 Spring 2012 Mid Q2}{\texttt{exec} then \texttt{fork}: a subtle shell bug}
You write a UNIX shell, but instead of calling \texttt{fork()} then \texttt{exec()} to launch a new job, you instead insert a subtle difference: the code first calls \texttt{exec()} and then calls \texttt{fork()}. What is the impact of this change to the shell, if any?
\end{Qn}
\begin{Sol}
\texttt{exec} \emph{replaces} the caller's address space. If it succeeds, the shell \emph{is} now the command, \texttt{fork} never runs, and there is no shell left to read the next line. If \texttt{exec} fails, \texttt{fork} may run, but then you have forked another copy of the (still broken) shell, not the user's program. The Lect~3 pattern is mandatory: fork first (so the shell survives), exec in the child.
\end{Sol}

\begin{Qn}{Medium}{Lect 3}{Mandar ISI 2010 End Q3(c)--(d)}{Shell: redirection, and \texttt{firefox} vs.\ \texttt{firefox \&}}
(c)~Clearly explain the steps by which a typical shell handles the following command:
\[
\texttt{a.out < input.txt > output.txt}
\]
(d)~Explain what happens when the following commands are given to a typical shell: (i)~\texttt{firefox}\quad (ii)~\texttt{firefox \&}
\end{Qn}
\begin{Sol}
(c)~The shell \texttt{fork}s. In the \emph{child}, before \texttt{exec}:
\begin{enumerate}
  \item close fd 0 (stdin); \texttt{open("input.txt", O\_RDONLY)} --- it lands in fd 0;
  \item close fd 1 (stdout); \texttt{open("output.txt", O\_CREAT|O\_WRONLY|O\_TRUNC, \ldots)} --- it lands in fd 1;
  \item \texttt{execv("a.out", \ldots)}.
\end{enumerate}
The parent \texttt{wait}s. This is Lect~3's redirection case study (\texttt{ls > a.text}): rewire the child's stdout/stdin, then exec. (\texttt{dup2} is the usual implementation; the lecture's story is ``rewire then exec''.)

(d)~(i)~\texttt{firefox}: fork, exec \texttt{firefox} in the child, parent \texttt{wait}s --- the shell is blocked until the browser exits (foreground). (ii)~\texttt{firefox \&}: fork, exec in the child, parent does \emph{not} wait --- the shell immediately prints the next prompt (background). The child is not reaped until later (or is adopted by \texttt{init} if the shell never waits).
\end{Sol}

\begin{Qn}{Medium}{Lect 3}{Mandar ISI 2016 Mid Q1(c)}{Terminal, Firefox, PDF viewer: \texttt{proc} tree, then the browser exits}
Suppose a user starts the Firefox browser by running \texttt{\$ firefox} from the command prompt in a terminal. While browsing, she clicks on a link to a PDF file, which is opened for viewing in a document viewer program.
\begin{itemize}
  \item[(i)] Draw a diagram showing the parent-child relationship between the \texttt{proc} structures corresponding to the terminal, browser and document viewer processes.
  \item[(ii)] Explain how the diagram will change if the user now exits the browser.
\end{itemize}
\end{Qn}
\begin{Sol}
(i)~Typical tree (Lect~3: \texttt{init} is ancestor of all; the shell is the terminal's child):
\[
\texttt{init} \to \text{terminal/shell} \to \text{firefox} \to \text{pdf-viewer}.
\]
(Firefox may fork a helper; the exam wants this three-level user picture.)

(ii)~When firefox \texttt{exit}s, its children are \emph{orphans} and are adopted by \texttt{init} (Lect~3). The pdf-viewer remains running, now with parent \texttt{init}. The shell, which was \texttt{wait}ing for firefox (foreground), is woken, reaps the firefox zombie, and reprints the prompt.
\end{Sol}

\begin{Qn}{Medium}{Lect 3}{Mandar ISI 2012 Mid Q1(c); also Mandar ISI 2017 Mid Q1(c)}{\texttt{exit}, zombies, orphans, when the PCB is reclaimed}
(a)~[2012 Mid] When a process $P$ calls \texttt{exit()}, how are the parent and children of $P$ affected?

(b)~[2017 Mid] When does the kernel reclaim the space that is allocated to hold the PCB (\texttt{proc} structure) of user-created processes? Clearly explain your answer considering all possible scenarios.
\end{Qn}
\begin{Sol}
(a)~$P$ becomes a \emph{zombie}: it no longer runs, but its PCB (exit status, pid) is kept so the parent can \texttt{wait()}. The parent, if waiting, is woken and reaps $P$. $P$'s children become \emph{orphans} and are adopted by \texttt{init}, which will \texttt{wait} for them.

(b)~The PCB cannot be freed at \texttt{exit}: the parent may still \texttt{wait()} for the status. It is reclaimed when
\begin{itemize}
  \item the parent calls \texttt{wait}/\texttt{waitpid} and reaps the zombie, or
  \item the parent itself has already exited, so \texttt{init} adopted the process and \texttt{init}'s \texttt{wait} reaps it.
\end{itemize}
A still-running process, and a zombie whose living parent has not yet waited, both occupy a PCB. (Do not confuse this with SIGCHLD details, which are not in Lect~1--7.)
\end{Sol}

\begin{Qn}{Medium}{Lect 3}{Mandar ISI 2015 Mid Q1}{Does \texttt{strcpy} / \texttt{strdup} / \texttt{stat} cause a syscall?}
For each of the standard C library functions listed below, explain in 1--2 lines whether the function may result in a system call.
\begin{itemize}
  \item[(a)] \texttt{strcpy(char *dest, const char *src)}: copies the string pointed to by \texttt{src}, including the terminating null byte, to the buffer pointed to by \texttt{dest};
  \item[(b)] \texttt{strdup(const char *s)}: returns a pointer to a new string which is a duplicate of \texttt{s} (memory for the new string is obtained with \texttt{malloc()}, and can be freed with \texttt{free()}).
  \item[(c)] \texttt{stat(const char *path, struct stat *buf)}: returns (in \texttt{buf}) information about a file (the size of the file, the time when it was last modified, etc.).
\end{itemize}
\end{Qn}
\begin{Sol}
(a)~\textbf{No} (not as such). Pure user-space copy; no kernel service. (A later fault if \texttt{dest}/\texttt{src} is a bad pointer is a trap, but \texttt{strcpy} itself is not a syscall wrapper.)

(b)~\textbf{Maybe.} \texttt{malloc} may be satisfied from the process's existing heap (no syscall) or may call \texttt{sbrk}/\texttt{mmap} to grow the heap (syscall). So \texttt{strdup} \emph{may} trap into the kernel.

(c)~\textbf{Yes.} File metadata lives in the kernel/FS; \texttt{stat} is a syscall (libc stub around it). This is Lect~3's ``libc hides syscalls'' example, dual to \texttt{printf}$\to$\texttt{write}.
\end{Sol}

\begin{Qn}{Easy-Medium}{Lect 3}{Mandar ISI 2014 Mid Q2(b)}{A libc function that is \emph{not} a syscall}
Give an example of a function from the standard C library that does not involve a system call. Justify your answer.
\end{Qn}
\begin{Sol}
Anything that is pure user-space: \texttt{strcpy}, \texttt{strlen}, \texttt{atoi}, \texttt{sin}, \ldots. They only manipulate the process's own memory or do arithmetic; they never need kernel mode. Contrast \texttt{read}/\texttt{write}/\texttt{fork}/\texttt{stat}/\texttt{open}.
\end{Sol}

\begin{Qn}{Medium}{Lect 3}{CS-537 Fall 2004 Mid1 Q2}{Why is a \texttt{fork}ed child a \emph{nearly} exact copy?}
When a Unix process calls \texttt{fork()}, the system creates (in the child) a nearly exact copy of the calling (parent) process. Why is the copy ``nearly exact'' and not an ``exact'' copy of the parent process?
\end{Qn}
\begin{Sol}
The child gets a copy of the parent's address space, register file, and open-file table, but a few PCB fields \emph{must} differ: a new pid; \texttt{fork}'s return value (0 in the child, child's pid in the parent); the parent-pointer; accounting/start time. If it were a bitwise-identical PCB, two processes would share a pid and \texttt{fork} could not tell parent from child.
\end{Sol}

\begin{Qn}{Medium}{Lect 2}{Mandar ISI 2012 Mid Q1(a)}{Seven constituents of a process's context}
List and briefly describe (in 2--3 sentences each) the 7 major constituents of a process's context. Which of these items need to be copied when saving a process's context? Justify your answer.
\end{Qn}
\begin{Sol}
A standard Unix-style list matching Lect~2's ``what constitutes a process'':
\begin{enumerate}
  \item \textbf{User address space} --- code, static data, heap, stack.
  \item \textbf{CPU registers} --- PC, SP, general-purpose registers.
  \item \textbf{Kernel stack / trap frame} --- saved user registers for return-from-trap.
  \item \textbf{PCB / scheduling state} --- pid, state (Ready/Running/Blocked), priority, parent pointer.
  \item \textbf{Open file descriptors} --- the fd table (stdin/stdout/stderr and any other files).
  \item \textbf{Environment / arguments} --- \texttt{argv}, \texttt{envp}.
  \item \textbf{Accounting / credentials} --- uid, resource usage, working directory.
\end{enumerate}
On a \emph{context switch} (Lect~4) you must save/restore the CPU registers and kernel SP (and MMU registers, later lectures). You do \emph{not} copy the address space or fd table --- they stay in memory, pointed to by the PCB. (A \texttt{fork} \emph{does} copy/share the address space and fd table; that is process creation, not a context switch.)
\end{Sol}

\begin{Qn}{Easy-Medium}{Lect 4}{Mandar ISI 2014 Mid Q2(a); also Mandar ISI 2024 Mid Q1(b)}{What the IDT stores, how it is accessed, and its size}
(a)~In the x86 family of processors, what does the Interrupt Descriptor Table store? How is this table accessed?

(b)~[2024 Mid] Recall that each entry of the Interrupt Descriptor Table is 8 bytes long. Compute the size of the Interrupt Descriptor Table.
\end{Qn}
\begin{Sol}
(a)~The IDT (Lect~4) is a table of handlers: for each interrupt/exception/syscall number, an entry holds the kernel address to jump to. It is built at boot. Hardware locates it via the privileged \texttt{idtr} register; a user process cannot load \texttt{idtr} (otherwise it could install its own trap table). On a trap, the hardware uses the vector to index the IDT and jumps to that handler.

(b)~x86 classic IDT has 256 vectors, 8 bytes each: $256\times 8 = \mathbf{2048}$ bytes (2\,KB). (The 2024 paper also asked whether \texttt{idtr} is part of a process's context: \textbf{no} --- it is system-wide, set at boot, not saved/restored on a context switch.)
\end{Sol}

\begin{Qn}{Easy-Medium}{Lect 4}{Mandar ISI 2011 Mid Q1(a)--(b)}{Synchronous vs.\ asynchronous; getting a syscall's return value}
(a)~Using system calls and interrupts as examples, clearly explain the difference between synchronous and asynchronous events.

(b)~Describe the mechanism by which a user's program gets back a return value from a kernel function at the end of a system call.
\end{Qn}
\begin{Sol}
(a)~A \textbf{syscall} is synchronous: it happens because the process executed a trap instruction at a well-defined program point. An \textbf{interrupt} is asynchronous: a device/timer raises a line at a moment unrelated to the current instruction.

(b)~The kernel places the return value in a calling-convention register (e.g.\ \texttt{eax}/\texttt{rax} on x86) in the process's trap frame on the kernel stack. Return-from-trap restores that register into the user-visible register file, then sets the PC to the instruction \emph{after} the trap. libc reads that register as the syscall's return value (negative values are turned into \texttt{errno}).
\end{Sol}

\begin{Qn}{Easy-Medium}{Lect 4}{Mandar ISI 2010 Mid Q1(a)--(b); also Mandar ISI 2014 Mid Q2(c)}{Two modes, and three situations for a context switch}
(a)~Processor architectures usually support at least two processor operating modes. Briefly explain the difference between these two modes.

(b)~List the three situations in which a context switch between processes can occur.
\end{Qn}
\begin{Sol}
(a)~Lect~3 dual-mode: \textbf{user mode} (mode bit 1) cannot execute privileged instructions (I/O, load IDTR, halt, \ldots); \textbf{kernel/monitor mode} (mode bit 0) can. A trap raises privilege; return-from-trap drops it.

(b)~The OS need not return to the same process after a trap (Lect~4). Three situations:
\begin{enumerate}
  \item The running process blocks (I/O syscall, \texttt{wait}, \ldots).
  \item The running process terminates (\texttt{exit}, fatal fault).
  \item Preemption: the timer interrupt (or a higher-priority wakeup) decides someone else should run.
\end{enumerate}
Cooperative systems have only (1)--(2); a timer adds (3).
\end{Sol}

\begin{Qn}{Easy-Medium}{Lect 4}{CS-537 Spring 2001 Mid Q5}{Cooperative scheduling: what, and why it is a bad idea}
What is a cooperative approach to scheduling, and why is it potentially a bad idea?
\end{Qn}
\begin{Sol}
The OS only gains control when the process \emph{volunteers}: a syscall, or an illegal operation (divide-by-zero, bad memory). There is no timer. Bad: a user-mode infinite loop never traps, so it monopolises the CPU until reboot (early Mac OS, Xerox Alto --- Lect~4).
\end{Sol}

\begin{Qn}{Medium}{Lect 4}{CS-537 Fall 2011 Mid Q1--Q3}{Syscall steps, installing the trap table, why a timer interrupt}
(1)~Most hardware provides user mode and kernel mode; applications are run in user mode, while the OS runs in kernel mode. The tricky part is transitioning from one mode to the other, which is done via a system call. What takes place during a system call?

(2)~To support trap handling, the OS has to set up a trap table of some kind; the trap table is used by the hardware to know which OS code to run when a particular trap occurs. How does the OS inform the hardware of the location of the trap table?

(3)~Timer interrupts are a useful mechanism for the OS. Why?
\end{Qn}
\begin{Sol}
(1)~User executes a trap instruction $\to$ hardware raises privilege, switches to the kernel stack, saves PC/registers, indexes the IDT, jumps to the handler $\to$ OS performs the syscall $\to$ return-from-trap restores registers, drops to user mode, resumes at the instruction after the trap. (The ``You Are The Grader'' paper had a student write this; the correct three beats are trap, run the designated handler, return-from-trap.)

(2)~At boot, the OS executes a \emph{privileged} instruction that loads the IDTR with the address of the IDT. The instruction must be privileged, otherwise any process could install a fake trap table.

(3)~Preemption: even a process that never makes a syscall is interrupted every few ms so the OS can schedule someone else, account CPU time, and implement RR/MLFQ. Also wall-clock time.
\end{Sol}

\begin{Qn}{Medium-Hard}{Lect 4}{CS-537 Fall 2004 Mid1, long Q1}{Timer interrupt vs.\ an instruction-counter interrupt}
The classic mechanism the OS uses is a timer interrupt.
\begin{itemize}
  \item[(a)] Describe how the OS uses a timer interrupt to gain control of the CPU, and why this is important.
  \item[(b)] Imagine a new hardware mechanism that counts the total number of instructions that have been executed and then raises an interrupt after some fixed number of instructions have passed. The OS uses this mechanism by configuring an instruction-counter interrupt register (ICIR); by setting the value of the ICIR to $x$, the OS makes sure that after $x$ instructions are executed, an interrupt is raised. Is the ICIR a good mechanism for the OS to use to gain control of the CPU? Why or why not?
  \item[(c)] Assume we are using the ICIR, but we wish to mimic the behaviour of the good old timer interrupt; that is, we want a clock tick every 10 milliseconds or so. How would you use the instruction-counter interrupt mechanism to achieve this effect?
  \item[(d)] Now imagine a new hardware mechanism that counts the number of load and store instructions that have been executed, and raises an interrupt after some (configurable) number of loads/stores have been executed. Is this a good mechanism for the OS to use to gain control of the CPU? Why or why not?
\end{itemize}
\end{Qn}
\begin{Sol}
(a)~At boot the OS arms a programmable timer. Every $T$\,ms the hardware interrupts, the trap path saves the user context on the kernel stack, and the scheduler may pick another process. This is how non-cooperative (preemptive) CPU virtualisation works.

(b)~Mostly yes: it guarantees the OS gets control after a bounded amount of \emph{work}, so a compute loop cannot run forever. Caveat: instruction time is not wall-clock time (halts, cache misses, I/O waits), so it is a poor real-time clock and a process that executes slowly (many stalls) holds the CPU longer in wall-clock terms.

(c)~Calibrate: measure (or know) roughly how many instructions this CPU executes in 10\,ms, and load that count into the ICIR at each tick. Recompute if frequency scaling changes. It will only be approximate.

(d)~No, not by itself. A process that only does arithmetic (no loads/stores), or that lives in registers, would never be interrupted. A process that streams memory would be interrupted constantly. The timer (or an instruction counter) does not have this workload bias.
\end{Sol}

\begin{Qn}{Easy-Medium}{Lect 4}{CS-537 Fall 2004 Mid1 Q1}{What is a context switch?}
What is a context switch? Describe what the OS must do to implement a context switch correctly.
\end{Qn}
\begin{Sol}
A context switch is the mechanism that stops one process and resumes another on the same CPU. The OS (typically from a trap/timer handler) saves the current process's general-purpose registers, PC, and kernel SP onto \emph{that} process's kernel stack / PCB, switches the kernel stack pointer to the next process, restores that process's registers, and return-from-traps into its user mode. It must also switch address-space registers (base/bounds or PTBR --- later lectures) so the next process sees its own memory. Lect~4's A$\to$B picture.
\end{Sol}

\begin{Qn}{Easy-Medium}{Lect 4}{CS-537 Fall 2004 Final Q12}{What does \emph{not} need to be saved on a context switch?}
A process has associated state which must be saved and restored when a context switch takes place. In the Broken OS, though, too much state is saved and restored. Which of the following state does \emph{not} need to be saved and restored during a context switch?
\begin{itemize}
  \item[(a)] The general purpose (integer) registers
  \item[(b)] The floating point registers
  \item[(c)] The contents of the TLB
  \item[(d)] The contents of the process's address space
\end{itemize}
\end{Qn}
\begin{Sol}
\textbf{(d)} --- and typically \textbf{(c)} as well, depending on the hardware. The address space \emph{stays in memory}; the OS only switches the pointer to it (base/bounds or page-table base). Copying the whole address space on every switch would make timesharing impossible. GPRs and FP registers \emph{are} per-process CPU state and must be saved. The TLB can be flushed (or tagged with an ASID --- Lect~13); one does not ``save the contents of the TLB'' into the PCB as if it were a register file, though flushing is a form of discarding, not saving. The intent of the Broken-OS question is that (d) is the absurd one.
\end{Sol}

\begin{Qn}{Easy-Medium}{Lect 5--6}{CS-537 Spring 2018 Mid Problem IX}{FIFO turnaround in $[10,20]$: which runtimes?}
Assume a FIFO scheduler and only one CPU. Each job has a ``required'' runtime.
Job A arrives at time $=0$, required runtime $=X$;
Job B arrives at time $=5$, required runtime $=Y$;
Job C arrives at time $=10$, required runtime $=Z$.
Assuming an average turnaround time between 10 and 20 time units (inclusive), which of the following run times for A, B, and C are possible? A $=$ possible, B $=$ not.
\begin{enumerate}
  \item[41.] $A=10$, $B=10$, $C=10$
  \item[42.] $A=20$, $B=20$, $C=20$
  \item[43.] $A=5$, $B=10$, $C=15$
  \item[44.] $A=20$, $B=30$, $C=40$
  \item[45.] $A=30$, $B=1$, $C=1$
\end{enumerate}
\end{Qn}
\begin{Sol}
FIFO, A then B then C. TAT $= T_{\mathrm{completion}}-T_{\mathrm{arrival}}$.
\begin{itemize}
  \item[41.] TAT $10$, $15$, $20$; avg $15$. \textbf{A.}
  \item[42.] $20$, $35$, $50$; avg $35$. \textbf{B.}
  \item[43.] $5$, $10$, $20$; avg $35/3\approx 11.67$. \textbf{A.}
  \item[44.] $20$, $45$, $80$; avg $\approx 48.3$. \textbf{B.}
  \item[45.] $30$, $26$, $22$; avg $26$. Convoy. \textbf{B.}
\end{itemize}
(Official key, CS-537 Sp18 Mid answers, Problem IX.)
\end{Sol}

\begin{Qn}{Easy-Medium}{Lect 5--7}{CS-537 Spring 2018 Mid Problem X}{Which discipline could produce this schedule?}
Assume the following schedule for three jobs A, B, and C:
A runs first (for 10 time units) but is not yet done;
B runs next (for 10) but is not yet done;
C runs next (for 10) and runs to completion;
A runs to completion (for 10);
B runs to completion (for 5).
Which scheduling disciplines could allow this schedule? A $=$ possible, B $=$ not.
\begin{enumerate}
  \item[46.] FIFO
  \item[47.] Round Robin
  \item[48.] STCF (Shortest Time to Completion First)
  \item[49.] Multi-level Feedback Queue
\end{enumerate}
\textit{(The original paper also listed lottery scheduling; that policy is not in the MTL458 syllabus and is omitted.)}
\end{Qn}
\begin{Sol}
Remaining work after the first three slices: A has 10 left, B has 5, C is done.
\begin{itemize}
  \item[46.] \textbf{B.} FIFO would run A to completion first.
  \item[47.] \textbf{A.} RR with a 10-unit slice, A-B-C order.
  \item[48.] \textbf{B.} After C finishes, STCF would run B (5 left) before A (10 left).
  \item[49.] \textbf{A.} MLFQ can behave like RR (same queue, or after they have all dropped).
\end{itemize}
(Official key, CS-537 Sp18 Mid answers, Problem X; lottery was A / ``anything is possible'' and is OOS.)
\end{Sol}

\begin{Qn}{Easy-Medium}{Lect 5--6}{CS-537 Fall 2021 Mid Q2--Q3}{FIFO convoy, and TAT of job B}
(2)~To keep scheduling simple in BeOS, Binus chose the FIFO policy. Which one of the following statements about FIFO scheduling is true?
\begin{itemize}
  \item[(a)] It reduces average turnaround time
  \item[(b)] It reduces average response time
  \item[(c)] It suffers from the ``convoy'' problem
  \item[(d)] It is a preemptive scheduling policy
  \item[(e)] It is different from the FCFS policy
\end{itemize}
(3)~Assuming FIFO: job A arrives at $T=0$ and needs 10\,s; job B arrives at $T=5$ and needs 20\,s; job C arrives at $T=10$ and needs 5\,s. No I/O. Turnaround time of job B?
\begin{itemize}
  \item[(a)] 10 \quad (b) 20 \quad (c) 30 \quad (d) 25 \quad (e) none of the above
\end{itemize}
\end{Qn}
\begin{Sol}
(2)~\textbf{(c).} FIFO $=$ FCFS, non-preemptive; a long job at the head of the queue delays everyone behind it (convoy). It does not specially reduce TAT or response time. Official key: C.

(3)~\textbf{(d) 25.} A runs $[0,10]$; B runs $[10,30]$; B's TAT $=30-5=25$. Official key: D.
\end{Sol}

\begin{Qn}{Easy-Medium}{Lect 5--6}{CS-537 Spring 2002 Mid Q4}{SJF vs.\ LJF turnaround}
Three jobs (A, B, and C) arrive to the job scheduler at time 0. Job A needs 10 seconds of CPU time, Job B needs 20 seconds, and Job C needs 30 seconds.
\begin{itemize}
  \item[(a)] What is the average turnaround time for the jobs, assuming a shortest-job-first (SJF) scheduling policy?
  \item[(b)] What is the average turnaround time assuming a longest-job-first (LJF) policy?
  \item[(c)] Which finishes first, Job C in SJF or Job A in LJF?
\end{itemize}
\end{Qn}
\begin{Sol}
(a)~Order A, B, C. Completions $10,30,60$. Avg TAT $100/3 \approx \mathbf{33.33}$\,s.\\
(b)~Order C, B, A. Completions $30,50,60$. Avg TAT $140/3 \approx \mathbf{46.67}$\,s.\\
(c)~C under SJF finishes at 60; A under LJF finishes at 60. \textbf{Tie.}
\end{Sol}

\begin{Qn}{Easy-Medium}{Lect 5--6}{CS-537 Spring 2001 Mid Q11}{Lowest \emph{total} run time of a mixed workload}
For a workload consisting of ten CPU-bound jobs of varying lengths (half run for 1 second, and the other half for ten seconds), which policy would result in the lowest total run time for the entire workload? Circle all that apply.
\begin{itemize}
  \item[(a)] Shortest Job First
  \item[(b)] Shortest-Time to Completion First
  \item[(c)] Round-robin with a 100 millisecond quantum
  \item[(d)] Multi-level Feedback Queue
\end{itemize}
\end{Qn}
\begin{Sol}
\textbf{All of them} (assuming no I/O and zero extra overhead). Total completion time of a pure CPU-bound, work-conserving schedule is just the sum of the runtimes, $5\cdot 1 + 5\cdot 10 = 55$\,s, independent of order. This is why CS-537 Fa11 Mid Q5's ``total completion time'' metric is not useful: it does not distinguish policies. (If context-switch cost is non-zero, RR/MLFQ pay more.)
\end{Sol}

\begin{Qn}{Easy-Medium}{Lect 5--6}{CS-537 Fall 2011 Mid Q5--Q6}{A useless metric, and why STCF is rarely implemented}
(5)~To compare scheduling policies, we often use certain metrics. Here we use a new metric: total completion time. The total completion time tells us when all of the jobs in a workload are done. What do you think about this metric?

(6)~One of the best scheduling policies out there is shortest-time-to-completion first, or STCF. What does STCF do? If it's so great, why do we rarely see it implemented?
\end{Qn}
\begin{Sol}
(5)~For CPU-only jobs and a work-conserving scheduler it is \emph{the same} for every policy (sum of burst lengths, after the last job). It does not capture fairness, response time, or convoy effects. Not a useful comparator.

(6)~STCF is preemptive SJF: always run the ready job with the least remaining time. Optimal for average turnaround (under the usual assumptions). Rarely used because the OS does \emph{not} know remaining runtimes (Lect~5--6 relaxes ``runtimes known''); MLFQ is the practical approximation that learns from behaviour.
\end{Sol}

\begin{Qn}{Easy-Medium}{Lect 5--6}{Mandar ISI 2012 Mid Q2(a)--(b)}{Starvation under FCFS vs.\ priority; TAT vs.\ response}
(a)~Explain whether starvation can happen under (i)~FCFS; (ii)~non-preemptive priority-based scheduling. Clearly state any assumptions.

(b)~What is the difference between turnaround time and response time?
\end{Qn}
\begin{Sol}
(a)~(i)~FCFS: \textbf{no} starvation if every job has finite burst (the queue is a FIFO; you eventually reach the head). (ii)~Non-preemptive priority: \textbf{yes} --- a low-priority job sits forever if a stream of higher-priority jobs keeps arriving. (Assumes priorities are static and a high-priority job can arrive while a low one is waiting; the currently running job is not preempted, but that does not help the waiter.)

(b)~$T_{\mathrm{turnaround}} = T_{\mathrm{completion}}-T_{\mathrm{arrival}}$; $T_{\mathrm{response}} = T_{\mathrm{firstschedule}}-T_{\mathrm{arrival}}$ (Lect~5--6). Response cares about the first slice (interactive feel); turnaround cares about finishing.
\end{Sol}

\begin{Qn}{Easy-Medium}{Lect 5--6}{Mandar ISI 2018 Mid Q3}{Define starvation; FCFS / priority / RR: enqueue, select, starve?}
(a)~Define starvation in the context of process scheduling.

(b)~Let $n$ denote the number of ready processes at any point. For each of FCFS, priority-based scheduling, and round-robin, write (i)~the asymptotic time complexity of adding a process to the ready queue; (ii)~the asymptotic time complexity of selecting a process during scheduling; (iii)~whether starvation is possible. Justify each answer in 1--2 lines.
\end{Qn}
\begin{Sol}
(a)~A ready process waits indefinitely --- it never gets the CPU --- even though the system keeps scheduling other processes.

(b)~Assume a linked-list ready queue unless noted.
\begin{center}
\small
\begin{tabularx}{\linewidth}{@{}lccX@{}}
\hline
Policy & Enqueue & Select & Starvation? \\
\hline
FCFS & $\Theta(1)$ (tail) & $\Theta(1)$ (head)
  & No (finite bursts). \\
Priority (unsorted)
  & $\Theta(1)$ & $\Theta(n)$ (scan)
  & \textbf{Yes}: a low-priority job can wait forever. \\
Priority (sorted)
  & $\Theta(n)$ (insert) & $\Theta(1)$
  & \textbf{Yes}, same. \\
RR & $\Theta(1)$ (tail) & $\Theta(1)$ (head/requeue)
  & No: every ready job runs every $n q$ time. \\
\hline
\end{tabularx}
\end{center}
\end{Sol}

\begin{Qn}{Easy-Medium}{Lect 5--6}{Mandar ISI 2010 Mid Q2(a)--(b)}{SJF vs.\ STCF; large RR quantum}
(a)~Using an example, explain the difference between the preemptive and non-preemptive versions of the Shortest Job First scheduling algorithm.

(b)~Explain one advantage and one disadvantage of having a large time quantum in round-robin scheduling.
\end{Qn}
\begin{Sol}
(a)~Non-preemptive SJF: once a job starts, it runs to the end of its CPU burst. Preemptive SJF $=$ STCF: if a new job arrives whose remaining time is shorter, the current job is preempted. Example: A (length 10) starts at $t=0$; B (length 2) arrives at $t=1$. SJF finishes A then B (B waits 9). STCF switches to B at $t=1$, B done at $t=3$, then A.

(b)~Advantage: amortises context-switch cost; with $q\to\infty$, RR $\to$ FIFO. Disadvantage: response time of short/interactive jobs collapses toward FIFO (a short job may wait a long slice behind a hog). Lect~5--6 quantum tradeoff.
\end{Sol}

\begin{Qn}{Medium}{Lect 5--6}{Mandar ISI 2011 Mid Q2}{Approximate SJF on a desktop; where true SJF works}
(a)~Recall that the Shortest Job First (SJF) scheduling algorithm cannot be used in the Operating Systems typically used in desktop computers. Describe an approximate version of this algorithm that can be used in such OSs.

(b)~Briefly describe an operating environment in which the original SJF scheduling algorithm can be used.
\end{Qn}
\begin{Sol}
(a)~You do not know burst lengths a priori. Approximate: exponential average of recent CPU bursts, $S_{n+1}=\alpha t_n+(1-\alpha)S_n$ (classic textbook), \emph{or} --- what MTL458 actually teaches --- \textbf{MLFQ}: new/interactive jobs stay high (look short); CPU hogs sink. That is Lect~7's ``approximate SJF''.

(b)~Batch systems where job lengths are known or declared (old batch processing; some HPC queues with user-supplied wall-time limits, if you trust them --- cf.\ CS-537 Fa08 Mid Q2's lying estimates).
\end{Sol}

\begin{Qn}{Medium}{Lect 5--6}{Mandar ISI 2014 Mid Q3(a)}{Petrol pump: SJF vs.\ RR vs.\ multi-level queue}
Consider a petrol pump that serves 2-wheelers, cars, and trucks. Assume that, on average,
\begin{itemize}
  \item it takes 3, 10, and 20 minutes to serve 2-wheelers, cars, and trucks, respectively;
  \item 2-wheelers, cars and trucks take a total of 2, 4, and 8 minutes respectively to move into position for fuelling, and vacating the spot afterwards;
  \item owners / drivers of 2-wheelers have the least patience, and will leave if they are made to wait for a long time, while owners / drivers of trucks have the greatest patience.
\end{itemize}
Clearly explain the advantages / disadvantages of (i)~Shortest-job first, (ii)~Round-robin, (iii)~Multi-level queue in this scenario. State any additional assumptions.
\end{Qn}
\begin{Sol}
Treat ``service time'' as the CPU burst and ``move in/out'' as context-switch cost. Patience $\approx$ response-time constraint.

(i)~\textbf{SJF.} Always serve the 2-wheeler (3\,min) before a car before a truck. Good average turnaround and matches 2-wheeler impatience. Trucks may starve if 2-wheelers keep arriving (same as SJF starvation of long jobs). Positioning cost is paid once per vehicle, which is fine.

(ii)~\textbf{RR.} Time-slicing a pump is disastrous: every quantum you pay 2--8 minutes of repositioning, so utilisation collapses (Lect~5--6: small $q$ $\Rightarrow$ CS cost dominates). Impatient 2-wheelers also hate being sent away mid-fill. One large quantum $\approx$ FIFO, which then has a convoy of trucks.

(iii)~\textbf{MLQ.} Three queues: 2-wheelers $>$ cars $>$ trucks, RR or FCFS within a queue. Matches patience, avoids RR's repositioning tax, and is the practical ``priority by class'' analogue of MLFQ's high-vs-low. Need a periodic boost (or aging) so a truck is not starved all day --- Lect~7 Rule~5 in this costume.
\end{Sol}

\begin{Qn}{Medium}{Lect 5--6}{CS-537 Fall 2009 Mid Q1}{RR at 100\,ms vs.\ 1\,ms vs.\ random}
According to the Old Gray Person, an old version of LostOS used a Round Robin scheduling policy with a 100-millisecond timeslice.
\begin{itemize}
  \item[(a)] What are the major strengths of the Round Robin approach?
  \item[(b)] What are its major weaknesses?
\end{itemize}
The creators later changed the timeslice to 1 millisecond (instead of 100).
\begin{itemize}
  \item[(c)] In what way did this change make the LostOS scheduler better? (if any)
  \item[(d)] In what way did this change make the LostOS scheduler worse? (if any)
\end{itemize}
Finally they replaced ``which process runs next'' with a Random approach (pick uniformly from the ready set).
\begin{itemize}
  \item[(e)] How does the Random approach compare to traditional Round Robin? Is it better, worse, or quite similar?
\end{itemize}
\end{Qn}
\begin{Sol}
Official key (CS-537 Fa09 Mid answers, Q1):
\begin{itemize}
  \item[(a)] Improves \emph{response time} (jobs get a slice quickly); simple; fair (no starvation).
  \item[(b)] Terrible \emph{turnaround} (stretches every job to the max); context-switch overhead.
  \item[(c)] Response time even better (a new job waits at most 1\,ms).
  \item[(d)] Context-switch cost can dominate; the machine spends its life in the OS.
  \item[(e)] If ``random'' means ``pick a random ready job every quantum'', behaviour is quite similar to RR (CS-537 Fa21 Mid Q9 makes the same point). It can be slightly less fair in the short run (a job might be skipped several times) but has the same flavour.
\end{itemize}
\end{Sol}

\begin{Qn}{Medium}{Lect 5--6}{CS-537 Spring 2012 Mid Q1}{New job at the front or the back of the RR queue?}
With the round robin (RR) scheduling policy, a question arises when a new job arrives in the system: should we put the job at the front of the RR queue, or the back? Does this subtle difference make a difference, or does RR behave pretty much the same way either way? (Explain)
\end{Qn}
\begin{Sol}
It \emph{does} make a difference, for \emph{response time} of the new job. Front: the new job runs at the next slice boundary (response $\approx 0$ after the current quantum). Back: it waits behind every already-ready job ($n q$). Turnaround of CPU-bound jobs is similar either way; interactivity of arrivals is not. (CS-537 Fa08 Mid assumed ``if X arrives just before Y, RR schedules X before Y'', i.e.\ arrival order at the tail.)
\end{Sol}

\begin{Qn}{Medium-Hard}{Lect 5--6}{CS-537 Spring 2009 Mid Q1}{SJF with an expensive sort and a switch cost}
Assume jobs are submitted in batches. Sorting $N$ jobs takes $\mathrm{Sort}(N)=N\times 10$ time. Switching between jobs takes $\mathrm{Switch}$ time. First, three jobs of lengths 30, 20, and 10. $\mathrm{Sort}(N)=N\times 10$, $\mathrm{Switch}=0$.
\begin{itemize}
  \item[(a)] How long will the system take to compute the schedule for an SJF policy?
  \item[(b)] How long will it take to run the jobs once the schedule has been computed?
  \item[(c)] Including scheduling and switching overhead, what is the average response time for each job?
  \item[(d)] Including scheduling and switching overhead, what is the average turnaround time for each job?
\end{itemize}
Now $\mathrm{Switch}=10$.
\begin{itemize}
  \item[(e)] Average response time?
  \item[(f)] Average turnaround time?
\end{itemize}
A different batch: five jobs of length 10. $\mathrm{Sort}(N)=N\times 10$, $\mathrm{Switch}=0$ again.
\begin{itemize}
  \item[(g)] SJF, including all overheads: average response time?
  \item[(h)] Average turnaround time?
  \item[(i)] What would a smarter policy be for this system?
\end{itemize}
\end{Qn}
\begin{Sol}
Official key (CS-537 Sp09 Mid answers, Q1). SJF order: 10, then 20, then 30.
\begin{itemize}
  \item[(a)] $3\times 10=\mathbf{30}$.
  \item[(b)] $10+20+30=\mathbf{60}$.
  \item[(c)] First job starts at 30, second at 40, third at 60. Avg response $(30+40+60)/3=\mathbf{43.33}$.
  \item[(d)] Completions 40, 60, 90. Avg TAT $190/3=\mathbf{63.33}$.
  \item[(e)] Counting a switch onto each job: starts 40, 60, 90. Avg response $\mathbf{63.33}$.
  \item[(f)] Completions 50, 80, 120. Avg TAT $\mathbf{83.33}$.
  \item[(g)] Sort costs 50. Response times 50, 60, 70, 80, 90. Avg $\mathbf{70}$.
  \item[(h)] Completions 60, 70, 80, 90, 100. Avg $\mathbf{80}$.
  \item[(i)] Any policy that \emph{skips the useless sort} --- FIFO or RR --- because all five lengths are equal, so SJF's order does not matter, and the sort is pure overhead.
\end{itemize}
\end{Sol}

\begin{Qn}{Medium}{Lect 5--6}{CS-537 Spring 2011 Mid, Problem 1}{Draw SJF and RR (with scheduler overhead); TAT and response}
(a)~Draw a picture of Shortest Job First (SJF) scheduling with three jobs A, B, and C, with run times of 5, 10, and 15 time units, respectively. Label the $x$-axis.

(b)~What is the average TURNAROUND TIME for jobs A, B, and C?

(c)~Draw a picture of ROUND-ROBIN scheduling for jobs A, B, and C, which each run for 6 time units, assuming a 2-time-unit time slice; also assume that the scheduler (S) takes 1 time unit to make a scheduling decision. Label the $x$-axis.

(d)~What is the average RESPONSE TIME for round robin for jobs A, B, and C?

(e)~What is the average TURNAROUND TIME for round robin for jobs A, B, and C?
\end{Qn}
\begin{Sol}
Assume arrivals at 0.

(a)~SJF: A $[0,5]$, B $[5,15]$, C $[15,30]$.

(b)~TAT $5$, $15$, $30$; avg $\mathbf{50/3 \approx 16.67}$.

(c)~Pattern \texttt{S A S B S C} repeating, each job-run 2 units, each S 1 unit. Each job needs three slices (6/2).
\begin{center}
\footnotesize
\begin{tabular}{c|ccccccccc}
$t$ & 0--1 & 1--3 & 3--4 & 4--6 & 6--7 & 7--9 & 9--10 & 10--12 & 12--13 \\
 & S & A & S & B & S & C & S & A & S \\
\end{tabular}\\[0.4em]
\begin{tabular}{c|ccccccccc}
$t$ & 13--15 & 15--16 & 16--18 & 18--19 & 19--21 & 21--22 & 22--24 & 24--25 & 25--27 \\
 & B & S & C & S & A & S & B & S & C \\
\end{tabular}
\end{center}
A done at 21, B at 24, C at 27.

(d)~First-scheduled times: A at 1, B at 4, C at 7. Response $1,4,7$; avg $\mathbf{4}$.

(e)~TAT $21,24,27$; avg $\mathbf{24}$.
\end{Sol}

\begin{Qn}{Easy-Medium}{Lect 5--6}{CS-537 Fall 2021 Mid Q9--Q11}{Random every 10\,ms $\approx$ RR; best response; worst TAT}
DOS's scheduler: every 10\,ms a timer interrupt fires and DOS picks a \emph{random} ready process.
\begin{enumerate}
  \item[9.] This approach is most similar to which policy?
  \begin{itemize}
    \item[(a)] FIFO \quad (b) SJF \quad (c) Round Robin \quad (d) STCF \quad (e) Rank-order Scheduling
  \end{itemize}
  \item[10.] Two jobs enter at the same time, each needs 30\,ms, timer 10\,ms. Best-case average \emph{response} time?
  \begin{itemize}
    \item[(a)] 0\,ms \quad (b) 5\,ms \quad (c) 10\,ms \quad (d) 15\,ms \quad (e) none
  \end{itemize}
  \item[11.] Same two jobs. Worst-case average \emph{turnaround} time?
  \begin{itemize}
    \item[(a)] 45\,ms \quad (b) 50\,ms \quad (c) 60\,ms \quad (d) 120\,ms \quad (e) none of the above
  \end{itemize}
\end{enumerate}
\end{Qn}
\begin{Sol}
9.\ \textbf{(c).} A random pick every tick is RR with a randomised (instead of deterministic) rotation. FIFO/SJF/STCF run to completion. Official key: C.

10.\ \textbf{(b) 5\,ms.} Best: job A runs immediately (response 0); at the first tick B runs (response 10); avg 5. Official key: B.

11.\ \textbf{(e).} Worst interleaving is ABABAB in 10\,ms chunks: TAT $50$ and $60$, avg $55$, which is not among (a)--(d). Official key: E.
\end{Sol}

\begin{Qn}{Easy-Medium}{Lect 7}{CS-537 Spring 2018 Mid Problem XI}{What is (and is not) true of MLFQ?}
The Multi-level Feedback Queue (MLFQ) is a fancy scheduler that does lots of things. Fill in A for things that are true about MLFQ, B for things that are not true about MLFQ.
\begin{enumerate}
  \item[51.] MLFQ learns things about running jobs
  \item[52.] MLFQ starves long running jobs
  \item[53.] MLFQ uses different length time slices for jobs
  \item[54.] MLFQ uses round robin
  \item[55.] MLFQ forgets what it has learned about running jobs sometimes
\end{enumerate}
\end{Qn}
\begin{Sol}
\begin{itemize}
  \item[51.] \textbf{A.} Demotion $=$ ``this job looks long''.
  \item[52.] \textbf{B.} Rule~5 (boost after $S$) prevents starvation.
  \item[53.] \textbf{A.} High queues: short slices (e.g.\ 10\,ms); low: long (e.g.\ 100\,ms).
  \item[54.] \textbf{A.} Rule~2: RR within a level.
  \item[55.] \textbf{A.} The boost puts everyone back at the top; history is forgotten (and that is how behaviour change is detected).
\end{itemize}
(Official key, CS-537 Sp18 Mid answers, Problem XI.)
\end{Sol}

\begin{Qn}{Easy-Medium}{Lect 7}{CS-537 Fall 2011 Mid Q7; also CS-537 Fall 2004 Mid1 Q12}{Purpose of MLFQ Rule 5 (boost after $S$)}
In the MLFQ scheduler, there is a rule that states the following: \emph{Rule 5: After some time period $S$, move all the jobs in the system to the top-most queue.} What is the purpose of this rule?
\end{Qn}
\begin{Sol}
Two purposes (Lect~7): (1)~\textbf{anti-starvation} --- long-running jobs that have sunk to the bottom get the CPU again; (2)~\textbf{behaviour change} --- a job that was a CPU hog and is now interactive is re-evaluated at high priority instead of remaining stuck at the bottom.
\end{Sol}

\begin{Qn}{Medium}{Lect 7}{CS-537 Spring 2012 Mid Q3}{Boost period $S$: 10\,s vs.\ 1\,s}
The multi-level feedback queue policy periodically moves all jobs back to the top-most queue. On a particular system, this is usually done every 10 seconds; the subtle difference we examine is that this value has been shortened to 1 second. How does this subtle difference affect the MLFQ scheduler? In general, what is the effect of shortening this value?
\end{Qn}
\begin{Sol}
Smaller $S$ $\Rightarrow$ more frequent boosts. Long jobs spend a larger fraction of time at high priority, so they look more like RR-at-the-top: better for them (less starvation), \emph{worse} for interactive jobs that now compete with recently-boosted hogs. In the limit $S\to 0$, MLFQ collapses toward ordinary RR and you lose the ``approximate SJF'' benefit. Larger $S$ risks starvation and slow reaction to behaviour change. $S$ is a tuning knob (Solaris TS: boost $\sim$ every 1\,s).
\end{Sol}

\begin{Qn}{Medium}{Lect 7}{CS-537 Fall 2021 Mid Q12}{MLFQ that inserts new jobs in a random queue}
DOS uses the MLFQ scheduler. However, it changes some rules. The biggest change: new processes are added to a \emph{random} queue (not the topmost one). What is the biggest effect this will have?
\begin{itemize}
  \item[(a)] Sometimes, short-running jobs won't get serviced quickly, thus decreasing interactivity
  \item[(b)] Over a long period of time, it will be unfair to long-running jobs
  \item[(c)] Jobs will generally finish more quickly
  \item[(d)] Jobs will be able to game the scheduler
  \item[(e)] None of the above
\end{itemize}
\end{Qn}
\begin{Sol}
\textbf{(a).} Rule~3 exists so a short/interactive job can run immediately at the highest queue. Starting at a random (possibly low) queue delays that first schedule --- McFlub's exact mistake on CS-537 Fa13 Mid Q3. Long jobs still round-robin at the bottom; no new gaming attack. Official key: A.
\end{Sol}

\begin{Qn}{Easy-Medium}{Lect 3}{CS-537 Fall 2011 Mid Q4}{Write \texttt{fork}/\texttt{exec}/\texttt{wait} to run \texttt{hello}}
A typical OS provides some APIs to create processes. In Unix-based systems, \texttt{fork()}, \texttt{exec()}, and \texttt{wait()} are used. Write some code that uses these system calls to launch a new child process, have the child execute a program named \texttt{"hello"} (with no arguments), and have the parent wait for the child to complete.
\end{Qn}
\begin{Sol}
\begin{lstlisting}[numbers=none]
int rc = fork();
if (rc < 0) { perror("fork"); exit(1); }
if (rc == 0) {
    char *argv[] = { "hello", NULL };
    execvp(argv[0], argv);
    perror("exec");   /* only on failure */
    exit(1);
} else {
    wait(NULL);
}
\end{lstlisting}
The \texttt{argv} array must be \texttt{NULL}-terminated. Do not ``wait'' by spinning on a user-space flag (the grader paper's student did that); \texttt{wait()} is the blocking syscall.
\end{Sol}

\clearpage

\section{Lectures 8--12 (Pre-Minor): Early Memory Virtualization (B\&B, Seg., Free Space \& Paging)}

\subsubsection*{Additional Challenging \& Analytical Problems}

\begin{Qn}{Hard}{Lect 9}{CS-537 Spring 2016 Mid, Q1}{Base-and-bounds blues: fill in the missing VA / PA / registers}
A system uses simple base/bounds. The bounds register holds the
\textbf{size} of the valid region (legal iff $0\le\mathrm{VA}<\mathrm{bounds}$).
All numbers are decimal. In each trace, fill in every missing virtual
address, physical address, base, or bounds. If a value is not unique,
give the tightest range.

\textbf{Trace 1.}
$0\to 1000$,\quad $100\to 1100$,\quad $1999\to 2999$,\quad
$2000\to$ fault.

\textbf{Trace 2.}
$0\to 1000$,\quad $100\to 1100$,\quad $1999\to 2999$,\quad
$2000\to 3000$.

\textbf{Trace 3.}
$0\to 6050$,\quad $100\to 6150$,\quad $2000\to 8050$,\quad
$2001\to$ fault.

\textbf{Trace 4.}
$9000\to 10001$,\quad $100\to 1101$,\quad $2000\to 3001$,\quad
$2001\to 3002$.
\end{Qn}
\begin{Sol}
Always $\mathrm{PA}=\mathrm{VA}+\mathrm{base}$ when the access is valid.

\textbf{Trace 1.} $0\to 1000$ forces $\mathrm{base}=1000$. The other
two valid pairs confirm it. VA $2000$ faults and $1999$ does not, so
$\mathrm{bounds}=2000$.

\textbf{Trace 2.} Same base $1000$, but VA $2000$ is now valid, so
$\mathrm{bounds}\ge 2001$. The trace gives no failing VA, so the
bounds cannot be pinned further.

\textbf{Trace 3.} $0\to 6050$ forces $\mathrm{base}=6050$. VA $2001$
faults and $2000$ does not, so $\mathrm{bounds}=2001$.

\textbf{Trace 4.} $100\to 1101$ forces $\mathrm{base}=1001$ (also
$9000+1001=10001$, $2000+1001=3001$, $2001+1001=3002$). Every listed
VA, including $9000$, is valid, so $\mathrm{bounds}\ge 9001$.
\end{Sol}

\begin{Qn}{Hard}{Lect 10}{CS-537 Spring 2016 Mid, Q9}{Set two-segment base/limit registers from a VA/PA trace}
Segmentation splits the virtual address space into two segments. The
top bit of the virtual address selects the segment.

Segment 0 is code/heap and grows toward higher addresses.
Segment 1 is a stack and grows backward. Follow the OSTEP convention:
the stack \textbf{base} points to the physical address one past the last
byte of the stack. In both segments the bounds (limit) register holds
the \textbf{size} (number of valid bytes).

Set the four registers so that every VA in the given trace translates
correctly and is valid, and so that every VA \emph{not} in the trace
generates a segmentation fault.

\textbf{Trace 1.} $16$-byte ($4$-bit) virtual address space.
Valid VAs: $0,1,2,3,15,14,13$.
VA $1$ translates to PA $101$.
VA $13$ translates to PA $998$.

\textbf{Trace 2.} $64$-byte ($6$-bit) virtual address space.
Valid VAs: $0,1,63$.
VA $1$ translates to PA $1001$.
VA $63$ translates to PA $899$.

\textbf{Trace 3.} $8$-byte ($3$-bit) virtual address space.
Valid VAs: $0,1,2,3$.
VA $3$ translates to PA $100$.
\end{Qn}
\begin{Sol}
Grow-positive: $\mathrm{PA}=\mathrm{base}_0+\mathrm{offset}$, valid iff
$0\le\mathrm{offset}<\mathrm{limit}_0$.
Grow-negative, VAS size $N$:
$\mathrm{PA}=\mathrm{base}_1-(N-\mathrm{VA})$, valid iff
$N-\mathrm{limit}_1\le\mathrm{VA}\le N-1$.

\textbf{Trace 1.} $N=16$, so SEG0 is VAs $0..7$ and SEG1 is $8..15$.
The valid SEG0 addresses are $0,1,2,3$, so $\mathrm{limit}_0=4$.
VA $1\to 101$ gives $\mathrm{base}_0=100$.
The valid SEG1 addresses are $13,14,15$ ($3$ bytes from the top), so
$\mathrm{limit}_1=3$.
VA $13\to 998$ gives $\mathrm{base}_1-(16-13)=998$, hence
$\mathrm{base}_1=1001$.

\textbf{Trace 2.} $N=64$, SEG0 is $0..31$, SEG1 is $32..63$.
Valid SEG0: $0,1$ so $\mathrm{limit}_0=2$, $\mathrm{base}_0=1000$.
Valid SEG1: only $63$, so $\mathrm{limit}_1=1$.
VA $63\to 899$ gives $\mathrm{base}_1-1=899$, hence
$\mathrm{base}_1=900$.

\textbf{Trace 3.} $N=8$, SEG0 is $0..3$, SEG1 is $4..7$.
All four valid VAs lie in SEG0, so $\mathrm{limit}_0=4$ and
$\mathrm{base}_0=100-3=97$. Segment 1 is never used:
$\mathrm{limit}_1=0$ and $\mathrm{base}_1$ is unconstrained
(unknown / unused).
\end{Sol}

\begin{Qn}{Hard}{Lect 10}{CS-537 Fall 2013 Mid, Q8}{McFlub's ``easy'' backward-stack translations are all faults}
McFlub: ``With segmentation, assume a $1$\,KB virtual address space.
Two segments: a code/heap segment that grows positive, and a
backwards-growing stack. The top bit of the VA selects the segment.

\begin{lstlisting}[numbers=none]
Segment 0 base (grows positive) : 0x00000000 (decimal 0)
Segment 0 limit                 : 1
Segment 1 base (grows negative) : 0x00000064 (decimal 100)
Segment 1 limit                 : 1
\end{lstlisting}

Because the base of segment 0 is $0$, translations are super easy:
\begin{lstlisting}[numbers=none]
VA 1: 0x00000065 (decimal 101) --> 0x65 (101)
VA 2: 0x00000169 (decimal 361) --> 0x169 (361)
VA 3: 0x000002ad (decimal 685) --> 0x2ad (685)
\end{lstlisting}
See, easy!''

Is McFlub right? If not, correct the translations.
\end{Qn}
\begin{Sol}
Each limit is \textbf{1 byte}. In a $1$\,KB ($1024$-byte) VAS the top
bit splits VAs $0..511$ (SEG0) from $512..1023$ (SEG1).

SEG0 is valid only for offset $0$ (VA $0$).
SEG1 is valid only for the single top-most byte, VA $1023$.

All three of McFlub's VAs are therefore \textbf{segmentation faults}:
\begin{itemize}
  \item $101$ and $361$ are in SEG0 with offset $\ge 1$;
  \item $685=0x2\mathrm{ad}$ has the top bit set (it is in SEG1) and
        is nowhere near VA $1023$.
\end{itemize}
The only legal VAs in this system are $0$ and $1023$. McFlub treated
the limit as if it were the whole address space.
\end{Sol}

\begin{Qn}{Hard}{Lect 10}{CS-537 Spring 2009 Mid, Q2}{Two-segment hardware: registers, then a looping pointer that crashes}
Segmentation uses several base/bounds pairs to virtualize memory.
\begin{enumerate}
  \item[(a)] What is a base register for?
  \item[(b)] What type of address is in the base register: physical, or
        virtual?
  \item[(c)] What is a bounds register for?
  \item[(d)] Why do we need more than one base/bounds pair?
  \item[(e)] How many segments should the hardware support, and why?
\end{enumerate}

Now the hardware supports two segments. Address spaces are $1$\,KB;
physical memory is $16$\,KB. Segment 0 has base $1$\,KB and bounds
(size) $300$ bytes, and grows upward. Segment 1 has base $5$\,KB and
bounds $300$, and grows downward (negative direction).

\begin{lstlisting}[numbers=none]
void *ptr = 20;
while (ptr <= 1024) {
    int x = (int *) *ptr;  /* LINE 1: read *ptr */
    ptr = ptr + 20;        /* LINE 2 */
}
\end{lstlisting}
\begin{enumerate}
  \item[(f)] What virtual addresses are generated at LINE 1 by
        dereferencing \texttt{ptr}, assuming the program runs to
        completion (ignore instruction fetches and the store into
        \texttt{x})?
  \item[(g)] How long will this program run before crashing (segmentation
        violation)?
  \item[(h)] What physical addresses are generated by those
        dereferences before the crash?
  \item[(i)] What legal physical addresses \emph{could} have been
        generated if the programmer had skipped the unmapped middle of
        the address space?
\end{enumerate}
\end{Qn}
\begin{Sol}
(a) The base holds the \textbf{physical} address at which a segment was
placed, so hardware can relocate it by
$\mathrm{PA}=\mathrm{base}+\mathrm{offset}$.

(b) \textbf{Physical}.

(c) The bounds is the segment \textbf{size} (or, in the other encoding,
the ending physical address). Hardware compares the offset against it
and raises a protection fault on overflow --- this is isolation.

(d) A single pair relocates the \emph{entire} address space as one
contiguous blob, so the unused virtual hole between heap and stack
still consumes physical memory. Several pairs let the OS place code,
heap, and stack independently (sparse VAS, independent growth,
per-segment protection/sharing).

(e) At least \textbf{two} (code+heap vs.\ backward stack) to support a
sparse space. \textbf{Three} is nicer: a separate read-execute code
segment can be shared. More than three is possible but not required.

(f) LINE 1 generates VA $20,40,60,\ldots,1020$ (start at $20$, step
$20$; the predicate is \texttt{ptr <= 1024}).

(g) The top bit of a $1$\,KB VAS splits SEG0 ($0..511$) from SEG1
($512..1023$). SEG0 size $300$ makes VAs $0..299$ legal. The loads
$20,40,\ldots,280$ are the $14$ legal SEG0 references; the next load,
VA $300$, exceeds the SEG0 bounds and \textbf{faults}. (SEG1 is never
reached.)

(h) Add the SEG0 base $1024$: PAs $1044,1064,\ldots,1304$
(i.e.\ $1024+20$ through $1024+280$).

(i) If the loop skipped the unmapped hole $300..723$, the last $300$
bytes of the VAS ($724..1023$) are a legal grow-negative stack.
Translation:
$\mathrm{PA}=5\,\mathrm{KB}-(1024-\mathrm{VA})=5120-(1024-\mathrm{VA})$.
VA $1020$ maps to $5120-4=5116$; the first legal stack VA $724$ maps
to $5120-300=4820$. Equivalently, every physical address in
$4820..5119$ is reachable from a careful loop. (The official key
quotes the slightly tighter example range $740\to 4836$ through
$1020\to 5116$, which is the subset obtained by keeping the step-$20$
stride.)
\end{Sol}

\begin{Qn}{Medium}{Lect 10}{CS-537 Spring 2018 Mid, Problem XIV (later appeared as MTL458 Minor Q1)}{Does segmentation actually beat base-and-bounds?}
Segmentation is a generalization of base-and-bounds. Which possible
advantages does segmentation have compared to base-and-bounds?

Fill in A if the statement is true about segmentation \emph{and} (as a
result) segmentation has a clear advantage over base-and-bounds; fill
in B otherwise.
\begin{enumerate}
  \item[66.] Faster translation
  \item[67.] Less physical memory waste
  \item[68.] Better sharing of code in memory
  \item[69.] More hardware support needed to implement it
  \item[70.] More OS issues to handle, such as compaction
\end{enumerate}
\end{Qn}
\begin{Sol}
Segmentation is several base/bounds pairs per process instead of one.
\begin{enumerate}
  \item[66.] \textbf{B.} Translation is not faster (still an add and a
        compare, plus a segment lookup).
  \item[67.] \textbf{A.} The virtual hole between heap and stack need
        not be backed by physical memory, so less waste (no giant
        internal hole).
  \item[68.] \textbf{A.} A read-execute code segment can be shared by
        pointing two processes' code-base registers at the same
        physical range (and copy-on-write follows the same idea).
  \item[69.] \textbf{B.} More hardware \emph{is} needed, but that is
        not an advantage.
  \item[70.] \textbf{B.} Compaction / external fragmentation \emph{are}
        extra OS issues, but that is not an advantage.
\end{enumerate}
\end{Sol}

\begin{Qn}{Medium-Hard}{Lect 10}{CS-537 Spring 2012 Mid, Q12}{Hardware and OS disagree by one byte on the stack base}
A system uses two-segment virtual memory. Segment 0 grows in the
increasing direction; segment 1 grows backwards. Unfortunately there is
confusion over the interpretation of the base register of segment 1:
the \textbf{hardware} thinks it should point to the physical address
\emph{one beyond} the bottom of the backwards-growing segment, while
the \textbf{OS} has been changed to assume it points to the
\emph{last byte} of that segment.

Describe what would happen when running processes on this
subtly-changed virtual memory system.
\end{Qn}
\begin{Sol}
Every stack translation is off by exactly one byte. If the OS writes
$B$ into the stack-base register, hardware computes
$\mathrm{PA}=B-(N-\mathrm{VA})$ as if $B$ were one-past-the-end, so it
actually refers to physical address $B-1-(N-\mathrm{VA})$ relative to
what the OS intended --- i.e.\ one byte \emph{lower} than the OS's
picture of the stack.

Consequences:
\begin{itemize}
  \item Loads/stores of stack slots hit the neighbour byte: saved
        return addresses, locals, and arguments are corrupted.
  \item The byte just below the intended stack (often belonging to
        another segment, another process, or unmapped physical memory)
        is read and written.
  \item The top-most intended stack byte is never touched by hardware.
  \item Bounds checks, which use the size and the same base convention,
        are also skewed, so a legal push may fault or an out-of-range
        access may sneak through.
\end{itemize}
The code/heap segment (grow-positive, no disagreement) continues to
translate correctly. The bug is a pure stack-base off-by-one between
software and MMU.
\end{Sol}

\begin{Qn}{Medium-Hard}{Lect 10}{CS-537 Spring 2011 Mid, Problem 4}{Draw the two-segment layout, then translate two VAs}
Hardware uses segmentation and splits the address space in two with the
\textbf{top bit} of the virtual address. Each segment is relocated
independently. Virtual address space size $16$ bytes; physical memory
$64$ bytes. Segment 1 grows backwards.

In pictures below, \texttt{0} marks a byte of segment 0, \texttt{1} a
byte of segment 1, and \texttt{F} free. (Spaces every 8 bytes are
for readability only.)
\begin{enumerate}
  \item[(a)] Draw physical memory for
        $\mathrm{seg0}$: base $12$, limit $6$;
        $\mathrm{seg1}$: base $10$, limit $3$.
  \item[(b)] Which physical byte is accessed by a load of virtual
        address $4$? Mark illegal with X.
  \item[(c)] Draw physical memory for
        $\mathrm{seg0}$: base $40$, limit $4$;
        $\mathrm{seg1}$: base $50$, limit $4$.
  \item[(d)] Load of virtual address $14$?
  \item[(e)] Load of virtual address $4$?
\end{enumerate}
\end{Qn}
\begin{Sol}
$N=16$, so SEG0 is VAs $0..7$ and SEG1 is $8..15$.
Grow-positive: PA $12..17$ hold SEG0 (limit $6$).
Grow-negative with base $10$ (one past the last stack byte) and limit
$3$: the three stack bytes are PA $7,8,9$, with VA $15\mapsto 9$,
VA $14\mapsto 8$, VA $13\mapsto 7$.

(a) Bytes $0..6$ free; $7..9$ are \texttt{1}; $10..11$ free;
$12..17$ are \texttt{0}; $18..63$ free.

(b) VA $4$ is SEG0, offset $4<6$, so PA $12+4=16$. Circle byte $16$.

(c) SEG0 occupies PA $40..43$. SEG1 with base $50$ and limit $4$
occupies PA $46..49$ (VA $15\mapsto 49$, \ldots, VA $12\mapsto 46$).

(d) VA $14$ is SEG1 and legal: PA $50-(16-14)=48$.

(e) VA $4$ is SEG0, offset $4\not< 4$. \textbf{Illegal} (X).
\end{Sol}

\begin{Qn}{Medium-Hard}{Lect 9}{CS-537 Fall 2013 Mid, Q5}{McFlub faults a legal base-and-bounds address}
McFlub: ``Base/bounds virtual memory is really easy. The base points to
the physical location where an address space is relocated; the bounds
tells us how big it can be.

\begin{lstlisting}[numbers=none]
Base   : 0x1000
Bounds : 0x10
\end{lstlisting}
Virtual references $0\mathrm{x}0$, $0\mathrm{x}4$, $0\mathrm{x}8$,
$0\mathrm{x}c$ then translate to $0\mathrm{x}1000$,
$0\mathrm{x}1004$, $0\mathrm{x}1008$, and FAULT (because
$0\mathrm{x}c$ is out of bounds). Make sense?''

Is McFlub right? If not, correct the translations.
\end{Qn}
\begin{Sol}
Bounds \texttt{0x10}$=16$ is a \textbf{size}: every VA $v$ with
$0\le v<16$ is legal. \texttt{0xc}$=12$ is strictly less than $16$,
so it is legal. McFlub is wrong.

The four translations are
\texttt{0x1000}, \texttt{0x1004}, \texttt{0x1008}, and
\texttt{0x100c} (no fault). The hardware does
$\mathrm{PA}=\mathrm{VA}+\mathrm{base}$ after a
$\mathrm{VA}<\mathrm{bounds}$ check.
\end{Sol}

\begin{Qn}{Hard}{Lect 12}{CS-537 Fall 2013 Mid, Q7}{McFlub's linear-page-table translation of \texttt{0x33ef}}
McFlub dumps this entire linear page table (entry $0$ first) and says
each page is $1$\,KB. High bit of a PTE is valid; the rest is the PFN.
\begin{lstlisting}[numbers=none]
0x8000002a
0x00000000
0x00000000
0x00000000
0x8000003e
0x00000000
0x80000017
0x00000000
0x00000000
0x8000000e
0x00000000
0x00000000
0x8000001f
0x00000000
0x8000001a
0x80000001
\end{lstlisting}
``This is the entire page table, so the VAS is $32$\,KB. Now translate
VA \texttt{0x33ef}: as you can quickly see, this is an INVALID ACCESS.
Stupid program!''

Is McFlub right? If not, show how the address really translates.
\end{Qn}
\begin{Sol}
There are $16$ PTEs, not $32$. With $1$\,KB pages the VAS is
$16$\,KB ($14$-bit VAs), not $32$\,KB.

\texttt{0x33ef}$=13{,}295<16{,}384$, so it \emph{is} a legal-width
virtual address. Split a $14$-bit VA into a $4$-bit VPN and a
$10$-bit offset:
\[
\texttt{0x33ef}=1100\,1111101111_2
\quad\Rightarrow\quad
\mathrm{VPN}=1100_2=12,\quad
\mathrm{offset}=1111101111_2=\texttt{0x3ef}.
\]
PTE $[12]$ is \texttt{0x8000001f}: valid, $\mathrm{PFN}=\texttt{0x1f}$.
\[
\mathrm{PA}=(\texttt{0x1f}\ll 10)\mid\texttt{0x3ef}
=\texttt{0x7c00}+\texttt{0x3ef}=\texttt{0x7fef}.
\]
McFlub both mis-sized the address space and then used the wrong VPN,
so he incorrectly declared a perfectly valid translation a fault.
\end{Sol}

\begin{Qn}{Hard}{Lect 12}{CS-537 Spring 2016 Mid, Q8}{Reverse-engineer a linear page table from a VA/PA trace}
Linear page table, format ``valid bit then PFN''. Parameters:
virtual address space $32$\,KB, page size $4$\,KB, physical memory
$64$\,KB. Reconstruct every PTE you can; mark the rest UNKNOWN.

\begin{lstlisting}[numbers=none]
VA 0x1063 --> PA 0x2063
VA 0x67b4 --> PA 0x67b4
VA 0x584a --> PA 0xe84a
VA 0x4dfe --> Invalid
VA 0x388a --> Invalid
VA 0x1c6b --> PA 0x2c6b
VA 0x50a9 --> PA 0xe0a9
VA 0x0bc6 --> Invalid
VA 0x2a9f --> PA 0x9a9f
VA 0x742b --> Invalid
VA 0x4b5e --> Invalid
VA 0x5597 --> PA 0xe597
\end{lstlisting}
\end{Qn}
\begin{Sol}
$32$\,KB VAS $\Rightarrow$ $15$-bit VA. $4$\,KB page $\Rightarrow$
$12$-bit offset, so a $3$-bit VPN and $8$ PTEs (VPN $0..7$).
$64$\,KB physical $\Rightarrow$ $4$-bit PFN. The VPN is the high hex
digit of these $15$-bit addresses; the offset is preserved by a
correct translation, which is a useful sanity check.

\begin{center}
\begin{tabular}{c|c|c|l}
VPN & Valid? & PFN & Evidence \\
\hline
0 & 0 & --- & VA \texttt{0x0bc6} invalid \\
1 & 1 & $2$ & \texttt{0x1063}$\to$\texttt{0x2063},
              \texttt{0x1c6b}$\to$\texttt{0x2c6b} \\
2 & 1 & $9$ & \texttt{0x2a9f}$\to$\texttt{0x9a9f} \\
3 & 0 & --- & VA \texttt{0x388a} invalid \\
4 & 0 & --- & VA \texttt{0x4dfe}, \texttt{0x4b5e} invalid \\
5 & 1 & \texttt{0xe} & \texttt{0x584a}$\to$\texttt{0xe84a} (and
                      \texttt{0x50a9}, \texttt{0x5597}) \\
6 & 1 & $6$ & \texttt{0x67b4}$\to$\texttt{0x67b4} \\
7 & 0 & --- & VA \texttt{0x742b} invalid \\
\end{tabular}
\end{center}

Every invalid VPN in the trace is enough to set that PTE's valid bit
to $0$ (the leftover bits of an invalid PTE are unused). No PTE is
UNKNOWN.
\end{Sol}

\begin{Qn}{Medium-Hard}{Lect 12}{CS-537 Fall 2008 Mid, Q1}{VPN/offset splits for two page sizes, then \texttt{find\_pte}}
A virtual address space of size $64$\,KB, paging, page size $8$\,KB.
\begin{enumerate}
  \item[(a)] How many bits in a virtual address?
  \item[(b)] How many bits in the VPN?
  \item[(c)] How many bits in the offset?
  \item[(d)] How many entries in a linear page table?
\end{enumerate}
Now the same $64$\,KB VAS, but page size $4$ \emph{bytes} (not KB).
Repeat (a)--(d).

Finally, the OS tracks the linear page table by its kernel-physical
base \texttt{pt\_base}. Complete
\begin{lstlisting}[numbers=none]
struct pte *find_pte(void *pt_base, int VPN)
{
    struct pte *p = ???;
    return p;
}
\end{lstlisting}
\end{Qn}
\begin{Sol}
$64$\,KB $=2^{16}$, so a virtual address is always $16$ bits.

\textbf{8\,KB pages.} $8$\,KB $=2^{13}$, so offset $=13$, VPN $=3$,
and a linear table has $2^3=8$ entries.

\textbf{4-byte pages.} Offset $=2$, VPN $=14$, linear table has
$2^{14}=16{,}384$ entries. Tiny pages explode the page table --- the
classic page-size trade-off of Lect 12 (less internal fragmentation,
far more PT overhead).

\textbf{\texttt{find\_pte}.} Index the array of PTEs:
\begin{lstlisting}[numbers=none]
struct pte *p = (struct pte *)pt_base + VPN;
\end{lstlisting}
Equivalently, if \texttt{pt\_base} is kept as \texttt{void *},
\begin{lstlisting}[numbers=none]
struct pte *p = pt_base + (VPN * sizeof(struct pte));
\end{lstlisting}
(the second form relies on byte-pointer arithmetic on
\texttt{void *}, which GCC allows). Hardware then reads
\texttt{*p}, checks Valid / Protection / Present / Dirty / Reference,
and forms $\mathrm{PA}=\mathrm{PFN}\times\mathrm{pagesize}+\mathrm{offset}$.
\end{Sol}

\begin{Qn}{Medium-Hard}{Lect 12}{CS-537 Fall 2004 Final, Q9}{The linear page tables are the wrong size --- fix them}
Each linear page table below is broken. Fix the stated quantity.
\begin{enumerate}
  \item[(a)] $32$-bit virtual address, $4$\,KB page, $32$-bit physical
        address, but the linear page table has $2^{16}$ entries. How
        many entries should it have?
  \item[(b)] Same VA/page/PA as (a), but each PTE has space for $2$
        bytes. How big should each entry be?
  \item[(c)] $32$-bit VA, $1$\,MB page, $16$-bit PA, but the linear
        page table has $2^{16}$ entries. How many entries should it
        have?
  \item[(d)] $32$-bit VA, $1$\,MB page, $22$-bit PA, but each PTE has
        space for $2$ bytes. How big should each entry be?
  \item[(e)] $32$-bit VA, $1$\,MB page, $20$-bit PA, but the linear
        page table has $2^{16}$ entries. How many entries should it
        have?
\end{enumerate}
\end{Qn}
\begin{Sol}
A linear table has one PTE per virtual page:
$\#\mathrm{PTEs}=2^{(\mathrm{VA\ bits}-\lg(\mathrm{page\ size}))}$.
A PTE must be wide enough for the PFN plus flags
(Valid / Protection / Present / Dirty / Reference).
$\mathrm{PFN\ bits}=\mathrm{PA\ bits}-\lg(\mathrm{page\ size})$.

(a) Offset $12$ bits $\Rightarrow$ VPN $20$ bits $\Rightarrow$
\textbf{$2^{20}$ entries} (not $2^{16}$).

(b) PFN is $32-12=20$ bits, plus at least a valid bit. Two bytes
($16$ bits) cannot hold a $20$-bit PFN. A typical fix is a
\textbf{$4$-byte} PTE (as in OSTEP / x86-32).

(c) $1$\,MB page $\Rightarrow$ $20$-bit offset, VPN $12$ bits
$\Rightarrow$ \textbf{$2^{12}=4096$ entries}.

(d) PFN is $22-20=2$ bits, plus a few flag bits --- well under a
byte. Two bytes is wasteful; \textbf{one byte} is already enough.
(Leaving the PTE at $2$ bytes still works, it just wastes PT space.)

(e) Same VPN width as (c): \textbf{$2^{12}$ entries}. The $20$-bit
PA equals the page size, i.e.\ there is only one physical frame, so
the PFN field is empty --- but the table is still indexed by VPN.
\end{Sol}

\subsubsection*{Conceptual Drills, MCQs \& Further Practice}

\begin{Qn}{Easy-Medium}{Lect 8}{CS-537 Fall 2004 Mid2, Q1}{Stack versus heap in a single-threaded process}
In a traditional single-threaded process we have both a stack and a heap.
\begin{enumerate}
  \item What is the stack used for?
  \item What is the heap used for?
  \item If you could only have one of the two, which would you pick, and
        why?
\end{enumerate}
\end{Qn}
\begin{Sol}
The virtual address space is the process's view of memory (code, heap,
stack), translated by the MMU.

The \textbf{stack} holds procedure activation records: return
addresses, saved registers, and local (automatic) variables. It grows
and shrinks at one end (typically toward lower VAs) in LIFO order, so
allocation is a pointer bump.

The \textbf{heap} holds dynamically allocated objects
(\texttt{malloc}/\texttt{new}) whose lifetime is not nested. It grows
toward higher VAs; the allocator must search a free list.

If forced to pick one: the \textbf{stack} is enough to run ordinary
code with locals and calls (and one can even overlay a bump-pointer
``heap'' on it). A heap without a stack cannot implement calls. In
practice both exist because their allocation patterns differ.
\end{Sol}

\begin{Qn}{Easy-Medium}{Lect 8}{Mandar ISI 2011 End-sem, Q3(a)}{Logical vs.\ physical addresses, and compile- / load- / run-time binding}
In the context of address-binding schemes, what is a (i) logical
address, (ii) physical address? Explain in 1--2 lines each whether the
two addresses are identical or different for compile-time, load-time,
and run-time address binding.
\end{Qn}
\begin{Sol}
A \textbf{logical} (virtual) address is the address generated by the
CPU / compiler, i.e.\ the process's view of the address space. A
\textbf{physical} address is the address that actually goes out on the
memory bus after translation.

\begin{itemize}
  \item \textbf{Compile-time binding:} the compiler emits absolute
        physical addresses. Logical $=$ physical. The program must be
        loaded at the compiled address.
  \item \textbf{Load-time binding:} the loader relocates the image as
        it is placed in memory. Logical $=$ physical \emph{after}
        loading, but the two differ until the load address is known.
  \item \textbf{Run-time binding} (dynamic relocation, segmentation,
        paging): the CPU always emits logical addresses; the MMU
        translates them on every reference. Logical $\ne$ physical in
        general. This is the scheme Lectures 9--12 implement.
\end{itemize}
\end{Sol}

\begin{Qn}{Easy-Medium}{Lect 9}{CS-537 Spring 2018 Mid, Problem XII}{Which physical addresses can this base-and-bounds process touch?}
Dynamic relocation (base-and-bounds). System:
$1$\,KB virtual address space, base $=10000$, bounds $=100$.

Which of the following \emph{physical} memory locations can be legally
accessed by the running program? A $=$ legally accessible, B $=$ not.
\begin{enumerate}
  \item[56.] $0$
  \item[57.] $1000$
  \item[58.] $10000$
  \item[59.] $10050$
  \item[60.] $10100$
\end{enumerate}
\end{Qn}
\begin{Sol}
Legal VAs are $0..99$ (bounds is a size). They map to physical
$10000..10099$. Thus:
56 B,\quad 57 B,\quad 58 A,\quad 59 A,\quad 60 B
($10100$ is the first illegal physical address).
\end{Sol}

\begin{Qn}{Easy-Medium}{Lect 9}{CS-537 Spring 2018 Mid, Problem XIII}{Which virtual addresses are in bounds?}
Same set-up ($1$\,KB VAS, base $=10000$, bounds $=100$). Which
\emph{virtual} addresses are valid? A $=$ valid, B $=$ not.
\begin{enumerate}
  \item[61.] $0$
  \item[62.] $1000$
  \item[63.] $10000$
  \item[64.] $10050$
  \item[65.] $10100$
\end{enumerate}
\end{Qn}
\begin{Sol}
Valid iff $0\le\mathrm{VA}<100$, independent of the base.
61 A,\quad 62 B,\quad 63 B,\quad 64 B,\quad 65 B.

(The $1$\,KB VAS would have allowed VAs up to $1023$, but the
\emph{bounds register} further restricts the process to the first
$100$ bytes. Base never makes a virtual address legal or illegal.)
\end{Sol}

\begin{Qn}{Medium}{Lect 9}{CS-537 Fall 2021 Mid, Q4--Q8}{BeOS: base-and-bounds protection, bounds-on-switch, static relocation}
BeOS originally used base-and-bounds.
\begin{enumerate}
  \item[4.] Which one is true about this approach?
        (a) It translates more slowly than most other approaches\quad
        (b) It prevents internal fragmentation\quad
        (c) It enables sharing of code between processes\quad
        (d) It provides protection between processes\quad
        (e) None of the above (all are false)
  \item[5.] VAS $64$ bytes, physical $512$ bytes, base $=10$,
        bounds $=20$. The process issues VAs $15,6,30,21,50$. How
        many of these are valid?
        (a) $1$\quad (b) $2$\quad (c) $3$\quad (d) $4$\quad (e) $5$
  \item[6.] A later BeOS has a $1$\,KB virtual address space on a
        machine with $16$\,KB of physical memory. What is the largest
        virtual address a process can \emph{generate}?
        (a) $1000$\quad (b) $1023$\quad (c) $1024$\quad
        (d) It depends on the base register\quad
        (e) It depends on the bounds register
  \item[7.] On a context switch BeOS updates the base but leaves the
        bounds register at whatever random value was first placed
        there. What could happen? (choose one)
        (a) A process could access another process's memory\quad
        (b) A process could fault on what should be valid memory\quad
        (c) A process could access physical memory assigned to nobody\quad
        (d) A process could be prevented from accessing valid code\quad
        (e) All of the above
  \item[8.] An even earlier BeOS ignored the MMU and \emph{statically
        rewrote} every address in the loaded image (e.g.\
        \texttt{load 10, R1} becomes \texttt{load 110, R1} if the
        image is placed at $100$). This approach:
        (a) Provides strong process isolation\quad
        (b) Makes programs run slowly\quad
        (c) Allows multiple processes to be active in the system\quad
        (d) Affects the CPU scheduling policy you can implement\quad
        (e) Helps realize a sparse virtual address space
\end{enumerate}
\end{Qn}
\begin{Sol}
4.\ \textbf{(d).} The bounds check is protection. Translation is
\emph{fast} (one add), so not (a). Internal fragmentation of the
heap--stack hole is a famous \emph{con}, so not (b). Sharing a code
segment needs more than one base/bounds pair, so not (c).

5.\ \textbf{(b) $2$.} Bounds $20$ makes VAs $0..19$ legal. Of
$15,6,30,21,50$ only $15$ and $6$ pass.

6.\ \textbf{(b) $1023$.} A $1$\,KB VAS can generate VAs $0..1023$.
Whether a generated VA is \emph{legal} depends on bounds; the
\emph{largest generable} VA does not.

7.\ \textbf{(e).} A too-large leftover bounds lets the process walk
off its slot onto a neighbour (or onto unassigned RAM). A too-small
leftover bounds faults legal VAs and can hide valid code. All four
symptoms follow. This is exactly why the OS must
save/restore base \emph{and} bounds in the PCB on every context
switch (Lect 9).

8.\ \textbf{(c).} Static relocation lets the loader place several
images at different physical origins. It does \emph{not} isolate them
(no bounds hardware), is not particularly slow (the rewrite is
once-per-load), does not change the scheduler, and does not create a
sparse VAS.
\end{Sol}

\begin{Qn}{Medium}{Lect 9}{CS-537 Fall 2022 Mid, Q1 (items 1--5, 7--8)}{x537 base-and-bounds: register count, \texttt{load}, and a forgotten bounds}
The x537 CPU uses base-and-bounds. VAS size $1$\,KB ($1024$ bytes).
\begin{enumerate}
  \item[1.] With one CPU, how many base registers are in the system
        (total)?
        (a) $1$\quad (b) $2$\quad (c) $4$\quad (d) $8$\quad
        (e) None of the above
  \item[2.] Base $=1000$, bounds $=100$. How many memory accesses does
        the instruction \texttt{load 10, r1} generate when it is
        fetched and then executed?
        (a) $0$\quad (b) $1$\quad (c) $2$\quad (d) $3$\quad
        (e) None of the above
  \item[3.] With base $=1000$ and bounds $=100$, to what physical
        address does virtual address $10$ translate?
        (a) $100$\quad (b) $1000$\quad (c) $1010$\quad (d) $10$\quad
        (e) None of the above
  \item[4.] With bounds $=100$, how many different legal virtual
        addresses can a process access?
        (a) $64$\quad (b) $99$\quad (c) $100$\quad (d) $1000$\quad
        (e) None of the above
  \item[5.] Process table:
        A (base $100$, bounds $10$),
        B (base $1000$, bounds $20$),
        C (base $500$, bounds $50$).
        After a switch from A to B, which of the following is a
        physical address B might legally refer to?
        (a) $0$\quad (b) $950$\quad (c) $500$\quad (d) $20$\quad
        (e) None of the above
  \item[7.] On the A$\to$B switch, xOS updates base from $100$ to
        $1000$ but forgets to update bounds (it stays $10$). What
        happens?
        (a) Everything works\quad
        (b) Process B might fault unexpectedly\quad
        (c) Process B might access memory it should not\quad
        (d) Process A might be mad\quad
        (e) None of the above
  \item[8.] The system is reconfigured with four CPUs. How many base
        registers in total?
        (a) $1$\quad (b) $2$\quad (c) $4$\quad (d) $8$\quad
        (e) None of the above
\end{enumerate}
\end{Qn}
\begin{Sol}
1.\ \textbf{(a) $1$.} Base (and bounds) live in the MMU of the CPU
that is running the process --- one pair per CPU.

2.\ \textbf{(c) $2$.} Fetch the instruction, then fetch the data at
VA $10$. Both go through the MMU. (No extra PTE walk: this is
base-and-bounds, not paging.)

3.\ \textbf{(c) $1010$.} $\mathrm{PA}=10+1000=1010$, and $10<100$.

4.\ \textbf{(c) $100$.} Legal VAs are $0..99$.

5.\ \textbf{(e).} B's legal physical range is $1000..1019$. None of
$0$, $950$, $500$, $20$ lies in that range.

7.\ \textbf{(b).} B's true size is $20$, but leftover bounds $10$
makes VAs $10..19$ fault even though they are B's memory. (The
symmetric ``too-big bounds'' leak is (c); here bounds is too
\emph{small}.)

8.\ \textbf{(c) $4$.} One base register per CPU, because each CPU
independently runs (and relocates) one process at a time.
\end{Sol}

\begin{Qn}{Medium}{Lect 9}{CS-537 Spring 2001 Mid, Q12}{Where does the new 2000-byte job go?}
A system performs dynamic relocation of processes in physical memory.
A process that has just started needs $2000$ total bytes. The OS
currently has two free regions: physical $1000$--$2000$, and physical
$5000$--$7000$. What value ends up in the base register, and what
value in the bounds register? Write down any assumption about the
encoding of bounds.
\end{Qn}
\begin{Sol}
The first hole is only $1000$ bytes and cannot hold a $2000$-byte
address space. The second hole is exactly $2000$ bytes, so
$\mathrm{base}=5000$.

If bounds holds the \textbf{size}, $\mathrm{bounds}=2000$
(legal iff $0\le\mathrm{VA}<2000$).
If bounds holds the \textbf{limit physical address},
$\mathrm{bounds}=7000$.
Either encoding is the Lect 9 pair; the exam asked you to state which
one you used. After termination the OS returns $[5000,7000)$ to the
free list.
\end{Sol}

\begin{Qn}{Medium}{Lect 9}{CS-537 Fall 2004 Mid2, Q3}{Lose the bounds register, or lose the base?}
With dynamic relocation the hardware has a base register and a bounds
register.
\begin{enumerate}
  \item If you had only a base register, what functionality do you
        lose?
  \item If you had only a bounds register, what functionality do you
        lose?
  \item In a multiprogrammed system, if you could keep only one of
        them, which would you choose, and why?
\end{enumerate}
\end{Qn}
\begin{Sol}
Base implements \textbf{relocation} ($\mathrm{PA}=\mathrm{VA}+\mathrm{base}$).
Bounds implements \textbf{protection} (fault if $\mathrm{VA}\ge\mathrm{bounds}$).

Without bounds, a process can emit any VA and, after adding base, walk
off its slot into another process (or the kernel). Dynamic relocation
without isolation.

Without base, every process must be loaded at physical address $0$
--- i.e.\ you lose the ability to host more than one process in
memory at a time (unless you do static rewrite, which is not dynamic
relocation).

In a multiprogrammed system, \textbf{bounds} is the one you cannot
give up if the goal is isolation: even a single process at a fixed
origin must be prevented from touching kernel / device memory. In
practice of course you need both; this is a forced choice.
\end{Sol}

\begin{Qn}{Medium}{Lect 9}{CS-537 Fall 2004 Final, Q2}{Broken Dynamic Relocation: which placements actually fit?}
The Broken Dynamic Relocation system places jobs in a $100$-byte
physical memory. Process A is loaded at address $0$ and needs $10$
bytes; process B is loaded at address $30$ and needs $20$ bytes.
For each proposed placement of a new process C, say whether it
\textbf{Works} or is \textbf{Broken}.
\begin{enumerate}
  \item[(a)] C needs $30$ bytes, placed at address $0$
  \item[(b)] C needs $20$ bytes, placed at address $10$
  \item[(c)] C needs $30$ bytes, placed at address $40$
  \item[(d)] C needs $30$ bytes, placed at address $80$
  \item[(e)] C needs $10$ bytes, placed at address $90$
\end{enumerate}
\end{Qn}
\begin{Sol}
Occupied: A occupies $[0,10)$, B occupies $[30,50)$. Free:
$[10,30)$ ($20$ bytes) and $[50,100)$ ($50$ bytes).

(a) \textbf{Broken.} Address $0$ overlaps A (and does not leave A
room).

(b) \textbf{Works.} $[10,30)$ is free and exactly $20$ bytes.

(c) \textbf{Broken.} Address $40$ sits inside B's $[30,50)$.

(d) \textbf{Broken.} $[80,110)$ runs off the end of the $100$-byte
memory.

(e) \textbf{Works.} $[90,100)$ is free and $10$ bytes.
\end{Sol}

\begin{Qn}{Medium}{Lect 9}{CS-537 Fall 2011 Mid, Q10}{Infer base and bounds from three translations}
A base-and-bounds pair virtualizes a tiny $1$\,KB address space.
\begin{lstlisting}[numbers=none]
VA 0x308 (decimal 776) --> VALID: PA 14611
VA 0x255 (decimal 597) --> VALID: PA 14432
VA 0x3a1 (decimal 929) --> SEGMENTATION VIOLATION
\end{lstlisting}
What can we say about the base register? About the bounds register?
\end{Qn}
\begin{Sol}
Two valid pairs give the same base:
$14611-776=13835$ and $14432-597=13835$. So
$\mathrm{base}=13835$.

VA $776$ is valid and VA $929$ faults, so (size encoding)
$777\le\mathrm{bounds}\le 929$. The trace does not pin a single
integer; any bounds in that closed range is consistent. (The VAS
being $1$\,KB only tells us generated VAs lie in $0..1023$.)
\end{Sol}

\begin{Qn}{Medium}{Lect 9}{CS-537 Spring 2012 Mid, Q8}{Bounds is checked only on writes}
In this subtly-different base-and-bounds, the bounds register is
checked only on \textbf{writes} to memory, not on reads. What is the
impact?
\end{Qn}
\begin{Sol}
A process may \textbf{read} any physical address of the form
$\mathrm{VA}+\mathrm{base}$, including kernel memory, other
processes, and MMIO --- isolation is gone for confidentiality. Writes
are still confined, so (most) corruption of other address spaces is
prevented, but a read gadget is enough to steal secrets. Instruction
fetches are reads, so the process can also \emph{execute} bytes
outside its slot if it can jump there. The hardware saved one compare
on the load path and in exchange lost the whole point of bounds.
\end{Sol}

\begin{Qn}{Medium-Hard}{Lect 10}{CS-537 Spring 2018 Mid, Problem XV}{Which VAs hit a 20-byte grow-negative stack?}
A simple segmentation system supports two segments: one
positive-growing for code/heap, one negative-growing for a stack.
VAS $128$ bytes; physical memory $512$ bytes.

\begin{lstlisting}[numbers=none]
Segment 0 base (grows positive) : 0
Segment 0 limit                 : 20
Segment 1 base (grows negative) : 0x200  (decimal 512)
Segment 1 limit                 : 20
\end{lstlisting}

Which of the following are valid virtual accesses? A $=$ valid,
B $=$ not.
\begin{enumerate}
  \item[71.] \texttt{0x1d} (decimal $29$)
  \item[72.] \texttt{0x7b} (decimal $123$)
  \item[73.] \texttt{0x10} (decimal $16$)
  \item[74.] \texttt{0x5a} (decimal $90$)
  \item[75.] \texttt{0x0a} (decimal $10$)
\end{enumerate}
\end{Qn}
\begin{Sol}
Top bit of a $128$-byte VAS: SEG0 is $0..63$, SEG1 is $64..127$.
SEG0 size $20$ $\Rightarrow$ VAs $0..19$ legal.
SEG1 size $20$ $\Rightarrow$ the last $20$ bytes, VAs $108..127$,
are legal. (Where they sit in physical memory does not affect
validity.)

71 B ($29$ is in SEG0 but $\ge 20$);\quad
72 A ($123\in[108,127]$);\quad
73 A ($16<20$);\quad
74 B ($90$ is in SEG1 but below $108$);\quad
75 A ($10<20$).
\end{Sol}

\begin{Qn}{Medium-Hard}{Lect 10}{CS-537 Fall 2011 Mid, Q13}{Translate three VAs in a 16-byte two-segment space}
Two segments; the first bit of the VA selects the segment. Segment 0
holds code/heap (grows positive); segment 1 holds a stack (grows
negative). VAS $16$ bytes; physical memory $64$ bytes.

\begin{lstlisting}[numbers=none]
Segment 0 base (grows positive) : 0x30 (decimal 48)
Segment 0 limit                 : 7
Segment 1 base (grows negative) : 0x09 (decimal  9)
Segment 1 limit                 : 4
\end{lstlisting}

Translate, or mark a segmentation violation:
\begin{lstlisting}[numbers=none]
VA 0x0e (decimal 14) --> ?
VA 0x06 (decimal  6) --> ?
VA 0x0a (decimal 10) --> ?
\end{lstlisting}
\end{Qn}
\begin{Sol}
SEG0 is VAs $0..7$, SEG1 is $8..15$.
Grow-negative:
$\mathrm{PA}=\mathrm{base}_1-(16-\mathrm{VA})$, legal iff
$12\le\mathrm{VA}\le 15$ (last $4$ bytes).

VA $14$: SEG1, legal, $\mathrm{PA}=9-(16-14)=7$.

VA $6$: SEG0, $6<7$, legal, $\mathrm{PA}=48+6=54$.

VA $10$: SEG1, but $10\notin[12,15]$. \textbf{Segmentation
violation.}
\end{Sol}

\begin{Qn}{Medium}{Lect 10}{CS-537 Fall 2013 Mid, Q9}{McFlub's advantages of paging over segmentation}
McFlub: ``Paging has a number of advantages over segmentation. They
are... well... paging sounds better, looks better, and in general
feels better. That's about it. Oh yeah, paging also makes address
translation much faster than segmentation. And, it removes
fragmentation of some kind. Yup, that's it.''

Correct him.
\end{Qn}
\begin{Sol}
Real advantages of paging (Lect 12): \textbf{no external
fragmentation} (all frames identical, so free-space management is a
simple list of frame numbers); fully flexible layout of the VAS (heap
and stack grow independently without compaction); sharing and
protection at page granularity.

McFlub is wrong that paging is \emph{faster}: a linear page table
lives in memory, so a translation costs an extra load of the PTE
(often slower than a pair of segment-register adds). Paging is also
more complex, and it \emph{introduces internal} fragmentation
(roughly half a page per mapped region on average). Segmentation, for
its part, still suffers external fragmentation and may need
compaction --- that is the cost of variable-sized segments, not
something paging ``sort of'' also has.
\end{Sol}

\begin{Qn}{Easy-Medium}{Lect 10}{Mandar ISI 2015 End-sem, Q2(a)--(b); cf.\ Mandar ISI 2018 End-sem Q4(b)}{Internal vs.\ external fragmentation under paging and segmentation}
\begin{enumerate}
  \item[(a)] Define internal and external fragmentation.
  \item[(b)] For which of the following memory-management schemes is
        internal / external fragmentation possible: (i) paging,
        (ii) segmentation? Justify each in 2--3 lines.
\end{enumerate}
\end{Qn}
\begin{Sol}
(a) \textbf{Internal} fragmentation is wasted space \emph{inside} an
allocated unit (the allocator gave you a block bigger than you
needed). \textbf{External} fragmentation is wasted space
\emph{between} allocated units: free memory exists, but not as one
hole large enough to satisfy the next request.

(b) \textbf{Paging:} no external fragmentation (every free frame is
interchangeable). Internal fragmentation yes: the last page of a
segment is on average half-empty, and the page size is the unit.

\textbf{Segmentation:} external fragmentation yes (variable-sized
segments leave odd holes; may need compaction). Internal
fragmentation is minor / typically said to be absent, because a
segment can be exactly as large as the code/heap/stack it backs
(the leftover is the unused \emph{virtual} hole, which is not
backed by physical memory --- that was the whole point of
segmentation versus one big base-and-bounds slot).
\end{Sol}

\begin{Qn}{Easy-Medium}{Lect 11}{CS-537 Spring 2001 Mid, Q3 and Q8}{Fixed-size chunks vs.\ best-fit: which fragmentation?}
\begin{enumerate}
  \item Assume the heap is divided into fixed-size units of size $S$,
        and no request is larger than $S$. Any request is given one
        block of size $S$ if a free block exists. What kind of
        fragmentation can occur, and why is it bad?
  \item Assume the heap is managed in a best-fit manner. What does
        best-fit try to do? What kind of fragmentation can occur, and
        why is it bad?
\end{enumerate}
\end{Qn}
\begin{Sol}
(Q3) This is the ``easy case'' of Lect 11 (and the reason paging is
easy). The only waste is \textbf{internal} fragmentation: a request
of $s\le S$ still occupies a whole $S$-byte slot, wasting $S-s$
bytes inside the unit. Bad because that wasted space cannot be used
by anyone else, and it grows with $S$ (the page-size trade-off of
Lect 12 is the same idea).

(Q8) Best-fit searches for the \textbf{smallest} free chunk that is
still $\ge$ the request, splits off the remainder, and returns it to
the free list. The typical failure mode is \textbf{external}
fragmentation: leftover slivers are too small to satisfy later
requests, even though the \emph{sum} of free bytes would be enough.
Without coalescing, freeing two adjacent objects does not even
recreate a useful hole.
\end{Sol}

\begin{Qn}{Easy-Medium}{Lect 11}{CS-537 Spring 2001 Mid, Q10}{Why is stack allocation faster than \texttt{malloc}?}
Why is stack-based memory management usually much faster than
heap-based memory management? If it is so much faster, why do we ever
use the heap?
\end{Qn}
\begin{Sol}
A stack allocate/free is a \textbf{pointer bump} of the stack
pointer (and a corresponding bump back on return). No search, no
split, no coalesce, no header. A heap \texttt{malloc} must find a
chunk (first/next/best/worst fit), possibly split it, write a
size+magic header, and later \texttt{free} must walk/coalesce an
embedded free list --- all on a structure that lives in the heap
itself (you cannot \texttt{malloc} to build the allocator).

We still need a heap because stack lifetime is LIFO and scoped to a
call. Objects that outlive the allocating frame, that are shared
across frames, or that have a size known only at run time, have to
live on the heap (and the heap is grown with \texttt{sbrk}/\texttt{mmap}
when the free list is empty).
\end{Sol}

\begin{Qn}{Medium}{Lect 11}{CS-537 Spring 2002 Mid, Q9}{A hybrid allocator: 4\,KB fixed units plus a best-fit half}
A new \texttt{malloc}/\texttt{free} divides available memory in two.
One half is split into fixed-size $4$\,KB units; the other half is a
best-fit free list. If a request is $\le 4$\,KB and the fixed-size
half still has a free unit, a whole $4$\,KB unit is returned;
otherwise the request is satisfied (at the requested size) from the
best-fit half.
\begin{enumerate}
  \item[(a)] Assuming $32$\,KB of total memory, what series of
        allocation requests will most quickly lead to all of memory
        getting allocated, while requesting the least total amount of
        memory?
  \item[(b)] What type(s) of fragmentation occur with this allocator?
\end{enumerate}
\end{Qn}
\begin{Sol}
(a) The fixed half is $16$\,KB $=4$ units. Four requests of $1$ byte
each consume that entire half (least payload that still takes a
$4$\,KB unit). The best-fit half is then consumed by a single
$16$\,KB request (requesting less would leave a leftover, so would
not fill ``all of memory''). Total requested:
$4+16384=16388$ bytes, all $32$\,KB allocated.

(b) \textbf{Internal} fragmentation on the fixed half (up to
$4095$ wasted bytes per $1$-byte request). \textbf{External}
fragmentation on the best-fit half, plus internal waste of the
size+magic header on every allocated chunk. The two policies do not
help each other: a $5$\,KB request can never use a free $4$\,KB unit,
and a $1$-byte request will not use a $16$\,KB best-fit hole while
fixed-size units remain.
\end{Sol}

\begin{Qn}{Easy-Medium}{Lect 11}{CS-537 Fall 2011 Mid, Q9}{When is a memory leak OK?}
A user allocates with \texttt{malloc()} and frees with
\texttt{free()}. When a user doesn't call \texttt{free()}, it is
called a ``memory leak''. When is it OK to have a memory leak in your
program?
\end{Qn}
\begin{Sol}
It is OK if the process is about to \textbf{exit} (or \texttt{exec},
which replaces the address space). The OS tears down the address
space on terminate and returns every frame/segment to the free list
--- user-level leak bookkeeping is then wasted work. A long-running
server that leaks on every request will still exhaust the heap, grow
via \texttt{sbrk}, and eventually die. ``Never'' is the wrong exam
answer; ``only if the process is dying'' is the right one.
\end{Sol}

\begin{Qn}{Medium}{Lect 11}{Mandar ISI 2010 End-sem, Q1(a); expanded Mandar ISI 2024 End-sem, Q4(a)--(c)}{Two ways to \texttt{malloc} a 2-D array, and how many pages they touch}
Consider two fragments that allocate an \texttt{int} matrix with
\texttt{rows=10}, \texttt{cols=1024}. The machine is $32$-bit, an
\texttt{int} is $4$ bytes, and the page size is $2^{12}=4096$ bytes.

Fragment I (one contiguous payload):
\begin{lstlisting}[numbers=none]
if (NULL == (temp = (int *) malloc(rows*cols*sizeof(int))) ||
    NULL == (A = (int **) malloc(rows * sizeof(int *))))
    exit(1);
for (ii = 0; ii < rows; temp += cols, ii++) {
    A[ii] = temp;
    printf("%d %lu\n", ii, ((unsigned long) A[ii]) % 4096);
}
\end{lstlisting}

Fragment II (one \texttt{malloc} per row):
\begin{lstlisting}[numbers=none]
if (NULL == (A = (int **) malloc(rows * sizeof(int *)))) exit(1);
for (ii = 0; ii < rows; ii++) {
    if (NULL == (A[ii] = (int *) malloc(cols * sizeof(int))))
        exit(1);
    printf("%d %lu ", ii, ((unsigned long) A[ii]) % 4096);
}
\end{lstlisting}

(2010) Which method is preferable, and why?

(2024, assuming \texttt{temp} in Fragment I points at the first byte
of a page.)
\begin{enumerate}
  \item[(a)] What is the output of Fragment I?
  \item[(b)] After Fragment I, the matrix is summed in row-major
        order. How many distinct pages are referenced?
  \item[(c)] Fragment II prints
        \texttt{0 1920  1 1936  ...  9 2064}
        (an arithmetic progression with difference $16$). The matrix
        is again summed in row-major order. How many distinct pages
        are referenced?
\end{enumerate}
\end{Qn}
\begin{Sol}
(2010) \textbf{Fragment I is preferable.} One payload
\texttt{malloc} plus a small pointer vector gives a contiguous
\texttt{rows*cols} block, so a row is a stride of \texttt{cols}
integers, the whole matrix is one cache-/DMA-/page-friendly region,
and a single \texttt{free(temp)} (plus \texttt{free(A)}) reclaims it.
Fragment II makes \texttt{rows} independent heap chunks: more
headers, more external fragmentation, worse spatial locality, and
\texttt{rows} calls to \texttt{free}.

(2024a) Each row is $1024\times 4=4096$ bytes $=$ one page, and
\texttt{temp += cols} advances exactly one page. If \texttt{temp}
starts at a page boundary, every \texttt{A[ii] \% 4096} is $0$:
\begin{lstlisting}[numbers=none]
0 0
1 0
...
9 0
\end{lstlisting}

(2024b) The payload is $10$ contiguous pages. Traversing
\texttt{A[i][j]} also reads the $10$ pointers in \texttt{A}, which
occupy $40$ bytes --- one extra page. \textbf{11 distinct pages}
(10 data + 1 pointer vector). Code/stack pages of the summing loop
are not asked for.

(2024c) Each row is still $4096$ bytes but starts at page offset
$1920,1936,\ldots,2064$. A $4096$-byte block that does not start at
offset $0$ spans \textbf{two} pages. Consecutive starts differ by
$16$, matching a $16$-byte allocator header between $4096$-byte
payloads (chunk size $4112=4096+16$). The rows therefore occupy a
chain of overlapping page pairs
$Q$--$Q+1$, $Q+1$--$Q+2$, \ldots, $Q+9$--$Q+10$:
\textbf{11 data pages}, plus the pointer-vector page, i.e.\
\textbf{12 distinct pages} referenced by the summation.
\end{Sol}

\begin{Qn}{Medium}{Lect 12}{CS-537 Spring 2018 Mid, Problem XVI}{Which VAs are valid in this 8-entry linear page table?}
Linear page table. VAS $128$ bytes, physical memory $1024$ bytes,
page size $16$ bytes. High-order bit of a PTE is VALID; if $1$, the
rest is the PFN.

\begin{lstlisting}[numbers=none]
[0] 0x80000034
[1] 0x00000000
[2] 0x00000000
[3] 0x00000000
[4] 0x8000001e
[5] 0x80000017
[6] 0x80000011
[7] 0x8000002e
\end{lstlisting}

Which virtual addresses are valid? A $=$ valid, B $=$ not.
\begin{enumerate}
  \item[76.] \texttt{0x34} (decimal $52$)
  \item[77.] \texttt{0x44} (decimal $68$)
  \item[78.] \texttt{0x57} (decimal $87$)
  \item[79.] \texttt{0x18} (decimal $24$)
  \item[80.] \texttt{0x46} (decimal $70$)
\end{enumerate}
\end{Qn}
\begin{Sol}
$128$-byte VAS $\Rightarrow$ $7$-bit VA. $16$-byte page $\Rightarrow$
$4$-bit offset, $3$-bit VPN. Valid VPNs (high nibble of these
hex VAs, essentially) are $0,4,5,6,7$.

76 B (VPN $3$);\quad
77 A (VPN $4$);\quad
78 A (VPN $5$);\quad
79 B (VPN $1$);\quad
80 A (VPN $4$).
\end{Sol}

\begin{Qn}{Medium-Hard}{Lect 12}{CS-537 Fall 2021 Mid, Q13--Q16}{DOS paging: PTE count, physical size, and translate \texttt{0x64}}
DOS uses a linear page table.
\begin{enumerate}
  \item[13.] VAS $128$ bytes, page size $2$ bytes. How many entries
        in each page table?
        (a) $128$\quad (b) $64$\quad (c) $32$\quad (d) $16$\quad
        (e) None of the above
  \item[14.] VAS still $128$ bytes, page size now $32$ bytes. The
        page table is
\begin{lstlisting}[numbers=none]
0x8000000c
0x00000000
0x00000000
0x80000006
\end{lstlisting}
        Leftmost bit $=$ valid, rightmost $4$ bits $=$ PFN. If $4$
        bits are needed for the PFN, how big is physical memory?
        (a) $128$ B\quad (b) $256$ B\quad (c) $512$ B\quad
        (d) $1024$ B\quad (e) None of the above
  \item[15.] Using the page table of Q14, translate VA \texttt{0x64}.
        (a) \texttt{0x46}\quad (b) \texttt{0x04}\quad
        (c) \texttt{0xc4}\quad (d) \texttt{0x64}\quad
        (e) None of the above
  \item[16.] When allocating a physical page, DOS selects a
        \emph{random} free frame (instead of the first free one).
        This approach will:
        (a) Make the page table larger\quad
        (b) Make physical memory more fragmented (externally)\quad
        (c) Slow down address translation\quad
        (d) Make heap allocation fail more often\quad
        (e) None of the above
\end{enumerate}
\end{Qn}
\begin{Sol}
13.\ \textbf{(b) $64$.} $128/2=64$ virtual pages $=$ $64$ linear
PTEs.

14.\ \textbf{(c) $512$ bytes.} $4$-bit PFN $\Rightarrow$ $16$
frames; $16\times 32=512$.

15.\ \textbf{(c) \texttt{0xc4}.}
\texttt{0x64}$=1100100_2$. Page size $32=2^5$ $\Rightarrow$ $5$-bit
offset, $2$-bit VPN. VPN $=11_2=3$, offset $=00100_2=4$.
PTE $[3]=\texttt{0x80000006}$ is valid, PFN $=6=0110_2$.
$\mathrm{PA}=0110\,00100_2=\texttt{0xc4}$.

16.\ \textbf{(e).} Under paging, frames are interchangeable, so
there is no external fragmentation to make worse; the linear table's
size is determined by the VAS and page size, not by which PFN is
chosen; translation is still $\mathrm{PA}=\mathrm{PFN}\times 32+\mathrm{offset}$;
user-level \texttt{malloc} is a VAS issue, not a frame-picking issue.
\end{Sol}

\begin{Qn}{Easy-Medium}{Lect 12}{CS-537 Spring 2002 Mid, Q3}{32-bit VAS, 4\,KB pages: split, then translate \texttt{0x00010033}}
Pure paging, $32$-bit address space, $4$\,KB page size.
\begin{enumerate}
  \item[(a)] How many bits specify the VPN, and how many specify the
        offset?
  \item[(b)] Translate logical address \texttt{0x00010033}, given
        that each logical page maps to the physical page one number
        higher (logical page $10$ $\to$ physical page $11$, etc.).
\end{enumerate}
\end{Qn}
\begin{Sol}
(a) $4$\,KB $=2^{12}$ $\Rightarrow$ $12$-bit offset, $20$-bit VPN.

(b) \texttt{0x00010033}: VPN \texttt{0x10}$=16$, offset
\texttt{0x033}. Logical page $16$ maps to physical page $17=\texttt{0x11}$,
so $\mathrm{PA}=\texttt{0x00011033}$.
\end{Sol}

\begin{Qn}{Medium}{Lect 12}{CS-537 Spring 2009 Mid, Q3(a)--(e)}{Linear page-table size for a 32-bit, 1\,KB-page system}
$32$-bit virtual address space, $1$\,KB page size. Each PTE is $4$
bytes. Linear page table:
\begin{enumerate}
  \item[(a)] Bits in the offset?
  \item[(b)] Bits in the VPN?
  \item[(c)] How many entries in the table?
  \item[(d)] Total size of the table?
  \item[(e)] In a live system, how much memory is occupied by page
        tables? What factors affect this?
\end{enumerate}
\end{Qn}
\begin{Sol}
(a) $1$\,KB $=2^{10}$ $\Rightarrow$ \textbf{10} offset bits.

(b) $32-10=$ \textbf{22} VPN bits.

(c) One PTE per virtual page: $2^{22}$ entries.

(d) $2^{22}\times 4=2^{24}=$ \textbf{16\,MB} per process.

(e) Roughly $16$\,MB $\times$ (number of processes), unless the OS
shares a table (or later uses a multi-level table, which is Lect 14
and not required here). A $100$-process system would carry about
$1.6$\,GB of linear page tables --- the Lect 12 motivation for not
stopping at a linear array.
\end{Sol}

\begin{Qn}{Medium}{Lect 12}{CS-537 Fall 2011 Mid, Q14}{100 processes, 30-bit VAS, 1\,KB pages: how much page-table memory?}
A system uses a simple linear page table. VAS $1$\,GB ($30$ bits);
page size $1$\,KB; each PTE holds only a valid bit and the PFN. The
system has at most $2^{15}$ physical pages ($32$\,MB of physical
memory). How much memory is used for page tables when $100$ processes
are running?
\end{Qn}
\begin{Sol}
Offset $10$ bits $\Rightarrow$ VPN $20$ bits $\Rightarrow$ $2^{20}$
PTEs per process. A PFN of $15$ bits plus a valid bit fits in $16$
bits, so a PTE is $2$ bytes.

One table: $2^{20}\times 2=2$\,MB.
One hundred processes: \textbf{200\,MB} of linear page tables ---
several times the $32$\,MB of DRAM they are describing. That is the
Lect 12 ``page tables can be bigger than the memory they map''
warning.
\end{Sol}

\begin{Qn}{Medium-Hard}{Lect 12}{CS-537 Spring 2012 Mid, Q11}{The valid bit moved to the low end of the PTE}
A simple PTE is usually a valid bit followed by a PFN. In this
subtly-different format the order is switched: PFN, then a valid bit.
VAS $128$ bytes, $32$-byte pages, linear table:
\begin{lstlisting}[numbers=none]
0xFE
0x0F
0xFE
0x09
\end{lstlisting}
For each VA, say whether it is valid and, if so, give the physical
address.
\begin{lstlisting}[numbers=none]
VA 0x065 (decimal 101) --> ?
VA 0x00c (decimal  12) --> ?
VA 0x026 (decimal  38) --> ?
VA 0x058 (decimal  88) --> ?
\end{lstlisting}
\end{Qn}
\begin{Sol}
$128$-byte VAS, $32$-byte pages: $2$-bit VPN, $5$-bit offset, $4$
PTEs. Each PTE is one byte: top $7$ bits $=$ PFN, bottom bit $=$
valid.

\begin{center}
\begin{tabular}{c|c|c}
PTE & valid bit & PFN \\
\hline
$[0]$ \texttt{0xFE} $=11111110_2$ & 0 & \texttt{0x7F} \\
$[1]$ \texttt{0x0F} $=00001111_2$ & 1 & $7$ \\
$[2]$ \texttt{0xFE} $=11111110_2$ & 0 & \texttt{0x7F} \\
$[3]$ \texttt{0x09} $=00001001_2$ & 1 & $4$ \\
\end{tabular}
\end{center}

VA $101=1100101_2$: VPN $3$, offset $5$. Valid, PFN $4$.
$\mathrm{PA}=4\times 32+5=133=\texttt{0x85}$.

VA $12=0001100_2$: VPN $0$, offset $12$. PTE $[0]$ invalid:
\textbf{fault}.

VA $38=0100110_2$: VPN $1$, offset $6$. Valid, PFN $7$.
$\mathrm{PA}=7\times 32+6=230=\texttt{0xE6}$.

VA $88=1011000_2$: VPN $2$, offset $24$. PTE $[2]$ invalid:
\textbf{fault}.
\end{Sol}

\begin{Qn}{Medium}{Lect 12}{CS-537 Fall 2004 Mid2, Q9}{PTE width follows physical memory, not the VAS}
A $20$-bit virtual address, $4$\,KB page size. Assume $2$ extra bits
of information are needed in each PTE beyond the PFN.
\begin{enumerate}
  \item The machine has a maximum of $1$\,GB of physical memory. How
        big is each page table?
  \item A different machine has a maximum of $64$\,KB of physical
        memory. How big is each page table?
  \item Which of the two cases is worse: having more physical memory
        than a process can address, or less? Why?
\end{enumerate}
\end{Qn}
\begin{Sol}
Offset $12$ bits $\Rightarrow$ VPN $8$ bits $\Rightarrow$ $256$ PTEs
(the table's \emph{length} depends only on the VAS).

(a) $1$\,GB $/\,4$\,KB $=2^{18}$ frames, so the PFN is $18$ bits.
Plus $2$ flag bits $=20$ bits/PTE. Packed into $4$-byte PTEs (the
natural rounding-up), the table is $256\times 4=$ \textbf{1\,KB}.
(If packed exactly, $256\times 20/8=640$ bytes.)

(b) $64$\,KB $/\,4$\,KB $=16$ frames, PFN $4$ bits $+2$ flags $=6$
bits/PTE. Table $\approx 256\times 6/8=$ \textbf{192 bytes}
(or $256$ bytes if rounded to a whole byte per PTE).

(c) \textbf{Less} physical memory than the VAS is worse for the
process: many valid virtual pages cannot be resident at once, so the
system must swap (Lect 15) or refuse allocations. Extra physical
memory the process cannot address is merely unused by \emph{this}
process; the linear table does not grow with DRAM --- only the PFN
field does, slightly.
\end{Sol}

\begin{Qn}{Easy-Medium}{Lect 12}{CS-537 Fall 2004 Mid2, Long Q2(a)}{Too-big vs.\ too-small a page}
If we have a single page size and it is quite big, what are the
possible negative consequences? In contrast, what negative
consequences arise if our page size is too small?
\end{Qn}
\begin{Sol}
\textbf{Too big:} more \textbf{internal fragmentation} (half a huge
page wasted, on average, at the end of each mapped region); a
sparse VAS still pulls in a lot of unused bytes with every mapped
page.

\textbf{Too small:} the linear page table explodes ($\#\mathrm{PTEs}=\mathrm{VAS}/\mathrm{pagesize}$),
and every I/O / copy moves less data per operation so overhead
dominates. Lect 12's trade-off is exactly this pair; TLB reach
(Lect 13) is the later, hardware-cache version of the same tension.
\end{Sol}

\begin{Qn}{Medium}{Lect 12}{Mandar ISI 2012 End-sem, Q3(a)--(b)}{Page-size disadvantages, and an optimum for a 64\,KB-max VAS}
\begin{enumerate}
  \item[(a)] What is the disadvantage of a large page/frame size? Of
        a small page/frame size?
  \item[(b)] Process address-space sizes are uniformly distributed
        between $1$ byte and $64$\,KB. The page table is the same
        size for all processes (linear). Each PTE is $4$ bytes. What
        is the optimum page size?
\end{enumerate}
\end{Qn}
\begin{Sol}
(a) Large pages: internal fragmentation. Small pages: a huge linear
page table, and more bookkeeping / I/O per byte of useful data.

(b) Let the page size be $p$. A linear table sized for the maximum
$64$\,KB VAS has $(2^{16}/p)$ PTEs and occupies
$4\cdot 2^{16}/p=2^{18}/p$ bytes. Average internal fragmentation is
$p/2$ (half a page). Total expected overhead
\[
E(p)=\frac{p}{2}+\frac{2^{18}}{p}.
\]
Set $E'(p)=0$: $\tfrac12=2^{18}/p^2$, so
$p=\sqrt{2^{19}}=\sqrt{524288}\approx 724$ bytes.
(A nearby power of two, $512$ or $1024$, is what hardware would
actually use. Using the \emph{average} process size $32$\,KB in
place of the max $64$\,KB yields the equally common textbook value
$p=\sqrt{2\cdot 4\cdot 2^{15}}=512$. The problem's wording ``the
size of the page table is the same for all processes'' prefers the
$64$\,KB / $p\approx 724$ reading.)
\end{Sol}

\begin{Qn}{Medium}{Lect 12}{Mandar ISI 2014 End-sem, Q2(a)}{Spare bits in a 4-byte PTE, and what they usually store}
A $32$-bit address bus, page/frame size $2$\,KB, PTE size $4$ bytes.
\begin{enumerate}
  \item[(i)] How many bits can be used to store additional information
        about a page (besides the physical frame number) in each PTE?
  \item[(ii)] List and briefly explain any three types of such
        additional information that are usually stored in PTEs.
\end{enumerate}
\end{Qn}
\begin{Sol}
(i) $2$\,KB $=2^{11}$ $\Rightarrow$ $11$-bit offset. A $32$-bit
physical address then has a $21$-bit PFN. A $4$-byte PTE is $32$
bits, so $32-21=$ \textbf{11 spare bits}.

(ii) Any three of the Lect 12 PTE flags:
\begin{itemize}
  \item \textbf{Valid:} this VPN is part of the process's address
        space (else segmentation fault, not a page fault).
  \item \textbf{Protection} (R/W/X): the allowed accesses; a write
        to a read-only page traps (used for CoW after \texttt{fork}
        as well).
  \item \textbf{Present:} the page is in a frame right now (else the
        PTE holds a swap-block number; the access is a page fault ---
        introduced here, policy in Lect 15).
  \item \textbf{Dirty:} the page has been written since it was
        brought in, so it must be written back if evicted.
  \item \textbf{Reference/Accessed:} the page has been touched
        recently; used later to approximate LRU.
\end{itemize}
\end{Sol}

\clearpage


\end{document}
